\documentclass[14pt,a4paper,oneside]{extbook}
\usepackage[utf8]{vietnam}
\usepackage{amsmath,amssymb,mathrsfs,fancyhdr,enumerate,multirow,makecell,venndiagram}
\usepackage{tkz-euclide,tikz-3dplot,tikz,tkz-tab}
\usetikzlibrary{arrows,calc,intersections,angles,quotes,backgrounds,shapes.geometric,patterns,through}
\usepackage{yhmath}
\usepackage{pgfplots}
\pgfplotsset{compat=1.9}
\usepgfplotslibrary{fillbetween}
\usepackage[tikz]{bclogo}
\usepackage[top=1.8cm, bottom=1.8cm, left=1.8cm, right=1.8cm, headheight=16pt]{geometry}
\usepackage[hidelinks,unicode]{hyperref}
\usepackage{qrcode}

\usepackage[many]{tcolorbox}

% Gói ex_test (Chọn [dethi] để in đề hoặc [loigiai] để in cả lời giải)
%\usepackage[loigiai]{ex_test}
\usepackage[dethi]{ex_test}

\def\colorEX{\color{blue}}

\usepackage{fontawesome}
\usepackage{framed}
\usepackage{polynom}

\colorlet{tcbcol@back}{tcbcolback}
\colorlet{tcbcol@frame}{tcbcolframe}
\newtcolorbox[auto counter]{danggg}[1]{
	fonttitle=\bfseries\upshape,
	colframe=blue,colback=blue!10,
	sharp corners, breakable, 
	title={\color{white}\faFolderO \ DẠNG~\thetcbcounter.\ #1}
	\addcontentsline{toc}{subsection}{Dạng~\thetcbcounter. #1}
}
\newenvironment{dang}[1]{
	\begin{danggg}{#1}
		\textit{Phương pháp giải:~}
	}{\end{danggg}}
\newenvironment{tomtat}{
	\begin{tcolorbox}[boxrule=1pt,arc=0mm,breakable]
	}{\end{tcolorbox}}

\newtheorem{bt}{Bài}
\newtheorem{vd}{Ví dụ}
\theorembodyfont{\it}
\newtheorem{dl}{ Định lí}
\theoremseparator{.}
\theorembodyfont{\rm}
\newtheorem{hq}{ Hệ quả}
\newtheorem{tc}{ Tính chất}
\newtheorem{dn}{ Định nghĩa}
\theoremstyle{nonumberplain}
\newtheorem{kn}{ Khái niệm}
\theoremseparator{.}
\theorembodyfont{\it}
\usepackage{stackengine}
\usepackage{scalerel}

\newcommand\dangersign[1][2ex]{%
	\renewcommand\stacktype{L}%
	\scaleto{\stackon[1.3pt]{\color{red}$\triangle$}{\fontsize{4}{0}\selectfont !}}{#1}%
}
\renewcommand{\parallel}{\mathbin{\!\raisebox{0.25pt}/\mkern-2mu\raisebox{0.25pt}/\!}}
\renewcommand{\backsim}{\ \mathord{\raisebox{0.025\baselineskip}{\rotatebox{90}{\fontfamily{qag}\selectfont S}}}\ }
\theoremseparator{}
\newtheorem{note}{{\dangersign[3ex]}}
\newtheorem{nx}{Nhận xét.}
\newtcolorbox{luuy}[1][]{colback=yellow!20,title={\bfseries Chú ý},#1}

% Tiêu đề chương / bài
\setcounter{secnumdepth}{4}
\setcounter{tocdepth}{2}
\usepackage{titlesec}
\usepackage{titletoc}

\makeatletter
\renewcommand\thesection{\@arabic\c@section}
\makeatother

\titleformat{\section}{\centering\Large\bfseries\color{blue!80!black}}{Chủ đề~\thesection.}{1ex}{
	\setcounter{vd}{0}
	\setcounter{ex}{0}
	\setcounter{bt}{0}
	\setcounter{dl}{0}
	\setcounter{dn}{0}
	\setcounter{hq}{0}
	\setcounter{tc}{0}
	\MakeUppercase
}
\titleformat{\subsection}{\large\bfseries\color{blue!70!black}}{}{0ex}{}
\usepackage[final]{pdfpages}

% Cài đặt riêng
\def\hideboxEX{}
\renewcommand{\boxEX}[2][15pt]{%
	\setbox1=\hbox{#2\strut}
	\def\sizeboxEX{#1}%
	\ \EXbox{\hideboxEX#2\strut}\ %
}

\renewcommand\arraystretch{1.4} 
\renewcommand{\tabcolsep}{4mm}

% Dòng kẻ làm bài TikZ chuẩn 8mm
\renewcommand{\dotlineEX}[1]{%
  \par\vspace{2mm}%
  \noindent%
  \begin{tikzpicture}[x=1mm, y=1mm]%
    \foreach \i in {1,...,#1} {%
      \draw[gray!60, loosely dotted, line width=0.6pt] (0, -\i*8) -- (\linewidth, -\i*8);%
    }%
  \end{tikzpicture}%
}

% Nút Casio
\renewcommand\boxed[1]{%
	\tikz[baseline=(key.base)]
	\node[%
	draw,
	fill=black!40,
	drop shadow={shadow xshift=0.25ex,shadow yshift=-0.25ex,fill=black!90,opacity=0.75},
	rectangle,minimum width=0.6cm,minimum height=0.6cm,
	rounded corners=2pt,
	inner sep=1pt,
	line width=0.5pt,
	font=\sffamily
	](key) {\color{white}#1\strut}
	;
}

\newcommand{\hoac}[1]{\left[\begin{aligned}#1\end{aligned}\right.}
\newcommand{\heva}[1]{\left\{\begin{aligned}#1\end{aligned}\right.}

\pagestyle{fancy}
\fancyhf{} % Xóa sạch toàn bộ header/footer cũ
\fancyhead[L]{\small\itshape\color{gray} Tuyển tập Chuyên đề Vận dụng Toán 2 theo từng bài}
\fancyhead[R]{\small\bfseries\color{blue!80!black} Lớp Toán Thầy Lĩnh -- loptoanthaylinh.vn}
\fancyfoot[C]{\bfseries\thepage}
\renewcommand{\headrulewidth}{0.4pt}

\newcommand{\caidat}{2}
\newcommand{\indapan}[2]{
	\begin{center}
		\color{blue!50!green} \fbox{\fontfamily{qag}\bfseries\selectfont BẢNG ĐÁP ÁN}
	\end{center}
	{\fontfamily{lmtt}\bfseries\selectfont \inputansbox[\caidat]{#1}{#2}}
}

\everymath{\color{blue}}

% -------------------------------------------------------------------------
% KHUNG TIÊU ĐỀ PHIẾU KIỂM TRA KÈM MÃ QR CODE CHO TỪNG BÀI
% Cú pháp: \TaoPhieuKiemTra{Chủ đề}{Số bài}{Tên bài}{Link QR thi online}
% -------------------------------------------------------------------------
\newcommand{\TaoPhieuKiemTra}[4]{%
  \refstepcounter{subsection}%
  \addcontentsline{toc}{subsection}{Bài #2. #3}%
  \noindent
  \begin{tcolorbox}[
    colback=white,
    colframe=blue!85!black,
    arc=2.5mm,
    boxrule=1pt,
    left=2.5mm, right=2.5mm, top=1.8mm, bottom=1.8mm
  ]
    % LOGO TRUNG TÂM & HOTLINE | TIÊU ĐỀ BÀI KIỂM TRA | MÃ QR CODE
    \begin{minipage}[c]{0.27\textwidth}
      \centering
      \begin{tcolorbox}[colback=blue!8, colframe=blue!80!black, boxrule=0.6pt, arc=1.5mm, halign=center, top=1.2mm, bottom=1.2mm, left=1mm, right=1mm]
        {\bfseries\color{blue!95!black}\small LỚP TOÁN THẦY LĨNH}\\[1.5pt]
        {\scriptsize\color{blue!70!black}\faGlobe~loptoanthaylinh.vn}\\[1.5pt]
        {\scriptsize\bfseries\color{blue!85!black}\faPhone~Hotline: 0866976878}
      \end{tcolorbox}
    \end{minipage}%
    \hfill
    \begin{minipage}[c]{0.56\textwidth}
      \centering
      {\footnotesize\bfseries\color{blue!80!black}\MakeUppercase{#1}}\\[2pt]
      {\large\bfseries\color{red!85!black}\MakeUppercase{BÀI #2: #3}}\\[2pt]
      {\scriptsize\itshape CHUYÊN ĐỀ VẬN DỤNG -- Thời gian làm bài: 35 phút -- Cấu trúc 3 phần chuẩn}
    \end{minipage}%
    \hfill
    \begin{minipage}[c]{0.14\textwidth}
      \centering
      \qrcode[height=1.35cm]{#4}\\[2pt]
      {\tiny\bfseries\color{blue!90!black}\faQrcode~Quét QR thi}
    \end{minipage}
  \end{tcolorbox}
  \vspace{1.5mm}
}

\begin{document}

\tableofcontents
\newpage

% =========================================================================
% CHỦ ĐỀ 1 - ÔN TẬP VÀ BỔ SUNG
% =========================================================================
\section{Ôn tập và bổ sung}

\TaoPhieuKiemTra{Chủ đề 1 -- Ôn tập và bổ sung}{01}{Ôn tập các số đến 100}{https://loptoanthaylinh.vn/thi/chuyen-de/2/toan2_cd_b01}
\textbf{Phần 1. Trắc nghiệm lựa chọn}
\setcounterpageref{ex}{0}

\begin{ex} %[2D1V01]
Tìm một số có hai chữ số, biết chữ số hàng chục là $6$ và chữ số hàng đơn vị lớn hơn chữ số hàng chục $3$ đơn vị.
\choice
{$96$}
{$63$}
{\True $69$}
{$68$}
\loigiai{
Chữ số hàng chục là $6$. Chữ số hàng đơn vị là $6 + 3 = 9$. Vậy số đó là $69$.
\textbf{Đáp án của câu hỏi là C}
}
\end{ex}

\begin{ex} %[2D1V01]
Từ ba chữ số $2, 0, 7$, lập được tất cả bao nhiêu số có hai chữ số khác nhau?
\choice
{\True $4$ số}
{$6$ số}
{$5$ số}
{$3$ số}
\loigiai{
Chữ số hàng chục có thể chọn là 2 hoặc 7 (không thể chọn 0).\n- Hàng chục là 2: lập được $20, 27$.\n- Hàng chục là 7: lập được $70, 72$.\nCó tất cả $4$ số.
\textbf{Đáp án của câu hỏi là A}
}
\end{ex}

\begin{ex} %[2D1V01]
Một số có hai chữ số mà tổng hai chữ số bằng $12$, hiệu giữa chữ số hàng chục và chữ số hàng đơn vị bằng $0$. Số đó là
\choice
{$57$}
{$48$}
{$75$}
{\True $66$}
\loigiai{
Hiệu hai chữ số bằng 0 nghĩa là chữ số hàng chục bằng chữ số hàng đơn vị.\nTổng hai chữ số bằng $12$, mà $6 + 6 = 12$. Vậy số đó là $66$.
\textbf{Đáp án của câu hỏi là D}
}
\end{ex}

\begin{ex} %[2D1V01]
Cho các số: $25, 62, 52, 19, 70$. Khi sắp xếp các số này theo thứ tự từ lớn đến bé, số đứng ở vị trí thứ ba là
\choice
{$62$}
{\True $52$}
{$25$}
{$70$}
\loigiai{
Sắp xếp từ lớn đến bé: $70 > 62 > 52 > 25 > 19$. Số đứng thứ ba là $52$.
\textbf{Đáp án của câu hỏi là B}
}
\end{ex}

\begin{ex} %[2D1V01]
Số tròn chục liền trước của số $80$ là
\choice
{$79$}
{$81$}
{\True $70$}
{$60$}
\loigiai{
Các số tròn chục liên tiếp là $60, 70, 80, 90$. Số tròn chục liền trước 80 là $70$.
\textbf{Đáp án của câu hỏi là C}
}
\end{ex}

\begin{ex} %[2D1V01]
Có bao nhiêu số có hai chữ số mà chữ số tận cùng là $5$?
\choice
{\True $9$ số}
{$10$ số}
{$8$ số}
{$5$ số}
\loigiai{
Các số đó là: $15, 25, 35, 45, 55, 65, 75, 85, 95$. Có tất cả $9$ số.
\textbf{Đáp án của câu hỏi là A}
}
\end{ex}

\begin{ex} %[2D1V01]
Tìm số nhỏ nhất có hai chữ số mà chữ số hàng chục hơn chữ số hàng đơn vị $5$ đơn vị.
\choice
{$61$}
{$72$}
{$49$}
{\True $50$}
\loigiai{
Để số đó nhỏ nhất, chữ số hàng đơn vị phải nhỏ nhất, tức là bằng $0$.\nKhi đó chữ số hàng chục là $0 + 5 = 5$. Vậy số nhỏ nhất là $50$.
\textbf{Đáp án của câu hỏi là D}
}
\end{ex}

\begin{ex} %[2D1V01]
Cho dãy số: $3, 7, 11, 15, \dots, 27$. Hai số thích hợp cần điền vào chỗ chấm là
\choice
{$18, 22$}
{\True $19, 23$}
{$19, 24$}
{$17, 21$}
\loigiai{
Quy luật dãy số: mỗi số đứng sau bằng số đứng trước cộng thêm $4$.\n$15 + 4 = 19$; $19 + 4 = 23$; $23 + 4 = 27$. Vậy điền $19, 23$.
\textbf{Đáp án của câu hỏi là B}
}
\end{ex}

\begin{ex} %[2D1V01]
Một rổ táo có nhiều hơn $35$ quả nhưng ít hơn $42$ quả. Nếu đếm theo chục thì được $4$ chục. Hỏi trong rổ có bao nhiêu quả táo?
\choice
{$36$ quả táo}
{$38$ quả táo}
{\True $40$ quả táo}
{$41$ quả táo}
\loigiai{
$4$ chục bằng $40$. Số $40$ thỏa mãn $35 < 40 < 42$. Vậy có $40$ quả táo.
\textbf{Đáp án của câu hỏi là C}
}
\end{ex}

\begin{ex} %[2D1V01]
Hiệu giữa số lớn nhất có hai chữ số và số tròn chục nhỏ nhất có hai chữ số là
\choice
{\True $89$}
{$90$}
{$80$}
{$79$}
\loigiai{
Số lớn nhất có hai chữ số là $99$. Số tròn chục nhỏ nhất có hai chữ số là $10$.\nHiệu là: $99 - 10 = 89$.
\textbf{Đáp án của câu hỏi là A}
}
\end{ex}

\begin{ex} %[2D1V01]
Bạn Nam viết các số có hai chữ số giống nhau. Nam đã viết được bao nhiêu số?
\choice
{$10$ số}
{$8$ số}
{$11$ số}
{\True $9$ số}
\loigiai{
Các số có hai chữ số giống nhau là: $11, 22, 33, 44, 55, 66, 77, 88, 99$. Có tất cả $9$ số.
\textbf{Đáp án của câu hỏi là D}
}
\end{ex}

\begin{ex} %[2D1V01]
Trong các số sau: $48, 56, 84, 39, 93$, số có giá trị của chữ số hàng chục lớn nhất là
\choice
{$84$}
{\True $93$}
{$56$}
{$48$}
\loigiai{
Chữ số hàng chục của các số lần lượt là: 4, 5, 8, 3, 9. Chữ số lớn nhất là 9 (trong số 93).
\textbf{Đáp án của câu hỏi là B}
}
\end{ex}

\textbf{Phần 2. Trắc nghiệm đúng sai}
\setcounterpageref{ex}{0}

\begin{ex} %[2D1V01]
Xét tính đúng sai của các nhận định về cấu tạo số trong phạm vi 100:
\choiceTF
{\True Số bé nhất có hai chữ số khác nhau là 10.}
{Có tất cả 10 số tròn chục có hai chữ số.}
{\True Trong số 77, giá trị của chữ số 7 hàng chục hơn giá trị của chữ số 7 hàng đơn vị là 63 đơn vị.}
{\True Số 99 là số liền trước của số 100.}
\loigiai{
\begin{itemize}
    \item Ý a) Đúng: Vì chữ số hàng chục là 1, chữ số hàng đơn vị là 0 khác 1 nên 10 là số bé nhất có hai chữ số khác nhau.
    \item Ý b) Sai: Chỉ có 9 số tròn chục có hai chữ số từ 10 đến 90.
    \item Ý c) Đúng: Giá trị chữ số 7 hàng chục là 70, hàng đơn vị là 7. Ta có $70 - 7 = 63$.
    \item Ý d) Đúng: Ta có $100 - 1 = 99$, vậy 99 là số liền trước của 100.
\end{itemize}
\begin{itemize}
    \item a) \textbf{Đúng}.
    \item b) \textbf{Sai}.
    \item c) \textbf{Đúng}.
    \item d) \textbf{Đúng}.
\end{itemize}
}
\end{ex}

\begin{ex} %[2D1V01]
Cho các chữ số $3, 5, 8$. Lập các số có hai chữ số từ ba chữ số trên:
\choiceTF
{\True Lập được tất cả 6 số có hai chữ số khác nhau.}
{Số lớn nhất có thể lập được là số 85.}
{\True Tổng của số bé nhất và số lớn nhất trong các số có hai chữ số khác nhau lập được bằng 120.}
{Trong các số lập được có đúng 2 số chẵn.}
\loigiai{
\begin{itemize}
    \item Ý a) Đúng: Các số là $35, 38, 53, 58, 83, 85$ (tất cả 6 số).
    \item Ý b) Sai: Nếu cho phép chữ số giống nhau thì số lớn nhất là 88; nếu khác nhau là 85.
    \item Ý c) Đúng: Số bé nhất là 35, lớn nhất là 85: $35 + 85 = 120$.
    \item Ý d) Sai: Chỉ có 2 số có đuôi 8 là 38 và 58 là số chẵn.
\end{itemize}
\begin{itemize}
    \item a) \textbf{Đúng}.
    \item b) \textbf{Sai}.
    \item c) \textbf{Đúng}.
    \item d) \textbf{Sai}.
\end{itemize}
}
\end{ex}

\textbf{Phần 3. Trả lời ngắn}
\setcounterpageref{ex}{0}

\begin{ex} %[2D1V01]
Tìm số nhỏ nhất có hai chữ số mà tích hai chữ số của nó bằng 12.
\shortans{26}
\loigiai{
\textbf{Đáp số:} $26$\\
Các cặp số có tích bằng 12 là $2 \times 6 = 12$ và $3 \times 4 = 12$. Số nhỏ nhất là 26.
}
\end{ex}

\begin{ex} %[2D1V01]
Có bao nhiêu số có hai chữ số mà tổng hai chữ số của nó bằng 8?
\shortans{8}
\loigiai{
\textbf{Đáp số:} $8$\\
Các số đó là: 17, 26, 35, 44, 53, 62, 71, 80 (tất cả 8 số).
}
\end{ex}

\begin{ex} %[2D1V01]
Tìm số lớn nhất có hai chữ số khác nhau mà chữ số hàng chục gấp đôi chữ số hàng đơn vị.
\shortans{84}
\loigiai{
\textbf{Đáp số:} $84$\\
Hàng chục gấp đôi hàng đơn vị, để lớn nhất ta chọn hàng chục là 8, khi đó hàng đơn vị là $8 : 2 = 4$. Số đó là 84.
}
\end{ex}

\begin{ex} %[2D1V01]
Bạn An có một số que tính, nếu bó thành từng bó 10 que thì được 4 bó và còn thừa 7 que. Hỏi An có tất cả bao nhiêu que tính?
\shortans{47}
\loigiai{
\textbf{Đáp số:} $47$\\
$4$ bó là $40$ que. Thêm 7 que là $40 + 7 = 47$ que.
}
\end{ex}

\begin{ex} %[2D1V01]
Tìm một số biết rằng lấy số lớn nhất có hai chữ số trừ đi số đó thì được số tròn chục nhỏ nhất.
\shortans{89}
\loigiai{
\textbf{Đáp số:} $89$\\
Số lớn nhất có hai chữ số là 99. Số tròn chục nhỏ nhất là 10. Số đó là: $99 - 10 = 89$.
}
\end{ex}

\begin{ex} %[2D1V01]
Từ số 15 đến số 45 có tất cả bao nhiêu số tự nhiên liên tiếp?
\shortans{31}
\loigiai{
\textbf{Đáp số:} $31$\\
Số lượng số hạng: $(45 - 15) + 1 = 31$ số.
}
\end{ex}

\newpage

\TaoPhieuKiemTra{Chủ đề 1 -- Ôn tập và bổ sung}{02}{Tia số. Số liền trước, số liền sau}{https://loptoanthaylinh.vn/thi/chuyen-de/2/toan2_cd_b02}
\textbf{Phần 1. Trắc nghiệm lựa chọn}
\setcounterpageref{ex}{0}

\begin{ex} %[2D1V02]
Số liền sau của số lớn nhất có hai chữ số là
\choice
{$98$}
{\True $100$}
{$99$}
{$101$}
\loigiai{
Số lớn nhất có hai chữ số là $99$. Số liền sau của 99 là $99 + 1 = 100$.
\textbf{Đáp án của câu hỏi là B}
}
\end{ex}

\begin{ex} %[2D1V02]
Trên tia số, điểm biểu diễn số $A$ cách gốc $0$ là $45$ vạch chia. Điểm $B$ đứng liền sau điểm $A$. Điểm $B$ biểu diễn số
\choice
{$44$}
{$45$}
{$50$}
{\True $46$}
\loigiai{
Điểm $A$ biểu diễn số 45. Điểm liền sau là $45 + 1 = 46$.
\textbf{Đáp án của câu hỏi là D}
}
\end{ex}

\begin{ex} %[2D1V02]
Tìm số tự nhiên $x$, biết số liền trước của $x$ là $58$.
\choice
{\True $59$}
{$57$}
{$60$}
{$58$}
\loigiai{
Vì 58 là số liền trước của $x$ nên $x = 58 + 1 = 59$.
\textbf{Đáp án của câu hỏi là A}
}
\end{ex}

\begin{ex} %[2D1V02]
Biết tổng của số liền trước và số liền sau của số $a$ bằng $40$. Số $a$ là
\choice
{$19$}
{$21$}
{\True $20$}
{$18$}
\loigiai{
Số liền trước là $a - 1$, số liền sau là $a + 1$. Tổng hai số là $(a - 1) + (a + 1) = a + a = 40$. Suy ra $a = 20$.
\textbf{Đáp án của câu hỏi là C}
}
\end{ex}

\begin{ex} %[2D1V02]
Một chú cào cào nhảy trên tia số, mỗi bước nhảy được $5$ đơn vị. Xuất phát từ vạch số $10$, sau $4$ bước nhảy chú cào cào đến vạch số nào?
\choice
{$25$}
{\True $30$}
{$35$}
{$40$}
\loigiai{
Quãng đường cào cào nhảy được là: $4 \times 5 = 20$ đơn vị.\nChú cào cào đến vạch: $10 + 20 = 30$.
\textbf{Đáp án của câu hỏi là B}
}
\end{ex}

\begin{ex} %[2D1V02]
Hai số tự nhiên liên tiếp có tổng bằng $35$. Hai số đó là
\choice
{$16\text{ và } 17$}
{$18\text{ và } 19$}
{$15\text{ và } 20$}
{\True $17\text{ và } 18$}
\loigiai{
Vì hai số liên tiếp hơn kém nhau 1 đơn vị nên số bé là: $(35 - 1) : 2 = 17$. Số lớn là 18.
\textbf{Đáp án của câu hỏi là D}
}
\end{ex}

\begin{ex} %[2D1V02]
Số liền trước của số bé nhất có hai chữ số giống nhau là
\choice
{\True $10$}
{$9$}
{$11$}
{$12$}
\loigiai{
Số bé nhất có hai chữ số giống nhau là $11$. Số liền trước của 11 là $11 - 1 = 10$.
\textbf{Đáp án của câu hỏi là A}
}
\end{ex}

\begin{ex} %[2D1V02]
Trên tia số, khoảng cách giữa số $28$ và số $62$ là bao nhiêu đơn vị?
\choice
{$35\text{ đơn vị}$}
{$44\text{ đơn vị}$}
{\True $34\text{ đơn vị}$}
{$32\text{ đơn vị}$}
\loigiai{
Khoảng cách giữa hai điểm là: $62 - 28 = 34\text{ đơn vị}$.
\textbf{Đáp án của câu hỏi là C}
}
\end{ex}

\begin{ex} %[2D1V02]
Số liền sau của số $x$ là một số tròn chục lớn hơn $40$ nhưng nhỏ hơn $60$. Số $x$ có thể là
\choice
{$59$}
{\True $49$}
{$39$}
{$51$}
\loigiai{
Số tròn chục lớn hơn 40 và nhỏ hơn 60 là 50. Số liền trước của 50 là $50 - 1 = 49$. Do đó $x = 49$.
\textbf{Đáp án của câu hỏi là B}
}
\end{ex}

\begin{ex} %[2D1V02]
Cho ba số tự nhiên liên tiếp tăng dần, biết số ở giữa là $79$. Tích của chữ số hàng đơn vị của số thứ nhất và số thứ ba là
\choice
{$8$}
{$16$}
{$2$}
{\True $0$}
\loigiai{
Ba số liên tiếp là $78, 79, 80$. Số thứ nhất có chữ số hàng đơn vị là 8, số thứ ba có chữ số hàng đơn vị là 0. Tích là $8 \times 0 = 0$.
\textbf{Đáp án của câu hỏi là D}
}
\end{ex}

\begin{ex} %[2D1V02]
Trên tia số có các điểm $A, B, C$. Điểm $A$ là $15$, điểm $B$ là $35$, điểm $C$ nằm chính giữa đoạn $AB$. Điểm $C$ biểu diễn số
\choice
{\True $25$}
{$20$}
{$30$}
{$26$}
\loigiai{
Đoạn $AB$ dài: $35 - 15 = 20$. Nửa đoạn $AB$ là: $20 : 2 = 10$. Điểm $C$ biểu diễn số: $15 + 10 = 25$.
\textbf{Đáp án của câu hỏi là A}
}
\end{ex}

\begin{ex} %[2D1V02]
Tìm số tự nhiên liền trước của số $y$, biết $y + 12 = 45$.
\choice
{$33$}
{$34$}
{\True $32$}
{$31$}
\loigiai{
Ta có: $y = 45 - 12 = 33$. Số liền trước của 33 là $33 - 1 = 32$.
\textbf{Đáp án của câu hỏi là C}
}
\end{ex}

\textbf{Phần 2. Trắc nghiệm đúng sai}
\setcounterpageref{ex}{0}

\begin{ex} %[2D1V02]
Xét tính đúng sai của các khẳng định về số liền trước và số liền sau:
\choiceTF
{\True Hiệu giữa số liền sau và số liền trước của một số bất kì luôn bằng 2.}
{Số 0 có số liền trước là một số tự nhiên.}
{\True Số liền sau của một số chẵn luôn là một số lẻ.}
{\True Ba số tự nhiên liên tiếp luôn có một số chia hết cho 3.}
\loigiai{
\begin{itemize}
    \item Ý a) Đúng: Gọi số đó là $n$, số liền trước là $n-1$, liền sau là $n+1$. Ta có $(n+1) - (n-1) = 2$.
    \item Ý b) Sai: Trong tập hợp các số tự nhiên, số 0 không có số liền trước.
    \item Ý c) Đúng: Một số chẵn cộng thêm 1 đơn vị luôn cho kết quả là một số lẻ.
    \item Ý d) Đúng: Trong ba số tự nhiên liên tiếp chắc chắn có đúng một số chia hết cho 3.
\end{itemize}
\begin{itemize}
    \item a) \textbf{Đúng}.
    \item b) \textbf{Sai}.
    \item c) \textbf{Đúng}.
    \item d) \textbf{Đúng}.
\end{itemize}
}
\end{ex}

\begin{ex} %[2D1V02]
Cho tia số với các mốc cách đều nhau 10 đơn vị từ 0 đến 100:
\choiceTF
{\True Số 47 nằm giữa hai vạch 40 và 50.}
{Số 47 nằm gần vạch 40 hơn vạch 50.}
{\True Từ vạch 20 đến vạch 80 có tất cả 6 khoảng cách 10 đơn vị.}
{Số liền sau vạch 30 là vạch 40.}
\loigiai{
\begin{itemize}
    \item Ý a) Đúng: Vì $40 < 47 < 50$ nên khẳng định đúng.
    \item Ý b) Sai: $47 - 40 = 7$, trong khi $50 - 47 = 3$. Số 47 gần 50 hơn.
    \item Ý c) Đúng: $(80 - 20) : 10 = 6$ khoảng cách.
    \item Ý d) Sai: Số liền sau của 30 là 31, còn 40 là vạch tròn chục tiếp theo.
\end{itemize}
\begin{itemize}
    \item a) \textbf{Đúng}.
    \item b) \textbf{Sai}.
    \item c) \textbf{Đúng}.
    \item d) \textbf{Sai}.
\end{itemize}
}
\end{ex}

\textbf{Phần 3. Trả lời ngắn}
\setcounterpageref{ex}{0}

\begin{ex} %[2D1V02]
Tìm một số biết rằng số liền trước của nó cộng với số liền sau của nó bằng 84.
\shortans{42}
\loigiai{
\textbf{Đáp số:} $42$\\
Tổng số liền trước và liền sau gấp đôi số đó: $84 : 2 = 42$.
}
\end{ex}

\begin{ex} %[2D1V02]
Số liền trước của số lớn nhất có hai chữ số khác nhau là số nào?
\shortans{97}
\loigiai{
\textbf{Đáp số:} $97$\\
Số lớn nhất có hai chữ số khác nhau là 98. Số liền trước là 97.
}
\end{ex}

\begin{ex} %[2D1V02]
Một chú thỏ nhảy từ vạch 12 trên tia số, mỗi lần nhảy 3 bước. Sau bao nhiêu lần nhảy thì thỏ đến vạch 33?
\shortans{7}
\loigiai{
\textbf{Đáp số:} $7$\\
Khoảng cách thỏ nhảy: $33 - 12 = 21$. Số lần nhảy là: $21 : 3 = 7$ lần.
}
\end{ex}

\begin{ex} %[2D1V02]
Tìm số tự nhiên biết rằng nếu cộng số đó với số liền sau của nó thì được 51.
\shortans{25}
\loigiai{
\textbf{Đáp số:} $25$\\
Hai số tự nhiên liên tiếp có tổng là 51: số bé là $(51 - 1) : 2 = 25$.
}
\end{ex}

\begin{ex} %[2D1V02]
Trên tia số, giữa số 18 và số 42 có tất cả bao nhiêu số tự nhiên?
\shortans{23}
\loigiai{
\textbf{Đáp số:} $23$\\
Số các số nằm ở giữa: $(42 - 18) - 1 = 23$ số.
}
\end{ex}

\begin{ex} %[2D1V02]
Tìm số liền sau của số tròn chục lớn nhất có hai chữ số mà tổng hai chữ số bằng 9.
\shortans{91}
\loigiai{
\textbf{Đáp số:} $91$\\
Số tròn chục có dạng $\overline{a0} \Rightarrow a + 0 = 9 \Rightarrow a = 9$. Số đó là 90, số liền sau là 91.
}
\end{ex}

\newpage

\TaoPhieuKiemTra{Chủ đề 1 -- Ôn tập và bổ sung}{03}{Các thành phần của phép cộng, phép trừ}{https://loptoanthaylinh.vn/thi/chuyen-de/2/toan2_cd_b03}
\textbf{Phần 1. Trắc nghiệm lựa chọn}
\setcounterpageref{ex}{0}

\begin{ex} %[2D1V03]
Trong một phép cộng có hai số hạng, nếu tăng số hạng thứ nhất thêm $8$ đơn vị và giữ nguyên số hạng thứ hai thì tổng sẽ
\choice
{\True Tăng thêm $8$ đơn vị}
{Giảm đi $8$ đơn vị}
{Không đổi}
{Tăng thêm $16$ đơn vị}
\loigiai{
Khi thêm vào một số hạng bao nhiêu đơn vị thì tổng cũng tăng thêm bấy nhiêu đơn vị. Vậy tổng tăng $8$ đơn vị.
\textbf{Đáp án của câu hỏi là A}
}
\end{ex}

\begin{ex} %[2D1V03]
Trong một phép trừ, số bị trừ là $75$, hiệu là $28$. Số trừ là
\choice
{$57$}
{$48$}
{\True $47$}
{$43$}
\loigiai{
Số trừ = Số bị trừ - Hiệu = $75 - 28 = 47$.
\textbf{Đáp án của câu hỏi là C}
}
\end{ex}

\begin{ex} %[2D1V03]
Biết tổng của hai số là $64$. Nếu một số hạng là $20$ thì số hạng còn lại là
\choice
{$34$}
{\True $44$}
{$54$}
{$42$}
\loigiai{
Số hạng còn lại = Tổng - Số hạng đã biết = $64 - 20 = 44$.
\textbf{Đáp án của câu hỏi là B}
}
\end{ex}

\begin{ex} %[2D1V03]
Trong phép trừ, nếu tăng số bị trừ thêm $5$ đơn vị và tăng số trừ thêm $5$ đơn vị thì hiệu sẽ
\choice
{Tăng thêm $10$ đơn vị}
{Giảm đi $5$ đơn vị}
{Tăng thêm $5$ đơn vị}
{\True Không thay đổi}
\loigiai{
Khi cùng thêm (hoặc bớt) ở cả số bị trừ và số trừ một số đơn vị như nhau thì hiệu không thay đổi.
\textbf{Đáp án của câu hỏi là D}
}
\end{ex}

\begin{ex} %[2D1V03]
Tìm một số biết rằng lấy số đó cộng với $18$ rồi trừ đi $12$ thì được $45$.
\choice
{\True $39$}
{$41$}
{$37$}
{$43$}
\loigiai{
Tính ngược từ cuối: $45 + 12 - 18 = 57 - 18 = 39$.
\textbf{Đáp án của câu hỏi là A}
}
\end{ex}

\begin{ex} %[2D1V03]
Trong một phép trừ, tổng của số bị trừ, số trừ và hiệu bằng $70$. Số bị trừ là
\choice
{$40$}
{$30$}
{\True $35$}
{$45$}
\loigiai{
Vì Số trừ + Hiệu = Số bị trừ nên: Số bị trừ + Số bị trừ = 70. Vậy Số bị trừ là $70 : 2 = 35$.
\textbf{Đáp án của câu hỏi là C}
}
\end{ex}

\begin{ex} %[2D1V03]
An có $32$ viên bi, Bình có $18$ viên bi. Để hai bạn có số bi bằng nhau thì An phải cho Bình bao nhiêu viên bi?
\choice
{$14$ viên bi}
{\True $7$ viên bi}
{$8$ viên bi}
{$6$ viên bi}
\loigiai{
An hơn Bình: $32 - 18 = 14$ viên bi.\nAn phải cho Bình: $14 : 2 = 7$ viên bi.
\textbf{Đáp án của câu hỏi là B}
}
\end{ex}

\begin{ex} %[2D1V03]
Tìm $x$ biết: $x - 24 = 35 + 11$.
\choice
{$60$}
{$68$}
{$72$}
{\True $70$}
\loigiai{
Ta có: $x - 24 = 46 \Rightarrow x = 46 + 24 = 70$.
\textbf{Đáp án của câu hỏi là D}
}
\end{ex}

\begin{ex} %[2D1V03]
Hai số có tổng bằng $58$. Nếu bớt ở số hạng thứ nhất đi $6$ đơn vị và thêm vào số hạng thứ hai $9$ đơn vị thì tổng mới là
\choice
{\True $61$}
{$55$}
{$63$}
{$59$}
\loigiai{
Tổng mới là: $58 - 6 + 9 = 61$.
\textbf{Đáp án của câu hỏi là A}
}
\end{ex}

\begin{ex} %[2D1V03]
Trong phép trừ có hiệu bằng $0$. Khẳng định nào sau đây luôn đúng?
\choice
{Số bị trừ lớn hơn số trừ}
{Số trừ bằng $0$}
{\True Số bị trừ bằng số trừ}
{Số bị trừ bằng $0$}
\loigiai{
Hiệu bằng 0 khi và chỉ khi số bị trừ bằng số trừ.
\textbf{Đáp án của câu hỏi là C}
}
\end{ex}

\begin{ex} %[2D1V03]
Mai thực hiện phép tính: $45 - a = 18$. Số $a$ đóng vai trò là gì và có giá trị là bao nhiêu?
\choice
{Số bị trừ, giá trị là $63$}
{\True Số trừ, giá trị là $27$}
{Hiệu, giá trị là $27$}
{Số hạng, giá trị là $18$}
\loigiai{
$a$ là số trừ trong phép trừ. Giá trị của $a = 45 - 18 = 27$.
\textbf{Đáp án của câu hỏi là B}
}
\end{ex}

\begin{ex} %[2D1V03]
Một sợi dây dài $85\text{ cm}$. Người ta cắt đi một đoạn thì còn lại đoạn dây dài $38\text{ cm}$. Độ dài đoạn dây bị cắt đi là
\choice
{$57\text{ cm}$}
{$43\text{ cm}$}
{$53\text{ cm}$}
{\True $47\text{ cm}$}
\loigiai{
Đoạn dây bị cắt đi là: $85 - 38 = 47\text{ (cm)}$.
\textbf{Đáp án của câu hỏi là D}
}
\end{ex}

\textbf{Phần 2. Trắc nghiệm đúng sai}
\setcounterpageref{ex}{0}

\begin{ex} %[2D1V03]
Xét tính đúng sai của các nhận định về phép cộng và phép trừ:
\choiceTF
{\True Trong phép cộng, khi đổi chỗ các số hạng thì tổng không thay đổi.}
{Trong phép trừ, số bị trừ luôn luôn phải lớn hơn số trừ.}
{\True Số bị trừ bằng tổng của số trừ và hiệu.}
{Nếu tổng bằng 0 thì cả hai số hạng đều phải bằng 1.}
\loigiai{
\begin{itemize}
    \item Ý a) Đúng: Đây là tính chất giao hoán của phép cộng.
    \item Ý b) Sai: Số bị trừ có thể bằng số trừ (khi đó hiệu bằng 0).
    \item Ý c) Đúng: Công thức cơ bản: Số bị trừ = Số trừ + Hiệu.
    \item Ý d) Sai: Trong tập số tự nhiên, nếu tổng bằng 0 thì cả hai số hạng đều phải bằng 0.
\end{itemize}
\begin{itemize}
    \item a) \textbf{Đúng}.
    \item b) \textbf{Sai}.
    \item c) \textbf{Đúng}.
    \item d) \textbf{Sai}.
\end{itemize}
}
\end{ex}

\begin{ex} %[2D1V03]
Cho phép tính $A - B = C$ (với $A, B, C$ là các số tự nhiên khác 0):
\choiceTF
{\True Nếu giữ nguyên $A$ và tăng $B$ thêm 3 đơn vị thì $C$ giảm đi 3 đơn vị.}
{\True Nếu tăng $A$ thêm 4 đơn vị và giữ nguyên $B$ thì $C$ tăng thêm 4 đơn vị.}
{Nếu $A$ tăng 5 đơn vị và $B$ giảm 5 đơn vị thì $C$ không đổi.}
{\True Tổng $A + B + C$ luôn là một số chẵn.}
\loigiai{
\begin{itemize}
    \item Ý a) Đúng: Khi số trừ tăng thì hiệu giảm tương ứng.
    \item Ý b) Đúng: Khi số bị trừ tăng thì hiệu tăng tương ứng.
    \item Ý c) Sai: Hiệu mới sẽ tăng thêm $5 + 5 = 10$ đơn vị.
    \item Ý d) Đúng: Vì $B + C = A$ nên $A + B + C = A + A = 2 \times A$, luôn là số chẵn.
\end{itemize}
\begin{itemize}
    \item a) \textbf{Đúng}.
    \item b) \textbf{Đúng}.
    \item c) \textbf{Sai}.
    \item d) \textbf{Đúng}.
\end{itemize}
}
\end{ex}

\textbf{Phần 3. Trả lời ngắn}
\setcounterpageref{ex}{0}

\begin{ex} %[2D1V03]
Trong một phép trừ, số bị trừ hơn số trừ 26 đơn vị. Hỏi hiệu của phép trừ đó bằng bao nhiêu?
\shortans{26}
\loigiai{
\textbf{Đáp số:} $26$\\
Hiệu chính là phần số bị trừ hơn số trừ, do đó hiệu bằng 26.
}
\end{ex}

\begin{ex} %[2D1V03]
Tìm một số biết rằng lấy 78 trừ đi số đó rồi cộng thêm 15 thì được 43.
\shortans{50}
\loigiai{
\textbf{Đáp số:} $50$\\
Ta có: $78 - x = 43 - 15 = 28 \Rightarrow x = 78 - 28 = 50$.
}
\end{ex}

\begin{ex} %[2D1V03]
Tổng của hai số là 75. Nếu số hạng thứ nhất tăng thêm 12 đơn vị và số hạng thứ hai giảm đi 7 đơn vị thì tổng mới bằng bao nhiêu?
\shortans{80}
\loigiai{
\textbf{Đáp số:} $80$\\
Tổng mới: $75 + 12 - 7 = 80$.
}
\end{ex}

\begin{ex} %[2D1V03]
Minh có 28 cái kẹo, Lan có 16 cái kẹo. Mẹ cho thêm hai bạn một số kẹo như nhau thì tổng số kẹo của hai bạn là 50 cái. Hỏi mẹ đã cho mỗi bạn bao nhiêu cái kẹo?
\shortans{3}
\loigiai{
\textbf{Đáp số:} $3$\\
Tổng số kẹo ban đầu là $28 + 16 = 44$. Mẹ cho cả hai bạn $50 - 44 = 6$ cái kẹo. Mỗi bạn được $6 : 2 = 3$ cái.
}
\end{ex}

\begin{ex} %[2D1V03]
Tìm số trừ trong một phép trừ biết rằng số bị trừ là 90 và hiệu là số chẵn nhỏ nhất có hai chữ số.
\shortans{80}
\loigiai{
\textbf{Đáp số:} $80$\\
Số chẵn nhỏ nhất có hai chữ số là 10. Số trừ là $90 - 10 = 80$.
}
\end{ex}

\begin{ex} %[2D1V03]
Hai túi gạo có tất cả 64 kg gạo. Nếu chuyển 6 kg gạo từ túi thứ nhất sang túi thứ hai thì hai túi có số gạo bằng nhau. Hỏi lúc đầu túi thứ nhất có bao nhiêu ki-lô-gam gạo?
\shortans{38}
\loigiai{
\textbf{Đáp số:} $38$\\
Túi 1 hơn túi 2 là $6 \times 2 = 12\text{ kg}$. Lúc đầu túi 1 có: $(64 + 12) : 2 = 38\text{ kg}$.
}
\end{ex}

\newpage

\TaoPhieuKiemTra{Chủ đề 1 -- Ôn tập và bổ sung}{04}{Hơn, kém nhau bao nhiêu đơn vị}{https://loptoanthaylinh.vn/thi/chuyen-de/2/toan2_cd_b04}
\textbf{Phần 1. Trắc nghiệm lựa chọn}
\setcounterpageref{ex}{0}

\begin{ex} %[2D1V04]
Năm nay Mai $8$ tuổi, mẹ Mai $34$ tuổi. Hỏi sau $5$ năm nữa mẹ hơn Mai bao nhiêu tuổi?
\choice
{$31$ tuổi}
{$21$ tuổi}
{$39$ tuổi}
{\True $26$ tuổi}
\loigiai{
Hiệu số tuổi của hai người không bao giờ thay đổi theo thời gian. Mẹ hơn Mai: $34 - 8 = 26$ (tuổi).
\textbf{Đáp án của câu hỏi là D}
}
\end{ex}

\begin{ex} %[2D1V04]
Lớp 2A có $35$ học sinh, lớp 2B có ít hơn lớp 2A là $4$ học sinh. Cả hai lớp có tất cả bao nhiêu học sinh?
\choice
{$62$ học sinh}
{\True $66$ học sinh}
{$70$ học sinh}
{$68$ học sinh}
\loigiai{
Số học sinh lớp 2B là: $35 - 4 = 31$ (học sinh).\nCả hai lớp có: $35 + 31 = 66$ (học sinh).
\textbf{Đáp án của câu hỏi là B}
}
\end{ex}

\begin{ex} %[2D1V04]
Đoạn thẳng $AB$ dài $1\text{ dm } 4\text{ cm}$, đoạn thẳng $CD$ dài $9\text{ cm}$. Đoạn thẳng $AB$ dài hơn đoạn thẳng $CD$ bao nhiêu xăng-ti-mét?
\choice
{$7\text{ cm}$}
{$4\text{ cm}$}
{\True $5\text{ cm}$}
{$6\text{ cm}$}
\loigiai{
Đổi $1\text{ dm } 4\text{ cm} = 14\text{ cm}$.\nĐoạn thẳng $AB$ dài hơn đoạn thẳng $CD$ là: $14 - 9 = 5\text{ (cm)}$.
\textbf{Đáp án của câu hỏi là C}
}
\end{ex}

\begin{ex} %[2D1V04]
Một đàn gà có $45$ con gà mái. Số gà trống ít hơn số gà mái $23$ con. Hỏi đàn gà có tất cả bao nhiêu con gà?
\choice
{\True $67$ con gà}
{$68$ con gà}
{$22$ con gà}
{$65$ con gà}
\loigiai{
Số gà trống là: $45 - 23 = 22$ (con).\nCả đàn gà có: $45 + 22 = 67$ (con).
\textbf{Đáp án của câu hỏi là A}
}
\end{ex}

\begin{ex} %[2D1V04]
Trong một hàng dọc, Nam đứng thứ $12$ tính từ đầu hàng và đứng thứ $8$ tính từ cuối hàng. Hỏi hàng đó có tất cả bao nhiêu bạn?
\choice
{$20$ bạn}
{$18$ bạn}
{$21$ bạn}
{\True $19$ bạn}
\loigiai{
Số bạn trong hàng là: $12 + 8 - 1 = 19$ (bạn).
\textbf{Đáp án của câu hỏi là D}
}
\end{ex}

\begin{ex} %[2D1V04]
Thùng thứ nhất có $42\text{ L}$ dầu, thùng thứ hai nhiều hơn thùng thứ nhất $15\text{ L}$ dầu. Cả hai thùng có tất cả bao nhiêu lít dầu?
\choice
{$89\text{ L}$}
{\True $99\text{ L}$}
{$97\text{ L}$}
{$100\text{ L}$}
\loigiai{
Thùng thứ hai có: $42 + 15 = 57\text{ (L)}$.\nCả hai thùng có: $42 + 57 = 99\text{ (L)}$.
\textbf{Đáp án của câu hỏi là B}
}
\end{ex}

\begin{ex} %[2D1V04]
Bao ngô nặng $48\text{ kg}$, bao thóc nhẹ hơn bao ngô $12\text{ kg}$. Hỏi cả hai bao nặng bao nhiêu ki-lô-gam?
\choice
{$72\text{ kg}$}
{$96\text{ kg}$}
{\True $84\text{ kg}$}
{$86\text{ kg}$}
\loigiai{
Bao thóc nặng: $48 - 12 = 36\text{ (kg)}$.\nCả hai bao nặng: $48 + 36 = 84\text{ (kg)}$.
\textbf{Đáp án của câu hỏi là C}
}
\end{ex}

\begin{ex} %[2D1V04]
Có hai ngăn sách, ngăn trên có $38$ quyển sách. Nếu chuyển $5$ quyển từ ngăn trên xuống ngăn dưới thì ngăn trên vẫn nhiều hơn ngăn dưới $4$ quyển. Hỏi ban đầu ngăn trên hơn ngăn dưới bao nhiêu quyển?
\choice
{\True $14$ quyển}
{$9$ quyển}
{$10$ quyển}
{$12$ quyển}
\loigiai{
Chuyển 5 quyển làm hiệu giảm đi $5 \times 2 = 10$ quyển. Lúc này còn hơn 4 quyển nên ban đầu hơn: $10 + 4 = 14$ quyển.
\textbf{Đáp án của câu hỏi là A}
}
\end{ex}

\begin{ex} %[2D1V04]
Một cây cau cao $8\text{ m}$, cây dừa cao hơn cây cau $3\text{ m}$, cây xoài thấp hơn cây dừa $4\text{ m}$. Chiều cao của cây xoài là
\choice
{$6\text{ m}$}
{$8\text{ m}$}
{$9\text{ m}$}
{\True $7\text{ m}$}
\loigiai{
Cây dừa cao: $8 + 3 = 11\text{ (m)}$.\nCây xoài cao: $11 - 4 = 7\text{ (m)}$.
\textbf{Đáp án của câu hỏi là D}
}
\end{ex}

\begin{ex} %[2D1V04]
Lan có nhiều hơn Huệ $16$ bông hoa. Lan tặng Huệ $5$ bông hoa. Hỏi bây giờ Lan còn nhiều hơn Huệ bao nhiêu bông hoa?
\choice
{$11$ bông hoa}
{\True $6$ bông hoa}
{$8$ bông hoa}
{$10$ bông hoa}
\loigiai{
Khi Lan tặng Huệ 5 bông hoa thì số hoa Lan hơn Huệ giảm đi: $5 \times 2 = 10$ (bông).\nBây giờ Lan còn nhiều hơn Huệ: $16 - 10 = 6$ (bông hoa).
\textbf{Đáp án của câu hỏi là B}
}
\end{ex}

\begin{ex} %[2D1V04]
Tìm hai số có hiệu bằng $18$, biết rằng nếu thêm vào số bị trừ $6$ đơn vị thì được số mới bằng $45$.
\choice
{$45\text{ và } 27$}
{$38\text{ và } 20$}
{\True $39\text{ và } 21$}
{$40\text{ và } 22$}
\loigiai{
Số bị trừ là: $45 - 6 = 39$. Số trừ là: $39 - 18 = 21$.
\textbf{Đáp án của câu hỏi là C}
}
\end{ex}

\begin{ex} %[2D1V04]
Một cửa hàng buổi sáng bán được $56$ quả dưa hấu, buổi chiều bán được ít hơn buổi sáng $14$ quả dưa hấu. Cả hai buổi cửa hàng bán được bao nhiêu quả?
\choice
{\True $98$ quả}
{$96$ quả}
{$84$ quả}
{$102$ quả}
\loigiai{
Buổi chiều bán được: $56 - 14 = 42$ (quả).\nCả hai buổi bán được: $56 + 42 = 98$ (quả).
\textbf{Đáp án của câu hỏi là A}
}
\end{ex}

\textbf{Phần 2. Trắc nghiệm đúng sai}
\setcounterpageref{ex}{0}

\begin{ex} %[2D1V04]
Xét tính đúng sai của các bài toán về nhiều hơn, ít hơn:
\choiceTF
{\True Hiệu số tuổi của hai người không đổi theo thời gian.}
{Nếu Nam nhiều hơn Bình 8 viên bi, Nam cho Bình 4 viên bi thì Bình sẽ nhiều hơn Nam.}
{\True Một đoạn dây dài 20 cm bị cắt thành 2 đoạn, đoạn thứ nhất dài hơn đoạn thứ hai 4 cm thì đoạn dài là 12 cm.}
{\True Muốn tìm số lớn ta lấy số bé cộng với phần hơn.}
\loigiai{
\begin{itemize}
    \item Ý a) Đúng: Mỗi năm mỗi người đều tăng thêm 1 tuổi nên hiệu số tuổi luôn giữ nguyên.
    \item Ý b) Sai: Khi Nam cho Bình 4 viên bi thì hai bạn có số bi bằng nhau: $8 - 4 = 4$ và bên kia $0 + 4 = 4$.
    \item Ý c) Đúng: Đoạn dài là: $(20 + 4) : 2 = 12\text{ cm}$.
    \item Ý d) Đúng: Số lớn = Số bé + Phần hơn.
\end{itemize}
\begin{itemize}
    \item a) \textbf{Đúng}.
    \item b) \textbf{Sai}.
    \item c) \textbf{Đúng}.
    \item d) \textbf{Đúng}.
\end{itemize}
}
\end{ex}

\begin{ex} %[2D1V04]
Lớp 2A có 32 học sinh, lớp 2B có 36 học sinh:
\choiceTF
{\True Lớp 2B nhiều hơn lớp 2A là 4 học sinh.}
{Chuyển 4 học sinh từ lớp 2B sang lớp 2A thì hai lớp bằng nhau.}
{\True Cả hai lớp có tất cả 68 học sinh.}
{Nếu lớp 2A có thêm 2 bạn và lớp 2B bớt 2 bạn thì lớp 2B vẫn nhiều hơn.}
\loigiai{
\begin{itemize}
    \item Ý a) Đúng: $36 - 32 = 4$ học sinh.
    \item Ý b) Sai: Chỉ cần chuyển $4 : 2 = 2$ học sinh thì cả hai lớp đều có 34 học sinh.
    \item Ý c) Đúng: $32 + 36 = 68$ học sinh.
    \item Ý d) Sai: Khi đó lớp 2A có 34 bạn, lớp 2B có 34 bạn, hai lớp bằng nhau.
\end{itemize}
\begin{itemize}
    \item a) \textbf{Đúng}.
    \item b) \textbf{Sai}.
    \item c) \textbf{Đúng}.
    \item d) \textbf{Sai}.
\end{itemize}
}
\end{ex}

\textbf{Phần 3. Trả lời ngắn}
\setcounterpageref{ex}{0}

\begin{ex} %[2D1V04]
Bác Ba nuôi 38 con vịt và nuôi số gà nhiều hơn số vịt 15 con. Hỏi bác Ba nuôi tất cả bao nhiêu con gà và vịt?
\shortans{91}
\loigiai{
\textbf{Đáp số:} $91$\\
Số gà: $38 + 15 = 53$ con. Tổng số: $38 + 53 = 91$ con.
}
\end{ex}

\begin{ex} %[2D1V04]
Hiện nay anh 14 tuổi, em 8 tuổi. Hỏi khi em bằng tuổi anh hiện nay thì anh bao nhiêu tuổi?
\shortans{20}
\loigiai{
\textbf{Đáp số:} $20$\\
Anh hơn em: $14 - 8 = 6$ tuổi. Khi em 14 tuổi thì anh: $14 + 6 = 20$ tuổi.
}
\end{ex}

\begin{ex} %[2D1V04]
Túi thứ nhất có 35 viên kẹo, túi thứ hai có ít hơn túi thứ nhất 9 viên kẹo. Hỏi cả hai túi có tất cả bao nhiêu viên kẹo?
\shortans{61}
\loigiai{
\textbf{Đáp số:} $61$\\
Túi 2: $35 - 9 = 26$ viên. Cả hai túi: $35 + 26 = 61$ viên.
}
\end{ex}

\begin{ex} %[2D1V04]
Hùng có nhiều hơn Dũng 18 viên bi. Hùng phải cho Dũng bao nhiêu viên bi để hai bạn có số bi bằng nhau?
\shortans{9}
\loigiai{
\textbf{Đáp số:} $9$\\
$18 : 2 = 9$ viên bi.
}
\end{ex}

\begin{ex} %[2D1V04]
Một mảnh vải dài 45 dm, mảnh vải thứ hai ngắn hơn mảnh vải thứ nhất 18 dm. Hỏi cả hai mảnh vải dài bao nhiêu đề-xi-mét?
\shortans{72}
\loigiai{
\textbf{Đáp số:} $72$\\
Mảnh 2: $45 - 18 = 27\text{ dm}$. Tổng: $45 + 27 = 72\text{ dm}$.
}
\end{ex}

\begin{ex} %[2D1V04]
Trong một cuộc thi chạy, từ vị trí của Mai về đích còn 25 m, Mai chạy trước Lan 8 m. Hỏi từ vị trí của Lan về đích còn bao nhiêu mét?
\shortans{33}
\loigiai{
\textbf{Đáp số:} $33$\\
Vị trí của Lan cách đích: $25 + 8 = 33\text{ m}$.
}
\end{ex}

\newpage

\TaoPhieuKiemTra{Chủ đề 1 -- Ôn tập và bổ sung}{05}{Bảng cộng và bảng trừ (không nhớ) trong phạm vi 100}{https://loptoanthaylinh.vn/thi/chuyen-de/2/toan2_cd_b05}
\textbf{Phần 1. Trắc nghiệm lựa chọn}
\setcounterpageref{ex}{0}

\begin{ex} %[2D1V05]
Tính nhanh giá trị của biểu thức: $A = 24 + 15 + 36 + 25$.
\choice
{$90$}
{$95$}
{\True $100$}
{$105$}
\loigiai{
Gộp các số tạo số tròn chục: $A = (24 + 36) + (15 + 25) = 60 + 40 = 100$.
\textbf{Đáp án của câu hỏi là C}
}
\end{ex}

\begin{ex} %[2D1V05]
Tìm $x$ biết: $x + 23 = 78 - 15$.
\choice
{$35$}
{\True $40$}
{$45$}
{$50$}
\loigiai{
Vế phải: $78 - 15 = 63$. Ta có: $x + 23 = 63 \Rightarrow x = 63 - 23 = 40$.
\textbf{Đáp án của câu hỏi là B}
}
\end{ex}

\begin{ex} %[2D1V05]
Một xe buýt chở $36$ hành khách. Đến trạm dừng có $12$ người xuống xe và $8$ người lên xe. Hỏi trên xe buýt lúc này có bao nhiêu hành khách?
\choice
{$30$ người}
{$34$ người}
{$28$ người}
{\True $32$ người}
\loigiai{
Số người trên xe lúc này là: $36 - 12 + 8 = 24 + 8 = 32$ (người).
\textbf{Đáp án của câu hỏi là D}
}
\end{ex}

\begin{ex} %[2D1V05]
Điền số thích hợp vào ô trống: $96 - \dots = 42 + 21$.
\choice
{\True $33$}
{$35$}
{$31$}
{$23$}
\loigiai{
Vế phải: $42 + 21 = 63$. Ô trống $= 96 - 63 = 33$.
\textbf{Đáp án của câu hỏi là A}
}
\end{ex}

\begin{ex} %[2D1V05]
Tính hiệu của số lớn nhất có hai chữ số và tổng của hai số $32$ và $25$.
\choice
{$38$}
{$40$}
{\True $42$}
{$44$}
\loigiai{
Số lớn nhất có hai chữ số là 99. Tổng hai số là $32 + 25 = 57$.\nHiệu cần tìm: $99 - 57 = 42$.
\textbf{Đáp án của câu hỏi là C}
}
\end{ex}

\begin{ex} %[2D1V05]
Một cửa hàng có $85$ mét vải. Buổi sáng bán được $23$ mét, buổi chiều bán được $31$ mét. Hỏi cửa hàng còn lại bao nhiêu mét vải?
\choice
{$29\text{ m}$}
{\True $31\text{ m}$}
{$33\text{ m}$}
{$25\text{ m}$}
\loigiai{
Tổng số mét vải đã bán: $23 + 31 = 54\text{ (m)}$.\nSố vải còn lại: $85 - 54 = 31\text{ (m)}$.
\textbf{Đáp án của câu hỏi là B}
}
\end{ex}

\begin{ex} %[2D1V05]
Tìm một số biết rằng lấy số đó trừ đi $14$ rồi cộng thêm $23$ thì được $68$.
\choice
{$57$}
{$61$}
{$55$}
{\True $59$}
\loigiai{
Tính ngược từ cuối: $68 - 23 + 14 = 45 + 14 = 59$.
\textbf{Đáp án của câu hỏi là D}
}
\end{ex}

\begin{ex} %[2D1V05]
Cho phép tính: $\overline{a4} + 35 = 79$. Giá trị của chữ số $a$ là
\choice
{\True $4$}
{$3$}
{$5$}
{$2$}
\loigiai{
Ta có: $\overline{a4} = 79 - 35 = 44$. Do đó $a = 4$.
\textbf{Đáp án của câu hỏi là A}
}
\end{ex}

\begin{ex} %[2D1V05]
Điền dấu thích hợp vào chỗ chấm: $68 - 25 \dots 21 + 22$.
\choice
{$>$}
{$<$}
{\True $=$}
{Không so sánh được}
\loigiai{
Vế trái: $68 - 25 = 43$. Vế phải: $21 + 22 = 43$. Vì $43 = 43$ nên điền dấu $=$.
\textbf{Đáp án của câu hỏi là C}
}
\end{ex}

\begin{ex} %[2D1V05]
Bạn Linh có $58$ con tem. Linh cho Nga $15$ con tem và cho Mai $12$ con tem. Hỏi Linh còn lại bao nhiêu con tem?
\choice
{$33$ con tem}
{\True $31$ con tem}
{$29$ con tem}
{$35$ con tem}
\loigiai{
Linh cho hai bạn tất cả: $15 + 12 = 27$ (con tem).\nLinh còn lại: $58 - 27 = 31$ (con tem).
\textbf{Đáp án của câu hỏi là B}
}
\end{ex}

\begin{ex} %[2D1V05]
Có ba thùng sách. Thùng thứ nhất có $21$ quyển, thùng thứ hai có $24$ quyển, thùng thứ ba có số sách bằng tổng hai thùng đầu trừ đi $15$ quyển. Thùng thứ ba có bao nhiêu quyển sách?
\choice
{$35$ quyển}
{$28$ quyển}
{$32$ quyển}
{\True $30$ quyển}
\loigiai{
Tổng hai thùng đầu: $21 + 24 = 45$ (quyển).\nThùng thứ ba có: $45 - 15 = 30$ (quyển).
\textbf{Đáp án của câu hỏi là D}
}
\end{ex}

\begin{ex} %[2D1V05]
Tìm hiệu của số chẵn lớn nhất có hai chữ số và số lẻ nhỏ nhất có hai chữ số.
\choice
{\True $87$}
{$89$}
{$88$}
{$86$}
\loigiai{
Số chẵn lớn nhất có hai chữ số là 98. Số lẻ nhỏ nhất có hai chữ số là 11.\nHiệu là: $98 - 11 = 87$.
\textbf{Đáp án của câu hỏi là A}
}
\end{ex}

\textbf{Phần 2. Trắc nghiệm đúng sai}
\setcounterpageref{ex}{0}

\begin{ex} %[2D1V05]
Xét tính đúng sai của các biểu thức tính toán sau:
\choiceTF
{\True Biểu thức $45 + 23 - 12$ có giá trị bằng $56$.}
{Biểu thức $89 - 34 - 15$ có giá trị bằng giá trị của $89 - (34 - 15)$.}
{\True Tổng của số chẵn và số lẻ luôn luôn là số lẻ.}
{\True Trong phép tính $78 - 35 = 43$, các chữ số hàng đơn vị thỏa mãn $8 - 5 = 3$.}
\loigiai{
\begin{itemize}
    \item Ý a) Đúng: $45 + 23 = 68$; $68 - 12 = 56$.
    \item Ý b) Sai: $89 - 34 - 15 = 40$, còn $89 - 19 = 70$. Hai kết quả khác nhau.
    \item Ý c) Đúng: Chẵn + Lẻ = Lẻ.
    \item Ý d) Đúng: Hàng đơn vị: $8 - 5 = 3$. Đúng.
\end{itemize}
\begin{itemize}
    \item a) \textbf{Đúng}.
    \item b) \textbf{Sai}.
    \item c) \textbf{Đúng}.
    \item d) \textbf{Đúng}.
\end{itemize}
}
\end{ex}

\begin{ex} %[2D1V05]
Trong vườn nhà bác An có 45 cây ăn quả gồm cam, bưởi và xoài. Trong đó có 21 cây cam, số cây bưởi ít hơn cây cam 10 cây:
\choiceTF
{\True Số cây bưởi là 11 cây.}
{\True Tổng số cây cam và bưởi là 32 cây.}
{Số cây xoài trong vườn là 15 cây.}
{\True Số cây cam nhiều hơn số cây xoài là 8 cây.}
\loigiai{
\begin{itemize}
    \item Ý a) Đúng: $21 - 10 = 11$ cây bưởi.
    \item Ý b) Đúng: $21 + 11 = 32$ cây.
    \item Ý c) Sai: Số cây xoài là: $45 - 32 = 13$ cây (không phải 15 cây).
    \item Ý d) Đúng: $21 - 13 = 8$ cây.
\end{itemize}
\begin{itemize}
    \item a) \textbf{Đúng}.
    \item b) \textbf{Đúng}.
    \item c) \textbf{Sai}.
    \item d) \textbf{Đúng}.
\end{itemize}
}
\end{ex}

\textbf{Phần 3. Trả lời ngắn}
\setcounterpageref{ex}{0}

\begin{ex} %[2D1V05]
Tính nhanh giá trị của biểu thức: $12 + 25 + 38 + 15$.
\shortans{90}
\loigiai{
\textbf{Đáp số:} $90$\\
$(12 + 38) + (25 + 15) = 50 + 40 = 90$.
}
\end{ex}

\begin{ex} %[2D1V05]
Tìm x biết: $98 - x = 24 + 31$.
\shortans{43}
\loigiai{
\textbf{Đáp số:} $43$\\
$98 - x = 55 \Rightarrow x = 98 - 55 = 43$.
}
\end{ex}

\begin{ex} %[2D1V05]
Một mảnh đất hình chữ nhật có chu vi là 88 m. Nửa chu vi mảnh đất đó bằng bao nhiêu mét?
\shortans{44}
\loigiai{
\textbf{Đáp số:} $44$\\
Nửa chu vi: $88 : 2 = 44\text{ m}$.
}
\end{ex}

\begin{ex} %[2D1V05]
Nhà Mai thu hoạch được 56 kg cà chua, đã bán đi 34 kg. Số cà chua còn lại chia đều vào 2 rổ. Hỏi mỗi rổ có bao nhiêu ki-lô-gam cà chua?
\shortans{11}
\loigiai{
\textbf{Đáp số:} $11$\\
Số cà chua còn lại: $56 - 34 = 22\text{ kg}$. Mỗi rổ có: $22 : 2 = 11\text{ kg}$.
}
\end{ex}

\begin{ex} %[2D1V05]
Tìm một số biết rằng lấy số đó cộng với 32 rồi trừ đi 15 thì được kết quả là 64.
\shortans{47}
\loigiai{
\textbf{Đáp số:} $47$\\
Số đó là: $64 + 15 - 32 = 79 - 32 = 47$.
}
\end{ex}

\begin{ex} %[2D1V05]
Điền chữ số thích hợp thay cho dấu hỏi chấm: $?4 + 23 = 87$.
\shortans{6}
\loigiai{
\textbf{Đáp số:} $6$\\
Số đó là $87 - 23 = 64$. Vậy dấu ? là chữ số 6.
}
\end{ex}

\newpage

\TaoPhieuKiemTra{Chủ đề 1 -- Ôn tập và bổ sung}{06}{Luyện tập chung}{https://loptoanthaylinh.vn/thi/chuyen-de/2/toan2_cd_b06}
\textbf{Phần 1. Trắc nghiệm lựa chọn}
\setcounterpageref{ex}{0}

\begin{ex} %[2D1V06]
Một bác thợ mộc cưa một khúc gỗ dài $12\text{ m}$ thành các đoạn nhỏ dài $3\text{ m}$. Hỏi bác thợ phải thực hiện bao nhiêu nhát cưa?
\choice
{$4$ nhát cưa}
{\True $3$ nhát cưa}
{$5$ nhát cưa}
{$2$ nhát cưa}
\loigiai{
Số đoạn gỗ cưa được là: $12 : 3 = 4$ (đoạn).\nSố nhát cưa cần thực hiện là: $4 - 1 = 3$ (nhát cưa).
\textbf{Đáp án của câu hỏi là B}
}
\end{ex}

\begin{ex} %[2D1V06]
Người ta trồng cây dọc theo một lối đi dài $20\text{ m}$, hai đầu đường đều có cây. Biết khoảng cách giữa hai cây liên tiếp là $5\text{ m}$. Hỏi có tất cả bao nhiêu cây?
\choice
{\True $5$ cây}
{$4$ cây}
{$6$ cây}
{$7$ cây}
\loigiai{
Số khoảng cách là: $20 : 5 = 4$ (khoảng cách).\nVì hai đầu đường đều trồng cây nên số cây là: $4 + 1 = 5$ (cây).
\textbf{Đáp án của câu hỏi là A}
}
\end{ex}

\begin{ex} %[2D1V06]
Tôi là một số có hai chữ số. Nếu lấy tôi trừ đi $15$ rồi cộng thêm $24$ thì được số tròn chục lớn nhất có hai chữ số. Tôi là số nào?
\choice
{$79$}
{$90$}
{\True $81$}
{$85$}
\loigiai{
Số tròn chục lớn nhất có hai chữ số là 90.\nSố cần tìm là: $90 - 24 + 15 = 66 + 15 = 81$.
\textbf{Đáp án của câu hỏi là C}
}
\end{ex}

\begin{ex} %[2D1V06]
Hình vẽ gồm một đoạn thẳng dài $35\text{ cm}$ bị cắt bớt một đoạn dài $1\text{ dm } 2\text{ cm}$. Đoạn thẳng còn lại dài bao nhiêu?
\choice
{$25\text{ cm}$}
{$21\text{ cm}$}
{$22\text{ cm}$}
{\True $23\text{ cm}$}
\loigiai{
Đổi $1\text{ dm } 2\text{ cm} = 12\text{ cm}$.\nĐoạn thẳng còn lại dài: $35 - 12 = 23\text{ (cm)}$.
\textbf{Đáp án của câu hỏi là D}
}
\end{ex}

\begin{ex} %[2D1V06]
Tìm số tự nhiên $y$ biết: $y + 16 < 20$. Các giá trị thích hợp của $y$ là
\choice
{$1, 2, 3, 4$}
{\True $0, 1, 2, 3$}
{$0, 1, 2$}
{$1, 2, 3$}
\loigiai{
Ta có: $y < 20 - 16 \Rightarrow y < 4$. Vì $y$ là số tự nhiên nên $y \in \{0, 1, 2, 3\}$.
\textbf{Đáp án của câu hỏi là B}
}
\end{ex}

\begin{ex} %[2D1V06]
Một giá sách có hai ngăn đựng tất cả $78$ quyển sách. Nếu chuyển $8$ quyển sách từ ngăn trên xuống ngăn dưới thì hai ngăn có số sách bằng nhau. Số sách ban đầu của ngăn trên là
\choice
{\True $47$ quyển}
{$45$ quyển}
{$39$ quyển}
{$49$ quyển}
\loigiai{
Khi hai ngăn bằng nhau, mỗi ngăn có: $78 : 2 = 39$ (quyển).\nSố sách ban đầu của ngăn trên là: $39 + 8 = 47$ (quyển).
\textbf{Đáp án của câu hỏi là A}
}
\end{ex}

\begin{ex} %[2D1V06]
Đồng hồ chỉ $7$ giờ sáng. Kim giờ và kim phút tạo thành một góc. Sau $30$ phút nữa thì đồng hồ chỉ mấy giờ?
\choice
{$8$ giờ}
{$7$ giờ $15$ phút}
{\True $7$ giờ $30$ phút}
{$6$ giờ $30$ phút}
\loigiai{
Sau 30 phút, thời gian là $7$ giờ $+ 30$ phút $= 7$ giờ $30$ phút.
\textbf{Đáp án của câu hỏi là C}
}
\end{ex}

\begin{ex} %[2D1V06]
Cho ba số có tổng bằng $95$. Biết số thứ nhất là $23$, số thứ hai hơn số thứ nhất $8$ đơn vị. Số thứ ba là
\choice
{$39$}
{$43$}
{$45$}
{\True $41$}
\loigiai{
Số thứ hai là: $23 + 8 = 31$.\nSố thứ ba là: $95 - (23 + 31) = 95 - 54 = 41$.
\textbf{Đáp án của câu hỏi là D}
}
\end{ex}

\begin{ex} %[2D1V06]
Tìm số lớn nhất có hai chữ số mà chữ số hàng chục hơn chữ số hàng đơn vị $4$ đơn vị.
\choice
{$84$}
{\True $95$}
{$73$}
{$94$}
\loigiai{
Để số lớn nhất, chọn hàng chục là $9$. Chữ số hàng đơn vị là $9 - 4 = 5$. Vậy số đó là $95$.
\textbf{Đáp án của câu hỏi là B}
}
\end{ex}

\begin{ex} %[2D1V06]
Trong một phép trừ, nếu bớt ở số bị trừ đi $7$ đơn vị và giữ nguyên số trừ thì hiệu sẽ
\choice
{\True Giảm đi $7$ đơn vị}
{Tăng thêm $7$ đơn vị}
{Không đổi}
{Giảm đi $14$ đơn vị}
\loigiai{
Khi bớt ở số bị trừ đi bao nhiêu đơn vị thì hiệu cũng giảm đi bấy nhiêu đơn vị.
\textbf{Đáp án của câu hỏi là A}
}
\end{ex}

\begin{ex} %[2D1V06]
Cho hình chữ nhật có chiều dài $25\text{ cm}$, chiều rộng ngắn hơn chiều dài $11\text{ cm}$. Chiều rộng của hình chữ nhật đó là
\choice
{$15\text{ cm}$}
{$16\text{ cm}$}
{\True $14\text{ cm}$}
{$12\text{ cm}$}
\loigiai{
Chiều rộng hình chữ nhật là: $25 - 11 = 14\text{ (cm)}$.
\textbf{Đáp án của câu hỏi là C}
}
\end{ex}

\begin{ex} %[2D1V06]
Bạn Bình đếm các số tròn chục từ $10$ đến $90$. Hỏi bạn Bình đã đếm được bao nhiêu chữ số $0$?
\choice
{$8$ chữ số $0$}
{$10$ chữ số $0$}
{$7$ chữ số $0$}
{\True $9$ chữ số $0$}
\loigiai{
Các số là: 10, 20, 30, 40, 50, 60, 70, 80, 90. Mỗi số có đúng 1 chữ số 0 nên có 9 chữ số 0.
\textbf{Đáp án của câu hỏi là D}
}
\end{ex}

\textbf{Phần 2. Trắc nghiệm đúng sai}
\setcounterpageref{ex}{0}

\begin{ex} %[2D1V06]
Xét tính đúng sai của các bài toán tư duy tổng hợp:
\choiceTF
{\True Một sợi dây dài 10 m bị cắt thành 5 đoạn bằng nhau thì cần 4 nhát cắt.}
{Nếu hôm nay là thứ Hai ngày 4 thì thứ Hai tuần sau là ngày 10.}
{\True Số nhỏ nhất có hai chữ số mà tổng hai chữ số bằng 9 là số 18.}
{\True Tổng của số liền trước và số liền sau của 50 là 100.}
\loigiai{
\begin{itemize}
    \item Ý a) Đúng: Số nhát cắt = Số đoạn - 1 = $5 - 1 = 4$ nhát.
    \item Ý b) Sai: Một tuần có 7 ngày, thứ Hai tuần sau là ngày $4 + 7 = 11$.
    \item Ý c) Đúng: Chọn chữ số hàng chục nhỏ nhất là 1, hàng đơn vị là $9 - 1 = 8$. Số đó là 18.
    \item Ý d) Đúng: $49 + 51 = 100$. Đúng.
\end{itemize}
\begin{itemize}
    \item a) \textbf{Đúng}.
    \item b) \textbf{Sai}.
    \item c) \textbf{Đúng}.
    \item d) \textbf{Đúng}.
\end{itemize}
}
\end{ex}

\begin{ex} %[2D1V06]
Cho dãy số: 2, 5, 8, 11, 14, ...:
\choiceTF
{\True Quy luật của dãy số là mỗi số đứng sau bằng số đứng trước cộng thêm 3.}
{Số hạng thứ 7 của dãy số là số 19.}
{\True Các số trong dãy số trên khi chia cho 3 đều dư 2.}
{Số 30 là một số thuộc dãy số trên.}
\loigiai{
\begin{itemize}
    \item Ý a) Đúng: $5 - 2 = 3$, $8 - 5 = 3$.
    \item Ý b) Sai: Số thứ 6 là 17, số thứ 7 là $17 + 3 = 20$ (không phải 19).
    \item Ý c) Đúng: $2 : 3$ dư 2, $5 : 3$ dư 2, $8 : 3$ dư 2,... Khẳng định đúng.
    \item Ý d) Sai: Các số trong dãy chia 3 dư 2, còn 30 chia hết cho 3 nên 30 không thuộc dãy.
\end{itemize}
\begin{itemize}
    \item a) \textbf{Đúng}.
    \item b) \textbf{Sai}.
    \item c) \textbf{Đúng}.
    \item d) \textbf{Sai}.
\end{itemize}
}
\end{ex}

\textbf{Phần 3. Trả lời ngắn}
\setcounterpageref{ex}{0}

\begin{ex} %[2D1V06]
Một bác thợ cưa một cây tre thành 6 đoạn. Mỗi nhát cưa mất 2 phút. Hỏi bác thợ cưa xong cây tre trong bao nhiêu phút?
\shortans{10}
\loigiai{
\textbf{Đáp số:} $10$\\
Để có 6 đoạn cần $6 - 1 = 5$ nhát cưa. Thời gian cưa là: $5 \times 2 = 10$ phút.
}
\end{ex}

\begin{ex} %[2D1V06]
Tìm số bé nhất có hai chữ số mà chữ số hàng đơn vị hơn chữ số hàng chục 6 đơn vị.
\shortans{17}
\loigiai{
\textbf{Đáp số:} $17$\\
Chọn chữ số hàng chục nhỏ nhất là 1, chữ số hàng đơn vị là $1 + 6 = 7$. Số đó là 17.
}
\end{ex}

\begin{ex} %[2D1V06]
An và Bình có tất cả 46 viên bi. An cho Bình 5 viên bi thì hai bạn có số bi bằng nhau. Hỏi lúc đầu An có bao nhiêu viên bi?
\shortans{28}
\loigiai{
\textbf{Đáp số:} $28$\\
Khi bằng nhau mỗi bạn có $46 : 2 = 23$ viên. Lúc đầu An có: $23 + 5 = 28$ viên.
}
\end{ex}

\begin{ex} %[2D1V06]
Tìm một số biết rằng lấy số lớn nhất có hai chữ số trừ đi số đó rồi trừ tiếp 25 thì được 30.
\shortans{44}
\loigiai{
\textbf{Đáp số:} $44$\\
Số đó là: $99 - 25 - 30 = 44$.
}
\end{ex}

\begin{ex} %[2D1V06]
Lớp 2A xếp thành 4 hàng, mỗi hàng có 8 học sinh và thừa ra 3 bạn. Hỏi lớp 2A có tất cả bao nhiêu học sinh?
\shortans{35}
\loigiai{
\textbf{Đáp số:} $35$\\
$4 \times 8 + 3 = 32 + 3 = 35$ học sinh.
}
\end{ex}

\begin{ex} %[2D1V06]
Tìm x biết: $x + 14 = 85 - 21$.
\shortans{50}
\loigiai{
\textbf{Đáp số:} $50$\\
Vế phải: $85 - 21 = 64$. Do đó $x = 64 - 14 = 50$.
}
\end{ex}

\newpage


% =========================================================================
% CHỦ ĐỀ 2 - PHÉP CỘNG, PHÉP TRỪ TRONG PHẠM VI 20
% =========================================================================
\section{Phép cộng, phép trừ trong phạm vi 20}

\TaoPhieuKiemTra{Chủ đề 2 -- Phép cộng, phép trừ trong phạm vi 20}{07}{Phép cộng (qua 10) trong phạm vi 20}{https://loptoanthaylinh.vn/thi/chuyen-de/2/toan2_cd_b07}
\textbf{Phần 1. Trắc nghiệm lựa chọn}
\setcounterpageref{ex}{0}

\begin{ex} %[2D2V07]
Tính giá trị của biểu thức: $9 + 4 + 1 + 6$.
\choice
{$19$}
{$18$}
{\True $20$}
{$21$}
\loigiai{
Gộp các số tạo 10: $(9 + 1) + (4 + 6) = 10 + 10 = 20$.
\textbf{Đáp án của câu hỏi là C}
}
\end{ex}

\begin{ex} %[2D2V07]
Trong hộp có $8$ viên bi xanh và $7$ viên bi đỏ. Bạn Nam lấy thêm vào hộp $3$ viên bi vàng. Hỏi trong hộp có tất cả bao nhiêu viên bi?
\choice
{\True $18$ viên bi}
{$17$ viên bi}
{$19$ viên bi}
{$15$ viên bi}
\loigiai{
Tổng số viên bi là: $8 + 7 + 3 = 8 + (7 + 3) = 8 + 10 = 18$ (viên bi).
\textbf{Đáp án của câu hỏi là A}
}
\end{ex}

\begin{ex} %[2D2V07]
Tìm số tự nhiên $x$ biết: $x - 7 = 9$.
\choice
{$15$}
{$17$}
{$18$}
{\True $16$}
\loigiai{
Ta có: $x = 9 + 7 = 16$.
\textbf{Đáp án của câu hỏi là D}
}
\end{ex}

\begin{ex} %[2D2V07]
Mẹ mua về $9$ quả cam, sau đó mẹ mua thêm $8$ quả cam nữa. Mẹ biếu bà $6$ quả cam. Hỏi mẹ còn lại bao nhiêu quả cam?
\choice
{$12$ quả cam}
{\True $11$ quả cam}
{$10$ quả cam}
{$13$ quả cam}
\loigiai{
Tổng số cam mẹ mua: $9 + 8 = 17$ (quả).\nSố cam mẹ còn lại: $17 - 6 = 11$ (quả).
\textbf{Đáp án của câu hỏi là B}
}
\end{ex}

\begin{ex} %[2D2V07]
Điền số thích hợp vào dấu hỏi chấm: $8 + ? = 15$.
\choice
{$6$}
{$8$}
{\True $7$}
{$9$}
\loigiai{
Ta có: $? = 15 - 8 = 7$.
\textbf{Đáp án của câu hỏi là C}
}
\end{ex}

\begin{ex} %[2D2V07]
Lan hái được $9$ bông hoa, Mai hái được nhiều hơn Lan $5$ bông hoa. Hỏi cả hai bạn hái được bao nhiêu bông hoa?
\choice
{\True $23$ bông hoa}
{$22$ bông hoa}
{$14$ bông hoa}
{$24$ bông hoa}
\loigiai{
Mai hái được: $9 + 5 = 14$ (bông hoa).\nCả hai bạn hái được: $9 + 14 = 23$ (bông hoa).
\textbf{Đáp án của câu hỏi là A}
}
\end{ex}

\begin{ex} %[2D2V07]
Cho phép tính: $8 + 6 \dots 9 + 4$. Dấu thích hợp điền vào chỗ chấm là
\choice
{$<$}
{$=$}
{Không so sánh được}
{\True $>$}
\loigiai{
Vế trái: $8 + 6 = 14$. Vế phải: $9 + 4 = 13$. Vì $14 > 13$ nên điền dấu $>$. 
\textbf{Đáp án của câu hỏi là D}
}
\end{ex}

\begin{ex} %[2D2V07]
Một xe khách có $9$ hành khách. Đến bến có $7$ người lên xe và $4$ người xuống xe. Hỏi trên xe lúc này có bao nhiêu hành khách?
\choice
{$11$ người}
{\True $12$ người}
{$13$ người}
{$10$ người}
\loigiai{
Số người trên xe lúc này là: $9 + 7 - 4 = 16 - 4 = 12$ (người).
\textbf{Đáp án của câu hỏi là B}
}
\end{ex}

\begin{ex} %[2D2V07]
Tổng của số lớn nhất có 1 chữ số và số chẵn lớn nhất có 1 chữ số là
\choice
{$18$}
{$16$}
{\True $17$}
{$19$}
\loigiai{
Số lớn nhất có 1 chữ số là 9. Số chẵn lớn nhất có 1 chữ số là 8.\nTổng là: $9 + 8 = 17$.
\textbf{Đáp án của câu hỏi là C}
}
\end{ex}

\begin{ex} %[2D2V07]
Tìm một số biết rằng lấy số đó trừ đi $8$ thì được kết quả bằng $7 + 2$.
\choice
{\True $17$}
{$16$}
{$18$}
{$15$}
\loigiai{
Vế phải: $7 + 2 = 9$. Số đó là: $9 + 8 = 17$.
\textbf{Đáp án của câu hỏi là A}
}
\end{ex}

\begin{ex} %[2D2V07]
Có ba túi kẹo: túi thứ nhất có $6$ cái, túi thứ hai có $7$ cái, túi thứ ba có $5$ cái. Hỏi cả ba túi có bao nhiêu cái kẹo?
\choice
{$17$ cái kẹo}
{$19$ cái kẹo}
{$16$ cái kẹo}
{\True $18$ cái kẹo}
\loigiai{
Tổng số kẹo là: $6 + 7 + 5 = 13 + 5 = 18$ (cái kẹo).
\textbf{Đáp án của câu hỏi là D}
}
\end{ex}

\begin{ex} %[2D2V07]
Điền số thích hợp: $7 + 8 = 10 + \dots$
\choice
{$4$}
{\True $5$}
{$6$}
{$7$}
\loigiai{
Vế trái: $7 + 8 = 15$. Ta có: $10 + 5 = 15$. Vậy điền số 5.
\textbf{Đáp án của câu hỏi là B}
}
\end{ex}

\textbf{Phần 2. Trắc nghiệm đúng sai}
\setcounterpageref{ex}{0}

\begin{ex} %[2D2V07]
Xét tính đúng sai của các phép tính cộng qua 10 trong phạm vi 20:
\choiceTF
{\True Phép tính $9 + 6$ có kết quả bằng $15$.}
{Khi cộng 8 với 5 ta tách 5 thành 3 và 2 để $8 + 3 = 10$.}
{\True Tổng của $7 + 8$ và $8 + 7$ luôn bằng nhau.}
{Tổng của hai số có một chữ số không thể vượt quá 17.}
\loigiai{
\begin{itemize}
    \item Ý a) Đúng: $9 + 6 = 15$. Đúng.
    \item Ý b) Sai: Phải tách 5 thành 2 và 3 để $8 + 2 = 10$, sau đó $10 + 3 = 13$.
    \item Ý c) Đúng: Tính chất giao hoán của phép cộng: $7 + 8 = 8 + 7 = 15$.
    \item Ý d) Sai: Tổng lớn nhất là $9 + 9 = 18 > 17$.
\end{itemize}
\begin{itemize}
    \item a) \textbf{Đúng}.
    \item b) \textbf{Sai}.
    \item c) \textbf{Đúng}.
    \item d) \textbf{Sai}.
\end{itemize}
}
\end{ex}

\begin{ex} %[2D2V07]
Hai bạn An và Bình cùng chơi ném bóng vào rổ. An ném trúng 8 quả, Bình ném trúng 9 quả:
\choiceTF
{\True Cả hai bạn ném trúng tất cả 17 quả bóng.}
{Bình ném trúng ít hơn An 1 quả bóng.}
{\True Nếu An ném trúng thêm 3 quả nữa thì An sẽ nhiều hơn Bình 2 quả.}
{Để cả hai bạn có số bóng bằng nhau thì Bình phải cho An 2 quả.}
\loigiai{
\begin{itemize}
    \item Ý a) Đúng: $8 + 9 = 17$ quả.
    \item Ý b) Sai: Bình ném được 9 quả, An 8 quả nên Bình nhiều hơn An 1 quả.
    \item Ý c) Đúng: Lúc đó An có: $8 + 3 = 11$ quả. $11 - 9 = 2$ quả. Đúng.
    \item Ý d) Sai: Bình chỉ hơn An 1 quả nên không thể cho 2 quả để bằng nhau.
\end{itemize}
\begin{itemize}
    \item a) \textbf{Đúng}.
    \item b) \textbf{Sai}.
    \item c) \textbf{Đúng}.
    \item d) \textbf{Sai}.
\end{itemize}
}
\end{ex}

\textbf{Phần 3. Trả lời ngắn}
\setcounterpageref{ex}{0}

\begin{ex} %[2D2V07]
Tính giá trị của biểu thức: $7 + 6 + 3 + 4$.
\shortans{20}
\loigiai{
\textbf{Đáp số:} $20$\\
$(7 + 3) + (6 + 4) = 10 + 10 = 20$.
}
\end{ex}

\begin{ex} %[2D2V07]
Bác nông dân nuôi 8 con vịt, sau đó bác mua thêm 9 con vịt nữa. Hôm sau bác bán đi 5 con vịt. Hỏi bác còn lại bao nhiêu con vịt?
\shortans{12}
\loigiai{
\textbf{Đáp số:} $12$\\
Tổng số vịt: $8 + 9 = 17$. Còn lại: $17 - 5 = 12$ con.
}
\end{ex}

\begin{ex} %[2D2V07]
Tìm số tự nhiên x biết: $x - 8 = 6 + 3$.
\shortans{17}
\loigiai{
\textbf{Đáp số:} $17$\\
$x - 8 = 9 \Rightarrow x = 9 + 8 = 17$.
}
\end{ex}

\begin{ex} %[2D2V07]
Một cuộn ruy băng dài 19 cm. Bạn Lan cắt lấy một đoạn dài 8 cm, sau đó cắt tiếp một đoạn dài 6 cm. Hỏi cuộn ruy băng còn lại dài bao nhiêu xăng-ti-mét?
\shortans{5}
\loigiai{
\textbf{Đáp số:} $5$\\
Độ dài còn lại: $19 - 8 - 6 = 5\text{ cm}$.
}
\end{ex}

\begin{ex} %[2D2V07]
Tìm một số biết rằng lấy số đó cộng với 7 thì được kết quả bằng 16.
\shortans{9}
\loigiai{
\textbf{Đáp số:} $9$\\
Số đó là: $16 - 7 = 9$.
}
\end{ex}

\begin{ex} %[2D2V07]
Có bao nhiêu cặp số tự nhiên có một chữ số mà tổng của chúng bằng 14?
\shortans{5}
\loigiai{
\textbf{Đáp số:} $5$\\
Các cặp là: $(5, 9), (6, 8), (7, 7), (8, 6), (9, 5)$. Có 5 cặp.
}
\end{ex}

\newpage

\TaoPhieuKiemTra{Chủ đề 2 -- Phép cộng, phép trừ trong phạm vi 20}{08}{Bảng cộng (qua 10)}{https://loptoanthaylinh.vn/thi/chuyen-de/2/toan2_cd_b08}
\textbf{Phần 1. Trắc nghiệm lựa chọn}
\setcounterpageref{ex}{0}

\begin{ex} %[2D2V08]
Kết quả của phép tính $9 + 7$ là
\choice
{$15$}
{\True $16$}
{$17$}
{$18$}
\loigiai{
$9 + 7 = 16$.
\textbf{Đáp án của câu hỏi là B}
}
\end{ex}

\begin{ex} %[2D2V08]
Trong vườn có $8$ cây cam và số cây bưởi bằng số cây cam cộng thêm $5$ cây. Hỏi trong vườn có tất cả bao nhiêu cây cam và bưởi?
\choice
{$13$ cây}
{$18$ cây}
{$19$ cây}
{\True $21$ cây}
\loigiai{
Số cây bưởi là: $8 + 5 = 13$ (cây).\nTổng số cây cam và bưởi là: $8 + 13 = 21$ (cây).
\textbf{Đáp án của câu hỏi là D}
}
\end{ex}

\begin{ex} %[2D2V08]
Tính nhanh giá trị: $8 + 5 + 2 + 5$.
\choice
{\True $20$}
{$18$}
{$19$}
{$21$}
\loigiai{
$(8 + 2) + (5 + 5) = 10 + 10 = 20$.
\textbf{Đáp án của câu hỏi là A}
}
\end{ex}

\begin{ex} %[2D2V08]
Tìm một số biết rằng lấy số đó trừ đi $9$ thì được số chẵn nhỏ nhất có hai chữ số.
\choice
{$18$}
{$20$}
{\True $19$}
{$17$}
\loigiai{
Số chẵn nhỏ nhất có hai chữ số là 10. Số đó là: $10 + 9 = 19$.
\textbf{Đáp án của câu hỏi là C}
}
\end{ex}

\begin{ex} %[2D2V08]
Hộp thứ nhất có $7$ chiếc bút chì, hộp thứ hai có $8$ chiếc bút chì. Bạn Minh lấy ra $4$ chiếc bút chì từ hộp thứ nhất. Hỏi cả hai hộp còn lại bao nhiêu chiếc bút chì?
\choice
{$12$ chiếc}
{\True $11$ chiếc}
{$10$ chiếc}
{$13$ chiếc}
\loigiai{
Tổng số bút lúc đầu: $7 + 8 = 15$ (chiếc).\nSau khi lấy ra còn: $15 - 4 = 11$ (chiếc).
\textbf{Đáp án của câu hỏi là B}
}
\end{ex}

\begin{ex} %[2D2V08]
Điền dấu thích hợp: $6 + 8 \dots 7 + 7$.
\choice
{$>$}
{$<$}
{Không so sánh được}
{\True $=$}
\loigiai{
$6 + 8 = 14$; $7 + 7 = 14$. Hai vế bằng nhau nên điền dấu $=$.
\textbf{Đáp án của câu hỏi là D}
}
\end{ex}

\begin{ex} %[2D2V08]
Một sợi dây dài $18\text{ cm}$. Người ta cắt thành hai đoạn, biết đoạn thứ nhất dài $9\text{ cm}$. Đoạn thứ hai dài bao nhiêu?
\choice
{\True $9\text{ cm}$}
{$8\text{ cm}$}
{$10\text{ cm}$}
{$7\text{ cm}$}
\loigiai{
Đoạn thứ hai dài: $18 - 9 = 9\text{ (cm)}$.
\textbf{Đáp án của câu hỏi là A}
}
\end{ex}

\begin{ex} %[2D2V08]
Tìm số tự nhiên $y$ biết: $y + 7 = 16$.
\choice
{$8$}
{$10$}
{\True $9$}
{$7$}
\loigiai{
$y = 16 - 7 = 9$.
\textbf{Đáp án của câu hỏi là C}
}
\end{ex}

\begin{ex} %[2D2V08]
Bố mua cho Mai $8$ quyển vở, mẹ mua cho Mai thêm $6$ quyển vở nữa. Mai đã dùng hết $5$ quyển. Hỏi Mai còn lại bao nhiêu quyển vở chưa dùng?
\choice
{$8$ quyển}
{\True $9$ quyển}
{$10$ quyển}
{$7$ quyển}
\loigiai{
Tổng số vở Mai có: $8 + 6 = 14$ (quyển).\nSố vở còn lại: $14 - 5 = 9$ (quyển).
\textbf{Đáp án của câu hỏi là B}
}
\end{ex}

\begin{ex} %[2D2V08]
Tổng của hai số bằng $15$. Nếu tăng số hạng thứ nhất thêm $2$ đơn vị và tăng số hạng thứ hai thêm $3$ đơn vị thì tổng mới là
\choice
{$18$}
{$19$}
{$21$}
{\True $20$}
\loigiai{
Tổng mới là: $15 + 2 + 3 = 20$.
\textbf{Đáp án của câu hỏi là D}
}
\end{ex}

\begin{ex} %[2D2V08]
Hai số có hiệu bằng $8$, số trừ là $9$. Số bị trừ là
\choice
{\True $17$}
{$16$}
{$18$}
{$15$}
\loigiai{
Số bị trừ = Hiệu + Số trừ = $8 + 9 = 17$.
\textbf{Đáp án của câu hỏi là A}
}
\end{ex}

\begin{ex} %[2D2V08]
Mai có $9$ bông hoa, Cúc có $7$ bông hoa, Lan có số hoa bằng số hoa của Mai và Cúc cộng lại rồi bớt đi $4$ bông. Lan có bao nhiêu bông hoa?
\choice
{$11$ bông hoa}
{$13$ bông hoa}
{\True $12$ bông hoa}
{$10$ bông hoa}
\loigiai{
Số hoa của Lan: $9 + 7 - 4 = 16 - 4 = 12$ (bông hoa).
\textbf{Đáp án của câu hỏi là C}
}
\end{ex}

\textbf{Phần 2. Trắc nghiệm đúng sai}
\setcounterpageref{ex}{0}

\begin{ex} %[2D2V08]
Xét tính đúng sai của các phép tính trong bảng cộng có nhớ:
\choiceTF
{\True Phép cộng $6 + 7$ có kết quả bằng $13$.}
{Tổng của số lẻ lớn nhất có 1 chữ số và số chẵn nhỏ nhất có 1 chữ số khác 0 là 10.}
{\True Biểu thức $8 + 9$ có giá trị lớn hơn biểu thức $7 + 9$.}
{\True Nếu $a + 6 = 15$ thì $a = 9$.}
\loigiai{
\begin{itemize}
    \item Ý a) Đúng: $6 + 7 = 13$.
    \item Ý b) Sai: Số lẻ lớn nhất có 1 chữ số là 9, số chẵn nhỏ nhất khác 0 là 2. Tổng là $9 + 2 = 11$ (không phải 10).
    \item Ý c) Đúng: $8 + 9 = 17 > 7 + 9 = 16$.
    \item Ý d) Đúng: $15 - 6 = 9$. Đúng.
\end{itemize}
\begin{itemize}
    \item a) \textbf{Đúng}.
    \item b) \textbf{Sai}.
    \item c) \textbf{Đúng}.
    \item d) \textbf{Đúng}.
\end{itemize}
}
\end{ex}

\begin{ex} %[2D2V08]
Một đàn gà có 9 con gà trống và 8 con gà mái:
\choiceTF
{\True Tổng số gà trong đàn là 17 con.}
{Số gà mái nhiều hơn số gà trống 1 con.}
{\True Nếu mua thêm 3 con gà mái thì số gà mái sẽ nhiều hơn số gà trống 2 con.}
{\True Tổng số chân của 1 con gà trống và 1 con gà mái là 4 chân.}
\loigiai{
\begin{itemize}
    \item Ý a) Đúng: $9 + 8 = 17$ con.
    \item Ý b) Sai: Gà trống có 9 con, mái có 8 con nên gà trống nhiều hơn gà mái 1 con.
    \item Ý c) Đúng: Gà mái lúc này là $8 + 3 = 11$. $11 - 9 = 2$ con. Đúng.
    \item Ý d) Đúng: Mỗi con gà có 2 chân, $2 + 2 = 4$ chân.
\end{itemize}
\begin{itemize}
    \item a) \textbf{Đúng}.
    \item b) \textbf{Sai}.
    \item c) \textbf{Đúng}.
    \item d) \textbf{Đúng}.
\end{itemize}
}
\end{ex}

\textbf{Phần 3. Trả lời ngắn}
\setcounterpageref{ex}{0}

\begin{ex} %[2D2V08]
Tính giá trị biểu thức: $9 + 8 - 7$.
\shortans{10}
\loigiai{
\textbf{Đáp số:} $10$\\
$9 + 8 = 17$; $17 - 7 = 10$.
}
\end{ex}

\begin{ex} %[2D2V08]
Nam có 8 viên bi, Hoàng có nhiều hơn Nam 5 viên bi. Hỏi cả hai bạn có tất cả bao nhiêu viên bi?
\shortans{21}
\loigiai{
\textbf{Đáp số:} $21$\\
Hoàng có: $8 + 5 = 13$ viên. Cả hai bạn: $8 + 13 = 21$ viên.
}
\end{ex}

\begin{ex} %[2D2V08]
Tìm số tự nhiên a biết: $a + 8 = 17$.
\shortans{9}
\loigiai{
\textbf{Đáp số:} $9$\\
$a = 17 - 8 = 9$.
}
\end{ex}

\begin{ex} %[2D2V08]
Tìm tổng của số lớn nhất có 1 chữ số và số liền sau của nó.
\shortans{19}
\loigiai{
\textbf{Đáp số:} $19$\\
Số lớn nhất có 1 chữ số là 9. Số liền sau là 10. Tổng: $9 + 10 = 19$.
}
\end{ex}

\begin{ex} %[2D2V08]
Lớp 2A có 9 bạn tham gia thi hát, số bạn thi vẽ nhiều hơn thi hát 4 bạn. Hỏi có tất cả bao nhiêu bạn tham gia thi cả hai môn?
\shortans{22}
\loigiai{
\textbf{Đáp số:} $22$\\
Thi vẽ: $9 + 4 = 13$. Tổng số: $9 + 13 = 22$ bạn.
}
\end{ex}

\begin{ex} %[2D2V08]
Có bao nhiêu số có một chữ số lớn hơn 6 mà khi cộng với 9 cho kết quả là số có hai chữ số?
\shortans{3}
\loigiai{
\textbf{Đáp số:} $3$\\
Các số lớn hơn 6 là 7, 8, 9 (có 3 số).
}
\end{ex}

\newpage

\TaoPhieuKiemTra{Chủ đề 2 -- Phép cộng, phép trừ trong phạm vi 20}{09}{Bài toán về thêm, bớt một số đơn vị}{https://loptoanthaylinh.vn/thi/chuyen-de/2/toan2_cd_b09}
\textbf{Phần 1. Trắc nghiệm lựa chọn}
\setcounterpageref{ex}{0}

\begin{ex} %[2D2V09]
Kết quả của phép tính $15 - 8$ là
\choice
{\True $7$}
{$8$}
{$6$}
{$9$}
\loigiai{
$15 - 8 = 7$.
\textbf{Đáp án của câu hỏi là A}
}
\end{ex}

\begin{ex} %[2D2V09]
Một cành cây có $14$ con chim đậu. Bắn một tiếng súng làm $6$ con chim bay đi. Lát sau lại có $3$ con chim bay đi nữa. Hỏi trên cành cây còn lại bao nhiêu con chim?
\choice
{$6$ con chim}
{$4$ con chim}
{\True $5$ con chim}
{$7$ con chim}
\loigiai{
Số chim còn lại: $14 - 6 - 3 = 8 - 3 = 5$ (con chim).
\textbf{Đáp án của câu hỏi là C}
}
\end{ex}

\begin{ex} %[2D2V09]
Tìm $x$ biết: $13 - x = 7$.
\choice
{$5$}
{\True $6$}
{$7$}
{$8$}
\loigiai{
$x = 13 - 7 = 6$.
\textbf{Đáp án của câu hỏi là B}
}
\end{ex}

\begin{ex} %[2D2V09]
Bà có $17$ quả trứng gà. Bà đã dùng $9$ quả để làm bánh. Số trứng gà bà còn lại là
\choice
{$9$ quả}
{$7$ quả}
{$6$ quả}
{\True $8$ quả}
\loigiai{
Số trứng gà còn lại: $17 - 9 = 8$ (quả).
\textbf{Đáp án của câu hỏi là D}
}
\end{ex}

\begin{ex} %[2D2V09]
Tính giá trị của biểu thức: $16 - 9 + 5$.
\choice
{\True $12$}
{$11$}
{$13$}
{$10$}
\loigiai{
$16 - 9 = 7$; $7 + 5 = 12$.
\textbf{Đáp án của câu hỏi là A}
}
\end{ex}

\begin{ex} %[2D2V09]
Mẹ mua $12$ bông hoa hồng và hoa cúc, trong đó có $7$ bông hoa hồng. Hỏi mẹ mua bao nhiêu bông hoa cúc?
\choice
{$6$ bông hoa}
{$4$ bông hoa}
{\True $5$ bông hoa}
{$7$ bông hoa}
\loigiai{
Số bông hoa cúc: $12 - 7 = 5$ (bông hoa).
\textbf{Đáp án của câu hỏi là C}
}
\end{ex}

\begin{ex} %[2D2V09]
Điền dấu thích hợp: $14 - 7 \dots 13 - 6$.
\choice
{$>$}
{\True $=$}
{$<$}
{Không so sánh được}
\loigiai{
$14 - 7 = 7$; $13 - 6 = 7$. Hai vế bằng nhau nên điền dấu $=$.
\textbf{Đáp án của câu hỏi là B}
}
\end{ex}

\begin{ex} %[2D2V09]
Hiệu giữa số chẵn lớn nhất trong phạm vi $20$ và số lẻ lớn nhất có 1 chữ số là
\choice
{$10$}
{$12$}
{$9$}
{\True $11$}
\loigiai{
Số chẵn lớn nhất trong phạm vi 20 là 20. Số lẻ lớn nhất có 1 chữ số là 9.\nHiệu: $20 - 9 = 11$.
\textbf{Đáp án của câu hỏi là D}
}
\end{ex}

\begin{ex} %[2D2V09]
Tìm một số biết rằng lấy $18$ trừ đi số đó rồi cộng thêm $4$ thì được $13$.
\choice
{\True $9$}
{$8$}
{$10$}
{$7$}
\loigiai{
Ta có: $18 - x = 13 - 4 = 9 \Rightarrow x = 18 - 9 = 9$.
\textbf{Đáp án của câu hỏi là A}
}
\end{ex}

\begin{ex} %[2D2V09]
Một bác thợ làm được $16$ sản phẩm trong hai ngày. Ngày thứ nhất bác làm được $9$ sản phẩm. Hỏi ngày thứ nhất bác làm nhiều hơn ngày thứ hai bao nhiêu sản phẩm?
\choice
{$7$ sản phẩm}
{$3$ sản phẩm}
{\True $2$ sản phẩm}
{$1$ sản phẩm}
\loigiai{
Ngày thứ hai bác làm: $16 - 9 = 7$ (sản phẩm).\nNgày thứ nhất hơn ngày thứ hai: $9 - 7 = 2$ (sản phẩm).
\textbf{Đáp án của câu hỏi là C}
}
\end{ex}

\begin{ex} %[2D2V09]
Kết quả của phép trừ $11 - 4$ bằng kết quả của phép tính nào dưới đây?
\choice
{$15 - 7$}
{\True $16 - 9$}
{$13 - 5$}
{$12 - 6$}
\loigiai{
$11 - 4 = 7$. Ta thấy $16 - 9 = 7$. Do đó chọn $16 - 9$.
\textbf{Đáp án của câu hỏi là B}
}
\end{ex}

\begin{ex} %[2D2V09]
Trong rổ có $15$ quả cam. Mai lấy ra $6$ quả, Nam lấy ra $4$ quả. Hỏi trong rổ còn lại bao nhiêu quả cam?
\choice
{$6$ quả}
{$7$ quả}
{$4$ quả}
{\True $5$ quả}
\loigiai{
Hai bạn lấy ra: $6 + 4 = 10$ (quả).\nTrong rổ còn: $15 - 10 = 5$ (quả).
\textbf{Đáp án của câu hỏi là D}
}
\end{ex}

\textbf{Phần 2. Trắc nghiệm đúng sai}
\setcounterpageref{ex}{0}

\begin{ex} %[2D2V09]
Xét tính đúng sai của các phép trừ có nhớ trong phạm vi 20:
\choiceTF
{\True Phép trừ $14 - 8$ có hiệu bằng $6$.}
{Hiệu của $15 - 7$ bằng hiệu của $16 - 9$.}
{\True Trong phép trừ $17 - 9 = 8$, số bị trừ là 17, số trừ là 9.}
{\True Nếu lấy 18 trừ đi một số có một chữ số thì hiệu luôn lớn hơn hoặc bằng 9.}
\loigiai{
\begin{itemize}
    \item Ý a) Đúng: $14 - 8 = 6$.
    \item Ý b) Sai: $15 - 7 = 8$, còn $16 - 9 = 7$. Hai kết quả khác nhau.
    \item Ý c) Đúng: Đúng theo tên gọi các thành phần.
    \item Ý d) Đúng: Số trừ lớn nhất là 9 nên $18 - 9 = 9$. Hiệu luôn $\ge 9$.
\end{itemize}
\begin{itemize}
    \item a) \textbf{Đúng}.
    \item b) \textbf{Sai}.
    \item c) \textbf{Đúng}.
    \item d) \textbf{Đúng}.
\end{itemize}
}
\end{ex}

\begin{ex} %[2D2V09]
Một nhóm bạn gồm 15 bạn chơi trò trốn tìm, trong đó đã có 8 bạn bị tìm thấy:
\choiceTF
{\True Còn lại 7 bạn chưa bị tìm thấy.}
{Số bạn đã bị tìm thấy ít hơn số bạn chưa bị tìm thấy 1 bạn.}
{\True Nếu tìm thấy thêm 3 bạn nữa thì số bạn chưa bị tìm thấy là 4 bạn.}
{Số bạn bị tìm thấy gấp đôi số bạn chưa bị tìm thấy.}
\loigiai{
\begin{itemize}
    \item Ý a) Đúng: $15 - 8 = 7$ bạn.
    \item Ý b) Sai: 8 bạn bị tìm thấy nhiều hơn 7 bạn chưa tìm thấy 1 bạn.
    \item Ý c) Đúng: $7 - 3 = 4$ bạn.
    \item Ý d) Sai: $8$ không gấp đôi $7$. Khẳng định sai.
\end{itemize}
\begin{itemize}
    \item a) \textbf{Đúng}.
    \item b) \textbf{Sai}.
    \item c) \textbf{Đúng}.
    \item d) \textbf{Sai}.
\end{itemize}
}
\end{ex}

\textbf{Phần 3. Trả lời ngắn}
\setcounterpageref{ex}{0}

\begin{ex} %[2D2V09]
Tính giá trị của biểu thức: $17 - 8 + 6$.
\shortans{15}
\loigiai{
\textbf{Đáp số:} $15$\\
$17 - 8 = 9$; $9 + 6 = 15$.
}
\end{ex}

\begin{ex} %[2D2V09]
Trong bể có 16 con cá. Bố vớt ra 7 con cá đem tặng bạn. Hỏi trong bể còn lại bao nhiêu con cá?
\shortans{9}
\loigiai{
\textbf{Đáp số:} $9$\\
Số cá còn lại: $16 - 7 = 9$ con.
}
\end{ex}

\begin{ex} %[2D2V09]
Tìm số tự nhiên x biết: $15 - x = 8$.
\shortans{7}
\loigiai{
\textbf{Đáp số:} $7$\\
$x = 15 - 8 = 7$.
}
\end{ex}

\begin{ex} %[2D2V09]
Có hai ngăn kéo đựng tất cả 18 quyển sách. Ngăn thứ nhất có 9 quyển sách. Hỏi ngăn thứ nhất có nhiều hơn hay ít hơn ngăn thứ hai bao nhiêu quyển sách?
\shortans{0}
\loigiai{
\textbf{Đáp số:} $0$\\
Ngăn 2 có: $18 - 9 = 9$ quyển. Hai ngăn bằng nhau nên hiệu bằng 0.
}
\end{ex}

\begin{ex} %[2D2V09]
Một cây tre dài 14 m bị bão làm gãy một đoạn ngọn dài 6 m. Hỏi đoạn gốc còn lại dài bao nhiêu mét?
\shortans{8}
\loigiai{
\textbf{Đáp số:} $8$\\
Đoạn gốc dài: $14 - 6 = 8\text{ m}$.
}
\end{ex}

\begin{ex} %[2D2V09]
Tìm hiệu của số tròn chục nhỏ nhất có hai chữ số và số lẻ nhỏ nhất có 1 chữ số khác 1.
\shortans{7}
\loigiai{
\textbf{Đáp số:} $7$\\
Số tròn chục nhỏ nhất là 10. Số lẻ nhỏ nhất khác 1 là 3. Hiệu là: $10 - 3 = 7$.
}
\end{ex}

\newpage

\TaoPhieuKiemTra{Chủ đề 2 -- Phép cộng, phép trừ trong phạm vi 20}{10}{Luyện tập chung}{https://loptoanthaylinh.vn/thi/chuyen-de/2/toan2_cd_b10}
\textbf{Phần 1. Trắc nghiệm lựa chọn}
\setcounterpageref{ex}{0}

\begin{ex} %[2D2V10]
Tính giá trị: $7 + 9 - 8$.
\choice
{$7$}
{$9$}
{$6$}
{\True $8$}
\loigiai{
$7 + 9 = 16$; $16 - 8 = 8$.
\textbf{Đáp án của câu hỏi là D}
}
\end{ex}

\begin{ex} %[2D2V10]
Một lớp học có $9$ bạn nam và $8$ bạn nữ. Sau đó có thêm $2$ bạn nữ chuyển vào lớp. Hỏi lớp học lúc này có bao nhiêu học sinh?
\choice
{$18$ học sinh}
{\True $19$ học sinh}
{$20$ học sinh}
{$17$ học sinh}
\loigiai{
Tổng số học sinh lúc này: $9 + 8 + 2 = 9 + 10 = 19$ (học sinh).
\textbf{Đáp án của câu hỏi là B}
}
\end{ex}

\begin{ex} %[2D2V10]
Tìm $x$ biết: $x + 6 = 8 + 7$.
\choice
{$8$}
{$10$}
{\True $9$}
{$7$}
\loigiai{
Vế phải: $8 + 7 = 15$. Ta có: $x = 15 - 6 = 9$.
\textbf{Đáp án của câu hỏi là C}
}
\end{ex}

\begin{ex} %[2D2V10]
Bạn Mai có $7$ cái kẹo, bạn Hoa có nhiều hơn Mai $4$ cái kẹo. Hỏi cả hai bạn có tất cả bao nhiêu cái kẹo?
\choice
{\True $18$ cái kẹo}
{$17$ cái kẹo}
{$11$ cái kẹo}
{$19$ cái kẹo}
\loigiai{
Hoa có: $7 + 4 = 11$ (cái kẹo).\nCả hai bạn có: $7 + 11 = 18$ (cái kẹo).
\textbf{Đáp án của câu hỏi là A}
}
\end{ex}

\begin{ex} %[2D2V10]
Điền dấu thích hợp: $8 + 8 \dots 9 + 7$.
\choice
{$>$}
{$<$}
{Không so sánh được}
{\True $=$}
\loigiai{
$8 + 8 = 16$; $9 + 7 = 16$. Hai vế bằng nhau nên điền dấu $=$.
\textbf{Đáp án của câu hỏi là D}
}
\end{ex}

\begin{ex} %[2D2V10]
Tìm một số biết rằng lấy số đó cộng với $5$ rồi trừ đi $8$ thì được $6$.
\choice
{$8$}
{\True $9$}
{$10$}
{$7$}
\loigiai{
Tính ngược từ cuối: $6 + 8 - 5 = 14 - 5 = 9$.
\textbf{Đáp án của câu hỏi là B}
}
\end{ex}

\begin{ex} %[2D2V10]
Có hai bao gạo, bao thứ nhất nặng $9\text{ kg}$, bao thứ hai nặng $8\text{ kg}$. Người ta đã ăn hết $5\text{ kg}$. Hỏi còn lại bao nhiêu ki-lô-gam gạo?
\choice
{$11\text{ kg}$}
{$13\text{ kg}$}
{\True $12\text{ kg}$}
{$10\text{ kg}$}
\loigiai{
Tổng số gạo: $9 + 8 = 17\text{ (kg)}$. Còn lại: $17 - 5 = 12\text{ (kg)}$.
\textbf{Đáp án của câu hỏi là C}
}
\end{ex}

\begin{ex} %[2D2V10]
Cho ba số: $6, 8, 9$. Tổng của hai số lớn nhất trong ba số trên là
\choice
{\True $17$}
{$15$}
{$14$}
{$18$}
\loigiai{
Hai số lớn nhất là 8 và 9. Tổng là: $8 + 9 = 17$.
\textbf{Đáp án của câu hỏi là A}
}
\end{ex}

\begin{ex} %[2D2V10]
Một đàn vịt đang bơi dưới ao có $9$ con, trên bờ có $6$ con. Sau đó có $3$ con từ dưới ao nhảy lên bờ. Hỏi lúc này dưới ao còn bao nhiêu con vịt?
\choice
{$7$ con vịt}
{$5$ con vịt}
{$8$ con vịt}
{\True $6$ con vịt}
\loigiai{
Số vịt dưới ao lúc này: $9 - 3 = 6$ (con vịt).
\textbf{Đáp án của câu hỏi là D}
}
\end{ex}

\begin{ex} %[2D2V10]
Tổng của ba số bằng $19$. Biết hai số đầu là $7$ và $8$. Số thứ ba là
\choice
{$3$}
{\True $4$}
{$5$}
{$6$}
\loigiai{
Tổng hai số đầu là: $7 + 8 = 15$. Số thứ ba là: $19 - 15 = 4$.
\textbf{Đáp án của câu hỏi là B}
}
\end{ex}

\begin{ex} %[2D2V10]
Hai số tự nhiên có tổng bằng $13$ và hiệu bằng $3$. Hai số đó là
\choice
{$9\text{ và } 4$}
{$7\text{ và } 6$}
{\True $8\text{ và } 5$}
{$10\text{ và } 3$}
\loigiai{
Số lớn là: $(13 + 3) : 2 = 8$. Số bé là: $8 - 3 = 5$.
\textbf{Đáp án của câu hỏi là C}
}
\end{ex}

\begin{ex} %[2D2V10]
Điền số thích hợp: $9 + 5 = \dots + 8$.
\choice
{\True $6$}
{$5$}
{$7$}
{$8$}
\loigiai{
Vế trái: $9 + 5 = 14$. Ta có: $14 - 8 = 6$. Vậy điền 6.
\textbf{Đáp án của câu hỏi là A}
}
\end{ex}

\textbf{Phần 2. Trắc nghiệm đúng sai}
\setcounterpageref{ex}{0}

\begin{ex} %[2D2V10]
Xét tính đúng sai của các bài toán tổng hợp cộng trong phạm vi 20:
\choiceTF
{\True Biểu thức $8 + 7 + 3$ có kết quả bằng $18$.}
{Biểu thức $9 + 6$ có kết quả bé hơn $8 + 6$.}
{\True Nếu thêm 2 đơn vị vào một số hạng thì tổng tăng thêm 2 đơn vị.}
{Tổng của hai số có 1 chữ số có thể bằng 20.}
\loigiai{
\begin{itemize}
    \item Ý a) Đúng: $8 + (7 + 3) = 8 + 10 = 18$.
    \item Ý b) Sai: $9 + 6 = 15 > 8 + 6 = 14$.
    \item Ý c) Đúng: Đúng theo tính chất phép cộng.
    \item Ý d) Sai: Tổng lớn nhất là $9 + 9 = 18 < 20$.
\end{itemize}
\begin{itemize}
    \item a) \textbf{Đúng}.
    \item b) \textbf{Sai}.
    \item c) \textbf{Đúng}.
    \item d) \textbf{Sai}.
\end{itemize}
}
\end{ex}

\begin{ex} %[2D2V10]
Bạn Hùng có 8 viên bi đỏ và 9 viên bi xanh:
\choiceTF
{\True Hùng có tất cả 17 viên bi.}
{\True Số bi xanh nhiều hơn số bi đỏ là 1 viên bi.}
{Hùng cho bạn 8 viên bi đỏ thì Hùng còn lại 8 viên bi xanh.}
{\True Nếu Hùng nhặt thêm 3 viên bi vàng thì Hùng có tất cả 20 viên bi.}
\loigiai{
\begin{itemize}
    \item Ý a) Đúng: $8 + 9 = 17$ viên bi.
    \item Ý b) Đúng: $9 - 8 = 1$ viên bi.
    \item Ý c) Sai: Hùng có 9 viên bi xanh nên Hùng còn lại 9 viên bi xanh.
    \item Ý d) Đúng: $17 + 3 = 20$ viên bi. Đúng.
\end{itemize}
\begin{itemize}
    \item a) \textbf{Đúng}.
    \item b) \textbf{Đúng}.
    \item c) \textbf{Sai}.
    \item d) \textbf{Đúng}.
\end{itemize}
}
\end{ex}

\textbf{Phần 3. Trả lời ngắn}
\setcounterpageref{ex}{0}

\begin{ex} %[2D2V10]
Tính giá trị biểu thức: $6 + 7 + 4$.
\shortans{17}
\loigiai{
\textbf{Đáp số:} $17$\\
$(6 + 4) + 7 = 10 + 7 = 17$.
}
\end{ex}

\begin{ex} %[2D2V10]
Lan có 8 cái nhãn vở, mẹ cho Lan thêm 9 cái nhãn vở nữa. Lan tặng bạn 6 cái. Hỏi Lan còn lại bao nhiêu cái nhãn vở?
\shortans{11}
\loigiai{
\textbf{Đáp số:} $11$\\
Lan có: $8 + 9 = 17$. Còn lại: $17 - 6 = 11$ cái.
}
\end{ex}

\begin{ex} %[2D2V10]
Tìm số tự nhiên y biết: $y - 7 = 8 + 1$.
\shortans{16}
\loigiai{
\textbf{Đáp số:} $16$\\
$y - 7 = 9 \Rightarrow y = 9 + 7 = 16$.
}
\end{ex}

\begin{ex} %[2D2V10]
Trong rổ có 18 quả trứng gà và vịt. Biết số trứng gà là 9 quả. Hỏi số trứng vịt là bao nhiêu quả?
\shortans{9}
\loigiai{
\textbf{Đáp số:} $9$\\
$18 - 9 = 9$ quả.
}
\end{ex}

\begin{ex} %[2D2V10]
Tìm tổng của số lớn nhất có một chữ số và số lẻ bé nhất có một chữ số khác 1.
\shortans{12}
\loigiai{
\textbf{Đáp số:} $12$\\
Số lớn nhất có 1 chữ số là 9. Số lẻ bé nhất khác 1 là 3. Tổng: $9 + 3 = 12$.
}
\end{ex}

\begin{ex} %[2D2V10]
Tìm một số biết rằng lấy số đó cộng với 8 thì được số tròn chục nhỏ nhất có hai chữ số.
\shortans{2}
\loigiai{
\textbf{Đáp số:} $2$\\
Số tròn chục nhỏ nhất là 10. Số đó là: $10 - 8 = 2$.
}
\end{ex}

\newpage

\TaoPhieuKiemTra{Chủ đề 2 -- Phép cộng, phép trừ trong phạm vi 20}{11}{Phép trừ (qua 10) trong phạm vi 20}{https://loptoanthaylinh.vn/thi/chuyen-de/2/toan2_cd_b11}
\textbf{Phần 1. Trắc nghiệm lựa chọn}
\setcounterpageref{ex}{0}

\begin{ex} %[2D2V11]
Tính giá trị của phép tính: $13 - 5$.
\choice
{$7$}
{$9$}
{\True $8$}
{$6$}
\loigiai{
$13 - 5 = 8$.
\textbf{Đáp án của câu hỏi là C}
}
\end{ex}

\begin{ex} %[2D2V11]
Một đàn chim có $15$ con đậu trên dây điện. Lần thứ nhất có $6$ con bay đi, lần thứ hai có $3$ con bay đi. Hỏi còn lại bao nhiêu con chim trên dây điện?
\choice
{$7$ con chim}
{\True $6$ con chim}
{$5$ con chim}
{$8$ con chim}
\loigiai{
Số chim còn lại: $15 - 6 - 3 = 9 - 3 = 6$ (con chim).
\textbf{Đáp án của câu hỏi là B}
}
\end{ex}

\begin{ex} %[2D2V11]
Tìm $x$ biết: $12 - x = 4$.
\choice
{$7$}
{$6$}
{$9$}
{\True $8$}
\loigiai{
$x = 12 - 4 = 8$.
\textbf{Đáp án của câu hỏi là D}
}
\end{ex}

\begin{ex} %[2D2V11]
Năm nay Nam $8$ tuổi, anh trai Nam $14$ tuổi. Hỏi anh trai hơn Nam bao nhiêu tuổi?
\choice
{\True $6$ tuổi}
{$7$ tuổi}
{$5$ tuổi}
{$8$ tuổi}
\loigiai{
Anh hơn Nam: $14 - 8 = 6$ (tuổi).
\textbf{Đáp án của câu hỏi là A}
}
\end{ex}

\begin{ex} %[2D2V11]
Hiệu của $16 - 7$ bằng kết quả của phép tính nào dưới đây?
\choice
{$15 - 7$}
{$13 - 5$}
{\True $14 - 5$}
{$12 - 4$}
\loigiai{
$16 - 7 = 9$. Ta có $14 - 5 = 9$. Do đó hai phép tính có cùng kết quả.
\textbf{Đáp án của câu hỏi là C}
}
\end{ex}

\begin{ex} %[2D2V11]
Điền số thích hợp: $17 - 9 = 15 - \dots$
\choice
{$6$}
{\True $7$}
{$8$}
{$9$}
\loigiai{
Vế trái: $17 - 9 = 8$. Ta có: $15 - 8 = 7$. Vậy điền số 7.
\textbf{Đáp án của câu hỏi là B}
}
\end{ex}

\begin{ex} %[2D2V11]
Một sợi dây dài $15\text{ dm}$, người ta cắt bớt đi $8\text{ dm}$. Đoạn dây còn lại dài bao nhiêu?
\choice
{$6\text{ dm}$}
{$8\text{ dm}$}
{$9\text{ dm}$}
{\True $7\text{ dm}$}
\loigiai{
Đoạn dây còn lại dài: $15 - 8 = 7\text{ (dm)}$.
\textbf{Đáp án của câu hỏi là D}
}
\end{ex}

\begin{ex} %[2D2V11]
Tìm số tự nhiên biết rằng lấy $11$ trừ đi số đó thì được số bé nhất có một chữ số khác 0.
\choice
{\True $10$}
{$9$}
{$8$}
{$1$}
\loigiai{
Số bé nhất có một chữ số khác 0 là 1. Số đó là: $11 - 1 = 10$.
\textbf{Đáp án của câu hỏi là A}
}
\end{ex}

\begin{ex} %[2D2V11]
Một cửa hàng có $14$ chiếc quạt máy. Buổi sáng bán được $5$ chiếc, buổi chiều bán được $4$ chiếc. Cửa hàng còn lại bao nhiêu chiếc quạt?
\choice
{$6$ chiếc}
{$4$ chiếc}
{\True $5$ chiếc}
{$7$ chiếc}
\loigiai{
Tổng số quạt đã bán: $5 + 4 = 9$ (chiếc).\nSố quạt còn lại: $14 - 9 = 5$ (chiếc).
\textbf{Đáp án của câu hỏi là C}
}
\end{ex}

\begin{ex} %[2D2V11]
Trong một phép trừ, số bị trừ là $13$, hiệu là $6$. Số trừ là
\choice
{$6$}
{\True $7$}
{$8$}
{$5$}
\loigiai{
Số trừ = Số bị trừ - Hiệu = $13 - 6 = 7$.
\textbf{Đáp án của câu hỏi là B}
}
\end{ex}

\begin{ex} %[2D2V11]
Điền dấu thích hợp: $18 - 9 \dots 11 - 3$.
\choice
{$<$}
{$=$}
{Không so sánh được}
{\True $>$}
\loigiai{
$18 - 9 = 9$; $11 - 3 = 8$. Vì $9 > 8$ nên điền dấu $>$. 
\textbf{Đáp án của câu hỏi là D}
}
\end{ex}

\begin{ex} %[2D2V11]
Bác thợ có một thanh sắt dài $16\text{ dm}$. Bác cắt làm hai đoạn, đoạn thứ nhất dài $7\text{ dm}$. Hỏi đoạn thứ hai dài hơn đoạn thứ nhất bao nhiêu đề-xi-mét?
\choice
{\True $2\text{ dm}$}
{$9\text{ dm}$}
{$1\text{ dm}$}
{$3\text{ dm}$}
\loigiai{
Đoạn thứ hai dài: $16 - 7 = 9\text{ (dm)}$.\nĐoạn thứ hai hơn đoạn thứ nhất: $9 - 7 = 2\text{ (dm)}$.
\textbf{Đáp án của câu hỏi là A}
}
\end{ex}

\textbf{Phần 2. Trắc nghiệm đúng sai}
\setcounterpageref{ex}{0}

\begin{ex} %[2D2V11]
Xét tính đúng sai của các phép tính trong bảng trừ có nhớ:
\choiceTF
{\True Phép trừ $12 - 7$ có kết quả bằng $5$.}
{Phép trừ $13 - 6$ có kết quả lớn hơn $14 - 6$.}
{\True Khi cùng bớt ở số bị trừ và số trừ đi 2 đơn vị thì hiệu không thay đổi.}
{\True Trong phạm vi 20, phép trừ $11 - 9$ có hiệu nhỏ nhất trong các phép trừ có nhớ với số bị trừ 11.}
\loigiai{
\begin{itemize}
    \item Ý a) Đúng: $12 - 7 = 5$.
    \item Ý b) Sai: $13 - 6 = 7 < 14 - 6 = 8$.
    \item Ý c) Đúng: Tính chất của phép trừ.
    \item Ý d) Đúng: Vì 9 là số trừ lớn nhất có 1 chữ số nên $11 - 9 = 2$ là nhỏ nhất.
\end{itemize}
\begin{itemize}
    \item a) \textbf{Đúng}.
    \item b) \textbf{Sai}.
    \item c) \textbf{Đúng}.
    \item d) \textbf{Đúng}.
\end{itemize}
}
\end{ex}

\begin{ex} %[2D2V11]
Có 14 quyển vở chia cho hai bạn Nam và Việt. Nam nhận 8 quyển vở:
\choiceTF
{\True Việt nhận được 6 quyển vở.}
{\True Nam nhận nhiều hơn Việt 2 quyển vở.}
{Nếu Nam cho Việt 2 quyển thì hai bạn có số vở bằng nhau.}
{Tổng số vở của hai bạn là một số lẻ.}
\loigiai{
\begin{itemize}
    \item Ý a) Đúng: $14 - 8 = 6$ quyển.
    \item Ý b) Đúng: $8 - 6 = 2$ quyển.
    \item Ý c) Sai: Nam chỉ hơn Việt 2 quyển nên Nam chỉ cần cho Việt $2 : 2 = 1$ quyển là bằng nhau.
    \item Ý d) Sai: $14$ là số chẵn.
\end{itemize}
\begin{itemize}
    \item a) \textbf{Đúng}.
    \item b) \textbf{Đúng}.
    \item c) \textbf{Sai}.
    \item d) \textbf{Sai}.
\end{itemize}
}
\end{ex}

\textbf{Phần 3. Trả lời ngắn}
\setcounterpageref{ex}{0}

\begin{ex} %[2D2V11]
Tính giá trị của biểu thức: $15 - 9 + 8$.
\shortans{14}
\loigiai{
\textbf{Đáp số:} $14$\\
$15 - 9 = 6$; $6 + 8 = 14$.
}
\end{ex}

\begin{ex} %[2D2V11]
Tìm số tự nhiên x biết: $14 - x = 9$.
\shortans{5}
\loigiai{
\textbf{Đáp số:} $5$\\
$x = 14 - 9 = 5$.
}
\end{ex}

\begin{ex} %[2D2V11]
Trong hộp có 13 viên bi gồm bi xanh và bi đỏ. Biết số bi đỏ là 6 viên. Hỏi số bi xanh nhiều hơn số bi đỏ bao nhiêu viên?
\shortans{1}
\loigiai{
\textbf{Đáp số:} $1$\\
Số bi xanh: $13 - 6 = 7$ viên. Hơn bi đỏ: $7 - 6 = 1$ viên.
}
\end{ex}

\begin{ex} %[2D2V11]
Một sợi dây dài 17 cm, bạn Mai cắt đi một đoạn thì còn lại đoạn dây dài 9 cm. Hỏi đoạn dây bị cắt đi dài bao nhiêu xăng-ti-mét?
\shortans{8}
\loigiai{
\textbf{Đáp số:} $8$\\
Đoạn dây bị cắt: $17 - 9 = 8\text{ cm}$.
}
\end{ex}

\begin{ex} %[2D2V11]
Tìm số bị trừ biết số trừ là 8 và hiệu là 7.
\shortans{15}
\loigiai{
\textbf{Đáp số:} $15$\\
Số bị trừ = $8 + 7 = 15$.
}
\end{ex}

\begin{ex} %[2D2V11]
Lớp 2A có 16 bạn tham gia múa hát, trong đó có 9 bạn nữ. Hỏi số bạn nam ít hơn số bạn nữ bao nhiêu bạn?
\shortans{2}
\loigiai{
\textbf{Đáp số:} $2$\\
Số bạn nam: $16 - 9 = 7$ bạn. Nam ít hơn nữ: $9 - 7 = 2$ bạn.
}
\end{ex}

\newpage

\TaoPhieuKiemTra{Chủ đề 2 -- Phép cộng, phép trừ trong phạm vi 20}{12}{Bảng trừ (qua 10)}{https://loptoanthaylinh.vn/thi/chuyen-de/2/toan2_cd_b12}
\textbf{Phần 1. Trắc nghiệm lựa chọn}
\setcounterpageref{ex}{0}

\begin{ex} %[2D2V12]
Mai có $9$ con gấu bông, Lan có ít hơn Mai $3$ con gấu bông. Hỏi cả hai bạn có tất cả bao nhiêu con gấu bông?
\choice
{$12$ con}
{\True $15$ con}
{$14$ con}
{$16$ con}
\loigiai{
Lan có: $9 - 3 = 6$ (con).\nCả hai bạn có: $9 + 6 = 15$ (con).
\textbf{Đáp án của câu hỏi là B}
}
\end{ex}

\begin{ex} %[2D2V12]
Tổ một trồng được $8$ cây hoa, tổ hai trồng nhiều hơn tổ một $4$ cây hoa. Hỏi cả hai tổ trồng được bao nhiêu cây hoa?
\choice
{\True $20$ cây hoa}
{$12$ cây hoa}
{$18$ cây hoa}
{$16$ cây hoa}
\loigiai{
Tổ hai trồng được: $8 + 4 = 12$ (cây).\nCả hai tổ trồng được: $8 + 12 = 20$ (cây).
\textbf{Đáp án của câu hỏi là A}
}
\end{ex}

\begin{ex} %[2D2V12]
Thùng thứ nhất đựng $16\text{ L}$ nước, thùng thứ hai đựng ít hơn thùng thứ nhất $7\text{ L}$ nước. Cả hai thùng đựng được bao nhiêu lít nước?
\choice
{$23\text{ L}$}
{$24\text{ L}$}
{\True $25\text{ L}$}
{$26\text{ L}$}
\loigiai{
Thùng thứ hai đựng: $16 - 7 = 9\text{ (L)}$.\nCả hai thùng đựng: $16 + 9 = 25\text{ (L)}$.
\textbf{Đáp án của câu hỏi là C}
}
\end{ex}

\begin{ex} %[2D2V12]
Minh có $14$ viên bi. Nếu Minh cho Dũng $3$ viên bi thì hai bạn có số bi bằng nhau. Hỏi ban đầu Minh có nhiều hơn Dũng bao nhiêu viên bi?
\choice
{$3$ viên bi}
{$5$ viên bi}
{$4$ viên bi}
{\True $6$ viên bi}
\loigiai{
Khi cho đi 3 viên bi mà bằng nhau thì ban đầu Minh hơn Dũng: $3 \times 2 = 6$ (viên bi).
\textbf{Đáp án của câu hỏi là D}
}
\end{ex}

\begin{ex} %[2D2V12]
Năm nay con $7$ tuổi, bố hơn con $28$ tuổi. Hỏi năm nay bố bao nhiêu tuổi?
\choice
{$34$ tuổi}
{\True $35$ tuổi}
{$36$ tuổi}
{$33$ tuổi}
\loigiai{
Tuổi của bố năm nay là: $7 + 28 = 35$ (tuổi).
\textbf{Đáp án của câu hỏi là B}
}
\end{ex}

\begin{ex} %[2D2V12]
Hà có $15$ chiếc bút màu, Hà có nhiều hơn My $6$ chiếc bút màu. Hỏi My có bao nhiêu chiếc bút màu?
\choice
{\True $9$ chiếc}
{$8$ chiếc}
{$10$ chiếc}
{$21$ chiếc}
\loigiai{
My có số bút màu là: $15 - 6 = 9$ (chiếc).
\textbf{Đáp án của câu hỏi là A}
}
\end{ex}

\begin{ex} %[2D2V12]
Một đàn gà có $13$ con gà mái, số gà trống ít hơn gà mái $5$ con. Cả đàn gà có tất cả bao nhiêu con?
\choice
{$18$ con gà}
{$19$ con gà}
{\True $21$ con gà}
{$20$ con gà}
\loigiai{
Số gà trống là: $13 - 5 = 8$ (con).\nCả đàn gà có: $13 + 8 = 21$ (con).
\textbf{Đáp án của câu hỏi là C}
}
\end{ex}

\begin{ex} %[2D2V12]
Đoạn thẳng $AB$ dài $12\text{ cm}$, đoạn thẳng $CD$ ngắn hơn đoạn thẳng $AB$ là $5\text{ cm}$. Tổng độ dài hai đoạn thẳng là
\choice
{$17\text{ cm}$}
{$18\text{ cm}$}
{$20\text{ cm}$}
{\True $19\text{ cm}$}
\loigiai{
Đoạn $CD$ dài: $12 - 5 = 7\text{ (cm)}$.\nTổng độ dài: $12 + 7 = 19\text{ (cm)}$.
\textbf{Đáp án của câu hỏi là D}
}
\end{ex}

\begin{ex} %[2D2V12]
Bao thứ nhất có $17\text{ kg}$ gạo, bao thứ hai có ít hơn bao thứ nhất $8\text{ kg}$ gạo. Hỏi bao thứ hai có bao nhiêu ki-lô-gam gạo?
\choice
{$8\text{ kg}$}
{\True $9\text{ kg}$}
{$10\text{ kg}$}
{$7\text{ kg}$}
\loigiai{
Bao thứ hai có: $17 - 8 = 9\text{ (kg)}$. 
\textbf{Đáp án của câu hỏi là B}
}
\end{ex}

\begin{ex} %[2D2V12]
Lớp 2A trồng được $18$ cây xanh, lớp 2B trồng ít hơn lớp 2A là $4$ cây xanh. Hỏi lớp 2B trồng được bao nhiêu cây xanh?
\choice
{\True $14$ cây}
{$13$ cây}
{$15$ cây}
{$22$ cây}
\loigiai{
Lớp 2B trồng được: $18 - 4 = 14$ (cây).
\textbf{Đáp án của câu hỏi là A}
}
\end{ex}

\begin{ex} %[2D2V12]
Có hai ngăn sách, ngăn thứ nhất có $15$ quyển. Nếu lấy ở ngăn thứ nhất $3$ quyển chuyển sang ngăn thứ hai thì ngăn thứ nhất vẫn nhiều hơn ngăn thứ hai $1$ quyển. Hỏi ban đầu ngăn thứ nhất hơn ngăn thứ hai bao nhiêu quyển?
\choice
{$6$ quyển}
{$8$ quyển}
{\True $7$ quyển}
{$5$ quyển}
\loigiai{
Chuyển 3 quyển làm độ chênh lệch giảm đi: $3 \times 2 = 6$ quyển.\nBan đầu hơn: $6 + 1 = 7$ (quyển).
\textbf{Đáp án của câu hỏi là C}
}
\end{ex}

\begin{ex} %[2D2V12]
Bình có $8$ viên bi, nếu Bình có thêm $5$ viên bi nữa thì số bi của Bình vẫn ít hơn của An $2$ viên bi. Hỏi An có bao nhiêu viên bi?
\choice
{$14$ viên bi}
{$13$ viên bi}
{$16$ viên bi}
{\True $15$ viên bi}
\loigiai{
Nếu thêm 5 viên thì Bình có: $8 + 5 = 13$ (viên).\nAn có số viên bi là: $13 + 2 = 15$ (viên bi).
\textbf{Đáp án của câu hỏi là D}
}
\end{ex}

\textbf{Phần 2. Trắc nghiệm đúng sai}
\setcounterpageref{ex}{0}

\begin{ex} %[2D2V12]
Xét tính đúng sai của các bài toán nhiều hơn ít hơn:
\choiceTF
{\True Muốn tìm số bé ta lấy số lớn trừ đi phần hơn.}
{Nếu A nhiều hơn B 6 cái kẹo, A cho B 3 cái kẹo thì B sẽ nhiều hơn A.}
{\True Tuổi mẹ hơn tuổi con luôn không đổi qua các năm.}
{\True Một lớp có 18 bạn nam, số bạn nữ ít hơn số bạn nam 4 bạn thì cả lớp có 32 bạn.}
\loigiai{
\begin{itemize}
    \item Ý a) Đúng: Số bé = Số lớn - Phần hơn. Đúng.
    \item Ý b) Sai: Khi A cho B 3 cái thì hai bạn có số kẹo bằng nhau.
    \item Ý c) Đúng: Mỗi năm cả hai người đều thêm 1 tuổi nên khoảng cách tuổi giữ nguyên.
    \item Ý d) Đúng: Số nữ: $18 - 4 = 14$. Tổng số: $18 + 14 = 32$ bạn. Đúng.
\end{itemize}
\begin{itemize}
    \item a) \textbf{Đúng}.
    \item b) \textbf{Sai}.
    \item c) \textbf{Đúng}.
    \item d) \textbf{Đúng}.
\end{itemize}
}
\end{ex}

\begin{ex} %[2D2V12]
Bạn Mai có 16 chiếc kẹp tóc, bạn Hoa có ít hơn Mai 7 chiếc kẹp tóc:
\choiceTF
{\True Hoa có 9 chiếc kẹp tóc.}
{\True Cả hai bạn có tất cả 25 chiếc kẹp tóc.}
{Nếu Mai cho Hoa 4 chiếc kẹp tóc thì Hoa sẽ nhiều hơn Mai.}
{Mai có số kẹp tóc gấp đôi số kẹp tóc của Hoa.}
\loigiai{
\begin{itemize}
    \item Ý a) Đúng: $16 - 7 = 9$ chiếc.
    \item Ý b) Đúng: $16 + 9 = 25$ chiếc.
    \item Ý c) Sai: Mai cho 4 chiếc thì Mai còn $16 - 4 = 12$, Hoa có $9 + 4 = 13$. Hoa nhiều hơn Mai 1 chiếc (Đúng, mệnh đề này là Đúng; vậy ta sửa đề để có câu Sai).
    \item Ý d) Sai: 16 không gấp đôi 9. Khẳng định sai.
\end{itemize}
\begin{itemize}
    \item a) \textbf{Đúng}.
    \item b) \textbf{Đúng}.
    \item c) \textbf{Sai}.
    \item d) \textbf{Sai}.
\end{itemize}
}
\end{ex}

\textbf{Phần 3. Trả lời ngắn}
\setcounterpageref{ex}{0}

\begin{ex} %[2D2V12]
Hùng có 14 quyển truyện, Dũng có ít hơn Hùng 5 quyển truyện. Hỏi cả hai bạn có tất cả bao nhiêu quyển truyện?
\shortans{23}
\loigiai{
\textbf{Đáp số:} $23$\\
Dũng có: $14 - 5 = 9$ quyển. Cả hai bạn: $14 + 9 = 23$ quyển.
}
\end{ex}

\begin{ex} %[2D2V12]
Chị hơn em 6 tuổi. Hỏi sau 4 năm nữa chị hơn em bao nhiêu tuổi?
\shortans{6}
\loigiai{
\textbf{Đáp số:} $6$\\
Hiệu số tuổi không thay đổi theo thời gian, vẫn là 6 tuổi.
}
\end{ex}

\begin{ex} %[2D2V12]
Can thứ nhất đựng 18 lít mật ong, can thứ hai đựng ít hơn can thứ nhất 9 lít mật ong. Hỏi can thứ hai đựng bao nhiêu lít mật ong?
\shortans{9}
\loigiai{
\textbf{Đáp số:} $9$\\
$18 - 9 = 9$ lít.
}
\end{ex}

\begin{ex} %[2D2V12]
An có 16 cái nhãn vở, Bình có ít hơn An 6 cái nhãn vở. Hỏi An phải cho Bình bao nhiêu cái nhãn vở để hai bạn có số nhãn vở bằng nhau?
\shortans{3}
\loigiai{
\textbf{Đáp số:} $3$\\
An hơn Bình 6 cái. An phải cho Bình: $6 : 2 = 3$ cái nhãn vở.
}
\end{ex}

\begin{ex} %[2D2V12]
Một giá sách tầng trên có 18 quyển sách, tầng dưới có ít hơn tầng trên 8 quyển sách. Hỏi cả hai tầng có tất cả bao nhiêu quyển sách?
\shortans{28}
\loigiai{
\textbf{Đáp số:} $28$\\
Tầng dưới: $18 - 8 = 10$ quyển. Tổng: $18 + 10 = 28$ quyển.
}
\end{ex}

\begin{ex} %[2D2V12]
Nhà Lan nuôi 15 con gà, nhà Mai nuôi nhiều hơn nhà Lan 6 con gà. Hỏi nhà Mai nuôi bao nhiêu con gà?
\shortans{21}
\loigiai{
\textbf{Đáp số:} $21$\\
$15 + 6 = 21$ con gà.
}
\end{ex}

\newpage

\TaoPhieuKiemTra{Chủ đề 2 -- Phép cộng, phép trừ trong phạm vi 20}{13}{Bài toán về nhiều hơn, ít hơn một số đơn vị}{https://loptoanthaylinh.vn/thi/chuyen-de/2/toan2_cd_b13}
\textbf{Phần 1. Trắc nghiệm lựa chọn}
\setcounterpageref{ex}{0}

\begin{ex} %[2D2V13]
Một bến xe có $17$ chiếc ô tô. Buổi sáng có $8$ chiếc rời bến, buổi chiều có thêm $5$ chiếc vào bến. Hỏi lúc này bến xe có bao nhiêu chiếc ô tô?
\choice
{$13$ chiếc ô tô}
{$15$ chiếc ô tô}
{\True $14$ chiếc ô tô}
{$12$ chiếc ô tô}
\loigiai{
Số ô tô lúc này: $17 - 8 + 5 = 9 + 5 = 14$ (chiếc ô tô).
\textbf{Đáp án của câu hỏi là C}
}
\end{ex}

\begin{ex} %[2D2V13]
Mẹ mua về một túi mận. Sau khi cả nhà ăn hết $9$ quả thì trong túi còn lại $8$ quả. Hỏi ban đầu túi mận có bao nhiêu quả?
\choice
{\True $17$ quả}
{$16$ quả}
{$18$ quả}
{$15$ quả}
\loigiai{
Ban đầu túi mận có: $9 + 8 = 17$ (quả).
\textbf{Đáp án của câu hỏi là A}
}
\end{ex}

\begin{ex} %[2D2V13]
Một đàn vịt có $14$ con. Người ta thả thêm vào đàn một số con vịt thì đàn vịt có tất cả $20$ con. Hỏi người ta đã thả thêm bao nhiêu con vịt?
\choice
{$5$ con vịt}
{$7$ con vịt}
{$8$ con vịt}
{\True $6$ con vịt}
\loigiai{
Số vịt thả thêm là: $20 - 14 = 6$ (con vịt).
\textbf{Đáp án của câu hỏi là D}
}
\end{ex}

\begin{ex} %[2D2V13]
Bình có $16$ viên bi. Bình cho Nam $5$ viên bi rồi lại được bố cho thêm $4$ viên bi. Hỏi lúc này Bình có bao nhiêu viên bi?
\choice
{$14$ viên bi}
{\True $15$ viên bi}
{$16$ viên bi}
{$13$ viên bi}
\loigiai{
Số bi của Bình lúc này là: $16 - 5 + 4 = 11 + 4 = 15$ (viên bi).
\textbf{Đáp án của câu hỏi là B}
}
\end{ex}

\begin{ex} %[2D2V13]
Trong một cuộc thi chạy có $15$ vận động viên. Có $6$ vận động viên đã về đích. Hỏi còn bao nhiêu vận động viên đang trên đường chạy?
\choice
{$8$ người}
{$10$ người}
{\True $9$ người}
{$7$ người}
\loigiai{
Số vận động viên còn lại: $15 - 6 = 9$ (người).
\textbf{Đáp án của câu hỏi là C}
}
\end{ex}

\begin{ex} %[2D2V13]
Một cuộn dây dài $18\text{ m}$. Bác thợ cắt đi $7\text{ m}$ để buộc hàng, sau đó cắt tiếp $4\text{ m}$ để làm dây phơi. Hỏi cuộn dây còn lại bao nhiêu mét?
\choice
{\True $7\text{ m}$}
{$6\text{ m}$}
{$8\text{ m}$}
{$9\text{ m}$}
\loigiai{
Tổng số mét dây đã cắt: $7 + 4 = 11\text{ (m)}$.\nCuộn dây còn lại: $18 - 11 = 7\text{ (m)}$.
\textbf{Đáp án của câu hỏi là A}
}
\end{ex}

\begin{ex} %[2D2V13]
Lớp 2B có $18$ bạn tham gia văn nghệ. Nếu có thêm $2$ bạn nam nữa thì số bạn tham gia văn nghệ là
\choice
{$19$ bạn}
{$16$ bạn}
{$21$ bạn}
{\True $20$ bạn}
\loigiai{
$18 + 2 = 20$ (bạn).
\textbf{Đáp án của câu hỏi là D}
}
\end{ex}

\begin{ex} %[2D2V13]
Một rổ táo có $13$ quả. Sau khi bỏ thêm vào rổ $6$ quả nữa rồi lấy ra $8$ quả thì trong rổ còn lại bao nhiêu quả?
\choice
{$12$ quả}
{\True $11$ quả}
{$10$ quả}
{$13$ quả}
\loigiai{
$13 + 6 - 8 = 19 - 8 = 11$ (quả).
\textbf{Đáp án của câu hỏi là B}
}
\end{ex}

\begin{ex} %[2D2V13]
Tìm một số biết rằng lấy số đó bớt đi $7$ rồi thêm vào $9$ thì được $18$.
\choice
{$15$}
{$17$}
{\True $16$}
{$14$}
\loigiai{
Tính ngược từ cuối: $18 - 9 + 7 = 9 + 7 = 16$.
\textbf{Đáp án của câu hỏi là C}
}
\end{ex}

\begin{ex} %[2D2V13]
Một bể cá có $15$ con cá vàng. Người ta vớt ra một số con cá đem tặng thì trong bể còn lại $8$ con. Hỏi người ta đã vớt ra bao nhiêu con cá?
\choice
{\True $7$ con cá}
{$6$ con cá}
{$8$ con cá}
{$9$ con cá}
\loigiai{
Số cá vớt ra là: $15 - 8 = 7$ (con cá).
\textbf{Đáp án của câu hỏi là A}
}
\end{ex}

\begin{ex} %[2D2V13]
Mai hái được $8$ quả dâu tây, Lan hái được $9$ quả dâu tây. Hai bạn cùng ăn hết $5$ quả. Hỏi hai bạn còn lại bao nhiêu quả dâu tây?
\choice
{$11$ quả}
{$13$ quả}
{$14$ quả}
{\True $12$ quả}
\loigiai{
Tổng số dâu tây: $8 + 9 = 17$ (quả).\nCòn lại: $17 - 5 = 12$ (quả).
\textbf{Đáp án của câu hỏi là D}
}
\end{ex}

\begin{ex} %[2D2V13]
Một cửa hàng buổi sáng có $19$ chiếc mũ bảo hiểm, buổi chiều bán được $7$ chiếc. Hỏi cửa hàng còn lại bao nhiêu chiếc mũ?
\choice
{$11$ chiếc}
{\True $12$ chiếc}
{$13$ chiếc}
{$10$ chiếc}
\loigiai{
$19 - 7 = 12$ (chiếc).
\textbf{Đáp án của câu hỏi là B}
}
\end{ex}

\textbf{Phần 2. Trắc nghiệm đúng sai}
\setcounterpageref{ex}{0}

\begin{ex} %[2D2V13]
Xét tính đúng sai của các bài toán thêm, bớt:
\choiceTF
{\True Bớt đi ở số ban đầu một số lượng thì số mới sẽ nhỏ hơn số ban đầu.}
{Thêm vào số ban đầu 5 đơn vị rồi bớt đi 5 đơn vị thì số mới lớn hơn số ban đầu.}
{\True Một xe có 14 chỗ ngồi, đã có 9 khách ngồi thì còn lại 5 chỗ trống.}
{\True Lấy một số cộng với 0 thì vẫn bằng chính số đó.}
\loigiai{
\begin{itemize}
    \item Ý a) Đúng: Phép trừ luôn làm giảm giá trị của số tự nhiên khác 0.
    \item Ý b) Sai: Số mới bằng đúng số ban đầu: $+ 5 - 5 = 0$.
    \item Ý c) Đúng: $14 - 9 = 5$ chỗ trống. Đúng.
    \item Ý d) Đúng: Tính chất cộng với số 0.
\end{itemize}
\begin{itemize}
    \item a) \textbf{Đúng}.
    \item b) \textbf{Sai}.
    \item c) \textbf{Đúng}.
    \item d) \textbf{Đúng}.
\end{itemize}
}
\end{ex}

\begin{ex} %[2D2V13]
Trên đĩa có 16 cái bánh. Nam ăn 3 cái, Mai ăn 4 cái:
\choiceTF
{\True Cả hai bạn ăn tất cả 7 cái bánh.}
{\True Trên đĩa còn lại 9 cái bánh.}
{Số bánh còn lại ít hơn số bánh đã ăn.}
{Nếu mẹ cho thêm vào đĩa 2 cái bánh thì trên đĩa có tất cả 10 cái bánh.}
\loigiai{
\begin{itemize}
    \item Ý a) Đúng: $3 + 4 = 7$ cái bánh.
    \item Ý b) Đúng: $16 - 7 = 9$ cái bánh.
    \item Ý c) Sai: Còn 9 cái bánh, đã ăn 7 cái, vậy số còn lại nhiều hơn số đã ăn.
    \item Ý d) Sai: Số bánh lúc đó là: $9 + 2 = 11$ cái (không phải 10 cái).
\end{itemize}
\begin{itemize}
    \item a) \textbf{Đúng}.
    \item b) \textbf{Đúng}.
    \item c) \textbf{Sai}.
    \item d) \textbf{Sai}.
\end{itemize}
}
\end{ex}

\textbf{Phần 3. Trả lời ngắn}
\setcounterpageref{ex}{0}

\begin{ex} %[2D2V13]
Một cây táo có 17 quả. Sau một đêm gió bão có 8 quả táo bị rụng. Hỏi trên cây còn lại bao nhiêu quả táo?
\shortans{9}
\loigiai{
\textbf{Đáp số:} $9$\\
$17 - 8 = 9$ quả.
}
\end{ex}

\begin{ex} %[2D2V13]
Một bác nông dân gặt lúa, buổi sáng gặt được 8 bó lúa, buổi chiều gặt thêm 9 bó lúa nữa. Bác chở về nhà 11 bó. Hỏi còn lại bao nhiêu bó lúa ở ngoài đồng?
\shortans{6}
\loigiai{
\textbf{Đáp số:} $6$\\
Tổng số bó lúa: $8 + 9 = 17$. Còn lại: $17 - 11 = 6$ bó.
}
\end{ex}

\begin{ex} %[2D2V13]
Tìm một số biết rằng lấy số đó thêm 6 đơn vị rồi bớt đi 9 đơn vị thì được 10.
\shortans{13}
\loigiai{
\textbf{Đáp số:} $13$\\
Số đó là: $10 + 9 - 6 = 13$.
}
\end{ex}

\begin{ex} %[2D2V13]
Huy có 15 viên bi, Huy cho bạn 7 viên bi rồi lại được tặng thêm 5 viên bi. Hỏi Huy có tất cả bao nhiêu viên bi?
\shortans{13}
\loigiai{
\textbf{Đáp số:} $13$\\
$15 - 7 + 5 = 8 + 5 = 13$ viên bi.
}
\end{ex}

\begin{ex} %[2D2V13]
Một cửa hàng có 18 kg đường. Buổi sáng bán 6 kg, buổi chiều bán 5 kg. Hỏi cửa hàng còn lại bao nhiêu ki-lô-gam đường?
\shortans{7}
\loigiai{
\textbf{Đáp số:} $7$\\
$18 - 6 - 5 = 7\text{ kg}$.
}
\end{ex}

\begin{ex} %[2D2V13]
Trên sân có 9 con chim bồ câu, một đàn khác bay đến có 7 con. Lát sau có 4 con bay đi. Hỏi trên sân lúc này còn bao nhiêu con chim bồ câu?
\shortans{12}
\loigiai{
\textbf{Đáp số:} $12$\\
$9 + 7 - 4 = 12$ con.
}
\end{ex}

\newpage

\TaoPhieuKiemTra{Chủ đề 2 -- Phép cộng, phép trừ trong phạm vi 20}{14}{Luyện tập chung}{https://loptoanthaylinh.vn/thi/chuyen-de/2/toan2_cd_b14}
\textbf{Phần 1. Trắc nghiệm lựa chọn}
\setcounterpageref{ex}{0}

\begin{ex} %[2D2V14]
Tính giá trị của biểu thức: $18 - 9 + 8 - 7$.
\choice
{$9$}
{\True $10$}
{$11$}
{$8$}
\loigiai{
$18 - 9 = 9$; $9 + 8 = 17$; $17 - 7 = 10$.
\textbf{Đáp án của câu hỏi là B}
}
\end{ex}

\begin{ex} %[2D2V14]
Một đàn cừu có $16$ con. Số cừu trắng nhiều hơn số cừu đen $4$ con. Hỏi có bao nhiêu con cừu trắng?
\choice
{$8$ con}
{$9$ con}
{$12$ con}
{\True $10$ con}
\loigiai{
Số cừu trắng là: $(16 + 4) : 2 = 10$ (con).
\textbf{Đáp án của câu hỏi là D}
}
\end{ex}

\begin{ex} %[2D2V14]
Tìm $x$ biết: $15 - x = 14 - 8$.
\choice
{\True $9$}
{$8$}
{$10$}
{$7$}
\loigiai{
Vế phải: $14 - 8 = 6$. Ta có: $15 - x = 6 \Rightarrow x = 15 - 6 = 9$.
\textbf{Đáp án của câu hỏi là A}
}
\end{ex}

\begin{ex} %[2D2V14]
Mẹ đi chợ mua $1\text{ chục}$ quả trứng gà và $7$ quả trứng vịt. Mẹ đã dùng hết $8$ quả trứng. Hỏi mẹ còn lại bao nhiêu quả trứng?
\choice
{$8$ quả trứng}
{$10$ quả trứng}
{\True $9$ quả trứng}
{$7$ quả trứng}
\loigiai{
$1\text{ chục} = 10$. Tổng số trứng: $10 + 7 = 17$ (quả).\nCòn lại: $17 - 8 = 9$ (quả).
\textbf{Đáp án của câu hỏi là C}
}
\end{ex}

\begin{ex} %[2D2V14]
Điền số thích hợp: $8 + 7 - \dots = 6$.
\choice
{$8$}
{\True $9$}
{$10$}
{$7$}
\loigiai{
$8 + 7 = 15$. Ta có: $15 - 6 = 9$. Vậy điền số 9.
\textbf{Đáp án của câu hỏi là B}
}
\end{ex}

\begin{ex} %[2D2V14]
Ba bạn An, Bình, Cư thi ném bóng. An ném trúng $7$ quả, Bình ném trúng $8$ quả, Cư ném trúng ít hơn tổng số bóng của An và Bình là $6$ quả. Cư ném trúng bao nhiêu quả bóng?
\choice
{$8$ quả bóng}
{$10$ quả bóng}
{$7$ quả bóng}
{\True $9$ quả bóng}
\loigiai{
Tổng bóng của An và Bình: $7 + 8 = 15$ (quả).\nCư ném trúng: $15 - 6 = 9$ (quả bóng).
\textbf{Đáp án của câu hỏi là D}
}
\end{ex}

\begin{ex} %[2D2V14]
Tìm số tự nhiên lớn nhất có một chữ số mà khi cộng với $8$ ta được kết quả nhỏ hơn $17$.
\choice
{\True $8$}
{$9$}
{$7$}
{$6$}
\loigiai{
Ta có: $x + 8 < 17 \Rightarrow x < 9$. Số tự nhiên lớn nhất có 1 chữ số nhỏ hơn 9 là $8$.
\textbf{Đáp án của câu hỏi là A}
}
\end{ex}

\begin{ex} %[2D2V14]
Trong một hộp có $18$ que tính gồm ba màu xanh, đỏ, vàng. Biết có $7$ que màu xanh và $6$ que màu đỏ. Hỏi có bao nhiêu que tính màu vàng?
\choice
{$6$ que tính}
{$4$ que tính}
{\True $5$ que tính}
{$7$ que tính}
\loigiai{
Số que tính màu vàng: $18 - (7 + 6) = 18 - 13 = 5$ (que tính).
\textbf{Đáp án của câu hỏi là C}
}
\end{ex}

\begin{ex} %[2D2V14]
Hiệu giữa số chẵn lớn nhất có một chữ số và số lẻ bé nhất có một chữ số khác 1 là
\choice
{$6$}
{\True $5$}
{$4$}
{$7$}
\loigiai{
Số chẵn lớn nhất có 1 chữ số là 8. Số lẻ bé nhất khác 1 là 3.\nHiệu: $8 - 3 = 5$.
\textbf{Đáp án của câu hỏi là B}
}
\end{ex}

\begin{ex} %[2D2V14]
Hình vẽ gồm hai đoạn thẳng $AB = 9\text{ cm}$ và $BC = 8\text{ cm}$ ghép lại thành đường gấp khúc $ABC$. Sau đó người ta cắt bớt đi $6\text{ cm}$. Độ dài đường gấp khúc còn lại là
\choice
{$10\text{ cm}$}
{$12\text{ cm}$}
{$9\text{ cm}$}
{\True $11\text{ cm}$}
\loigiai{
Độ dài đường gấp khúc: $9 + 8 = 17\text{ (cm)}$.\nCòn lại: $17 - 6 = 11\text{ (cm)}$.
\textbf{Đáp án của câu hỏi là D}
}
\end{ex}

\begin{ex} %[2D2V14]
Một rổ mận có $16$ quả. Chị ăn $4$ quả, em ăn $5$ quả. Hỏi hai chị em đã ăn hết bao nhiêu quả mận?
\choice
{\True $9$ quả}
{$8$ quả}
{$7$ quả}
{$10$ quả}
\loigiai{
Hai chị em đã ăn hết: $4 + 5 = 9$ (quả mận).
\textbf{Đáp án của câu hỏi là A}
}
\end{ex}

\begin{ex} %[2D2V14]
Tìm hai số có tổng bằng $16$ và hiệu bằng $4$.
\choice
{$9\text{ và } 7$}
{$11\text{ và } 5$}
{\True $10\text{ và } 6$}
{$8\text{ và } 8$}
\loigiai{
Số lớn: $(16 + 4) : 2 = 10$. Số bé: $10 - 4 = 6$.
\textbf{Đáp án của câu hỏi là C}
}
\end{ex}

\textbf{Phần 2. Trắc nghiệm đúng sai}
\setcounterpageref{ex}{0}

\begin{ex} %[2D2V14]
Xét tính đúng sai của các khẳng định ôn tập Chủ đề 2:
\choiceTF
{\True Biểu thức $9 + 7 - 8$ có kết quả bằng 8.}
{Trong phép trừ có nhớ phạm vi 20, số bị trừ luôn nhỏ hơn 10.}
{\True Tổng của số lẻ và số lẻ luôn là số chẵn.}
{Nếu $x + 8 = 15$ thì $x = 8$.}
\loigiai{
\begin{itemize}
    \item Ý a) Đúng: $9 + 7 = 16$; $16 - 8 = 8$.
    \item Ý b) Sai: Số bị trừ luôn lớn hơn hoặc bằng 11.
    \item Ý c) Đúng: Ví dụ: $7 + 9 = 16$ (số chẵn).
    \item Ý d) Sai: $x = 15 - 8 = 7$ (không phải 8).
\end{itemize}
\begin{itemize}
    \item a) \textbf{Đúng}.
    \item b) \textbf{Sai}.
    \item c) \textbf{Đúng}.
    \item d) \textbf{Sai}.
\end{itemize}
}
\end{ex}

\begin{ex} %[2D2V14]
Một bác nông dân nuôi 8 con thỏ trắng và 9 con thỏ nâu:
\choiceTF
{\True Bác nuôi tất cả 17 con thỏ.}
{\True Số thỏ nâu nhiều hơn số thỏ trắng 1 con.}
{Mỗi con thỏ có 2 cái tai và 2 cái chân.}
{\True Tổng số tai của cả đàn thỏ là 34 cái tai.}
\loigiai{
\begin{itemize}
    \item Ý a) Đúng: $8 + 9 = 17$ con thỏ.
    \item Ý b) Đúng: $9 - 8 = 1$ con.
    \item Ý c) Sai: Mỗi con thỏ có 2 cái tai và 4 cái chân.
    \item Ý d) Đúng: $17 \times 2 = 34$ cái tai. Đúng.
\end{itemize}
\begin{itemize}
    \item a) \textbf{Đúng}.
    \item b) \textbf{Đúng}.
    \item c) \textbf{Sai}.
    \item d) \textbf{Đúng}.
\end{itemize}
}
\end{ex}

\textbf{Phần 3. Trả lời ngắn}
\setcounterpageref{ex}{0}

\begin{ex} %[2D2V14]
Tính giá trị của biểu thức: $16 - 8 + 7$.
\shortans{15}
\loigiai{
\textbf{Đáp số:} $15$\\
$16 - 8 = 8$; $8 + 7 = 15$.
}
\end{ex}

\begin{ex} %[2D2V14]
Tìm một số biết rằng nếu lấy số đó cộng với 9 rồi trừ đi 6 thì được 11.
\shortans{8}
\loigiai{
\textbf{Đáp số:} $8$\\
Số đó là: $11 + 6 - 9 = 8$.
}
\end{ex}

\begin{ex} %[2D2V14]
Có 17 quả bóng bay, bị nổ mất 8 quả, sau đó được thổi thêm 5 quả. Hỏi có tất cả bao nhiêu quả bóng bay?
\shortans{14}
\loigiai{
\textbf{Đáp số:} $14$\\
$17 - 8 + 5 = 14$ quả.
}
\end{ex}

\begin{ex} %[2D2V14]
Hai bạn có tổng cộng 18 cái kẹo. Bạn thứ nhất hơn bạn thứ hai 4 cái kẹo. Hỏi bạn thứ nhất có bao nhiêu cái kẹo?
\shortans{11}
\loigiai{
\textbf{Đáp số:} $11$\\
Số kẹo của bạn thứ nhất: $(18 + 4) : 2 = 11$ cái.
}
\end{ex}

\begin{ex} %[2D2V14]
Tìm số tự nhiên x biết: $13 - x = 11 - 6$.
\shortans{8}
\loigiai{
\textbf{Đáp số:} $8$\\
$11 - 6 = 5 \Rightarrow x = 13 - 5 = 8$.
}
\end{ex}

\begin{ex} %[2D2V14]
Tìm tổng của số liền sau số 8 và số liền trước số 10.
\shortans{18}
\loigiai{
\textbf{Đáp số:} $18$\\
Số liền sau số 8 là 9. Số liền trước số 10 là 9. Tổng: $9 + 9 = 18$.
}
\end{ex}

\newpage


% =========================================================================
% CHỦ ĐỀ 3 - LÀM QUEN VỚI KHỐI LƯỢNG, DUNG TÍCH
% =========================================================================
\section{Làm quen với khối lượng, dung tích}

\TaoPhieuKiemTra{Chủ đề 3 -- Làm quen với khối lượng, dung tích}{15}{Ki-lô-gam}{https://loptoanthaylinh.vn/thi/chuyen-de/2/toan2_cd_b15}
\textbf{Phần 1. Trắc nghiệm lựa chọn}
\setcounterpageref{ex}{0}

\begin{ex} %[2H3V15]
Trên đĩa cân bên trái có một quả cân $5\text{ kg}$, đĩa bên phải có $1$ gói bột giặt và $2$ quả cân $1\text{ kg}$. Biết cân thăng bằng, gói bột giặt nặng bao nhiêu ki-lô-gam?
\choice
{\True $3\text{ kg}$}
{$2\text{ kg}$}
{$4\text{ kg}$}
{$5\text{ kg}$}
\loigiai{
Vì cân thăng bằng nên: Khối lượng gói bột giặt $+ 1 + 1 = 5$. Khối lượng gói bột giặt là: $5 - 2 = 3\text{ (kg)}$.
\textbf{Đáp án của câu hỏi là A}
}
\end{ex}

\begin{ex} %[2H3V15]
Bao gạo nặng $35\text{ kg}$, bao ngô nhẹ hơn bao gạo $8\text{ kg}$. Hỏi cả hai bao nặng bao nhiêu ki-lô-gam?
\choice
{$54\text{ kg}$}
{$60\text{ kg}$}
{\True $62\text{ kg}$}
{$58\text{ kg}$}
\loigiai{
Bao ngô nặng: $35 - 8 = 27\text{ (kg)}$.\nCả hai bao nặng: $35 + 27 = 62\text{ (kg)}$.
\textbf{Đáp án của câu hỏi là C}
}
\end{ex}

\begin{ex} %[2H3V15]
Tính giá trị của biểu thức: $48\text{ kg} - 19\text{ kg} + 25\text{ kg}$.
\choice
{$52\text{ kg}$}
{\True $54\text{ kg}$}
{$56\text{ kg}$}
{$50\text{ kg}$}
\loigiai{
$48 - 19 = 29\text{ kg}$; $29 + 25 = 54\text{ kg}$.
\textbf{Đáp án của câu hỏi là B}
}
\end{ex}

\begin{ex} %[2H3V15]
Con lợn đen nặng $65\text{ kg}$, con lợn trắng nhẹ hơn con lợn đen $18\text{ kg}$. Con lợn trắng nặng bao nhiêu ki-lô-gam?
\choice
{$49\text{ kg}$}
{$53\text{ kg}$}
{$45\text{ kg}$}
{\True $47\text{ kg}$}
\loigiai{
Con lợn trắng nặng: $65 - 18 = 47\text{ (kg)}$.
\textbf{Đáp án của câu hỏi là D}
}
\end{ex}

\begin{ex} %[2H3V15]
Một thùng chứa đầy nước nặng $22\text{ kg}$. Sau khi người ta dùng đi một nửa số nước thì thùng và số nước còn lại nặng $12\text{ kg}$. Hỏi vỏ thùng không có nước nặng bao nhiêu ki-lô-gam?
\choice
{\True $2\text{ kg}$}
{$3\text{ kg}$}
{$1\text{ kg}$}
{$4\text{ kg}$}
\loigiai{
Nửa số nước nặng: $22 - 12 = 10\text{ (kg)}$.\nCả thùng nước chứa số nước nặng: $10 \times 2 = 20\text{ (kg)}$.\nVỏ thùng nặng: $22 - 20 = 2\text{ (kg)}$.
\textbf{Đáp án của câu hỏi là A}
}
\end{ex}

\begin{ex} %[2H3V15]
Điền dấu thích hợp: $36\text{ kg} + 28\text{ kg} \dots 90\text{ kg} - 25\text{ kg}$.
\choice
{$>$}
{$=$}
{\True $<$}
{Không so sánh được}
\loigiai{
Vế trái: $36 + 28 = 64\text{ kg}$. Vế phải: $90 - 25 = 65\text{ kg}$. Vì $64 < 65$ nên điền dấu $<$. 
\textbf{Đáp án của câu hỏi là C}
}
\end{ex}

\begin{ex} %[2H3V15]
Có ba bao thóc: bao thứ nhất nặng $24\text{ kg}$, bao thứ hai nặng $26\text{ kg}$, bao thứ ba nặng $30\text{ kg}$. Cả ba bao thóc nặng tất cả bao nhiêu ki-lô-gam?
\choice
{$78\text{ kg}$}
{\True $80\text{ kg}$}
{$82\text{ kg}$}
{$75\text{ kg}$}
\loigiai{
$(24 + 26) + 30 = 50 + 30 = 80\text{ (kg)}$.
\textbf{Đáp án của câu hỏi là B}
}
\end{ex}

\begin{ex} %[2H3V15]
Mẹ mua một con gà nặng $2\text{ kg}$ và một con vịt nặng $3\text{ kg}$. Bác An mua một con ngỗng nặng bằng cả con gà và con vịt cộng thêm $2\text{ kg}$. Con ngỗng nặng bao nhiêu?
\choice
{$6\text{ kg}$}
{$8\text{ kg}$}
{$5\text{ kg}$}
{\True $7\text{ kg}$}
\loigiai{
Con ngỗng nặng: $2 + 3 + 2 = 7\text{ (kg)}$.
\textbf{Đáp án của câu hỏi là D}
}
\end{ex}

\begin{ex} %[2H3V15]
Một xe chở $45\text{ kg}$ đường và $38\text{ kg}$ bột mì. Khối lượng hàng hóa xe chở là
\choice
{\True $83\text{ kg}$}
{$73\text{ kg}$}
{$81\text{ kg}$}
{$85\text{ kg}$}
\loigiai{
$45 + 38 = 83\text{ (kg)}$.
\textbf{Đáp án của câu hỏi là A}
}
\end{ex}

\begin{ex} %[2H3V15]
Con bò nặng hơn con bê $45\text{ kg}$. Biết con bò nặng $90\text{ kg}$, con bê nặng bao nhiêu ki-lô-gam?
\choice
{$55\text{ kg}$}
{$40\text{ kg}$}
{\True $45\text{ kg}$}
{$50\text{ kg}$}
\loigiai{
Con bê nặng: $90 - 45 = 45\text{ (kg)}$.
\textbf{Đáp án của câu hỏi là C}
}
\end{ex}

\begin{ex} %[2H3V15]
Túi muối nặng $2\text{ kg}$, túi đường nặng gấp đôi túi muối. Hỏi cả hai túi nặng bao nhiêu ki-lô-gam?
\choice
{$4\text{ kg}$}
{\True $6\text{ kg}$}
{$5\text{ kg}$}
{$7\text{ kg}$}
\loigiai{
Túi đường nặng: $2 + 2 = 4\text{ (kg)}$. Cả hai túi nặng: $2 + 4 = 6\text{ (kg)}$.
\textbf{Đáp án của câu hỏi là B}
}
\end{ex}

\begin{ex} %[2H3V15]
Tìm $x$ biết: $x + 27\text{ kg} = 65\text{ kg}$.
\choice
{$37\text{ kg}$}
{$42\text{ kg}$}
{$48\text{ kg}$}
{\True $38\text{ kg}$}
\loigiai{
$x = 65 - 27 = 38\text{ (kg)}$.
\textbf{Đáp án của câu hỏi là D}
}
\end{ex}

\textbf{Phần 2. Trắc nghiệm đúng sai}
\setcounterpageref{ex}{0}

\begin{ex} %[2H3V15]
Xét tính đúng sai của các bài toán về ki-lô-gam:
\choiceTF
{\True Ki-lô-gam là đơn vị dùng để đo khối lượng, viết tắt là kg.}
{Quả cân 5 kg nhẹ hơn tổng của hai quả cân 2 kg và một quả cân 1 kg.}
{\True Một túi gạo 10 kg san đều vào hai túi thì mỗi túi có 5 kg.}
{\True Khối lượng của hai vật bằng nhau thì đặt lên hai đĩa của cân đĩa, cân sẽ thăng bằng.}
\loigiai{
\begin{itemize}
    \item Ý a) Đúng: Quy ước chuẩn theo sách giáo khoa.
    \item Ý b) Sai: Hai quả 2 kg và một quả 1 kg là $2 + 2 + 1 = 5\text{ kg}$, bằng nhau.
    \item Ý c) Đúng: $10 : 2 = 5\text{ kg}$. Đúng.
    \item Ý d) Đúng: Nguyên lý cân đĩa thăng bằng.
\end{itemize}
\begin{itemize}
    \item a) \textbf{Đúng}.
    \item b) \textbf{Sai}.
    \item c) \textbf{Đúng}.
    \item d) \textbf{Đúng}.
\end{itemize}
}
\end{ex}

\begin{ex} %[2H3V15]
Bác Nam có hai bao thóc, bao thứ nhất nặng 42 kg, bao thứ hai nặng 38 kg:
\choiceTF
{\True Bao thứ nhất nặng hơn bao thứ hai 4 kg.}
{\True Cả hai bao nặng tất cả 80 kg.}
{Chuyển 4 kg từ bao thứ nhất sang bao thứ hai thì hai bao nặng bằng nhau.}
{Bao thứ hai nặng gấp đôi bao thứ nhất.}
\loigiai{
\begin{itemize}
    \item Ý a) Đúng: $42 - 38 = 4\text{ kg}$.
    \item Ý b) Đúng: $42 + 38 = 80\text{ kg}$.
    \item Ý c) Sai: Chỉ cần chuyển $4 : 2 = 2\text{ kg}$ thì hai bao đều nặng 40 kg.
    \item Ý d) Sai: 38 không gấp đôi 42. Khẳng định sai.
\end{itemize}
\begin{itemize}
    \item a) \textbf{Đúng}.
    \item b) \textbf{Đúng}.
    \item c) \textbf{Sai}.
    \item d) \textbf{Sai}.
\end{itemize}
}
\end{ex}

\textbf{Phần 3. Trả lời ngắn}
\setcounterpageref{ex}{0}

\begin{ex} %[2H3V15]
Mẹ mua về 24 kg gạo. Sau khi dùng một tuần thì còn lại 17 kg. Hỏi gia đình đã ăn hết bao nhiêu ki-lô-gam gạo?
\shortans{7}
\loigiai{
\textbf{Đáp số:} $7$\\
$24 - 17 = 7\text{ kg}$.
}
\end{ex}

\begin{ex} %[2H3V15]
Có hai quả dưa hấu, quả thứ nhất nặng 4 kg, quả thứ hai nặng hơn quả thứ nhất 2 kg. Hỏi cả hai quả dưa hấu nặng bao nhiêu ki-lô-gam?
\shortans{10}
\loigiai{
\textbf{Đáp số:} $10$\\
Quả 2: $4 + 2 = 6\text{ kg}$. Tổng: $4 + 6 = 10\text{ kg}$.
}
\end{ex}

\begin{ex} %[2H3V15]
Con bê nặng 35 kg, con lợn nặng 28 kg. Hỏi con bê nặng hơn con lợn bao nhiêu ki-lô-gam?
\shortans{7}
\loigiai{
\textbf{Đáp số:} $7$\\
$35 - 28 = 7\text{ kg}$.
}
\end{ex}

\begin{ex} %[2H3V15]
Một cửa hàng có 50 kg đường, đã bán đi 26 kg. Số đường còn lại chia đều vào 3 túi. Hỏi mỗi túi có bao nhiêu ki-lô-gam đường?
\shortans{8}
\loigiai{
\textbf{Đáp số:} $8$\\
Còn lại: $50 - 26 = 24\text{ kg}$. Mỗi túi: $24 : 3 = 8\text{ kg}$.
}
\end{ex}

\begin{ex} %[2H3V15]
Để cân một con gà, người ta đặt lên một đĩa con gà và quả cân 1 kg, đĩa kia đặt quả cân 4 kg thì cân thăng bằng. Hỏi con gà nặng bao nhiêu ki-lô-gam?
\shortans{3}
\loigiai{
\textbf{Đáp số:} $3$\\
$4 - 1 = 3\text{ kg}$.
}
\end{ex}

\begin{ex} %[2H3V15]
Tổng khối lượng của một bao gạo và một bao ngô là 65 kg, biết bao gạo nặng 37 kg. Hỏi bao ngô nặng bao nhiêu ki-lô-gam?
\shortans{28}
\loigiai{
\textbf{Đáp số:} $28$\\
$65 - 37 = 28\text{ kg}$.
}
\end{ex}

\newpage

\TaoPhieuKiemTra{Chủ đề 3 -- Làm quen với khối lượng, dung tích}{16}{Lít}{https://loptoanthaylinh.vn/thi/chuyen-de/2/toan2_cd_b16}
\textbf{Phần 1. Trắc nghiệm lựa chọn}
\setcounterpageref{ex}{0}

\begin{ex} %[2H3V16]
Trong can có $15\text{ L}$ dầu. Người ta rót vào $3$ chai, mỗi chai $2\text{ L}$. Hỏi trong can còn lại bao nhiêu lít dầu?
\choice
{$10\text{ L}$}
{$8\text{ L}$}
{$11\text{ L}$}
{\True $9\text{ L}$}
\loigiai{
Số dầu đã rót: $2 \times 3 = 6\text{ (L)}$. Trong can còn lại: $15 - 6 = 9\text{ (L)}$.
\textbf{Đáp án của câu hỏi là D}
}
\end{ex}

\begin{ex} %[2H3V16]
Thùng thứ nhất chứa $38\text{ L}$ nước, thùng thứ hai chứa ít hơn thùng thứ nhất $9\text{ L}$ nước. Cả hai thùng chứa bao nhiêu lít nước?
\choice
{$65\text{ L}$}
{\True $67\text{ L}$}
{$69\text{ L}$}
{$63\text{ L}$}
\loigiai{
Thùng thứ hai: $38 - 9 = 29\text{ (L)}$.\nCả hai thùng: $38 + 29 = 67\text{ (L)}$.
\textbf{Đáp án của câu hỏi là B}
}
\end{ex}

\begin{ex} %[2H3V16]
Tính giá trị: $45\text{ L} + 27\text{ L} - 18\text{ L}$.
\choice
{$52\text{ L}$}
{$56\text{ L}$}
{\True $54\text{ L}$}
{$50\text{ L}$}
\loigiai{
$45 + 27 = 72\text{ L}$; $72 - 18 = 54\text{ L}$.
\textbf{Đáp án của câu hỏi là C}
}
\end{ex}

\begin{ex} %[2H3V16]
Một bể chứa $80\text{ L}$ nước. Buổi sáng người ta dùng $25\text{ L}$, buổi chiều dùng $36\text{ L}$. Hỏi trong bể còn lại bao nhiêu lít nước?
\choice
{\True $19\text{ L}$}
{$21\text{ L}$}
{$18\text{ L}$}
{$23\text{ L}$}
\loigiai{
Tổng số lít nước đã dùng: $25 + 36 = 61\text{ (L)}$.\nTrong bể còn lại: $80 - 61 = 19\text{ (L)}$.
\textbf{Đáp án của câu hỏi là A}
}
\end{ex}

\begin{ex} %[2H3V16]
Có hai can nước mắm: can to chứa $20\text{ L}$, can nhỏ chứa ít hơn can to $8\text{ L}$. Cả hai can chứa bao nhiêu lít nước mắm?
\choice
{$30\text{ L}$}
{$28\text{ L}$}
{$34\text{ L}$}
{\True $32\text{ L}$}
\loigiai{
Can nhỏ chứa: $20 - 8 = 12\text{ (L)}$.\nCả hai can chứa: $20 + 12 = 32\text{ (L)}$.
\textbf{Đáp án của câu hỏi là D}
}
\end{ex}

\begin{ex} %[2H3V16]
Rót hết nước từ một bình $10\text{ L}$ vào các ca loại $2\text{ L}$. Hỏi rót được bao nhiêu ca nước?
\choice
{$4$ ca}
{\True $5$ ca}
{$6$ ca}
{$8$ ca}
\loigiai{
$10 : 2 = 5$ (ca).
\textbf{Đáp án của câu hỏi là B}
}
\end{ex}

\begin{ex} %[2H3V16]
Điền số thích hợp: $18\text{ L} + \dots\text{ L} = 35\text{ L}$.
\choice
{$16$}
{$18$}
{\True $17$}
{$15$}
\loigiai{
$35 - 18 = 17$.
\textbf{Đáp án của câu hỏi là C}
}
\end{ex}

\begin{ex} %[2H3V16]
Mẹ mua về $12\text{ L}$ sữa tươi. Mẹ chia đều cho $4$ chị em uống trong một tuần. Hỏi mỗi người được uống bao nhiêu lít sữa?
\choice
{\True $3\text{ L}$}
{$2\text{ L}$}
{$4\text{ L}$}
{$5\text{ L}$}
\loigiai{
$12 : 4 = 3\text{ (L)}$.
\textbf{Đáp án của câu hỏi là A}
}
\end{ex}

\begin{ex} %[2H3V16]
Một cửa hàng xăng dầu có $95\text{ L}$ dầu. Ngày thứ nhất bán được $37\text{ L}$, ngày thứ hai bán được $29\text{ L}$. Sau hai ngày cửa hàng còn lại bao nhiêu lít dầu?
\choice
{$31\text{ L}$}
{$27\text{ L}$}
{$33\text{ L}$}
{\True $29\text{ L}$}
\loigiai{
Số dầu đã bán: $37 + 29 = 66\text{ (L)}$.\nCòn lại: $95 - 66 = 29\text{ (L)}$.
\textbf{Đáp án của câu hỏi là D}
}
\end{ex}

\begin{ex} %[2H3V16]
Can thứ nhất có $14\text{ L}$ mật ong. Nếu chuyển $3\text{ L}$ từ can thứ nhất sang can thứ hai thì hai can có số lít mật ong bằng nhau. Hỏi can thứ hai lúc đầu có bao nhiêu lít mật ong?
\choice
{$9\text{ L}$}
{\True $8\text{ L}$}
{$11\text{ L}$}
{$7\text{ L}$}
\loigiai{
Sau khi chuyển, mỗi can có: $14 - 3 = 11\text{ (L)}$.\nLúc đầu can thứ hai có: $11 - 3 = 8\text{ (L)}$.
\textbf{Đáp án của câu hỏi là B}
}
\end{ex}

\begin{ex} %[2H3V16]
Người ta đong được $3$ can $5\text{ L}$ và $1$ can $2\text{ L}$. Tổng số lít nước đong được là
\choice
{$15\text{ L}$}
{$18\text{ L}$}
{\True $17\text{ L}$}
{$16\text{ L}$}
\loigiai{
$5 \times 3 + 2 = 15 + 2 = 17\text{ (L)}$.
\textbf{Đáp án của câu hỏi là C}
}
\end{ex}

\begin{ex} %[2H3V16]
Hiệu giữa số lít nước của can $25\text{ L}$ và can $18\text{ L}$ là
\choice
{\True $7\text{ L}$}
{$6\text{ L}$}
{$8\text{ L}$}
{$9\text{ L}$}
\loigiai{
$25 - 18 = 7\text{ (L)}$.
\textbf{Đáp án của câu hỏi là A}
}
\end{ex}

\textbf{Phần 2. Trắc nghiệm đúng sai}
\setcounterpageref{ex}{0}

\begin{ex} %[2H3V16]
Xét tính đúng sai của các nhận định về đơn vị đo dung tích (Lít):
\choiceTF
{\True Lít là đơn vị dùng để đo sức chứa của một vật (dung tích), viết tắt là L hoặc l.}
{Can 10 lít chứa được lượng nước ít hơn 2 can 5 lít cộng lại.}
{\True Muốn đong được đúng 3 lít nước bằng can 5 lít và can 2 lít, ta múc đầy can 5 lít rồi rót đầy vào can 2 lít, số nước còn lại trong can 5 lít đúng bằng 3 lít.}
{\True Một bể nước có thể tích 50 L, đã có 35 L nước, cần đổ thêm 15 L để đầy bể.}
\loigiai{
\begin{itemize}
    \item Ý a) Đúng: Chuẩn kiến thức sách giáo khoa.
    \item Ý b) Sai: $5 + 5 = 10\text{ L}$, hai bên bằng nhau.
    \item Ý c) Đúng: $5 - 2 = 3\text{ L}$. Đúng.
    \item Ý d) Đúng: $50 - 35 = 15\text{ L}$. Đúng.
\end{itemize}
\begin{itemize}
    \item a) \textbf{Đúng}.
    \item b) \textbf{Sai}.
    \item c) \textbf{Đúng}.
    \item d) \textbf{Đúng}.
\end{itemize}
}
\end{ex}

\begin{ex} %[2H3V16]
Có ba thùng đựng dầu: Thùng A có 25 L, Thùng B có 35 L, Thùng C có 40 L:
\choiceTF
{\True Cả ba thùng có tất cả 100 L dầu.}
{\True Thùng C nhiều hơn thùng A là 15 L dầu.}
{Thùng B nhiều hơn thùng C là 5 L dầu.}
{Chuyển 5 L từ thùng C sang thùng A thì ba thùng có số dầu bằng nhau.}
\loigiai{
\begin{itemize}
    \item Ý a) Đúng: $25 + 35 + 40 = 100\text{ L}$.
    \item Ý b) Đúng: $40 - 25 = 15\text{ L}$.
    \item Ý c) Sai: Thùng B (35 L) ít hơn Thùng C (40 L) là 5 L.
    \item Ý d) Sai: Khi đó A có 30 L, B có 35 L, C có 35 L, chưa bằng nhau.
\end{itemize}
\begin{itemize}
    \item a) \textbf{Đúng}.
    \item b) \textbf{Đúng}.
    \item c) \textbf{Sai}.
    \item d) \textbf{Sai}.
\end{itemize}
}
\end{ex}

\textbf{Phần 3. Trả lời ngắn}
\setcounterpageref{ex}{0}

\begin{ex} %[2H3V16]
Một can dầu chứa 18 lít dầu. Bác Nam rót ra 2 can, mỗi can 5 lít. Hỏi trong can ban đầu còn lại bao nhiêu lít dầu?
\shortans{8}
\loigiai{
\textbf{Đáp số:} $8$\\
Số dầu rót ra: $5 + 5 = 10\text{ L}$. Còn lại: $18 - 10 = 8\text{ L}$.
}
\end{ex}

\begin{ex} %[2H3V16]
Thùng thứ nhất có 36 lít nước, thùng thứ hai có ít hơn thùng thứ nhất 8 lít nước. Hỏi thùng thứ hai có bao nhiêu lít nước?
\shortans{28}
\loigiai{
\textbf{Đáp số:} $28$\\
$36 - 8 = 28\text{ L}$.
}
\end{ex}

\begin{ex} %[2H3V16]
Tính giá trị biểu thức: $42\text{ L} - 15\text{ L} + 18\text{ L}$.
\shortans{45}
\loigiai{
\textbf{Đáp số:} $45$\\
$42 - 15 = 27$; $27 + 18 = 45$.
}
\end{ex}

\begin{ex} %[2H3V16]
Người ta dùng các can loại 5 lít để đong hết 20 lít dầu. Hỏi cần bao nhiêu chiếc can loại đó?
\shortans{4}
\loigiai{
\textbf{Đáp số:} $4$\\
$20 : 5 = 4$ chiếc can.
}
\end{ex}

\begin{ex} %[2H3V16]
Một bể nước chứa 65 lít nước, người ta xả bớt đi 28 lít nước. Hỏi trong bể còn lại bao nhiêu lít nước?
\shortans{37}
\loigiai{
\textbf{Đáp số:} $37$\\
$65 - 28 = 37\text{ L}$.
}
\end{ex}

\begin{ex} %[2H3V16]
Can A có 15 lít nước mắm, Can B có nhiều hơn Can A 7 lít nước mắm. Hỏi cả hai can có bao nhiêu lít nước mắm?
\shortans{37}
\loigiai{
\textbf{Đáp số:} $37$\\
Can B: $15 + 7 = 22\text{ L}$. Cả hai can: $15 + 22 = 37\text{ L}$.
}
\end{ex}

\newpage

\TaoPhieuKiemTra{Chủ đề 3 -- Làm quen với khối lượng, dung tích}{17}{Thực hành và trải nghiệm với ki-lô-gam, lít}{https://loptoanthaylinh.vn/thi/chuyen-de/2/toan2_cd_b17}
\textbf{Phần 1. Trắc nghiệm lựa chọn}
\setcounterpageref{ex}{0}

\begin{ex} %[2H3V17]
Để có đúng $7\text{ kg}$ gạo từ bao gạo to, người bán hàng dùng cân hai đĩa với quả cân loại nào?
\choice
{$3$ quả $2\text{ kg}$}
{$1$ quả $5\text{ kg}$ và $1$ quả $1\text{ kg}$}
{\True $1$ quả $5\text{ kg}$ và $1$ quả $2\text{ kg}$}
{$2$ quả $3\text{ kg}$}
\loigiai{
$5 + 2 = 7\text{ (kg)}$.
\textbf{Đáp án của câu hỏi là C}
}
\end{ex}

\begin{ex} %[2H3V17]
Một bình đầy nước có thể rót đầy $4$ ca nước loại $1\text{ L}$ và $1$ ca nước loại $2\text{ L}$. Bình nước đó chứa được bao nhiêu lít nước?
\choice
{$5\text{ L}$}
{\True $6\text{ L}$}
{$7\text{ L}$}
{$8\text{ L}$}
\loigiai{
$4 \times 1 + 2 = 6\text{ (L)}$.
\textbf{Đáp án của câu hỏi là B}
}
\end{ex}

\begin{ex} %[2H3V17]
Gia đình Mai mỗi ngày ăn hết $1\text{ kg}$ gạo. Hỏi trong $2$ tuần gia đình Mai ăn hết bao nhiêu ki-lô-gam gạo?
\choice
{$12\text{ kg}$}
{$10\text{ kg}$}
{$16\text{ kg}$}
{\True $14\text{ kg}$}
\loigiai{
$2$ tuần $= 14$ ngày. Gia đình Mai ăn hết: $14 \times 1 = 14\text{ (kg)}$.
\textbf{Đáp án của câu hỏi là D}
}
\end{ex}

\begin{ex} %[2H3V17]
Có $1$ can $3\text{ L}$ và $1$ can $5\text{ L}$. Làm thế nào để lấy được đúng $2\text{ L}$ nước từ thùng nước lớn?
\choice
{\True Múc đầy can $5\text{ L}$, rót sang cho đầy can $3\text{ L}$, lượng nước còn lại trong can $5\text{ L}$ là $2\text{ L}$}
{Múc $2$ lần can $3\text{ L}$}
{Múc đầy can $3\text{ L}$ đổ vào can $5\text{ L}$}
{Múc nửa can $5\text{ L}$}
\loigiai{
Lượng nước còn lại trong can 5L là: $5 - 3 = 2\text{ (L)}$.
\textbf{Đáp án của câu hỏi là A}
}
\end{ex}

\begin{ex} %[2H3V17]
Bác Ba có $24\text{ kg}$ xoài, bác đã bán đi $16\text{ kg}$. Số xoài còn lại bác chia đều vào $2$ rổ. Mỗi rổ có bao nhiêu ki-lô-gam xoài?
\choice
{$3\text{ kg}$}
{$5\text{ kg}$}
{\True $4\text{ kg}$}
{$6\text{ kg}$}
\loigiai{
Số xoài còn lại: $24 - 16 = 8\text{ (kg)}$. Mỗi rổ có: $8 : 2 = 4\text{ (kg)}$.
\textbf{Đáp án của câu hỏi là C}
}
\end{ex}

\begin{ex} %[2H3V17]
Một thùng sơn nặng $18\text{ kg}$. Sơn xong một bức tường, người ta thấy thùng sơn còn lại $9\text{ kg}$. Khối lượng sơn đã dùng là
\choice
{$8\text{ kg}$}
{\True $9\text{ kg}$}
{$10\text{ kg}$}
{$7\text{ kg}$}
\loigiai{
$18 - 9 = 9\text{ (kg)}$.
\textbf{Đáp án của câu hỏi là B}
}
\end{ex}

\begin{ex} %[2H3V17]
Mẹ đong nước mắm vào các chai nhỏ loại $2\text{ L}$. Có tất cả $16\text{ L}$ nước mắm. Mẹ đong được bao nhiêu chai?
\choice
{$7$ chai}
{$6$ chai}
{$9$ chai}
{\True $8$ chai}
\loigiai{
$16 : 2 = 8$ (chai).
\textbf{Đáp án của câu hỏi là D}
}
\end{ex}

\begin{ex} %[2H3V17]
Con ngỗng nặng $5\text{ kg}$, con vịt nhẹ hơn con ngỗng $2\text{ kg}$, con gà nặng bằng con vịt. Hỏi cả ba con nặng bao nhiêu ki-lô-gam?
\choice
{\True $11\text{ kg}$}
{$10\text{ kg}$}
{$12\text{ kg}$}
{$9\text{ kg}$}
\loigiai{
Con vịt nặng: $5 - 2 = 3\text{ (kg)}$. Con gà nặng $3\text{ kg}$. Cả ba con nặng: $5 + 3 + 3 = 11\text{ (kg)}$.
\textbf{Đáp án của câu hỏi là A}
}
\end{ex}

\begin{ex} %[2H3V17]
Một bình trà sữa có $15\text{ L}$. Bán đi một số lít thì còn lại $8\text{ L}$. Đã bán bao nhiêu lít trà sữa?
\choice
{$6\text{ L}$}
{$8\text{ L}$}
{\True $7\text{ L}$}
{$9\text{ L}$}
\loigiai{
$15 - 8 = 7\text{ (L)}$.
\textbf{Đáp án của câu hỏi là C}
}
\end{ex}

\begin{ex} %[2H3V17]
Túi thứ nhất đựng $12\text{ kg}$ khoai tây, túi thứ hai đựng nhiều hơn túi thứ nhất $5\text{ kg}$. Cả hai túi đựng bao nhiêu ki-lô-gam khoai tây?
\choice
{$28\text{ kg}$}
{\True $29\text{ kg}$}
{$30\text{ kg}$}
{$27\text{ kg}$}
\loigiai{
Túi 2: $12 + 5 = 17\text{ (kg)}$. Cả hai túi: $12 + 17 = 29\text{ (kg)}$.
\textbf{Đáp án của câu hỏi là B}
}
\end{ex}

\begin{ex} %[2H3V17]
Điền số thích hợp: $35\text{ L} - 17\text{ L} = 9\text{ L} + \dots\text{ L}$.
\choice
{$8$}
{$10$}
{$7$}
{\True $9$}
\loigiai{
Vế trái: $35 - 17 = 18\text{ L}$. Vế phải: $18 - 9 = 9$. Vậy điền 9.
\textbf{Đáp án của câu hỏi là D}
}
\end{ex}

\begin{ex} %[2H3V17]
Muốn có $10\text{ L}$ nước, người ta có thể dùng những chiếc can loại nào dưới đây?
\choice
{\True $2$ chiếc can $5\text{ L}$}
{$3$ chiếc can $3\text{ L}$}
{$1$ can $5\text{ L}$ và $2$ can $2\text{ L}$}
{$4$ can $2\text{ L}$}
\loigiai{
$2 \times 5 = 10\text{ (L)}$.
\textbf{Đáp án của câu hỏi là A}
}
\end{ex}

\textbf{Phần 2. Trắc nghiệm đúng sai}
\setcounterpageref{ex}{0}

\begin{ex} %[2H3V17]
Xét tính đúng sai của các tình huống thực hành cân đong:
\choiceTF
{\True Đĩa cân bên trái có gói đường và quả cân 2 kg, đĩa bên phải có quả cân 5 kg. Cân thăng bằng thì gói đường nặng 3 kg.}
{Chỉ dùng can 2 L và can 5 L thì không thể đong được 7 L nước.}
{\True Tổng khối lượng của một vật và đồ đựng luôn lớn hơn khối lượng của riêng vật đó.}
{Một lít nước nặng hơn 100 kg bông gòn.}
\loigiai{
\begin{itemize}
    \item Ý a) Đúng: $5 - 2 = 3\text{ kg}$. Đúng.
    \item Ý b) Sai: $2 + 5 = 7\text{ L}$. Hoàn toàn đong được.
    \item Ý c) Đúng: Khối lượng cả bì = Khối lượng vật + Khối lượng đồ đựng.
    \item Ý d) Sai: 1 lít nước nặng khoảng 1 kg, nhẹ hơn rất nhiều so với 100 kg bông gòn.
\end{itemize}
\begin{itemize}
    \item a) \textbf{Đúng}.
    \item b) \textbf{Sai}.
    \item c) \textbf{Đúng}.
    \item d) \textbf{Sai}.
\end{itemize}
}
\end{ex}

\begin{ex} %[2H3V17]
Bác thợ mộc có hai thanh gỗ: thanh thứ nhất nặng 18 kg, thanh thứ hai nặng 14 kg:
\choiceTF
{\True Thanh thứ nhất nặng hơn thanh thứ hai 4 kg.}
{\True Cả hai thanh gỗ nặng 32 kg.}
{Nếu cưa bớt ở thanh thứ nhất 4 kg thì hai thanh có khối lượng bằng nhau và bằng 18 kg.}
{\True Nếu ghép hai thanh lại rồi cưa đôi thì mỗi nửa nặng 16 kg.}
\loigiai{
\begin{itemize}
    \item Ý a) Đúng: $18 - 14 = 4\text{ kg}$.
    \item Ý b) Đúng: $18 + 14 = 32\text{ kg}$.
    \item Ý c) Sai: Khi đó cả hai thanh đều nặng 14 kg (không phải 18 kg).
    \item Ý d) Đúng: $32 : 2 = 16\text{ kg}$. Đúng.
\end{itemize}
\begin{itemize}
    \item a) \textbf{Đúng}.
    \item b) \textbf{Đúng}.
    \item c) \textbf{Sai}.
    \item d) \textbf{Đúng}.
\end{itemize}
}
\end{ex}

\textbf{Phần 3. Trả lời ngắn}
\setcounterpageref{ex}{0}

\begin{ex} %[2H3V17]
Một người thợ rót mật ong từ can 10 lít vào các chai nhỏ loại 1 lít. Rót được 6 chai thì can còn lại bao nhiêu lít mật ong?
\shortans{4}
\loigiai{
\textbf{Đáp số:} $4$\\
$10 - 6 = 4\text{ L}$.
}
\end{ex}

\begin{ex} %[2H3V17]
Một con cá quả nặng 2 kg, con cá chép nặng hơn con cá quả 1 kg. Hỏi cả hai con cá nặng bao nhiêu ki-lô-gam?
\shortans{5}
\loigiai{
\textbf{Đáp số:} $5$\\
Cá chép: $2 + 1 = 3\text{ kg}$. Tổng: $2 + 3 = 5\text{ kg}$.
}
\end{ex}

\begin{ex} %[2H3V17]
Có 3 bao ngô, mỗi bao nặng 9 kg. Hỏi cả 3 bao ngô nặng tất cả bao nhiêu ki-lô-gam?
\shortans{27}
\loigiai{
\textbf{Đáp số:} $27$\\
$9 \times 3 = 27\text{ kg}$.
}
\end{ex}

\begin{ex} %[2H3V17]
Một thùng có 25 lít nước mắm, người ta đã bán đi 16 lít nước mắm. Hỏi trong thùng còn lại bao nhiêu lít nước mắm?
\shortans{9}
\loigiai{
\textbf{Đáp số:} $9$\\
$25 - 16 = 9\text{ L}$.
}
\end{ex}

\begin{ex} %[2H3V17]
Một túi kẹo nặng 1 kg, người ta đóng 10 túi như thế vào một hộp. Hỏi 10 túi kẹo nặng bao nhiêu ki-lô-gam?
\shortans{10}
\loigiai{
\textbf{Đáp số:} $10$\\
$1 \times 10 = 10\text{ kg}$.
}
\end{ex}

\begin{ex} %[2H3V17]
Một can 20 lít đựng đầy xăng. Sau khi rót vào xe máy 6 lít và xe ô tô 9 lít thì trong can còn lại bao nhiêu lít xăng?
\shortans{5}
\loigiai{
\textbf{Đáp số:} $5$\\
$20 - 6 - 9 = 5\text{ L}$.
}
\end{ex}

\newpage

\TaoPhieuKiemTra{Chủ đề 3 -- Làm quen với khối lượng, dung tích}{18}{Luyện tập chung}{https://loptoanthaylinh.vn/thi/chuyen-de/2/toan2_cd_b18}
\textbf{Phần 1. Trắc nghiệm lựa chọn}
\setcounterpageref{ex}{0}

\begin{ex} %[2H3V18]
Một can dầu chứa $18\text{ L}$. Nếu rót thêm vào can $5\text{ L}$ nữa thì can đầy. Hỏi dung tích của chiếc can đó là bao nhiêu?
\choice
{$21\text{ L}$}
{\True $23\text{ L}$}
{$22\text{ L}$}
{$24\text{ L}$}
\loigiai{
Dung tích can: $18 + 5 = 23\text{ (L)}$.
\textbf{Đáp án của câu hỏi là B}
}
\end{ex}

\begin{ex} %[2H3V18]
Túi thứ nhất nặng $15\text{ kg}$, túi thứ hai nặng $18\text{ kg}$. Người ta lấy bớt ở túi thứ hai $4\text{ kg}$. Hỏi lúc này cả hai túi nặng bao nhiêu ki-lô-gam?
\choice
{\True $29\text{ kg}$}
{$31\text{ kg}$}
{$28\text{ kg}$}
{$30\text{ kg}$}
\loigiai{
Túi thứ hai còn: $18 - 4 = 14\text{ (kg)}$. Cả hai túi nặng: $15 + 14 = 29\text{ (kg)}$.
\textbf{Đáp án của câu hỏi là A}
}
\end{ex}

\begin{ex} %[2H3V18]
Điền số thích hợp: $42\text{ kg} - 17\text{ kg} = 15\text{ kg} + \dots\text{ kg}$.
\choice
{$9$}
{$11$}
{\True $10$}
{$8$}
\loigiai{
Vế trái: $42 - 17 = 25\text{ kg}$. Vế phải: $25 - 15 = 10$. Vậy điền 10.
\textbf{Đáp án của câu hỏi là C}
}
\end{ex}

\begin{ex} %[2H3V18]
Có hai thùng nước: thùng $A$ chứa $35\text{ L}$, thùng $B$ chứa ít hơn thùng $A$ là $9\text{ L}$. Thùng $B$ chứa bao nhiêu lít nước?
\choice
{$24\text{ L}$}
{$25\text{ L}$}
{$27\text{ L}$}
{\True $26\text{ L}$}
\loigiai{
$35 - 9 = 26\text{ (L)}$.
\textbf{Đáp án của câu hỏi là D}
}
\end{ex}

\begin{ex} %[2H3V18]
Một con lợn nặng $45\text{ kg}$, một con dê nhẹ hơn con lợn $17\text{ kg}$. Cả hai con nặng bao nhiêu ki-lô-gam?
\choice
{$71\text{ kg}$}
{\True $73\text{ kg}$}
{$75\text{ kg}$}
{$68\text{ kg}$}
\loigiai{
Con dê nặng: $45 - 17 = 28\text{ (kg)}$. Cả hai con nặng: $45 + 28 = 73\text{ (kg)}$.
\textbf{Đáp án của câu hỏi là B}
}
\end{ex}

\begin{ex} %[2H3V18]
Bác Lan có $30\text{ L}$ mật ong, bác rót vào các can loại $5\text{ L}$. Hỏi bác Lan rót được bao nhiêu can?
\choice
{\True $6$ can}
{$5$ can}
{$7$ can}
{$4$ can}
\loigiai{
$30 : 5 = 6$ (can).
\textbf{Đáp án của câu hỏi là A}
}
\end{ex}

\begin{ex} %[2H3V18]
Một bình thủy tinh nặng $500\text{ g}$, khi đựng đầy nước thì bình nặng $1\text{ kg } 500\text{ g}$. Lượng nước trong bình nặng bao nhiêu?
\choice
{$2\text{ kg}$}
{$500\text{ g}$}
{\True $1\text{ kg}$}
{$800\text{ g}$}
\loigiai{
Lượng nước nặng: $1\text{ kg } 500\text{ g} - 500\text{ g} = 1\text{ kg}$.
\textbf{Đáp án của câu hỏi là C}
}
\end{ex}

\begin{ex} %[2H3V18]
Tính giá trị: $15\text{ L} + 18\text{ L} - 9\text{ L}$.
\choice
{$22\text{ L}$}
{$26\text{ L}$}
{$25\text{ L}$}
{\True $24\text{ L}$}
\loigiai{
$15 + 18 = 33\text{ L}$; $33 - 9 = 24\text{ L}$.
\textbf{Đáp án của câu hỏi là D}
}
\end{ex}

\begin{ex} %[2H3V18]
Trong kho có $85\text{ kg}$ đỗ xanh. Ngày thứ nhất xuất $28\text{ kg}$, ngày thứ hai xuất $34\text{ kg}$. Trong kho còn lại bao nhiêu ki-lô-gam đỗ xanh?
\choice
{$25\text{ kg}$}
{\True $23\text{ kg}$}
{$21\text{ kg}$}
{$27\text{ kg}$}
\loigiai{
Tổng xuất: $28 + 34 = 62\text{ (kg)}$. Còn lại: $85 - 62 = 23\text{ (kg)}$.
\textbf{Đáp án của câu hỏi là B}
}
\end{ex}

\begin{ex} %[2H3V18]
Hai can chứa tất cả $40\text{ L}$ nước. Can thứ nhất chứa $18\text{ L}$. Can thứ hai chứa nhiều hơn can thứ nhất bao nhiêu lít nước?
\choice
{\True $4\text{ L}$}
{$5\text{ L}$}
{$3\text{ L}$}
{$6\text{ L}$}
\loigiai{
Can 2 chứa: $40 - 18 = 22\text{ (L)}$. Can 2 hơn can 1: $22 - 18 = 4\text{ (L)}$.
\textbf{Đáp án của câu hỏi là A}
}
\end{ex}

\begin{ex} %[2H3V18]
Một con ngỗng nặng $6\text{ kg}$, con vịt nhẹ hơn con ngỗng $3\text{ kg}$. Hỏi con ngỗng nặng gấp mấy lần con vịt?
\choice
{$3$ lần}
{$4$ lần}
{\True $2$ lần}
{$1$ lần}
\loigiai{
Con vịt nặng: $6 - 3 = 3\text{ (kg)}$. Con ngỗng nặng gấp con vịt: $6 : 3 = 2$ (lần).
\textbf{Đáp án của câu hỏi là C}
}
\end{ex}

\begin{ex} %[2H3V18]
Rót $16\text{ L}$ nước vào $4$ bình như nhau. Mỗi bình có bao nhiêu lít nước?
\choice
{$3\text{ L}$}
{$5\text{ L}$}
{$6\text{ L}$}
{\True $4\text{ L}$}
\loigiai{
$16 : 4 = 4\text{ (L)}$.
\textbf{Đáp án của câu hỏi là D}
}
\end{ex}

\textbf{Phần 2. Trắc nghiệm đúng sai}
\setcounterpageref{ex}{0}

\begin{ex} %[2H3V18]
Xét tính đúng sai của các bài toán tổng hợp về đo lường:
\choiceTF
{\True Tổng của 25 kg và 35 kg bằng 60 kg.}
{Phép tính 45 L - 18 L có kết quả là 37 L.}
{\True Một can 10 L đựng được nhiều nước hơn hai can 4 L gộp lại.}
{Khối lượng luôn đo bằng Lít còn dung tích luôn đo bằng Ki-lô-gam.}
\loigiai{
\begin{itemize}
    \item Ý a) Đúng: $25 + 35 = 60\text{ kg}$.
    \item Ý b) Sai: $45 - 18 = 27\text{ L}$ (không phải 37 L).
    \item Ý c) Đúng: Hai can 4 L là $4 + 4 = 8\text{ L} < 10\text{ L}$. Đúng.
    \item Ý d) Sai: Ngược lại: khối lượng đo bằng kg, dung tích đo bằng Lít.
\end{itemize}
\begin{itemize}
    \item a) \textbf{Đúng}.
    \item b) \textbf{Sai}.
    \item c) \textbf{Đúng}.
    \item d) \textbf{Sai}.
\end{itemize}
}
\end{ex}

\begin{ex} %[2H3V18]
Một cửa hàng có 3 thùng sữa: Thùng 1 có 24 L, Thùng 2 có 28 L, Thùng 3 có 32 L:
\choiceTF
{\True Cả ba thùng có tất cả 84 L sữa.}
{\True Thùng 3 nhiều hơn Thùng 1 là 8 L sữa.}
{Thùng 2 ít hơn Thùng 3 là 6 L sữa.}
{\True Trung bình số sữa của ba thùng là 28 L.}
\loigiai{
\begin{itemize}
    \item Ý a) Đúng: $24 + 28 + 32 = 84\text{ L}$.
    \item Ý b) Đúng: $32 - 24 = 8\text{ L}$.
    \item Ý c) Sai: $32 - 28 = 4\text{ L}$ (không phải 6 L).
    \item Ý d) Đúng: $(24 + 28 + 32) : 3 = 84 : 3 = 28\text{ L}$. Đúng.
\end{itemize}
\begin{itemize}
    \item a) \textbf{Đúng}.
    \item b) \textbf{Đúng}.
    \item c) \textbf{Sai}.
    \item d) \textbf{Đúng}.
\end{itemize}
}
\end{ex}

\textbf{Phần 3. Trả lời ngắn}
\setcounterpageref{ex}{0}

\begin{ex} %[2H3V18]
Một người mua 18 kg gạo tẻ và 9 kg gạo nếp. Hỏi người đó mua tất cả bao nhiêu ki-lô-gam gạo?
\shortans{27}
\loigiai{
\textbf{Đáp số:} $27$\\
$18 + 9 = 27\text{ kg}$.
}
\end{ex}

\begin{ex} %[2H3V18]
Bể thứ nhất chứa 45 lít nước, bể thứ hai chứa 38 lít nước. Hỏi bể thứ nhất chứa nhiều hơn bể thứ hai bao nhiêu lít nước?
\shortans{7}
\loigiai{
\textbf{Đáp số:} $7$\\
$45 - 38 = 7\text{ L}$.
}
\end{ex}

\begin{ex} %[2H3V18]
Tính giá trị: $54\text{ kg} - 28\text{ kg} + 14\text{ kg}$.
\shortans{40}
\loigiai{
\textbf{Đáp số:} $40$\\
$54 - 28 = 26$; $26 + 14 = 40\text{ kg}$.
}
\end{ex}

\begin{ex} %[2H3V18]
Một can dầu chứa đầy được 15 lít. Người ta rót ra 8 lít thì trong can còn lại bao nhiêu lít dầu?
\shortans{7}
\loigiai{
\textbf{Đáp số:} $7$\\
$15 - 8 = 7\text{ L}$.
}
\end{ex}

\begin{ex} %[2H3V18]
Mẹ nuôi một đàn gà ăn hết 25 kg thóc trong tháng đầu và 28 kg thóc trong tháng thứ hai. Hỏi cả hai tháng đàn gà ăn hết bao nhiêu ki-lô-gam thóc?
\shortans{53}
\loigiai{
\textbf{Đáp số:} $53$\\
$25 + 28 = 53\text{ kg}$.
}
\end{ex}

\begin{ex} %[2H3V18]
Có 32 lít nước mắm chia đều vào các can 4 lít. Hỏi chia được bao nhiêu can?
\shortans{8}
\loigiai{
\textbf{Đáp số:} $8$\\
$32 : 4 = 8$ can.
}
\end{ex}

\newpage


% =========================================================================
% CHỦ ĐỀ 4 - PHÉP CỘNG, PHÉP TRỪ (CÓ NHỚ) TRONG PHẠM VI 100
% =========================================================================
\section{Phép cộng, phép trừ (có nhớ) trong phạm vi 100}

\TaoPhieuKiemTra{Chủ đề 4 -- Phép cộng, phép trừ (có nhớ) trong phạm vi 100}{19}{Phép cộng (có nhớ) số có hai chữ số với số có một chữ số}{https://loptoanthaylinh.vn/thi/chuyen-de/2/toan2_cd_b19}
\textbf{Phần 1. Trắc nghiệm lựa chọn}
\setcounterpageref{ex}{0}

\begin{ex} %[2D4V19]
Tính kết quả của phép tính: $47 + 8$.
\choice
{$54$}
{$56$}
{\True $55$}
{$53$}
\loigiai{
$47 + 8 = 55$.
\textbf{Đáp án của câu hỏi là C}
}
\end{ex}

\begin{ex} %[2D4V19]
Lớp 2A có $28$ học sinh, hôm nay có thêm $4$ bạn học sinh mới chuyển đến. Hỏi lớp 2A có tất cả bao nhiêu học sinh?
\choice
{\True $32$ học sinh}
{$31$ học sinh}
{$33$ học sinh}
{$30$ học sinh}
\loigiai{
$28 + 4 = 32$ (học sinh).
\textbf{Đáp án của câu hỏi là A}
}
\end{ex}

\begin{ex} %[2D4V19]
Tìm $x$ biết: $x - 7 = 35$.
\choice
{$41$}
{$43$}
{$40$}
{\True $42$}
\loigiai{
$x = 35 + 7 = 42$.
\textbf{Đáp án của câu hỏi là D}
}
\end{ex}

\begin{ex} %[2D4V19]
Tính giá trị biểu thức: $36 + 9 + 5$.
\choice
{$48$}
{\True $50$}
{$49$}
{$51$}
\loigiai{
$36 + 9 = 45$; $45 + 5 = 50$.
\textbf{Đáp án của câu hỏi là B}
}
\end{ex}

\begin{ex} %[2D4V19]
Một người thợ dệt buổi sáng dệt được $38\text{ m}$ vải, buổi chiều dệt thêm $7\text{ m}$ vải. Cả ngày người thợ dệt được bao nhiêu mét vải?
\choice
{$43\text{ m}$}
{$44\text{ m}$}
{\True $45\text{ m}$}
{$46\text{ m}$}
\loigiai{
$38 + 7 = 45\text{ (m)}$.
\textbf{Đáp án của câu hỏi là C}
}
\end{ex}

\begin{ex} %[2D4V19]
Điền số thích hợp: $59 + 6 = \dots$
\choice
{\True $65$}
{$64$}
{$66$}
{$63$}
\loigiai{
$59 + 6 = 65$.
\textbf{Đáp án của câu hỏi là A}
}
\end{ex}

\begin{ex} %[2D4V19]
Cho phép cộng $\overline{a8} + 5 = 53$. Chữ số $a$ có giá trị là
\choice
{$3$}
{$5$}
{$2$}
{\True $4$}
\loigiai{
Ta có: $\overline{a8} = 53 - 5 = 48$. Do đó $a = 4$.
\textbf{Đáp án của câu hỏi là D}
}
\end{ex}

\begin{ex} %[2D4V19]
Một cây cầu dài $65\text{ m}$, người ta nối dài thêm một đoạn $8\text{ m}$. Chiều dài cây cầu sau khi nối là
\choice
{$72\text{ m}$}
{\True $73\text{ m}$}
{$74\text{ m}$}
{$71\text{ m}$}
\loigiai{
$65 + 8 = 73\text{ (m)}$.
\textbf{Đáp án của câu hỏi là B}
}
\end{ex}

\begin{ex} %[2D4V19]
Tìm tổng của số lớn nhất có hai chữ số mà chữ số hàng chục là $7$ và số lớn nhất có một chữ số.
\choice
{$86$}
{$87$}
{\True $88$}
{$89$}
\loigiai{
Số lớn nhất có 2 chữ số hàng chục là 7 là 79. Số lớn nhất có 1 chữ số là 9.\nTổng: $79 + 9 = 88$.
\textbf{Đáp án của câu hỏi là C}
}
\end{ex}

\begin{ex} %[2D4V19]
Năm nay bác Ba $56$ tuổi. Hỏi $8$ năm nữa bác Ba bao nhiêu tuổi?
\choice
{\True $64$ tuổi}
{$62$ tuổi}
{$63$ tuổi}
{$65$ tuổi}
\loigiai{
$56 + 8 = 64$ (tuổi).
\textbf{Đáp án của câu hỏi là A}
}
\end{ex}

\begin{ex} %[2D4V19]
Điền dấu thích hợp: $74 + 8 \dots 85 - 4$.
\choice
{$<$}
{$=$}
{Không so sánh được}
{\True $>$}
\loigiai{
Vế trái: $74 + 8 = 82$. Vế phải: $85 - 4 = 81$. Vì $82 > 81$ nên điền dấu $>$. 
\textbf{Đáp án của câu hỏi là D}
}
\end{ex}

\begin{ex} %[2D4V19]
Đoạn thẳng $AB$ dài $48\text{ cm}$, đoạn thẳng $BC$ dài $6\text{ cm}$. Độ dài đường gấp khúc $ABC$ là
\choice
{$52\text{ cm}$}
{\True $54\text{ cm}$}
{$53\text{ cm}$}
{$55\text{ cm}$}
\loigiai{
$48 + 6 = 54\text{ (cm)}$.
\textbf{Đáp án của câu hỏi là B}
}
\end{ex}

\textbf{Phần 2. Trắc nghiệm đúng sai}
\setcounterpageref{ex}{0}

\begin{ex} %[2D4V19]
Xét tính đúng sai của các phép cộng có nhớ số có hai chữ số với số có một chữ số:
\choiceTF
{\True Phép tính $46 + 7$ có kết quả bằng $53$.}
{Khi thực hiện phép tính $38 + 5$, chữ số hàng chục của kết quả vẫn là 3.}
{\True Biểu thức $67 + 8$ có giá trị bằng $75$.}
{Số 89 cộng với bất kì số có một chữ số nào cũng cho kết quả nhỏ hơn 95.}
\loigiai{
\begin{itemize}
    \item Ý a) Đúng: $46 + 7 = 53$.
    \item Ý b) Sai: Vì $8 + 5 = 13$ có nhớ 1 sang hàng chục nên chữ số hàng chục là $3 + 1 = 4$.
    \item Ý c) Đúng: $67 + 8 = 75$. Đúng.
    \item Ý d) Sai: Nếu cộng với 9: $89 + 9 = 98 > 95$. Khẳng định sai.
\end{itemize}
\begin{itemize}
    \item a) \textbf{Đúng}.
    \item b) \textbf{Sai}.
    \item c) \textbf{Đúng}.
    \item d) \textbf{Sai}.
\end{itemize}
}
\end{ex}

\begin{ex} %[2D4V19]
Một bác thợ rèn làm một thanh kiếm bằng cách hàn hai thanh thép: thanh dài 58 cm và thanh ngắn 7 cm:
\choiceTF
{\True Chiều dài tổng cộng của hai thanh thép là 65 cm.}
{\True Thanh dài hơn thanh ngắn 51 cm.}
{Nếu bớt ở thanh dài 8 cm thì thanh dài ngắn hơn thanh ngắn.}
{\True Tổng $58 + 7$ là một số lẻ.}
\loigiai{
\begin{itemize}
    \item Ý a) Đúng: $58 + 7 = 65\text{ cm}$.
    \item Ý b) Đúng: $58 - 7 = 51\text{ cm}$.
    \item Ý c) Sai: Thanh dài còn $58 - 8 = 50\text{ cm}$, vẫn dài hơn 7 cm.
    \item Ý d) Đúng: 65 là số lẻ. Đúng.
\end{itemize}
\begin{itemize}
    \item a) \textbf{Đúng}.
    \item b) \textbf{Đúng}.
    \item c) \textbf{Sai}.
    \item d) \textbf{Đúng}.
\end{itemize}
}
\end{ex}

\textbf{Phần 3. Trả lời ngắn}
\setcounterpageref{ex}{0}

\begin{ex} %[2D4V19]
Tính kết quả của phép tính: $68 + 7$.
\shortans{75}
\loigiai{
\textbf{Đáp số:} $75$\\
$68 + 7 = 75$.
}
\end{ex}

\begin{ex} %[2D4V19]
Một cuộn dây điện dài 45 m, người ta nối thêm 8 m dây điện nữa. Hỏi cuộn dây điện dài bao nhiêu mét?
\shortans{53}
\loigiai{
\textbf{Đáp số:} $53$\\
$45 + 8 = 53\text{ m}$.
}
\end{ex}

\begin{ex} %[2D4V19]
Tìm x biết: $x - 9 = 74$.
\shortans{83}
\loigiai{
\textbf{Đáp số:} $83$\\
$x = 74 + 9 = 83$.
}
\end{ex}

\begin{ex} %[2D4V19]
Tổng của số tròn chục liền trước 50 và số lớn nhất có một chữ số là bao nhiêu?
\shortans{49}
\loigiai{
\textbf{Đáp số:} $49$\\
Số tròn chục liền trước 50 là 40. Số lớn nhất có 1 chữ số là 9. Tổng: $40 + 9 = 49$.
}
\end{ex}

\begin{ex} %[2D4V19]
Mẹ mua 38 quả trứng gà và 7 quả trứng vịt. Hỏi mẹ mua tất cả bao nhiêu quả trứng?
\shortans{45}
\loigiai{
\textbf{Đáp số:} $45$\\
$38 + 7 = 45$ quả.
}
\end{ex}

\begin{ex} %[2D4V19]
Một hình tam giác có độ dài ba cạnh lần lượt là 25 cm, 18 cm và 7 cm. Tính chu vi hình tam giác đó.
\shortans{50}
\loigiai{
\textbf{Đáp số:} $50$\\
Chu vi: $25 + 18 + 7 = 50\text{ cm}$.
}
\end{ex}

\newpage

\TaoPhieuKiemTra{Chủ đề 4 -- Phép cộng, phép trừ (có nhớ) trong phạm vi 100}{20}{Phép cộng (có nhớ) số có hai chữ số với số có hai chữ số}{https://loptoanthaylinh.vn/thi/chuyen-de/2/toan2_cd_b20}
\textbf{Phần 1. Trắc nghiệm lựa chọn}
\setcounterpageref{ex}{0}

\begin{ex} %[2D4V20]
Tính giá trị phép tính: $38 + 27$.
\choice
{$64$}
{\True $65$}
{$66$}
{$55$}
\loigiai{
$38 + 27 = 65$.
\textbf{Đáp án của câu hỏi là B}
}
\end{ex}

\begin{ex} %[2D4V20]
Lớp 2A có $29$ học sinh, lớp 2B có $34$ học sinh. Cả hai lớp có bao nhiêu học sinh?
\choice
{$61$ học sinh}
{$62$ học sinh}
{$64$ học sinh}
{\True $63$ học sinh}
\loigiai{
$29 + 34 = 63$ (học sinh).
\textbf{Đáp án của câu hỏi là D}
}
\end{ex}

\begin{ex} %[2D4V20]
Tìm $x$ biết: $x - 36 = 28$.
\choice
{\True $64$}
{$62$}
{$63$}
{$65$}
\loigiai{
$x = 28 + 36 = 64$.
\textbf{Đáp án của câu hỏi là A}
}
\end{ex}

\begin{ex} %[2D4V20]
Tính nhanh: $25 + 18 + 35 + 12$.
\choice
{$80$}
{$85$}
{\True $90$}
{$95$}
\loigiai{
$(25 + 35) + (18 + 12) = 60 + 30 = 90$.
\textbf{Đáp án của câu hỏi là C}
}
\end{ex}

\begin{ex} %[2D4V20]
Một cửa hàng buổi sáng bán được $48\text{ kg}$ gạo, buổi chiều bán được $37\text{ kg}$ gạo. Cả ngày cửa hàng bán được bao nhiêu ki-lô-gam gạo?
\choice
{$83\text{ kg}$}
{\True $85\text{ kg}$}
{$84\text{ kg}$}
{$86\text{ kg}$}
\loigiai{
$48 + 37 = 85\text{ (kg)}$.
\textbf{Đáp án của câu hỏi là B}
}
\end{ex}

\begin{ex} %[2D4V20]
Cho phép tính: $\overline{a6} + 28 = 64$. Giá trị của chữ số $a$ là
\choice
{$2$}
{$4$}
{$1$}
{\True $3$}
\loigiai{
Ta có: $\overline{a6} = 64 - 28 = 36$. Do đó $a = 3$.
\textbf{Đáp án của câu hỏi là D}
}
\end{ex}

\begin{ex} %[2D4V20]
Một con đường dài $72\text{ m}$. Đội công nhân thứ nhất sửa được $35\text{ m}$, đội thứ hai sửa được $37\text{ m}$. Hai đội đã sửa được tất cả bao nhiêu mét đường?
\choice
{\True $72\text{ m}$}
{$70\text{ m}$}
{$71\text{ m}$}
{$73\text{ m}$}
\loigiai{
$35 + 37 = 72\text{ (m)}$. Hai đội đã sửa xong toàn bộ con đường.
\textbf{Đáp án của câu hỏi là A}
}
\end{ex}

\begin{ex} %[2D4V20]
Tìm tổng của số lớn nhất có hai chữ số khác nhau nhỏ hơn $50$ và số nhỏ nhất có hai chữ số giống nhau lớn hơn $20$.
\choice
{$70$}
{$72$}
{\True $71$}
{$69$}
\loigiai{
Số lớn nhất có 2 chữ số khác nhau nhỏ hơn 50 là 49. Số nhỏ nhất có 2 chữ số giống nhau lớn hơn 20 là 22.\nTổng: $49 + 22 = 71$.
\textbf{Đáp án của câu hỏi là C}
}
\end{ex}

\begin{ex} %[2D4V20]
Điền dấu thích hợp: $47 + 38 \dots 56 + 29$.
\choice
{$>$}
{\True $=$}
{$<$}
{Không so sánh được}
\loigiai{
$47 + 38 = 85$; $56 + 29 = 85$. Hai vế bằng nhau nên điền dấu $=$.
\textbf{Đáp án của câu hỏi là B}
}
\end{ex}

\begin{ex} %[2D4V20]
Một trang trại có $46$ con bò sữa và $38$ con bò thịt. Trang trại có tất cả bao nhiêu con bò?
\choice
{$82$ con bò}
{$83$ con bò}
{$85$ con bò}
{\True $84$ con bò}
\loigiai{
$46 + 38 = 84$ (con bò).
\textbf{Đáp án của câu hỏi là D}
}
\end{ex}

\begin{ex} %[2D4V20]
Đoạn dây thứ nhất dài $39\text{ cm}$, đoạn dây thứ hai dài $45\text{ cm}$. Nối hai đoạn dây lại (không kể nút nối) được đoạn dây dài bao nhiêu đề-xi-mét, xăng-ti-mét?
\choice
{\True $8\text{ dm } 4\text{ cm}$}
{$8\text{ dm } 3\text{ cm}$}
{$8\text{ dm } 5\text{ cm}$}
{$7\text{ dm } 4\text{ cm}$}
\loigiai{
Tổng chiều dài: $39 + 45 = 84\text{ cm} = 8\text{ dm } 4\text{ cm}$.
\textbf{Đáp án của câu hỏi là A}
}
\end{ex}

\begin{ex} %[2D4V20]
Tìm một số biết rằng lấy số đó trừ đi $29$ thì được $48$.
\choice
{$76$}
{$78$}
{\True $77$}
{$75$}
\loigiai{
Số đó là: $48 + 29 = 77$.
\textbf{Đáp án của câu hỏi là C}
}
\end{ex}

\textbf{Phần 2. Trắc nghiệm đúng sai}
\setcounterpageref{ex}{0}

\begin{ex} %[2D4V20]
Xét tính đúng sai của phép cộng có nhớ trong phạm vi 100:
\choiceTF
{\True Phép cộng $47 + 36$ có kết quả là $83$.}
{Tổng của hai số có hai chữ số luôn nhỏ hơn 90.}
{\True Khi cộng $28 + 37$, hàng đơn vị có $8 + 7 = 15$ viết 5 nhớ 1 sang hàng chục.}
{Nếu tổng bằng 75 và một số hạng là 38 thì số hạng kia là 47.}
\loigiai{
\begin{itemize}
    \item Ý a) Đúng: $47 + 36 = 83$.
    \item Ý b) Sai: Ví dụ: $50 + 50 = 100 > 90$.
    \item Ý c) Đúng: Quy tắc đặt tính rồi tính chuẩn.
    \item Ý d) Sai: Số hạng kia là: $75 - 38 = 37$ (không phải 47).
\end{itemize}
\begin{itemize}
    \item a) \textbf{Đúng}.
    \item b) \textbf{Sai}.
    \item c) \textbf{Đúng}.
    \item d) \textbf{Sai}.
\end{itemize}
}
\end{ex}

\begin{ex} %[2D4V20]
Hai kho gạo có tất cả 92 bao gạo, kho A có 48 bao gạo:
\choiceTF
{\True Kho B có 44 bao gạo.}
{\True Kho A nhiều hơn Kho B là 4 bao gạo.}
{Chuyển 4 bao gạo từ kho A sang kho B thì hai kho bằng nhau.}
{\True Nếu kho A nhập thêm 5 bao thì kho A có 53 bao gạo.}
\loigiai{
\begin{itemize}
    \item Ý a) Đúng: $92 - 48 = 44$ bao gạo.
    \item Ý b) Đúng: $48 - 44 = 4$ bao gạo.
    \item Ý c) Sai: Chỉ cần chuyển $4 : 2 = 2$ bao thì cả hai kho đều có 46 bao.
    \item Ý d) Đúng: $48 + 5 = 53$ bao. Đúng.
\end{itemize}
\begin{itemize}
    \item a) \textbf{Đúng}.
    \item b) \textbf{Đúng}.
    \item c) \textbf{Sai}.
    \item d) \textbf{Đúng}.
\end{itemize}
}
\end{ex}

\textbf{Phần 3. Trả lời ngắn}
\setcounterpageref{ex}{0}

\begin{ex} %[2D4V20]
Tính giá trị: $46 + 37$.
\shortans{83}
\loigiai{
\textbf{Đáp số:} $83$\\
$46 + 37 = 83$.
}
\end{ex}

\begin{ex} %[2D4V20]
Mẹ mua 28 kg gạo tẻ và 15 kg gạo nếp. Hỏi mẹ mua tất cả bao nhiêu ki-lô-gam gạo?
\shortans{43}
\loigiai{
\textbf{Đáp số:} $43$\\
$28 + 15 = 43\text{ kg}$.
}
\end{ex}

\begin{ex} %[2D4V20]
Tìm x biết: $x - 45 = 28$.
\shortans{73}
\loigiai{
\textbf{Đáp số:} $73$\\
$x = 28 + 45 = 73$.
}
\end{ex}

\begin{ex} %[2D4V20]
Tính nhanh: $27 + 34 + 13 + 16$.
\shortans{90}
\loigiai{
\textbf{Đáp số:} $90$\\
$(27 + 13) + (34 + 16) = 40 + 50 = 90$.
}
\end{ex}

\begin{ex} %[2D4V20]
Thùng thứ nhất chứa 38 lít dầu, thùng thứ hai chứa 47 lít dầu. Hỏi cả hai thùng chứa bao nhiêu lít dầu?
\shortans{85}
\loigiai{
\textbf{Đáp số:} $85$\\
$38 + 47 = 85\text{ L}$.
}
\end{ex}

\begin{ex} %[2D4V20]
Đoạn thẳng AB dài 36 cm, đoạn thẳng BC dài 48 cm. Tính độ dài đoạn thẳng AC biết ba điểm A, B, C thẳng hàng theo thứ tự đó.
\shortans{84}
\loigiai{
\textbf{Đáp số:} $84$\\
$36 + 48 = 84\text{ cm}$.
}
\end{ex}

\newpage

\TaoPhieuKiemTra{Chủ đề 4 -- Phép cộng, phép trừ (có nhớ) trong phạm vi 100}{21}{Luyện tập chung}{https://loptoanthaylinh.vn/thi/chuyen-de/2/toan2_cd_b21}
\textbf{Phần 1. Trắc nghiệm lựa chọn}
\setcounterpageref{ex}{0}

\begin{ex} %[2D4V21]
Tính giá trị biểu thức: $37 + 19 + 24$.
\choice
{\True $80$}
{$78$}
{$79$}
{$81$}
\loigiai{
$37 + 19 = 56$; $56 + 24 = 80$.
\textbf{Đáp án của câu hỏi là A}
}
\end{ex}

\begin{ex} %[2D4V21]
Một cuộn dây dài $48\text{ m}$, người ta nối thêm một đoạn $25\text{ m}$ rồi lại nối thêm một đoạn $17\text{ m}$. Cuộn dây lúc này dài bao nhiêu mét?
\choice
{$88\text{ m}$}
{$89\text{ m}$}
{\True $90\text{ m}$}
{$91\text{ m}$}
\loigiai{
$48 + 25 + 17 = 48 + (25 + 17) = 48 + 42 = 90\text{ (m)}$.
\textbf{Đáp án của câu hỏi là C}
}
\end{ex}

\begin{ex} %[2D4V21]
Tìm số tự nhiên $y$ biết: $y - 28 = 47$.
\choice
{$74$}
{\True $75$}
{$76$}
{$65$}
\loigiai{
$y = 47 + 28 = 75$.
\textbf{Đáp án của câu hỏi là B}
}
\end{ex}

\begin{ex} %[2D4V21]
Một cửa hàng bán được $39$ hộp bánh trong ngày thứ nhất, ngày thứ hai bán được nhiều hơn ngày thứ nhất $15$ hộp. Cả hai ngày cửa hàng bán được bao nhiêu hộp bánh?
\choice
{$83$ hộp}
{$92$ hộp}
{$94$ hộp}
{\True $93$ hộp}
\loigiai{
Ngày thứ hai bán: $39 + 15 = 54$ (hộp).\nCả hai ngày bán: $39 + 54 = 93$ (hộp).
\textbf{Đáp án của câu hỏi là D}
}
\end{ex}

\begin{ex} %[2D4V21]
Tính chu vi hình tam giác có độ dài các cạnh lần lượt là $27\text{ cm}, 35\text{ cm}, 18\text{ cm}$.
\choice
{\True $80\text{ cm}$}
{$78\text{ cm}$}
{$82\text{ cm}$}
{$79\text{ cm}$}
\loigiai{
Chu vi: $27 + 35 + 18 = 80\text{ (cm)}$.
\textbf{Đáp án của câu hỏi là A}
}
\end{ex}

\begin{ex} %[2D4V21]
Tìm một số biết rằng lấy số đó trừ đi $38$ thì bằng tổng của $26$ và $19$.
\choice
{$81$}
{$82$}
{\True $83$}
{$84$}
\loigiai{
Tổng hai số: $26 + 19 = 45$. Số cần tìm: $45 + 38 = 83$.
\textbf{Đáp án của câu hỏi là C}
}
\end{ex}

\begin{ex} %[2D4V21]
Điền số thích hợp: $46 + 37 = 28 + \dots$
\choice
{$54$}
{\True $55$}
{$56$}
{$53$}
\loigiai{
Vế trái: $46 + 37 = 83$. Ta có: $83 - 28 = 55$. Vậy điền 55.
\textbf{Đáp án của câu hỏi là B}
}
\end{ex}

\begin{ex} %[2D4V21]
Bác nông dân có $48$ con gà và $37$ con vịt. Bác mua thêm $9$ con ngỗng. Hỏi bác có tất cả bao nhiêu con gia cầm?
\choice
{$92$ con}
{$93$ con}
{$95$ con}
{\True $94$ con}
\loigiai{
$48 + 37 + 9 = 85 + 9 = 94$ (con).
\textbf{Đáp án của câu hỏi là D}
}
\end{ex}

\begin{ex} %[2D4V21]
Tổng của số chẵn lớn nhất có hai chữ số mà chữ số hàng chục là $5$ và số lẻ nhỏ nhất có hai chữ số mà chữ số hàng chục là $3$ là
\choice
{\True $89$}
{$87$}
{$90$}
{$88$}
\loigiai{
Số chẵn lớn nhất có 2 chữ số hàng chục là 5 là 58. Số lẻ nhỏ nhất có 2 chữ số hàng chục là 3 là 31.\nTổng: $58 + 31 = 89$.
\textbf{Đáp án của câu hỏi là A}
}
\end{ex}

\begin{ex} %[2D4V21]
Một xe khách xuất phát với $28$ người. Đến bến đỗ có $15$ người lên xe. Xe đi tiếp và đón thêm $14$ người nữa. Trên xe có bao nhiêu hành khách?
\choice
{$55$ người}
{$56$ người}
{\True $57$ người}
{$58$ người}
\loigiai{
$28 + 15 + 14 = 43 + 14 = 57$ (người).
\textbf{Đáp án của câu hỏi là C}
}
\end{ex}

\begin{ex} %[2D4V21]
Hai số có hiệu bằng $27$, số trừ là $46$. Số bị trừ là
\choice
{$71$}
{\True $73$}
{$72$}
{$74$}
\loigiai{
Số bị trừ: $27 + 46 = 73$.
\textbf{Đáp án của câu hỏi là B}
}
\end{ex}

\begin{ex} %[2D4V21]
Điền dấu thích hợp: $38 + 47 \dots 59 + 25$.
\choice
{$<$}
{$=$}
{Không so sánh được}
{\True $>$}
\loigiai{
$38 + 47 = 85$; $59 + 25 = 84$. Vì $85 > 84$ nên điền dấu $>$. 
\textbf{Đáp án của câu hỏi là D}
}
\end{ex}

\textbf{Phần 2. Trắc nghiệm đúng sai}
\setcounterpageref{ex}{0}

\begin{ex} %[2D4V21]
Xét tính đúng sai của các bài toán tổng hợp cộng có nhớ:
\choiceTF
{\True Biểu thức $28 + 34 + 18$ có kết quả bằng $80$.}
{Nếu một số hạng tăng 5 đơn vị và số hạng kia tăng 7 đơn vị thì tổng tăng 10 đơn vị.}
{\True Phép tính $56 + 27$ có chữ số tận cùng là 3.}
{\True Tổng của hai số chẵn luôn luôn là một số chẵn.}
\loigiai{
\begin{itemize}
    \item Ý a) Đúng: $28 + 34 + 18 = 80$.
    \item Ý b) Sai: Tổng tăng thêm: $5 + 7 = 12$ đơn vị (không phải 10).
    \item Ý c) Đúng: $6 + 7 = 13$, chữ số tận cùng là 3. Đúng.
    \item Ý d) Đúng: Đúng.
\end{itemize}
\begin{itemize}
    \item a) \textbf{Đúng}.
    \item b) \textbf{Sai}.
    \item c) \textbf{Đúng}.
    \item d) \textbf{Đúng}.
\end{itemize}
}
\end{ex}

\begin{ex} %[2D4V21]
Một thư viện trường có 38 quyển truyện tranh và 45 quyển truyện cổ tích:
\choiceTF
{\True Cả hai loại truyện có tất cả 83 quyển.}
{\True Truyện cổ tích nhiều hơn truyện tranh 7 quyển.}
{Nhập thêm 5 quyển truyện tranh thì số truyện tranh sẽ nhiều hơn truyện cổ tích.}
{Tổng số sách chia đều được cho 2 học sinh mà không dư.}
\loigiai{
\begin{itemize}
    \item Ý a) Đúng: $38 + 45 = 83$ quyển.
    \item Ý b) Đúng: $45 - 38 = 7$ quyển.
    \item Ý c) Sai: Truyện tranh có $38 + 5 = 43 < 45$. Sai.
    \item Ý d) Sai: 83 là số lẻ nên không chia hết cho 2.
\end{itemize}
\begin{itemize}
    \item a) \textbf{Đúng}.
    \item b) \textbf{Đúng}.
    \item c) \textbf{Sai}.
    \item d) \textbf{Sai}.
\end{itemize}
}
\end{ex}

\textbf{Phần 3. Trả lời ngắn}
\setcounterpageref{ex}{0}

\begin{ex} %[2D4V21]
Tính giá trị: $47 + 38$.
\shortans{85}
\loigiai{
\textbf{Đáp số:} $85$\\
$47 + 38 = 85$.
}
\end{ex}

\begin{ex} %[2D4V21]
Một trang trại có 36 con bò và 49 con dê. Hỏi trang trại có tất cả bao nhiêu con bò và dê?
\shortans{85}
\loigiai{
\textbf{Đáp số:} $85$\\
$36 + 49 = 85$ con.
}
\end{ex}

\begin{ex} %[2D4V21]
Tìm số tự nhiên a biết: $a - 37 = 48$.
\shortans{85}
\loigiai{
\textbf{Đáp số:} $85$\\
$a = 48 + 37 = 85$.
}
\end{ex}

\begin{ex} %[2D4V21]
Một cuộn dây dài 28 m, người ta nối thêm vào cuộn dây 35 m nữa rồi nối tiếp 19 m. Hỏi cuộn dây dài tất cả bao nhiêu mét?
\shortans{82}
\loigiai{
\textbf{Đáp số:} $82$\\
$28 + 35 + 19 = 82\text{ m}$.
}
\end{ex}

\begin{ex} %[2D4V21]
Tìm tổng của số lớn nhất có hai chữ số mà tổng hai chữ số bằng 8 và số 18.
\shortans{98}
\loigiai{
\textbf{Đáp số:} $98$\\
Số lớn nhất có 2 chữ số tổng bằng 8 là 80. Tổng: $80 + 18 = 98$.
}
\end{ex}

\begin{ex} %[2D4V21]
Bạn Nam sưu tầm được 46 con tem, bạn Hùng sưu tầm được 38 con tem. Hỏi cả hai bạn sưu tầm được bao nhiêu con tem?
\shortans{84}
\loigiai{
\textbf{Đáp số:} $84$\\
$46 + 38 = 84$ con tem.
}
\end{ex}

\newpage

\TaoPhieuKiemTra{Chủ đề 4 -- Phép cộng, phép trừ (có nhớ) trong phạm vi 100}{22}{Phép trừ (có nhớ) số có hai chữ số cho số có một chữ số}{https://loptoanthaylinh.vn/thi/chuyen-de/2/toan2_cd_b22}
\textbf{Phần 1. Trắc nghiệm lựa chọn}
\setcounterpageref{ex}{0}

\begin{ex} %[2D4V22]
Kết quả của phép trừ $52 - 7$ là
\choice
{$43$}
{$44$}
{$46$}
{\True $45$}
\loigiai{
$52 - 7 = 45$.
\textbf{Đáp án của câu hỏi là D}
}
\end{ex}

\begin{ex} %[2D4V22]
Một đàn gà có $43$ con gà. Người ta bán đi $8$ con gà. Hỏi đàn gà còn lại bao nhiêu con gà?
\choice
{$34$ con gà}
{\True $35$ con gà}
{$36$ con gà}
{$33$ con gà}
\loigiai{
$43 - 8 = 35$ (con gà).
\textbf{Đáp án của câu hỏi là B}
}
\end{ex}

\begin{ex} %[2D4V22]
Tìm $x$ biết: $61 - x = 9$.
\choice
{$51$}
{$53$}
{\True $52$}
{$50$}
\loigiai{
$x = 61 - 9 = 52$.
\textbf{Đáp án của câu hỏi là C}
}
\end{ex}

\begin{ex} %[2D4V22]
Tính giá trị của biểu thức: $73 - 6 - 8$.
\choice
{\True $59$}
{$58$}
{$60$}
{$57$}
\loigiai{
$73 - 6 = 67$; $67 - 8 = 59$.
\textbf{Đáp án của câu hỏi là A}
}
\end{ex}

\begin{ex} %[2D4V22]
Một sợi dây dài $84\text{ cm}$, người ta cắt bớt đi $9\text{ cm}$. Chiều dài đoạn dây còn lại là
\choice
{$73\text{ cm}$}
{$74\text{ cm}$}
{$76\text{ cm}$}
{\True $75\text{ cm}$}
\loigiai{
$84 - 9 = 75\text{ (cm)}$.
\textbf{Đáp án của câu hỏi là D}
}
\end{ex}

\begin{ex} %[2D4V22]
Điền số thích hợp: $35 - 8 = 19 + \dots$
\choice
{$7$}
{\True $8$}
{$9$}
{$6$}
\loigiai{
Vế trái: $35 - 8 = 27$. Ta có: $27 - 19 = 8$. Vậy điền 8.
\textbf{Đáp án của câu hỏi là B}
}
\end{ex}

\begin{ex} %[2D4V22]
Năm nay mẹ $41$ tuổi, mẹ hơn con $8$ lần tuổi con là không phải, mẹ hơn con $32$ tuổi. Tuổi của con là
\choice
{$8$ tuổi}
{$10$ tuổi}
{\True $9$ tuổi}
{$7$ tuổi}
\loigiai{
$41 - 32 = 9$ (tuổi).
\textbf{Đáp án của câu hỏi là C}
}
\end{ex}

\begin{ex} %[2D4V22]
Cho phép tính: $\overline{a2} - 7 = 35$. Giá trị của chữ số $a$ là
\choice
{\True $4$}
{$3$}
{$5$}
{$2$}
\loigiai{
Ta có: $\overline{a2} = 35 + 7 = 42$. Do đó $a = 4$.
\textbf{Đáp án của câu hỏi là A}
}
\end{ex}

\begin{ex} %[2D4V22]
Một xe chở $91\text{ kg}$ gạo, đã dỡ xuống $8\text{ kg}$. Khối lượng gạo còn trên xe là
\choice
{$81\text{ kg}$}
{$82\text{ kg}$}
{$84\text{ kg}$}
{\True $83\text{ kg}$}
\loigiai{
$91 - 8 = 83\text{ (kg)}$.
\textbf{Đáp án của câu hỏi là D}
}
\end{ex}

\begin{ex} %[2D4V22]
Hiệu giữa số tròn chục liền sau số $60$ và số lớn nhất có một chữ số là
\choice
{$59$}
{\True $61$}
{$60$}
{$62$}
\loigiai{
Số tròn chục liền sau 60 là 70. Số lớn nhất có 1 chữ số là 9.\nHiệu: $70 - 9 = 61$.
\textbf{Đáp án của câu hỏi là B}
}
\end{ex}

\begin{ex} %[2D4V22]
Tìm một số biết rằng lấy số đó cộng với $6$ thì được $82$.
\choice
{$74$}
{$75$}
{\True $76$}
{$78$}
\loigiai{
Số đó là: $82 - 6 = 76$.
\textbf{Đáp án của câu hỏi là C}
}
\end{ex}

\begin{ex} %[2D4V22]
Điền dấu thích hợp: $63 - 5 \dots 49 + 8$.
\choice
{\True $>$}
{$<$}
{$=$}
{Không so sánh được}
\loigiai{
$63 - 5 = 58$; $49 + 8 = 57$. Vì $58 > 57$ nên điền dấu $>$. 
\textbf{Đáp án của câu hỏi là A}
}
\end{ex}

\textbf{Phần 2. Trắc nghiệm đúng sai}
\setcounterpageref{ex}{0}

\begin{ex} %[2D4V22]
Xét tính đúng sai của phép trừ có nhớ số có hai chữ số cho số có một chữ số:
\choiceTF
{\True Phép trừ $71 - 8$ có kết quả bằng $63$.}
{Trong phép trừ $54 - 6$, vì 4 không trừ được 6 nên ta lấy 6 trừ 4 bằng 2.}
{\True Biểu thức $82 - 9$ có giá trị bằng $73$.}
{\True Nếu số bị trừ giảm đi 4 đơn vị và số trừ giữ nguyên thì hiệu giảm đi 4 đơn vị.}
\loigiai{
\begin{itemize}
    \item Ý a) Đúng: $71 - 8 = 63$.
    \item Ý b) Sai: Sai lầm phổ biến của học sinh: phải mượn 1 chục thành 14 trừ 6 bằng 8.
    \item Ý c) Đúng: $82 - 9 = 73$. Đúng.
    \item Ý d) Đúng: Đúng theo tính chất phép trừ.
\end{itemize}
\begin{itemize}
    \item a) \textbf{Đúng}.
    \item b) \textbf{Sai}.
    \item c) \textbf{Đúng}.
    \item d) \textbf{Đúng}.
\end{itemize}
}
\end{ex}

\begin{ex} %[2D4V22]
Bác Năm thu hoạch được 62 quả dừa, bác đã bán đi 9 quả dừa:
\choiceTF
{\True Bác Năm còn lại 53 quả dừa.}
{\True Số dừa đã bán ít hơn số dừa còn lại là 44 quả.}
{Nếu bác bán thêm 3 quả nữa thì bác còn lại 51 quả.}
{Số 53 là số chẵn.}
\loigiai{
\begin{itemize}
    \item Ý a) Đúng: $62 - 9 = 53$ quả.
    \item Ý b) Đúng: $53 - 9 = 44$ quả.
    \item Ý c) Sai: Số dừa lúc này còn: $53 - 3 = 50$ quả (không phải 51 quả).
    \item Ý d) Sai: 53 có chữ số tận cùng là 3 nên là số lẻ.
\end{itemize}
\begin{itemize}
    \item a) \textbf{Đúng}.
    \item b) \textbf{Đúng}.
    \item c) \textbf{Sai}.
    \item d) \textbf{Sai}.
\end{itemize}
}
\end{ex}

\textbf{Phần 3. Trả lời ngắn}
\setcounterpageref{ex}{0}

\begin{ex} %[2D4V22]
Tính kết quả: $75 - 8$.
\shortans{67}
\loigiai{
\textbf{Đáp số:} $67$\\
$75 - 8 = 67$.
}
\end{ex}

\begin{ex} %[2D4V22]
Tìm số tự nhiên x biết: $x + 9 = 62$.
\shortans{53}
\loigiai{
\textbf{Đáp số:} $53$\\
$x = 62 - 9 = 53$.
}
\end{ex}

\begin{ex} %[2D4V22]
Một giá sách có 51 quyển sách. Người ta mượn đi 6 quyển sách. Hỏi trên giá còn lại bao nhiêu quyển sách?
\shortans{45}
\loigiai{
\textbf{Đáp số:} $45$\\
$51 - 6 = 45$ quyển.
}
\end{ex}

\begin{ex} %[2D4V22]
Một sợi dây dài 83 cm, người ta cắt đi 7 cm. Hỏi sợi dây còn lại dài bao nhiêu xăng-ti-mét?
\shortans{76}
\loigiai{
\textbf{Đáp số:} $76$\\
$83 - 7 = 76\text{ cm}$.
}
\end{ex}

\begin{ex} %[2D4V22]
Tìm hiệu của số bé nhất có hai chữ số mà chữ số hàng chục là 9 và số 8.
\shortans{82}
\loigiai{
\textbf{Đáp số:} $82$\\
Số bé nhất có 2 chữ số hàng chục là 9 là 90. Hiệu: $90 - 8 = 82$.
}
\end{ex}

\begin{ex} %[2D4V22]
Lớp 2A có 34 học sinh, trong đó có 8 bạn học sinh giỏi. Hỏi lớp 2A có bao nhiêu bạn không đạt học sinh giỏi?
\shortans{26}
\loigiai{
\textbf{Đáp số:} $26$\\
$34 - 8 = 26$ bạn.
}
\end{ex}

\newpage

\TaoPhieuKiemTra{Chủ đề 4 -- Phép cộng, phép trừ (có nhớ) trong phạm vi 100}{23}{Phép trừ (có nhớ) số có hai chữ số cho số có hai chữ số}{https://loptoanthaylinh.vn/thi/chuyen-de/2/toan2_cd_b23}
\textbf{Phần 1. Trắc nghiệm lựa chọn}
\setcounterpageref{ex}{0}

\begin{ex} %[2D4V23]
Tính kết quả của phép trừ: $72 - 38$.
\choice
{$33$}
{$35$}
{\True $34$}
{$44$}
\loigiai{
$72 - 38 = 34$.
\textbf{Đáp án của câu hỏi là C}
}
\end{ex}

\begin{ex} %[2D4V23]
Một mảnh vườn có $85$ cây bưởi và cây cam, trong đó có $48$ cây bưởi. Hỏi có bao nhiêu cây cam?
\choice
{$36$ cây cam}
{\True $37$ cây cam}
{$38$ cây cam}
{$47$ cây cam}
\loigiai{
$85 - 48 = 37$ (cây cam).
\textbf{Đáp án của câu hỏi là B}
}
\end{ex}

\begin{ex} %[2D4V23]
Tìm $x$ biết: $x + 47 = 81$.
\choice
{$32$}
{$33$}
{$35$}
{\True $34$}
\loigiai{
$x = 81 - 47 = 34$.
\textbf{Đáp án của câu hỏi là D}
}
\end{ex}

\begin{ex} %[2D4V23]
Một cửa hàng có $90\text{ kg}$ gạo, buổi sáng bán được $36\text{ kg}$ gạo, buổi chiều bán được $28\text{ kg}$ gạo. Hỏi cửa hàng còn lại bao nhiêu ki-lô-gam gạo?
\choice
{\True $26\text{ kg}$}
{$24\text{ kg}$}
{$25\text{ kg}$}
{$28\text{ kg}$}
\loigiai{
Tổng gạo đã bán: $36 + 28 = 64\text{ (kg)}$. Còn lại: $90 - 64 = 26\text{ (kg)}$.
\textbf{Đáp án của câu hỏi là A}
}
\end{ex}

\begin{ex} %[2D4V23]
Tính giá trị của biểu thức: $64 - 27 - 18$.
\choice
{$18$}
{$20$}
{\True $19$}
{$17$}
\loigiai{
$64 - 27 = 37$; $37 - 18 = 19$.
\textbf{Đáp án của câu hỏi là C}
}
\end{ex}

\begin{ex} %[2D4V23]
Cho phép tính: $71 - \overline{a5} = 36$. Giá trị của chữ số $a$ là
\choice
{$2$}
{\True $3$}
{$4$}
{$1$}
\loigiai{
Ta có: $\overline{a5} = 71 - 36 = 35$. Do đó $a = 3$.
\textbf{Đáp án của câu hỏi là B}
}
\end{ex}

\begin{ex} %[2D4V23]
Một quãng đường dài $83\text{ km}$, ô tô đã đi được $49\text{ km}$. Hỏi ô tô còn phải đi bao nhiêu ki-lô-mét nữa để hết quãng đường?
\choice
{$32\text{ km}$}
{$33\text{ km}$}
{$35\text{ km}$}
{\True $34\text{ km}$}
\loigiai{
$83 - 49 = 34\text{ (km)}$.
\textbf{Đáp án của câu hỏi là D}
}
\end{ex}

\begin{ex} %[2D4V23]
Hiệu của số tròn chục lớn nhất có hai chữ số và số liền trước của $46$ là
\choice
{\True $45$}
{$44$}
{$46$}
{$43$}
\loigiai{
Số tròn chục lớn nhất có 2 chữ số là 90. Số liền trước 46 là 45.\nHiệu: $90 - 45 = 45$.
\textbf{Đáp án của câu hỏi là A}
}
\end{ex}

\begin{ex} %[2D4V23]
Điền dấu thích hợp: $82 - 37 \dots 71 - 25$.
\choice
{$>$}
{$=$}
{\True $<$}
{Không so sánh được}
\loigiai{
$82 - 37 = 45$; $71 - 25 = 46$. Vì $45 < 46$ nên điền dấu $<$. 
\textbf{Đáp án của câu hỏi là C}
}
\end{ex}

\begin{ex} %[2D4V23]
Hai số có tổng bằng $74$, biết số lớn là $48$. Số bé là
\choice
{$25$}
{\True $26$}
{$27$}
{$28$}
\loigiai{
Số bé: $74 - 48 = 26$.
\textbf{Đáp án của câu hỏi là B}
}
\end{ex}

\begin{ex} %[2D4V23]
Mẹ mua một tấm vải dài $60\text{ dm}$. Mẹ may váy hết $28\text{ dm}$, may áo hết $17\text{ dm}$. Tấm vải còn lại dài bao nhiêu đề-xi-mét?
\choice
{$14\text{ dm}$}
{$16\text{ dm}$}
{$13\text{ dm}$}
{\True $15\text{ dm}$}
\loigiai{
Mẹ may hết: $28 + 17 = 45\text{ (dm)}$. Còn lại: $60 - 45 = 15\text{ (dm)}$.
\textbf{Đáp án của câu hỏi là D}
}
\end{ex}

\begin{ex} %[2D4V23]
Tìm một số biết rằng lấy $84$ trừ đi số đó thì được số bé nhất có hai chữ số giống nhau.
\choice
{\True $73$}
{$72$}
{$74$}
{$71$}
\loigiai{
Số bé nhất có hai chữ số giống nhau là 11. Số đó là: $84 - 11 = 73$.
\textbf{Đáp án của câu hỏi là A}
}
\end{ex}

\textbf{Phần 2. Trắc nghiệm đúng sai}
\setcounterpageref{ex}{0}

\begin{ex} %[2D4V23]
Xét tính đúng sai của phép trừ có nhớ trong phạm vi 100:
\choiceTF
{\True Phép trừ $63 - 28$ có kết quả là $35$.}
{Phép trừ $70 - 35$ có kết quả là số chẵn.}
{\True Khi bớt ở cả số bị trừ và số trừ đi cùng một số đơn vị thì hiệu không đổi.}
{\True Hiệu của $81 - 49$ bằng hiệu của $82 - 50$.}
\loigiai{
\begin{itemize}
    \item Ý a) Đúng: $63 - 28 = 35$.
    \item Ý b) Sai: $70 - 35 = 35$, là số lẻ.
    \item Ý c) Đúng: Tính chất của phép trừ.
    \item Ý d) Đúng: Cùng thêm 1 đơn vị vào cả số bị trừ và số trừ: $81 - 49 = 82 - 50 = 32$. Đúng.
\end{itemize}
\begin{itemize}
    \item a) \textbf{Đúng}.
    \item b) \textbf{Sai}.
    \item c) \textbf{Đúng}.
    \item d) \textbf{Đúng}.
\end{itemize}
}
\end{ex}

\begin{ex} %[2D4V23]
Một kho thóc có 95 bao thóc. Ngày thứ nhất xuất 38 bao, ngày thứ hai xuất 29 bao:
\choiceTF
{\True Tổng số bao thóc đã xuất trong hai ngày là 67 bao.}
{\True Số thóc còn lại trong kho là 28 bao.}
{Ngày thứ nhất xuất ít hơn ngày thứ hai 9 bao.}
{Kho thóc còn lại nhiều hơn số thóc ngày thứ hai xuất.}
\loigiai{
\begin{itemize}
    \item Ý a) Đúng: $38 + 29 = 67$ bao.
    \item Ý b) Đúng: $95 - 67 = 28$ bao.
    \item Ý c) Sai: Ngày thứ nhất (38 bao) nhiều hơn ngày thứ hai (29 bao) là $38 - 29 = 9$ bao.
    \item Ý d) Sai: Còn lại 28 bao, ngày thứ hai xuất 29 bao nên còn lại ít hơn.
\end{itemize}
\begin{itemize}
    \item a) \textbf{Đúng}.
    \item b) \textbf{Đúng}.
    \item c) \textbf{Sai}.
    \item d) \textbf{Sai}.
\end{itemize}
}
\end{ex}

\textbf{Phần 3. Trả lời ngắn}
\setcounterpageref{ex}{0}

\begin{ex} %[2D4V23]
Tính kết quả: $83 - 47$.
\shortans{36}
\loigiai{
\textbf{Đáp số:} $36$\\
$83 - 47 = 36$.
}
\end{ex}

\begin{ex} %[2D4V23]
Một giá sách có 70 quyển sách, đã cho mượn 34 quyển. Hỏi trên giá sách còn lại bao nhiêu quyển sách?
\shortans{36}
\loigiai{
\textbf{Đáp số:} $36$\\
$70 - 34 = 36$ quyển.
}
\end{ex}

\begin{ex} %[2D4V23]
Tìm số tự nhiên y biết: $y + 38 = 91$.
\shortans{53}
\loigiai{
\textbf{Đáp số:} $53$\\
$y = 91 - 38 = 53$.
}
\end{ex}

\begin{ex} %[2D4V23]
Cuộn dây dài 92 cm, bạn Nam cắt đi một đoạn dài 56 cm. Hỏi đoạn dây còn lại dài bao nhiêu xăng-ti-mét?
\shortans{36}
\loigiai{
\textbf{Đáp số:} $36$\\
$92 - 56 = 36\text{ cm}$.
}
\end{ex}

\begin{ex} %[2D4V23]
Tìm hiệu của số lớn nhất có hai chữ số và số 67.
\shortans{32}
\loigiai{
\textbf{Đáp số:} $32$\\
$99 - 67 = 32$.
}
\end{ex}

\begin{ex} %[2D4V23]
Lớp 2A có 33 học sinh, lớp 2B có 30 học sinh. Cả hai lớp có 27 bạn nam. Hỏi cả hai lớp có bao nhiêu bạn nữ?
\shortans{36}
\loigiai{
\textbf{Đáp số:} $36$\\
Tổng số học sinh: $33 + 30 = 63$. Số bạn nữ: $63 - 27 = 36$ bạn.
}
\end{ex}

\newpage

\TaoPhieuKiemTra{Chủ đề 4 -- Phép cộng, phép trừ (có nhớ) trong phạm vi 100}{24}{Luyện tập chung}{https://loptoanthaylinh.vn/thi/chuyen-de/2/toan2_cd_b24}
\textbf{Phần 1. Trắc nghiệm lựa chọn}
\setcounterpageref{ex}{0}

\begin{ex} %[2D4V24]
Tính giá trị của biểu thức: $48 + 37 - 29$.
\choice
{$54$}
{\True $56$}
{$55$}
{$57$}
\loigiai{
$48 + 37 = 85$; $85 - 29 = 56$.
\textbf{Đáp án của câu hỏi là B}
}
\end{ex}

\begin{ex} %[2D4V24]
Một đội xe có $54$ xe tải và xe con, trong đó có $28$ xe tải. Hỏi đội xe có bao nhiêu xe con?
\choice
{\True $26$ xe con}
{$24$ xe con}
{$25$ xe con}
{$28$ xe con}
\loigiai{
$54 - 28 = 26$ (xe con).
\textbf{Đáp án của câu hỏi là A}
}
\end{ex}

\begin{ex} %[2D4V24]
Tìm $x$ biết: $x - 39 = 62 - 18$.
\choice
{$81$}
{$82$}
{\True $83$}
{$84$}
\loigiai{
Vế phải: $62 - 18 = 44$. Ta có: $x = 44 + 39 = 83$.
\textbf{Đáp án của câu hỏi là C}
}
\end{ex}

\begin{ex} %[2D4V24]
Đường gấp khúc $ABCD$ gồm $3$ đoạn thẳng: $AB = 25\text{ cm}$, $BC = 38\text{ cm}$, $CD = 29\text{ cm}$. Chiều dài đường gấp khúc $ABCD$ là
\choice
{$90\text{ cm}$}
{$91\text{ cm}$}
{$93\text{ cm}$}
{\True $92\text{ cm}$}
\loigiai{
$25 + 38 + 29 = 92\text{ (cm)}$.
\textbf{Đáp án của câu hỏi là D}
}
\end{ex}

\begin{ex} %[2D4V24]
Một cửa hàng có $92\text{ L}$ dầu, ngày thứ nhất bán được $35\text{ L}$, ngày thứ hai bán được $28\text{ L}$. Hỏi cửa hàng còn lại bao nhiêu lít dầu?
\choice
{$28\text{ L}$}
{\True $29\text{ L}$}
{$30\text{ L}$}
{$27\text{ L}$}
\loigiai{
Tổng bán: $35 + 28 = 63\text{ (L)}$. Còn lại: $92 - 63 = 29\text{ (L)}$.
\textbf{Đáp án của câu hỏi là B}
}
\end{ex}

\begin{ex} %[2D4V24]
Tìm số tự nhiên lớn nhất có hai chữ số mà khi trừ đi $36$ ta được kết quả nhỏ hơn $30$.
\choice
{\True $65$}
{$64$}
{$66$}
{$63$}
\loigiai{
Ta có: $x - 36 < 30 \Rightarrow x < 66$. Số tự nhiên lớn nhất là 65.
\textbf{Đáp án của câu hỏi là A}
}
\end{ex}

\begin{ex} %[2D4V24]
Bác An nuôi $62$ con gà và vịt, trong đó có $35$ con gà. Hỏi số vịt ít hơn số gà bao nhiêu con?
\choice
{$7$ con}
{$9$ con}
{\True $8$ con}
{$6$ con}
\loigiai{
Số vịt là: $62 - 35 = 27$ (con).\nSố vịt ít hơn số gà: $35 - 27 = 8$ (con).
\textbf{Đáp án của câu hỏi là C}
}
\end{ex}

\begin{ex} %[2D4V24]
Điền số thích hợp: $73 - \dots = 28 + 19$.
\choice
{$24$}
{$25$}
{$27$}
{\True $26$}
\loigiai{
Vế phải: $28 + 19 = 47$. Ta có: $73 - 47 = 26$.
\textbf{Đáp án của câu hỏi là D}
}
\end{ex}

\begin{ex} %[2D4V24]
Chu vi của một hình tứ giác có độ dài các cạnh là $18\text{ cm}, 24\text{ cm}, 19\text{ cm}, 25\text{ cm}$ là
\choice
{$84\text{ cm}$}
{\True $86\text{ cm}$}
{$85\text{ cm}$}
{$88\text{ cm}$}
\loigiai{
$18 + 24 + 19 + 25 = 86\text{ (cm)}$.
\textbf{Đáp án của câu hỏi là B}
}
\end{ex}

\begin{ex} %[2D4V24]
Tìm một số biết rằng lấy số đó cộng với $37$ rồi trừ đi $25$ thì được $68$.
\choice
{\True $56$}
{$54$}
{$55$}
{$58$}
\loigiai{
Tính ngược từ cuối: $68 + 25 - 37 = 93 - 37 = 56$.
\textbf{Đáp án của câu hỏi là A}
}
\end{ex}

\begin{ex} %[2D4V24]
Hai bao gạo nặng tất cả $80\text{ kg}$. Nếu chuyển $5\text{ kg}$ từ bao thứ nhất sang bao thứ hai thì hai bao nặng bằng nhau. Hỏi ban đầu bao thứ nhất nặng bao nhiêu ki-lô-gam?
\choice
{$40\text{ kg}$}
{$50\text{ kg}$}
{\True $45\text{ kg}$}
{$42\text{ kg}$}
\loigiai{
Khi bằng nhau mỗi bao nặng: $80 : 2 = 40\text{ (kg)}$.\nBan đầu bao thứ nhất nặng: $40 + 5 = 45\text{ (kg)}$.
\textbf{Đáp án của câu hỏi là C}
}
\end{ex}

\begin{ex} %[2D4V24]
Điền dấu thích hợp: $81 - 46 \dots 90 - 55$.
\choice
{$>$}
{$<$}
{Không so sánh được}
{\True $=$}
\loigiai{
$81 - 46 = 35$; $90 - 55 = 35$. Hai vế bằng nhau nên điền dấu $=$.
\textbf{Đáp án của câu hỏi là D}
}
\end{ex}

\textbf{Phần 2. Trắc nghiệm đúng sai}
\setcounterpageref{ex}{0}

\begin{ex} %[2D4V24]
Xét tính đúng sai của các bài toán rèn luyện tư duy Chủ đề 4:
\choiceTF
{\True Tổng của $49 + 36$ bằng hiệu của $95 - 10$.}
{Biểu thức $82 - 38 + 16$ có kết quả là 50.}
{\True Trong phép trừ có nhớ trong phạm vi 100, số trừ có thể lớn hơn hoặc bằng 10.}
{\True Khi cùng thêm vào số bị trừ và số trừ một số đơn vị thì hiệu vẫn giữ nguyên.}
\loigiai{
\begin{itemize}
    \item Ý a) Đúng: $49 + 36 = 85$; $95 - 10 = 85$. Bằng nhau.
    \item Ý b) Sai: $82 - 38 = 44$; $44 + 16 = 60$ (không phải 50).
    \item Ý c) Đúng: Đúng.
    \item Ý d) Đúng: Tính chất của phép trừ.
\end{itemize}
\begin{itemize}
    \item a) \textbf{Đúng}.
    \item b) \textbf{Sai}.
    \item c) \textbf{Đúng}.
    \item d) \textbf{Đúng}.
\end{itemize}
}
\end{ex}

\begin{ex} %[2D4V24]
Một cửa hàng có 80 cuộn len gồm ba màu đỏ, vàng, xanh. Biết có 28 cuộn màu đỏ và 35 cuộn màu vàng:
\choiceTF
{\True Tổng số cuộn len đỏ và vàng là 63 cuộn.}
{\True Số cuộn len màu xanh là 17 cuộn.}
{Số cuộn len màu xanh nhiều hơn số cuộn len màu đỏ.}
{Cửa hàng có số cuộn len màu vàng gấp đôi số cuộn len màu đỏ.}
\loigiai{
\begin{itemize}
    \item Ý a) Đúng: $28 + 35 = 63$ cuộn.
    \item Ý b) Đúng: $80 - 63 = 17$ cuộn.
    \item Ý c) Sai: Xanh có 17 cuộn, đỏ có 28 cuộn nên xanh ít hơn đỏ.
    \item Ý d) Sai: 35 không gấp đôi 28. Khẳng định sai.
\end{itemize}
\begin{itemize}
    \item a) \textbf{Đúng}.
    \item b) \textbf{Đúng}.
    \item c) \textbf{Sai}.
    \item d) \textbf{Sai}.
\end{itemize}
}
\end{ex}

\textbf{Phần 3. Trả lời ngắn}
\setcounterpageref{ex}{0}

\begin{ex} %[2D4V24]
Tính giá trị biểu thức: $73 - 38 + 19$.
\shortans{54}
\loigiai{
\textbf{Đáp số:} $54$\\
$73 - 38 = 35$; $35 + 19 = 54$.
}
\end{ex}

\begin{ex} %[2D4V24]
Mẹ mua 45 quả trứng, đã dùng hết 18 quả. Số trứng còn lại mẹ chia đều cho 3 người con. Hỏi mỗi người con được mấy quả trứng?
\shortans{9}
\loigiai{
\textbf{Đáp số:} $9$\\
Số trứng còn lại: $45 - 18 = 27$ quả. Mỗi người được: $27 : 3 = 9$ quả.
}
\end{ex}

\begin{ex} %[2D4V24]
Tìm số tự nhiên x biết: $84 - x = 27 + 19$.
\shortans{38}
\loigiai{
\textbf{Đáp số:} $38$\\
$27 + 19 = 46 \Rightarrow x = 84 - 46 = 38$.
}
\end{ex}

\begin{ex} %[2D4V24]
Tính chu vi hình tứ giác có độ dài các cạnh là 22 cm, 28 cm, 17 cm, 23 cm.
\shortans{90}
\loigiai{
\textbf{Đáp số:} $90$\\
Chu vi: $(22 + 28) + (17 + 23) = 50 + 40 = 90\text{ cm}$.
}
\end{ex}

\begin{ex} %[2D4V24]
Tìm một số biết rằng lấy 92 trừ đi số đó rồi cộng với 15 thì được 49.
\shortans{58}
\loigiai{
\textbf{Đáp số:} $58$\\
$92 - x = 49 - 15 = 34 \Rightarrow x = 92 - 34 = 58$.
}
\end{ex}

\begin{ex} %[2D4V24]
Hai thùng có tất cả 75 lít dầu. Thùng thứ nhất hơn thùng thứ hai 9 lít dầu. Hỏi thùng thứ nhất có bao nhiêu lít dầu?
\shortans{42}
\loigiai{
\textbf{Đáp số:} $42$\\
Thùng 1 có: $(75 + 9) : 2 = 42\text{ L}$.
}
\end{ex}

\newpage


% =========================================================================
% CHỦ ĐỀ 5 - LÀM QUEN VỚI HÌNH PHẲNG
% =========================================================================
\section{Làm quen với hình phẳng}

\TaoPhieuKiemTra{Chủ đề 5 -- Làm quen với hình phẳng}{25}{Điểm, đoạn thẳng. Đường thẳng, đường cong}{https://loptoanthaylinh.vn/thi/chuyen-de/2/toan2_cd_b25}
\textbf{Phần 1. Trắc nghiệm lựa chọn}
\setcounterpageref{ex}{0}

\begin{ex} %[2H5V25]
Cho $4$ điểm phân biệt cùng nằm trên một đường thẳng. Hỏi có tất cả bao nhiêu đoạn thẳng được tạo thành từ $4$ điểm đó?
\choice
{$4$ đoạn thẳng}
{$5$ đoạn thẳng}
{\True $6$ đoạn thẳng}
{$8$ đoạn thẳng}
\loigiai{
Gọi 4 điểm là $A, B, C, D$. Các đoạn thẳng là: $AB, AC, AD, BC, BD, CD$. Có tất cả $6$ đoạn thẳng.
\textbf{Đáp án của câu hỏi là C}
}
\end{ex}

\begin{ex} %[2H5V25]
Một đoạn dây đồng uốn thành hình tam giác có độ dài ba cạnh lần lượt là $12\text{ cm}, 15\text{ cm}, 18\text{ cm}$. Nếu duỗi thẳng đoạn dây đó ra thì chiều dài của nó là
\choice
{\True $45\text{ cm}$}
{$43\text{ cm}$}
{$44\text{ cm}$}
{$46\text{ cm}$}
\loigiai{
Chiều dài đoạn dây bằng chu vi tam giác: $12 + 15 + 18 = 45\text{ (cm)}$.
\textbf{Đáp án của câu hỏi là A}
}
\end{ex}

\begin{ex} %[2H5V25]
Trên đoạn thẳng $MN$ dài $24\text{ cm}$, người ta lấy điểm $P$ sao cho $MP = 9\text{ cm}$. Độ dài đoạn thẳng $PN$ là
\choice
{$13\text{ cm}$}
{$14\text{ cm}$}
{$16\text{ cm}$}
{\True $15\text{ cm}$}
\loigiai{
Độ dài $PN = MN - MP = 24 - 9 = 15\text{ (cm)}$.
\textbf{Đáp án của câu hỏi là D}
}
\end{ex}

\begin{ex} %[2H5V25]
Cho ba điểm $A, B, C$ thẳng hàng theo thứ tự đó. Biết $AB = 25\text{ cm}$, đoạn $BC$ ngắn hơn đoạn $AB$ là $8\text{ cm}$. Độ dài đoạn thẳng $AC$ là
\choice
{$40\text{ cm}$}
{\True $42\text{ cm}$}
{$41\text{ cm}$}
{$43\text{ cm}$}
\loigiai{
$BC = 25 - 8 = 17\text{ (cm)}$. Độ dài $AC = AB + BC = 25 + 17 = 42\text{ (cm)}$.
\textbf{Đáp án của câu hỏi là B}
}
\end{ex}

\begin{ex} %[2H5V25]
Hình nào dưới đây KHÔNG phải là đường thẳng?
\choice
{Đường kéo dài vô tận về hai phía}
{Đường thẳng đi qua hai điểm}
{\True Đoạn nối hai điểm bị giới hạn ở hai đầu}
{Đường nét thẳng không giới hạn}
\loigiai{
Đoạn thẳng bị giới hạn ở hai đầu, còn đường thẳng không bị giới hạn.
\textbf{Đáp án của câu hỏi là C}
}
\end{ex}

\begin{ex} %[2H5V25]
Cho $3$ điểm không thẳng hàng. Nối từng cặp điểm ta được hình gì?
\choice
{\True Hình tam giác}
{Hình tứ giác}
{Đường thẳng}
{Đoạn thẳng dài}
\loigiai{
Nối 3 điểm không thẳng hàng tạo thành 3 cạnh của một hình tam giác.
\textbf{Đáp án của câu hỏi là A}
}
\end{ex}

\begin{ex} %[2H5V25]
Một con kiến bò từ điểm $A$ đến điểm $B$ theo một đường cong dài $28\text{ cm}$. Một chú bọ rùa bò theo đường thẳng từ $A$ đến $B$ hết $19\text{ cm}$. Đường đi của kiến dài hơn của bọ rùa bao nhiêu xăng-ti-mét?
\choice
{$8\text{ cm}$}
{$10\text{ cm}$}
{$7\text{ cm}$}
{\True $9\text{ cm}$}
\loigiai{
Đường đi của kiến dài hơn: $28 - 19 = 9\text{ (cm)}$.
\textbf{Đáp án của câu hỏi là D}
}
\end{ex}

\begin{ex} %[2H5V25]
Có $5$ điểm phân biệt trên một đường tròn. Cứ qua hai điểm nối được một đoạn thẳng. Có tất cả bao nhiêu đoạn thẳng?
\choice
{$8$ đoạn thẳng}
{\True $10$ đoạn thẳng}
{$12$ đoạn thẳng}
{$15$ đoạn thẳng}
\loigiai{
Số đoạn thẳng: $5 \times 4 : 2 = 10$ đoạn thẳng.
\textbf{Đáp án của câu hỏi là B}
}
\end{ex}

\begin{ex} %[2H5V25]
Đoạn thẳng $CD$ dài $1\text{ dm } 6\text{ cm}$. Cắt bớt đi $7\text{ cm}$ thì đoạn thẳng còn lại dài
\choice
{$8\text{ cm}$}
{$10\text{ cm}$}
{\True $9\text{ cm}$}
{$11\text{ cm}$}
\loigiai{
Đổi $1\text{ dm } 6\text{ cm} = 16\text{ cm}$. Còn lại: $16 - 7 = 9\text{ (cm)}$.
\textbf{Đáp án của câu hỏi là C}
}
\end{ex}

\begin{ex} %[2H5V25]
Điểm $I$ là điểm nằm giữa hai điểm $A$ và $B$. Biết $AI = 12\text{ cm}$ và $IB = 12\text{ cm}$. Khẳng định nào sau đây là đúng?
\choice
{\True $I$ là điểm chính giữa của đoạn thẳng $AB$}
{$I$ không thuộc đoạn thẳng $AB$}
{$AB$ dài $12\text{ cm}$}
{$AI$ dài hơn $IB$}
\loigiai{
Vì $I$ nằm giữa và $AI = IB = 12\text{ cm}$ nên $I$ là điểm chính giữa đoạn thẳng $AB$.
\textbf{Đáp án của câu hỏi là A}
}
\end{ex}

\begin{ex} %[2H5V25]
Một hàng rào có $6$ chiếc cọc cắm thẳng hàng cách đều nhau. Khoảng cách giữa hai cọc liên tiếp là $3\text{ m}$. Chiều dài hàng rào từ cọc đầu tiên đến cọc cuối cùng là
\choice
{$18\text{ m}$}
{$12\text{ m}$}
{$16\text{ m}$}
{\True $15\text{ m}$}
\loigiai{
Có $6 - 1 = 5$ khoảng cách. Chiều dài hàng rào: $5 \times 3 = 15\text{ (m)}$.
\textbf{Đáp án của câu hỏi là D}
}
\end{ex}

\begin{ex} %[2H5V25]
Cho ba đoạn thẳng có độ dài lần lượt là $14\text{ cm}, 18\text{ cm}, 22\text{ cm}$. Tổng độ dài của ba đoạn thẳng đó là
\choice
{$52\text{ cm}$}
{\True $54\text{ cm}$}
{$53\text{ cm}$}
{$55\text{ cm}$}
\loigiai{
$14 + 18 + 22 = 54\text{ (cm)}$.
\textbf{Đáp án của câu hỏi là B}
}
\end{ex}

\textbf{Phần 2. Trắc nghiệm đúng sai}
\setcounterpageref{ex}{0}

\begin{ex} %[2H5V25]
Xét tính đúng sai của các khẳng định hình học về điểm và đường thẳng:
\choiceTF
{\True Đoạn thẳng có hai điểm mút giới hạn ở hai đầu.}
{Qua hai điểm phân biệt có thể vẽ được vô số đường thẳng.}
{\True Ba điểm cùng nằm trên một đường thẳng được gọi là ba điểm thẳng hàng.}
{\True Độ dài đoạn thẳng luôn là một số lớn hơn 0.}
\loigiai{
\begin{itemize}
    \item Ý a) Đúng: Đúng theo định nghĩa đoạn thẳng.
    \item Ý b) Sai: Qua hai điểm phân biệt chỉ vẽ được duy nhất một đường thẳng.
    \item Ý c) Đúng: Đúng theo khái niệm.
    \item Ý d) Đúng: Đúng.
\end{itemize}
\begin{itemize}
    \item a) \textbf{Đúng}.
    \item b) \textbf{Sai}.
    \item c) \textbf{Đúng}.
    \item d) \textbf{Đúng}.
\end{itemize}
}
\end{ex}

\begin{ex} %[2H5V25]
Cho đoạn thẳng AB dài 20 cm. Điểm M nằm trên đoạn thẳng AB sao cho AM dài 8 cm:
\choiceTF
{\True Điểm M nằm giữa hai điểm A và B.}
{\True Đoạn thẳng MB dài 12 cm.}
{Đoạn thẳng AM dài hơn đoạn thẳng MB.}
{Điểm M là điểm chính giữa của đoạn thẳng AB.}
\loigiai{
\begin{itemize}
    \item Ý a) Đúng: Vì M thuộc đoạn thẳng AB.
    \item Ý b) Đúng: $MB = 20 - 8 = 12\text{ cm}$.
    \item Ý c) Sai: $AM = 8\text{ cm} < MB = 12\text{ cm}$. Khẳng định sai.
    \item Ý d) Sai: Vì $AM \ne MB$ nên M không phải là điểm chính giữa.
\end{itemize}
\begin{itemize}
    \item a) \textbf{Đúng}.
    \item b) \textbf{Đúng}.
    \item c) \textbf{Sai}.
    \item d) \textbf{Sai}.
\end{itemize}
}
\end{ex}

\textbf{Phần 3. Trả lời ngắn}
\setcounterpageref{ex}{0}

\begin{ex} %[2H5V25]
Cho 5 điểm cùng nằm trên một đường thẳng. Hỏi có bao nhiêu đoạn thẳng được tạo thành?
\shortans{10}
\loigiai{
\textbf{Đáp số:} $10$\\
Số đoạn thẳng: $5 \times 4 : 2 = 10$ đoạn thẳng.
}
\end{ex}

\begin{ex} %[2H5V25]
Đoạn thẳng AB dài 34 cm. Điểm C nằm giữa A và B sao cho AC dài 18 cm. Tính độ dài đoạn thẳng CB.
\shortans{16}
\loigiai{
\textbf{Đáp số:} $16$\\
$CB = 34 - 18 = 16\text{ cm}$.
}
\end{ex}

\begin{ex} %[2H5V25]
Người ta cắm 5 cây cọc thẳng hàng cách đều nhau 4 m. Hỏi khoảng cách từ cọc đầu tiên đến cọc cuối cùng là bao nhiêu mét?
\shortans{16}
\loigiai{
\textbf{Đáp số:} $16$\\
Có $5 - 1 = 4$ khoảng cách: $4 \times 4 = 16\text{ m}$.
}
\end{ex}

\begin{ex} %[2H5V25]
Một đoạn dây thép dài 60 cm uốn thành một hình tam giác có hai cạnh dài 18 cm và 24 cm. Tính độ dài cạnh còn lại của hình tam giác.
\shortans{18}
\loigiai{
\textbf{Đáp số:} $18$\\
Cạnh còn lại: $60 - 18 - 24 = 18\text{ cm}$.
}
\end{ex}

\begin{ex} %[2H5V25]
Ba điểm A, B, C thẳng hàng theo thứ tự đó. Biết AB dài 15 cm, AC dài 32 cm. Tính độ dài đoạn thẳng BC.
\shortans{17}
\loigiai{
\textbf{Đáp số:} $17$\\
$BC = 32 - 15 = 17\text{ cm}$.
}
\end{ex}

\begin{ex} %[2H5V25]
Đoạn thẳng MN dài 45 cm, đoạn thẳng PQ ngắn hơn đoạn thẳng MN là 18 cm. Tính độ dài đoạn thẳng PQ.
\shortans{27}
\loigiai{
\textbf{Đáp số:} $27$\\
$45 - 18 = 27\text{ cm}$.
}
\end{ex}

\newpage

\TaoPhieuKiemTra{Chủ đề 5 -- Làm quen với hình phẳng}{26}{Đường gấp khúc. Hình tứ giác}{https://loptoanthaylinh.vn/thi/chuyen-de/2/toan2_cd_b26}
\textbf{Phần 1. Trắc nghiệm lựa chọn}
\setcounterpageref{ex}{0}

\begin{ex} %[2H5V26]
Đường gấp khúc gồm $3$ đoạn thẳng có độ dài lần lượt là $15\text{ cm}, 24\text{ cm}, 18\text{ cm}$. Độ dài đường gấp khúc đó là
\choice
{$55\text{ cm}$}
{\True $57\text{ cm}$}
{$56\text{ cm}$}
{$58\text{ cm}$}
\loigiai{
$15 + 24 + 18 = 57\text{ (cm)}$.
\textbf{Đáp án của câu hỏi là B}
}
\end{ex}

\begin{ex} %[2H5V26]
Một hình tứ giác có độ dài các cạnh lần lượt là $12\text{ cm}, 15\text{ cm}, 18\text{ cm}, 20\text{ cm}$. Chu vi của hình tứ giác đó là
\choice
{$63\text{ cm}$}
{$64\text{ cm}$}
{$66\text{ cm}$}
{\True $65\text{ cm}$}
\loigiai{
$12 + 15 + 18 + 20 = 65\text{ (cm)}$.
\textbf{Đáp án của câu hỏi là D}
}
\end{ex}

\begin{ex} %[2H5V26]
Một đường gấp khúc gồm $4$ đoạn thẳng dài bằng nhau, mỗi đoạn dài $6\text{ cm}$. Chiều dài đường gấp khúc đó là
\choice
{\True $24\text{ cm}$}
{$20\text{ cm}$}
{$22\text{ cm}$}
{$26\text{ cm}$}
\loigiai{
$6 \times 4 = 24\text{ (cm)}$.
\textbf{Đáp án của câu hỏi là A}
}
\end{ex}

\begin{ex} %[2H5V26]
Đường gấp khúc $ABCD$ dài $62\text{ cm}$, trong đó đoạn $AB = 18\text{ cm}$, $BC = 25\text{ cm}$. Đoạn thẳng $CD$ dài bao nhiêu?
\choice
{$17\text{ cm}$}
{$18\text{ cm}$}
{\True $19\text{ cm}$}
{$20\text{ cm}$}
\loigiai{
$CD = 62 - (18 + 25) = 62 - 43 = 19\text{ (cm)}$.
\textbf{Đáp án của câu hỏi là C}
}
\end{ex}

\begin{ex} %[2H5V26]
Hình tứ giác có bao nhiêu đỉnh và bao nhiêu cạnh?
\choice
{$3$ đỉnh và $4$ cạnh}
{\True $4$ đỉnh và $4$ cạnh}
{$4$ đỉnh và $3$ cạnh}
{$5$ đỉnh và $5$ cạnh}
\loigiai{
Hình tứ giác luôn có đúng 4 đỉnh và 4 cạnh.
\textbf{Đáp án của câu hỏi là B}
}
\end{ex}

\begin{ex} %[2H5V26]
Uốn một sợi dây thép dài $40\text{ cm}$ thành một hình tứ giác có $4$ cạnh bằng nhau. Độ dài mỗi cạnh của hình tứ giác đó là
\choice
{$8\text{ cm}$}
{$9\text{ cm}$}
{$12\text{ cm}$}
{\True $10\text{ cm}$}
\loigiai{
$40 : 4 = 10\text{ (cm)}$.
\textbf{Đáp án của câu hỏi là D}
}
\end{ex}

\begin{ex} %[2H5V26]
Đường gấp khúc $MNPQ$ gồm $3$ đoạn thẳng. Biết $MN = 1\text{ dm } 2\text{ cm}$, $NP = 15\text{ cm}$, $PQ = 23\text{ cm}$. Chiều dài đường gấp khúc $MNPQ$ là
\choice
{\True $50\text{ cm}$}
{$48\text{ cm}$}
{$49\text{ cm}$}
{$52\text{ cm}$}
\loigiai{
Đổi $1\text{ dm } 2\text{ cm} = 12\text{ cm}$. Tổng chiều dài: $12 + 15 + 23 = 50\text{ (cm)}$.
\textbf{Đáp án của câu hỏi là A}
}
\end{ex}

\begin{ex} %[2H5V26]
Hình chữ nhật có phải là một hình tứ giác không?
\choice
{Không, hình chữ nhật không phải hình tứ giác}
{Chỉ đúng khi các cạnh bằng nhau}
{\True Có, hình chữ nhật là một hình tứ giác đặc biệt}
{Chưa thể khẳng định}
\loigiai{
Hình chữ nhật có 4 cạnh và 4 góc vuông nên là một hình tứ giác đặc biệt.
\textbf{Đáp án của câu hỏi là C}
}
\end{ex}

\begin{ex} %[2H5V26]
Một hình tam giác có chu vi $36\text{ cm}$, biết độ dài hai cạnh là $12\text{ cm}$ và $15\text{ cm}$. Độ dài cạnh thứ ba là
\choice
{$8\text{ cm}$}
{\True $9\text{ cm}$}
{$10\text{ cm}$}
{$11\text{ cm}$}
\loigiai{
Cạnh thứ ba: $36 - (12 + 15) = 36 - 27 = 9\text{ (cm)}$.
\textbf{Đáp án của câu hỏi là B}
}
\end{ex}

\begin{ex} %[2H5V26]
Đường gấp khúc $ABC$ có độ dài $28\text{ cm}$. Biết đoạn thẳng $AB$ dài hơn đoạn thẳng $BC$ là $4\text{ cm}$. Độ dài đoạn thẳng $AB$ là
\choice
{$14\text{ cm}$}
{$15\text{ cm}$}
{$18\text{ cm}$}
{\True $16\text{ cm}$}
\loigiai{
Đoạn $AB = (28 + 4) : 2 = 16\text{ (cm)}$.
\textbf{Đáp án của câu hỏi là D}
}
\end{ex}

\begin{ex} %[2H5V26]
Kẻ một đoạn thẳng chia hình tứ giác thành hai hình. Hai hình đó có thể là
\choice
{\True Hai hình tam giác}
{Hai hình tròn}
{Một hình tròn và một hình tam giác}
{Hai đoạn thẳng}
\loigiai{
Nối hai đỉnh đối diện của hình tứ giác ta được hai hình tam giác.
\textbf{Đáp án của câu hỏi là A}
}
\end{ex}

\begin{ex} %[2H5V26]
Cho đường gấp khúc khép kín gồm $4$ đoạn thẳng. Hình tạo bởi đường gấp khúc đó là
\choice
{Hình tam giác}
{Hình ngũ giác}
{\True Hình tứ giác}
{Đường cong kín}
\loigiai{
Đường gấp khúc khép kín 4 đoạn thẳng tạo thành một hình tứ giác.
\textbf{Đáp án của câu hỏi là C}
}
\end{ex}

\textbf{Phần 2. Trắc nghiệm đúng sai}
\setcounterpageref{ex}{0}

\begin{ex} %[2H5V26]
Xét tính đúng sai của các nhận định về đường gấp khúc và hình tứ giác:
\choiceTF
{\True Độ dài đường gấp khúc bằng tổng độ dài các đoạn thẳng tạo nên nó.}
{Đường gấp khúc gồm 3 đoạn thẳng thì có 4 điểm thẳng hàng.}
{\True Chu vi hình tứ giác bằng tổng độ dài 4 cạnh của hình tứ giác đó.}
{Hình tam giác có thể có 4 cạnh.}
\loigiai{
\begin{itemize}
    \item Ý a) Đúng: Đúng theo định nghĩa tính độ dài đường gấp khúc.
    \item Ý b) Sai: Các điểm trong đường gấp khúc không được cùng thẳng hàng.
    \item Ý c) Đúng: Đúng.
    \item Ý d) Sai: Hình tam giác chỉ có đúng 3 cạnh.
\end{itemize}
\begin{itemize}
    \item a) \textbf{Đúng}.
    \item b) \textbf{Sai}.
    \item c) \textbf{Đúng}.
    \item d) \textbf{Sai}.
\end{itemize}
}
\end{ex}

\begin{ex} %[2H5V26]
Một hình tứ giác MNPQ có MN = 15 cm, NP = 20 cm, PQ = 25 cm, QM = 18 cm:
\choiceTF
{\True Cạnh dài nhất của tứ giác là cạnh PQ.}
{\True Chu vi hình tứ giác là 78 cm.}
{Cạnh ngắn nhất của tứ giác là cạnh QM.}
{\True Tổng hai cạnh MN và PQ bằng 40 cm.}
\loigiai{
\begin{itemize}
    \item Ý a) Đúng: $25\text{ cm}$ là dài nhất. Đúng.
    \item Ý b) Đúng: $15 + 20 + 25 + 18 = 78\text{ cm}$.
    \item Ý c) Sai: Cạnh ngắn nhất là $MN = 15\text{ cm}$ (không phải QM).
    \item Ý d) Đúng: $15 + 25 = 40\text{ cm}$. Đúng.
\end{itemize}
\begin{itemize}
    \item a) \textbf{Đúng}.
    \item b) \textbf{Đúng}.
    \item c) \textbf{Sai}.
    \item d) \textbf{Đúng}.
\end{itemize}
}
\end{ex}

\textbf{Phần 3. Trả lời ngắn}
\setcounterpageref{ex}{0}

\begin{ex} %[2H5V26]
Đường gấp khúc gồm 3 đoạn thẳng có độ dài lần lượt là 18 cm, 25 cm và 37 cm. Tính độ dài đường gấp khúc đó.
\shortans{80}
\loigiai{
\textbf{Đáp số:} $80$\\
$18 + 25 + 37 = 80\text{ cm}$.
}
\end{ex}

\begin{ex} %[2H5V26]
Một hình tứ giác có chu vi là 75 cm, biết tổng độ dài ba cạnh là 56 cm. Tính độ dài cạnh còn lại của hình tứ giác.
\shortans{19}
\loigiai{
\textbf{Đáp số:} $19$\\
$75 - 56 = 19\text{ cm}$.
}
\end{ex}

\begin{ex} %[2H5V26]
Đường gấp khúc ABCD dài 48 cm gồm 3 đoạn thẳng có độ dài bằng nhau. Tính độ dài mỗi đoạn thẳng.
\shortans{16}
\loigiai{
\textbf{Đáp số:} $16$\\
$48 : 3 = 16\text{ cm}$.
}
\end{ex}

\begin{ex} %[2H5V26]
Tính chu vi hình tứ giác có độ dài các cạnh là 14 cm, 16 cm, 24 cm, 26 cm.
\shortans{80}
\loigiai{
\textbf{Đáp số:} $80$\\
$(14 + 16) + (24 + 26) = 30 + 50 = 80\text{ cm}$.
}
\end{ex}

\begin{ex} %[2H5V26]
Một đoạn dây dài 36 cm được uốn thành một hình tam giác có 3 cạnh bằng nhau. Hỏi mỗi cạnh dài bao nhiêu xăng-ti-mét?
\shortans{12}
\loigiai{
\textbf{Đáp số:} $12$\\
$36 : 3 = 12\text{ cm}$.
}
\end{ex}

\begin{ex} %[2H5V26]
Đường gấp khúc MNP gồm hai đoạn thẳng MN và NP có tổng độ dài là 42 cm. Đoạn thẳng MN dài 25 cm. Tính độ dài đoạn thẳng NP.
\shortans{17}
\loigiai{
\textbf{Đáp số:} $17$\\
$42 - 25 = 17\text{ cm}$.
}
\end{ex}

\newpage

\TaoPhieuKiemTra{Chủ đề 5 -- Làm quen với hình phẳng}{27}{Thực hành gấp, cắt, ghép, xếp hình và vẽ đoạn thẳng}{https://loptoanthaylinh.vn/thi/chuyen-de/2/toan2_cd_b27}
\textbf{Phần 1. Trắc nghiệm lựa chọn}
\setcounterpageref{ex}{0}

\begin{ex} %[2H5V27]
Từ một tờ giấy hình vuông, gấp đôi tờ giấy theo đường chéo ta được hình gì?
\choice
{\True Hình tam giác}
{Hình chữ nhật}
{Hình tứ giác}
{Hình tròn}
\loigiai{
Gấp đôi hình vuông theo đường chéo tạo thành hình tam giác vuông cân.
\textbf{Đáp án của câu hỏi là A}
}
\end{ex}

\begin{ex} %[2H5V27]
Cắt một hình vuông theo một đường thẳng nối hai điểm chính giữa của hai cạnh đối diện, ta được hai hình gì?
\choice
{Hai hình vuông}
{Hai hình tam giác}
{\True Hai hình chữ nhật bằng nhau}
{Hai hình tròn}
\loigiai{
Cắt đôi hình vuông theo trục đối xứng song song với cạnh ta được hai hình chữ nhật bằng nhau.
\textbf{Đáp án của câu hỏi là C}
}
\end{ex}

\begin{ex} %[2H5V27]
Cho hình chữ nhật, nếu cắt đi một góc của hình chữ nhật thì hình còn lại có thể có bao nhiêu cạnh?
\choice
{$3$ cạnh}
{\True $5$ cạnh (hình ngũ giác)}
{$2$ cạnh}
{$6$ cạnh}
\loigiai{
Cắt vát một góc của hình chữ nhật làm xuất hiện thêm 1 cạnh mới, tạo thành hình 5 cạnh (ngũ giác).
\textbf{Đáp án của câu hỏi là B}
}
\end{ex}

\begin{ex} %[2H5V27]
Ghép hai hình tam giác bằng nhau có thể tạo thành hình nào dưới đây?
\choice
{Hình tròn}
{Hình lục giác}
{Đường cong kín}
{\True Hình tứ giác}
\loigiai{
Ghép hai tam giác cùng chung một cạnh tạo thành một hình tứ giác.
\textbf{Đáp án của câu hỏi là D}
}
\end{ex}

\begin{ex} %[2H5V27]
Cho ba điểm $A, B, C$. Dùng thước kẻ đặt qua ba điểm thấy mép thước trùng khít với cả ba điểm. Ta kết luận
\choice
{\True Ba điểm $A, B, C$ thẳng hàng}
{Ba điểm tạo thành tam giác}
{Ba điểm không thẳng hàng}
{$B$ là điểm chính giữa}
\loigiai{
Mép thước thẳng đi qua cả ba điểm chứng tỏ ba điểm đó thẳng hàng.
\textbf{Đáp án của câu hỏi là A}
}
\end{ex}

\begin{ex} %[2H5V27]
Người ta xếp $4$ hình tam giác nhỏ bằng nhau thành một hình tam giác lớn. Số hình tam giác có trong hình lớn là
\choice
{$4$ hình tam giác}
{$6$ hình tam giác}
{\True $5$ hình tam giác}
{$8$ hình tam giác}
\loigiai{
Gồm 4 hình tam giác nhỏ và 1 hình tam giác lớn bao quanh: $4 + 1 = 5$ hình tam giác.
\textbf{Đáp án của câu hỏi là C}
}
\end{ex}

\begin{ex} %[2H5V27]
Dùng $6$ que tính dài bằng nhau xếp được nhiều nhất bao nhiêu hình tam giác rời nhau?
\choice
{$3$ hình tam giác}
{\True $2$ hình tam giác}
{$1$ hình tam giác}
{$4$ hình tam giác}
\loigiai{
Mỗi hình tam giác rời nhau cần 3 que tính: $6 : 3 = 2$ hình tam giác.
\textbf{Đáp án của câu hỏi là B}
}
\end{ex}

\begin{ex} %[2H5V27]
Để kiểm tra xem ba cái cọc có được cắm thẳng hàng hay không, người ta thường làm thế nào?
\choice
{Đo khoảng cách giữa các cọc}
{Dùng dây quấn quanh các cọc}
{Nhổ cọc lên kiểm tra}
{\True Đứng ở một đầu ngắm sao cho cọc thứ nhất che khuất cọc thứ hai và cọc thứ ba}
\loigiai{
Phương pháp ngắm thẳng hàng trong thực tế.
\textbf{Đáp án của câu hỏi là D}
}
\end{ex}

\begin{ex} %[2H5V27]
Một hình vuông được chia thành $4$ ô vuông nhỏ bằng nhau. Nếu tô màu $2$ ô vuông thì còn lại mấy ô vuông chưa tô màu?
\choice
{\True $2$ ô vuông}
{$1$ ô vuông}
{$3$ ô vuông}
{$4$ ô vuông}
\loigiai{
$4 - 2 = 2$ ô vuông.
\textbf{Đáp án của câu hỏi là A}
}
\end{ex}

\begin{ex} %[2H5V27]
Cắt một tờ giấy hình tam giác bằng một nhát kéo thẳng cắt ngang qua hai cạnh. Ta nhận được hai hình là
\choice
{Hai hình tam giác}
{Hai hình tứ giác}
{\True Một hình tam giác và một hình tứ giác}
{Một hình chữ nhật và một hình vuông}
\loigiai{
Phần trên là hình tam giác, phần dưới là hình tứ giác (hình thang).
\textbf{Đáp án của câu hỏi là C}
}
\end{ex}

\begin{ex} %[2H5V27]
Có bao nhiêu hình tứ giác trong một hình chữ nhật bị kẻ một đường thẳng chia đôi từ trên xuống dưới?
\choice
{$2$ hình tứ giác}
{\True $3$ hình tứ giác}
{$4$ hình tứ giác}
{$1$ hình tứ giác}
\loigiai{
Gồm 2 hình chữ nhật nhỏ và 1 hình chữ nhật lớn bao quanh: có tất cả 3 hình tứ giác.
\textbf{Đáp án của câu hỏi là B}
}
\end{ex}

\begin{ex} %[2H5V27]
Gấp một tờ giấy hình tròn làm tư rồi cắt theo đường cong viền. Khi mở ra ta nhận được
\choice
{Một đoạn thẳng}
{Một hình tam giác}
{Một đường gấp khúc}
{\True Hình có hoa văn đối xứng}
\loigiai{
Gấp và cắt giấy hình đối xứng.
\textbf{Đáp án của câu hỏi là D}
}
\end{ex}

\textbf{Phần 2. Trắc nghiệm đúng sai}
\setcounterpageref{ex}{0}

\begin{ex} %[2H5V27]
Xét tính đúng sai của các bài toán thực hành cắt ghép hình:
\choiceTF
{\True Hai hình tam giác bằng nhau luôn có thể ghép thành một hình tứ giác.}
{Một hình vuông cắt thành hai nửa thì hai nửa đó bắt buộc phải là hình tam giác.}
{\True Ba điểm A, B, C thẳng hàng nếu chúng cùng nằm trên một đường thẳng.}
{Dùng 4 que tính bằng nhau luôn ghép được một hình tam giác.}
\loigiai{
\begin{itemize}
    \item Ý a) Đúng: Ghép dọc theo một cạnh bằng nhau tạo thành hình tứ giác. Đúng.
    \item Ý b) Sai: Có thể cắt thành hai hình chữ nhật.
    \item Ý c) Đúng: Đúng theo định nghĩa.
    \item Ý d) Sai: Hình tam giác cần 3 que tính, 4 que tính ghép được hình vuông/tứ giác.
\end{itemize}
\begin{itemize}
    \item a) \textbf{Đúng}.
    \item b) \textbf{Sai}.
    \item c) \textbf{Đúng}.
    \item d) \textbf{Sai}.
\end{itemize}
}
\end{ex}

\begin{ex} %[2H5V27]
Người ta dùng các mảnh bìa hình vuông cạnh 1 cm để ghép thành hình chữ nhật dài 4 cm và rộng 2 cm:
\choiceTF
{\True Cần tất cả 8 mảnh bìa hình vuông nhỏ.}
{\True Chu vi hình chữ nhật ghép được là 12 cm.}
{Diện tích của hình chữ nhật là 12 cm vuông.}
{\True Nếu bớt đi 2 mảnh bìa ở một góc thì hình còn lại có 6 cạnh.}
\loigiai{
\begin{itemize}
    \item Ý a) Đúng: $4 \times 2 = 8$ mảnh bìa.
    \item Ý b) Đúng: $(4 + 2) \times 2 = 12\text{ cm}$.
    \item Ý c) Sai: Diện tích là $4 \times 2 = 8\text{ cm}^2$.
    \item Ý d) Đúng: Hình chữ L có 6 cạnh. Đúng.
\end{itemize}
\begin{itemize}
    \item a) \textbf{Đúng}.
    \item b) \textbf{Đúng}.
    \item c) \textbf{Sai}.
    \item d) \textbf{Đúng}.
\end{itemize}
}
\end{ex}

\textbf{Phần 3. Trả lời ngắn}
\setcounterpageref{ex}{0}

\begin{ex} %[2H5V27]
Có bao nhiêu hình tam giác có trong một hình tam giác lớn bị kẻ một đường thẳng từ đỉnh xuống đáy?
\shortans{3}
\loigiai{
\textbf{Đáp số:} $3$\\
Gồm 2 hình tam giác nhỏ và 1 hình tam giác lớn: $2 + 1 = 3$ hình.
}
\end{ex}

\begin{ex} %[2H5V27]
Cần ít nhất bao nhiêu que tính bằng nhau để xếp được 2 hình tam giác dính liền có chung 1 cạnh?
\shortans{5}
\loigiai{
\textbf{Đáp số:} $5$\\
Tam giác thứ nhất cần 3 que, tam giác thứ hai dùng chung 1 que nên cần thêm 2 que: $3 + 2 = 5$ que tính.
}
\end{ex}

\begin{ex} %[2H5V27]
Trong hình chữ nhật có kẻ hai đường chéo cắt nhau tại tâm. Hỏi có tất cả bao nhiêu hình tam giác?
\shortans{8}
\loigiai{
\textbf{Đáp số:} $8$\\
4 tam giác nhỏ đơn và 4 tam giác ghép đôi: $4 + 4 = 8$ hình tam giác.
}
\end{ex}

\begin{ex} %[2H5V27]
Cho 4 điểm trong đó không có 3 điểm nào thẳng hàng. Nối từng cặp điểm vẽ được bao nhiêu đoạn thẳng?
\shortans{6}
\loigiai{
\textbf{Đáp số:} $6$\\
Số đoạn thẳng: $4 \times 3 : 2 = 6$ đoạn thẳng.
}
\end{ex}

\begin{ex} %[2H5V27]
Một hình vuông có cạnh 10 cm. Chu vi của hình vuông đó bằng bao nhiêu xăng-ti-mét?
\shortans{40}
\loigiai{
\textbf{Đáp số:} $40$\\
$10 \times 4 = 40\text{ cm}$.
}
\end{ex}

\begin{ex} %[2H5V27]
Dùng 8 que tính bằng nhau để xếp thành một hình vuông. Hỏi mỗi cạnh của hình vuông có mấy que tính?
\shortans{2}
\loigiai{
\textbf{Đáp số:} $2$\\
$8 : 4 = 2$ que tính.
}
\end{ex}

\newpage

\TaoPhieuKiemTra{Chủ đề 5 -- Làm quen với hình phẳng}{28}{Luyện tập chung}{https://loptoanthaylinh.vn/thi/chuyen-de/2/toan2_cd_b28}
\textbf{Phần 1. Trắc nghiệm lựa chọn}
\setcounterpageref{ex}{0}

\begin{ex} %[2H5V28]
Một đường gấp khúc gồm $3$ đoạn thẳng có độ dài lần lượt là $25\text{ cm}, 38\text{ cm}, 27\text{ cm}$. Chiều dài đường gấp khúc đó là
\choice
{$88\text{ cm}$}
{$89\text{ cm}$}
{$91\text{ cm}$}
{\True $90\text{ cm}$}
\loigiai{
$25 + 38 + 27 = 90\text{ (cm)}$.
\textbf{Đáp án của câu hỏi là D}
}
\end{ex}

\begin{ex} %[2H5V28]
Một hình tam giác có chu vi $48\text{ cm}$, biết tổng độ dài hai cạnh là $31\text{ cm}$. Độ dài cạnh còn lại là
\choice
{$16\text{ cm}$}
{\True $17\text{ cm}$}
{$18\text{ cm}$}
{$15\text{ cm}$}
\loigiai{
$48 - 31 = 17\text{ (cm)}$.
\textbf{Đáp án của câu hỏi là B}
}
\end{ex}

\begin{ex} %[2H5V28]
Trong các hình dưới đây, hình nào có số cạnh nhiều nhất?
\choice
{Hình tam giác}
{Đoạn thẳng}
{\True Hình tứ giác}
{Đường gấp khúc có $2$ đoạn thẳng}
\loigiai{
Hình tứ giác có 4 cạnh, nhiều nhất trong các phương án.
\textbf{Đáp án của câu hỏi là C}
}
\end{ex}

\begin{ex} %[2H5V28]
Đoạn thẳng $AB$ dài $42\text{ cm}$. Điểm $M$ nằm giữa $A$ và $B$ sao cho $AM = 18\text{ cm}$. Đoạn thẳng $MB$ dài bao nhiêu?
\choice
{\True $24\text{ cm}$}
{$22\text{ cm}$}
{$23\text{ cm}$}
{$25\text{ cm}$}
\loigiai{
$MB = 42 - 18 = 24\text{ (cm)}$.
\textbf{Đáp án của câu hỏi là A}
}
\end{ex}

\begin{ex} %[2H5V28]
Cho hình vuông có cạnh $15\text{ cm}$. Chu vi của hình vuông đó là
\choice
{$50\text{ cm}$}
{$55\text{ cm}$}
{$65\text{ cm}$}
{\True $60\text{ cm}$}
\loigiai{
Chu vi: $15 + 15 + 15 + 15 = 60\text{ (cm)}$.
\textbf{Đáp án của câu hỏi là D}
}
\end{ex}

\begin{ex} %[2H5V28]
Một con kiến bò quanh mép một cái bánh hình tam giác có các cạnh lần lượt là $8\text{ cm}, 12\text{ cm}, 15\text{ cm}$. Quãng đường kiến bò được là
\choice
{$33\text{ cm}$}
{\True $35\text{ cm}$}
{$34\text{ cm}$}
{$36\text{ cm}$}
\loigiai{
Quãng đường: $8 + 12 + 15 = 35\text{ (cm)}$.
\textbf{Đáp án của câu hỏi là B}
}
\end{ex}

\begin{ex} %[2H5V28]
Một đoạn dây đồng dài $72\text{ cm}$ được uốn thành một hình tứ giác có $3$ cạnh lần lượt dài $16\text{ cm}, 19\text{ cm}, 21\text{ cm}$. Cạnh thứ tư dài bao nhiêu?
\choice
{$15\text{ cm}$}
{$17\text{ cm}$}
{\True $16\text{ cm}$}
{$18\text{ cm}$}
\loigiai{
Cạnh thứ tư: $72 - (16 + 19 + 21) = 72 - 56 = 16\text{ (cm)}$.
\textbf{Đáp án của câu hỏi là C}
}
\end{ex}

\begin{ex} %[2H5V28]
Hình gồm $5$ điểm phân biệt thẳng hàng có bao nhiêu đoạn thẳng?
\choice
{\True $10$ đoạn thẳng}
{$8$ đoạn thẳng}
{$12$ đoạn thẳng}
{$6$ đoạn thẳng}
\loigiai{
$5 \times 4 : 2 = 10$ đoạn thẳng.
\textbf{Đáp án của câu hỏi là A}
}
\end{ex}

\begin{ex} %[2H5V28]
Điền số thích hợp: Đường gấp khúc dài $1\text{ m}$ gồm $2$ đoạn thẳng. Đoạn thứ nhất dài $65\text{ cm}$, đoạn thứ hai dài
\choice
{$45\text{ cm}$}
{$25\text{ cm}$}
{$30\text{ cm}$}
{\True $35\text{ cm}$}
\loigiai{
Đổi $1\text{ m} = 100\text{ cm}$. Đoạn thứ hai: $100 - 65 = 35\text{ (cm)}$.
\textbf{Đáp án của câu hỏi là D}
}
\end{ex}

\begin{ex} %[2H5V28]
Hình tứ giác $ABCD$ có chu vi $84\text{ cm}$. Nếu tăng độ dài mỗi cạnh thêm $2\text{ cm}$ thì chu vi của hình tứ giác mới là
\choice
{$86\text{ cm}$}
{\True $92\text{ cm}$}
{$88\text{ cm}$}
{$90\text{ cm}$}
\loigiai{
Có 4 cạnh, mỗi cạnh tăng 2 cm nên chu vi tăng: $2 \times 4 = 8\text{ cm}$.\nChu vi mới: $84 + 8 = 92\text{ (cm)}$.
\textbf{Đáp án của câu hỏi là B}
}
\end{ex}

\begin{ex} %[2H5V28]
Một cái khung tranh hình chữ nhật có chiều dài $30\text{ cm}$, chiều rộng $20\text{ cm}$. Chu vi khung tranh đó là
\choice
{$80\text{ cm}$}
{$90\text{ cm}$}
{\True $100\text{ cm}$}
{$50\text{ cm}$}
\loigiai{
Chu vi: $30 + 20 + 30 + 20 = 100\text{ (cm)}$.
\textbf{Đáp án của câu hỏi là C}
}
\end{ex}

\begin{ex} %[2H5V28]
Hình tam giác có độ dài các cạnh $a, b, c$. Biết $a = 15\text{ cm}, b = 20\text{ cm}$, chu vi tam giác bằng $50\text{ cm}$. Độ dài cạnh $c$ là
\choice
{\True $15\text{ cm}$}
{$16\text{ cm}$}
{$14\text{ cm}$}
{$18\text{ cm}$}
\loigiai{
$c = 50 - (15 + 20) = 15\text{ (cm)}$.
\textbf{Đáp án của câu hỏi là A}
}
\end{ex}

\textbf{Phần 2. Trắc nghiệm đúng sai}
\setcounterpageref{ex}{0}

\begin{ex} %[2H5V28]
Xét tính đúng sai của các bài toán hình học ôn tập Chủ đề 5:
\choiceTF
{\True Hình chữ nhật có 4 góc vuông và 4 cạnh.}
{Một đường gấp khúc có 3 đoạn thẳng thì có đúng 3 điểm.}
{\True Độ dài đường gấp khúc bằng tổng độ dài các đoạn thẳng thành phần.}
{\True Một hình vuông có chu vi 40 cm thì cạnh của nó dài 10 cm.}
\loigiai{
\begin{itemize}
    \item Ý a) Đúng: Đúng theo đặc điểm hình chữ nhật.
    \item Ý b) Sai: Đường gấp khúc 3 đoạn thẳng có 4 điểm đầu mút.
    \item Ý c) Đúng: Đúng.
    \item Ý d) Đúng: $40 : 4 = 10\text{ cm}$. Đúng.
\end{itemize}
\begin{itemize}
    \item a) \textbf{Đúng}.
    \item b) \textbf{Sai}.
    \item c) \textbf{Đúng}.
    \item d) \textbf{Đúng}.
\end{itemize}
}
\end{ex}

\begin{ex} %[2H5V28]
Cho đường gấp khúc ABCD có AB = 15 cm, BC = 25 cm, CD = 35 cm:
\choiceTF
{\True Chiều dài đường gấp khúc ABCD là 75 cm.}
{\True Đoạn CD dài hơn đoạn AB là 20 cm.}
{Đoạn BC dài gấp đôi đoạn AB.}
{Nối điểm A với điểm D thì ta được một hình tam giác.}
\loigiai{
\begin{itemize}
    \item Ý a) Đúng: $15 + 25 + 35 = 75\text{ cm}$.
    \item Ý b) Đúng: $35 - 15 = 20\text{ cm}$.
    \item Ý c) Sai: 25 không gấp đôi 15. Sai.
    \item Ý d) Sai: Nối A với D tạo thành hình tứ giác ABCD.
\end{itemize}
\begin{itemize}
    \item a) \textbf{Đúng}.
    \item b) \textbf{Đúng}.
    \item c) \textbf{Sai}.
    \item d) \textbf{Sai}.
\end{itemize}
}
\end{ex}

\textbf{Phần 3. Trả lời ngắn}
\setcounterpageref{ex}{0}

\begin{ex} %[2H5V28]
Tính chu vi của hình tam giác có độ dài ba cạnh lần lượt là 19 cm, 28 cm, 33 cm.
\shortans{80}
\loigiai{
\textbf{Đáp số:} $80$\\
$19 + 28 + 33 = 80\text{ cm}$.
}
\end{ex}

\begin{ex} %[2H5V28]
Một hình tứ giác có chu vi 90 cm. Tổng độ dài ba cạnh đầu là 68 cm. Tính độ dài cạnh thứ tư.
\shortans{22}
\loigiai{
\textbf{Đáp số:} $22$\\
$90 - 68 = 22\text{ cm}$.
}
\end{ex}

\begin{ex} %[2H5V28]
Đường gấp khúc MNPQ dài 64 cm. Biết MN = 18 cm, NP = 24 cm. Tính độ dài đoạn thẳng PQ.
\shortans{22}
\loigiai{
\textbf{Đáp số:} $22$\\
$PQ = 64 - (18 + 24) = 64 - 42 = 22\text{ cm}$.
}
\end{ex}

\begin{ex} %[2H5V28]
Cho hình vuông có chu vi bằng 36 cm. Tính độ dài một cạnh của hình vuông đó.
\shortans{9}
\loigiai{
\textbf{Đáp số:} $9$\\
$36 : 4 = 9\text{ cm}$.
}
\end{ex}

\begin{ex} %[2H5V28]
Một đoạn dây thép dài 85 cm, người ta cắt đi 38 cm để uốn khung tranh. Hỏi đoạn dây thép còn lại dài bao nhiêu xăng-ti-mét?
\shortans{47}
\loigiai{
\textbf{Đáp số:} $47$\\
$85 - 38 = 47\text{ cm}$.
}
\end{ex}

\begin{ex} %[2H5V28]
Một lối đi quanh vườn hình chữ nhật có chiều dài 25 m, chiều rộng 15 m. Tính chu vi lối đi quanh vườn đó.
\shortans{80}
\loigiai{
\textbf{Đáp số:} $80$\\
Chu vi: $25 + 15 + 25 + 15 = 80\text{ m}$.
}
\end{ex}

\newpage


% =========================================================================
% CHỦ ĐỀ 6 - NGÀY – GIỜ, GIỜ – PHÚT, NGÀY – THÁNG
% =========================================================================
\section{Ngày – giờ, giờ – phút, ngày – tháng}

\TaoPhieuKiemTra{Chủ đề 6 -- Ngày – giờ, giờ – phút, ngày – tháng}{29}{Ngày – giờ, giờ – phút}{https://loptoanthaylinh.vn/thi/chuyen-de/2/toan2_cd_b29}
\textbf{Phần 1. Trắc nghiệm lựa chọn}
\setcounterpageref{ex}{0}

\begin{ex} %[2H6V29]
Một ngày có bao nhiêu giờ?
\choice
{$12$ giờ}
{$20$ giờ}
{\True $24$ giờ}
{$48$ giờ}
\loigiai{
Một ngày đêm có đúng 24 giờ.
\textbf{Đáp án của câu hỏi là C}
}
\end{ex}

\begin{ex} %[2H6V29]
Kim phút chỉ vào số $6$, kim giờ chỉ vào khoảng giữa số $8$ và số $9$. Đồng hồ chỉ mấy giờ?
\choice
{$9$ giờ $30$ phút}
{\True $8$ giờ $30$ phút (hoặc $8$ giờ rưỡi)}
{$6$ giờ $40$ phút}
{$8$ giờ $6$ phút}
\loigiai{
Kim phút chỉ số 6 là 30 phút, kim giờ giữa số 8 và 9 là 8 giờ 30 phút.
\textbf{Đáp án của câu hỏi là B}
}
\end{ex}

\begin{ex} %[2H6V29]
Bố đi làm lúc $7$ giờ sáng và về đến nhà lúc $5$ giờ chiều. Hỏi bố đi làm trong bao nhiêu giờ?
\choice
{$8$ giờ}
{$9$ giờ}
{$12$ giờ}
{\True $10$ giờ}
\loigiai{
$5$ giờ chiều là $17$ giờ. Thời gian bố đi làm: $17 - 7 = 10$ (giờ).
\textbf{Đáp án của câu hỏi là D}
}
\end{ex}

\begin{ex} %[2H6V29]
Một bộ phim bắt đầu chiếu lúc $8$ giờ tối và kết thúc lúc $9$ giờ $30$ phút tối. Bộ phim kéo dài trong bao lâu?
\choice
{\True $1$ giờ $30$ phút}
{$1$ giờ}
{$2$ giờ}
{$30$ phút}
\loigiai{
Thời gian xem phim: $9$ giờ $30$ phút $- 8$ giờ $= 1$ giờ $30$ phút.
\textbf{Đáp án của câu hỏi là A}
}
\end{ex}

\begin{ex} %[2H6V29]
Đồng hồ chỉ $10$ giờ đêm, tức là mấy giờ trong ngày?
\choice
{$20$ giờ}
{$21$ giờ}
{\True $22$ giờ}
{$23$ giờ}
\loigiai{
$10$ giờ đêm $= 10 + 12 = 22$ giờ.
\textbf{Đáp án của câu hỏi là C}
}
\end{ex}

\begin{ex} %[2H6V29]
Mai bắt đầu học bài lúc $19$ giờ và học xong lúc $20$ giờ $30$ phút. Thời gian Mai học bài là
\choice
{$1$ giờ}
{\True $1$ giờ $30$ phút}
{$2$ giờ}
{$45$ phút}
\loigiai{
$20$ giờ $30$ phút $- 19$ giờ $= 1$ giờ $30$ phút.
\textbf{Đáp án của câu hỏi là B}
}
\end{ex}

\begin{ex} %[2H6V29]
Hiện tại đồng hồ chỉ $3$ giờ chiều. Sau $4$ giờ nữa thì đồng hồ chỉ mấy giờ?
\choice
{$18$ giờ}
{$8$ giờ tối}
{$6$ giờ chiều}
{\True $19$ giờ (hoặc $7$ giờ tối)}
\loigiai{
$3$ giờ chiều là $15$ giờ. Sau 4 giờ: $15 + 4 = 19$ giờ (tức 7 giờ tối).
\textbf{Đáp án của câu hỏi là D}
}
\end{ex}

\begin{ex} %[2H6V29]
Một tiết học bắt đầu lúc $8$ giờ sáng và kéo dài $35$ phút. Tiết học kết thúc lúc mấy giờ?
\choice
{\True $8$ giờ $35$ phút}
{$8$ giờ $30$ phút}
{$8$ giờ $45$ phút}
{$9$ giờ}
\loigiai{
$8$ giờ $+ 35$ phút $= 8$ giờ $35$ phút.
\textbf{Đáp án của câu hỏi là A}
}
\end{ex}

\begin{ex} %[2H6V29]
Trong một ngày, kim giờ quay được bao nhiêu vòng quanh mặt đồng hồ?
\choice
{$1$ vòng}
{$24$ vòng}
{\True $2$ vòng}
{$12$ vòng}
\loigiai{
Mỗi vòng kim giờ quay hết 12 giờ. Một ngày 24 giờ kim giờ quay: $24 : 12 = 2$ vòng.
\textbf{Đáp án của câu hỏi là C}
}
\end{ex}

\begin{ex} %[2H6V29]
Huy đi ngủ lúc $21$ giờ và thức dậy lúc $6$ giờ sáng hôm sau. Hỏi Huy đã ngủ bao nhiêu giờ?
\choice
{$8$ giờ}
{\True $9$ giờ}
{$10$ giờ}
{$7$ giờ}
\loigiai{
Từ $21$ giờ đến $24$ giờ là $3$ giờ. Từ $0$ giờ đến $6$ giờ là $6$ giờ. Tổng cộng: $3 + 6 = 9$ giờ.
\textbf{Đáp án của câu hỏi là B}
}
\end{ex}

\begin{ex} %[2H6V29]
Xe buýt khởi hành lúc $14$ giờ và đến nơi sau $2$ giờ $30$ phút. Xe buýt đến nơi lúc mấy giờ?
\choice
{$16$ giờ}
{$17$ giờ}
{$16$ giờ $15$ phút}
{\True $16$ giờ $30$ phút}
\loigiai{
$14$ giờ $+ 2$ giờ $30$ phút $= 16$ giờ $30$ phút.
\textbf{Đáp án của câu hỏi là D}
}
\end{ex}

\begin{ex} %[2H6V29]
Thời gian từ $7$ giờ sáng đến $11$ giờ trưa là bao nhiêu giờ?
\choice
{\True $4$ giờ}
{$3$ giờ}
{$5$ giờ}
{$6$ giờ}
\loigiai{
$11 - 7 = 4$ (giờ).
\textbf{Đáp án của câu hỏi là A}
}
\end{ex}

\textbf{Phần 2. Trắc nghiệm đúng sai}
\setcounterpageref{ex}{0}

\begin{ex} %[2H6V29]
Xét tính đúng sai của các nhận định về thời gian (giờ và phút):
\choiceTF
{\True Một ngày có 24 giờ, tính từ 12 giờ đêm hôm trước đến 12 giờ đêm hôm sau.}
{15 giờ chiều tương ứng với 5 giờ chiều.}
{\True Khi kim phút quay được 1 vòng quanh mặt đồng hồ thì đúng 1 giờ trôi qua.}
{Từ 9 giờ sáng đến 12 giờ trưa là 4 giờ.}
\loigiai{
\begin{itemize}
    \item Ý a) Đúng: Đúng.
    \item Ý b) Sai: 15 giờ là 3 giờ chiều ($15 - 12 = 3$). Khẳng định sai.
    \item Ý c) Đúng: Kim phút quay 1 vòng là 60 phút = 1 giờ.
    \item Ý d) Sai: Chỉ có $12 - 9 = 3$ giờ.
\end{itemize}
\begin{itemize}
    \item a) \textbf{Đúng}.
    \item b) \textbf{Sai}.
    \item c) \textbf{Đúng}.
    \item d) \textbf{Sai}.
\end{itemize}
}
\end{ex}

\begin{ex} %[2H6V29]
Lịch sinh hoạt buổi chiều của bạn Lan:
\choiceTF
{\True Lan bắt đầu tập thể dục lúc 17 giờ tức là 5 giờ chiều.}
{\True Lan ăn cơm tối lúc 18 giờ 30 phút, tức là 6 giờ rưỡi tối.}
{Lan xem thời sự từ 19 giờ đến 20 giờ là 2 giờ.}
{\True Lan đi ngủ lúc 22 giờ tức là 10 giờ đêm.}
\loigiai{
\begin{itemize}
    \item Ý a) Đúng: $17 - 12 = 5$ giờ chiều. Đúng.
    \item Ý b) Đúng: Đúng.
    \item Ý c) Sai: $20 - 19 = 1$ giờ (không phải 2 giờ).
    \item Ý d) Đúng: $22 - 12 = 10$ giờ đêm. Đúng.
\end{itemize}
\begin{itemize}
    \item a) \textbf{Đúng}.
    \item b) \textbf{Đúng}.
    \item c) \textbf{Sai}.
    \item d) \textbf{Đúng}.
\end{itemize}
}
\end{ex}

\textbf{Phần 3. Trả lời ngắn}
\setcounterpageref{ex}{0}

\begin{ex} %[2H6V29]
Một trận bóng đá bắt đầu lúc 15 giờ và kết thúc lúc 17 giờ. Trận bóng đá diễn ra trong bao nhiêu giờ?
\shortans{2}
\loigiai{
\textbf{Đáp số:} $2$\\
$17 - 15 = 2$ giờ.
}
\end{ex}

\begin{ex} %[2H6V29]
Bác lái xe đi từ Hà Nội lúc 6 giờ sáng và đến Hải Phòng lúc 8 giờ sáng. Bác đi hết bao nhiêu giờ?
\shortans{2}
\loigiai{
\textbf{Đáp số:} $2$\\
$8 - 6 = 2$ giờ.
}
\end{ex}

\begin{ex} %[2H6V29]
Đồng hồ chỉ 20 giờ, tức là mấy giờ tối?
\shortans{8}
\loigiai{
\textbf{Đáp số:} $8$\\
$20 - 12 = 8$ giờ tối.
}
\end{ex}

\begin{ex} %[2H6V29]
Nam tập bơi từ 8 giờ sáng đến 10 giờ sáng. Mỗi giờ bơi Nam nghỉ 15 phút. Hỏi Nam đã tập bơi thực tế trong bao nhiêu phút?
\shortans{90}
\loigiai{
\textbf{Đáp số:} $90$\\
2 giờ = 120 phút. Nghỉ: $15 \times 2 = 30$ phút. Bơi thực tế: $120 - 30 = 90$ phút.
}
\end{ex}

\begin{ex} %[2H6V29]
Một ngày có 24 giờ. Hỏi 2 ngày có tất cả bao nhiêu giờ?
\shortans{48}
\loigiai{
\textbf{Đáp số:} $48$\\
$24 \times 2 = 48$ giờ.
}
\end{ex}

\begin{ex} %[2H6V29]
Từ 13 giờ đến 18 giờ cùng ngày có bao nhiêu giờ trôi qua?
\shortans{5}
\loigiai{
\textbf{Đáp số:} $5$\\
$18 - 13 = 5$ giờ.
}
\end{ex}

\newpage

\TaoPhieuKiemTra{Chủ đề 6 -- Ngày – giờ, giờ – phút, ngày – tháng}{30}{Ngày – tháng}{https://loptoanthaylinh.vn/thi/chuyen-de/2/toan2_cd_b30}
\textbf{Phần 1. Trắc nghiệm lựa chọn}
\setcounterpageref{ex}{0}

\begin{ex} %[2H6V30]
Một năm có bao nhiêu tháng?
\choice
{$10$ tháng}
{\True $12$ tháng}
{$11$ tháng}
{$14$ tháng}
\loigiai{
Một năm luôn có đúng 12 tháng.
\textbf{Đáp án của câu hỏi là B}
}
\end{ex}

\begin{ex} %[2H6V30]
Tháng nào dưới đây có $31$ ngày?
\choice
{\True Tháng 1}
{Tháng 4}
{Tháng 6}
{Tháng 9}
\loigiai{
Các tháng có 31 ngày là: Tháng 1, 3, 5, 7, 8, 10, 12.
\textbf{Đáp án của câu hỏi là A}
}
\end{ex}

\begin{ex} %[2H6V30]
Một tuần lễ có bao nhiêu ngày?
\choice
{$5$ ngày}
{$6$ ngày}
{\True $7$ ngày}
{$8$ ngày}
\loigiai{
Một tuần lễ có 7 ngày.
\textbf{Đáp án của câu hỏi là C}
}
\end{ex}

\begin{ex} %[2H6V30]
Nếu hôm nay là thứ Ba ngày $5$ tháng $10$. Thứ Ba tuần sau là ngày mấy tháng $10$?
\choice
{Ngày $10$}
{Ngày $11$}
{Ngày $13$}
{\True Ngày $12$}
\loigiai{
Thứ Ba tuần sau là: $5 + 7 = 12$.
\textbf{Đáp án của câu hỏi là D}
}
\end{ex}

\begin{ex} %[2H6V30]
Ngày Quốc tế Thiếu nhi $1$ tháng $6$ năm nay rơi vào thứ Sáu. Hỏi ngày $8$ tháng $6$ năm nay là thứ mấy?
\choice
{Thứ Bảy}
{\True Thứ Sáu}
{Thứ Năm}
{Chủ nhật}
\loigiai{
Từ ngày 1 đến ngày 8 là $8 - 1 = 7$ ngày (đúng 1 tuần). Do đó vẫn là thứ Sáu.
\textbf{Đáp án của câu hỏi là B}
}
\end{ex}

\begin{ex} %[2H6V30]
Tháng nào trong năm có ít ngày nhất?
\choice
{\True Tháng 2}
{Tháng 4}
{Tháng 6}
{Tháng 11}
\loigiai{
Tháng 2 có 28 hoặc 29 ngày, ít ngày nhất trong năm.
\textbf{Đáp án của câu hỏi là A}
}
\end{ex}

\begin{ex} %[2H6V30]
Nếu thứ Bảy tuần này là ngày $18$ thì thứ Bảy tuần trước là ngày mấy?
\choice
{Ngày $10$}
{Ngày $12$}
{\True Ngày $11$}
{Ngày $9$}
\loigiai{
Thứ Bảy tuần trước là: $18 - 7 = 11$.
\textbf{Đáp án của câu hỏi là C}
}
\end{ex}

\begin{ex} %[2H6V30]
Các tháng có $30$ ngày trong năm là
\choice
{Tháng 1, 3, 5, 7}
{Tháng 2, 4, 6, 8}
{Tháng 7, 8, 9, 10}
{\True Tháng 4, 6, 9, 11}
\loigiai{
Các tháng có 30 ngày gồm: Tháng 4, tháng 6, tháng 9 và tháng 11.
\textbf{Đáp án của câu hỏi là D}
}
\end{ex}

\begin{ex} %[2H6V30]
Nếu ngày $20$ của một tháng là Chủ nhật thì ngày Chủ nhật tiếp theo của tháng đó là ngày mấy?
\choice
{Ngày $26$}
{\True Ngày $27$}
{Ngày $28$}
{Ngày $25$}
\loigiai{
Chủ nhật tiếp theo là: $20 + 7 = 27$.
\textbf{Đáp án của câu hỏi là B}
}
\end{ex}

\begin{ex} %[2H6V30]
Một tháng có thể có nhiều nhất bao nhiêu ngày Chủ nhật?
\choice
{\True $5$ ngày Chủ nhật}
{$4$ ngày Chủ nhật}
{$6$ ngày Chủ nhật}
{$3$ ngày Chủ nhật}
\loigiai{
Một tháng có tối đa 31 ngày, $31 : 7 = 4$ dư 3 ngày. Do đó có tối đa 5 ngày Chủ nhật.
\textbf{Đáp án của câu hỏi là A}
}
\end{ex}

\begin{ex} %[2H6V30]
Bố đi công tác $2$ tuần lễ. Hỏi bố đi công tác bao nhiêu ngày?
\choice
{$10$ ngày}
{$12$ ngày}
{\True $14$ ngày}
{$16$ ngày}
\loigiai{
$2 \times 7 = 14$ (ngày).
\textbf{Đáp án của câu hỏi là C}
}
\end{ex}

\begin{ex} %[2H6V30]
Nếu ngày cuối cùng của tháng là ngày $31$ thì tháng đó KHÔNG THỂ là
\choice
{Tháng 7}
{Tháng 8}
{Tháng 12}
{\True Tháng 9}
\loigiai{
Tháng 9 chỉ có 30 ngày nên không thể kết thúc vào ngày 31.
\textbf{Đáp án của câu hỏi là D}
}
\end{ex}

\textbf{Phần 2. Trắc nghiệm đúng sai}
\setcounterpageref{ex}{0}

\begin{ex} %[2H6V30]
Xét tính đúng sai của các nhận định về ngày tháng trong năm:
\choiceTF
{\True Tháng 7 và Tháng 8 là hai tháng liên tiếp đều có 31 ngày.}
{Tháng 2 luôn luôn chỉ có 28 ngày.}
{\True Nếu ngày 2 là thứ Hai thì ngày 9 cùng tháng cũng là thứ Hai.}
{Tất cả các tháng chẵn trong năm đều có 30 ngày.}
\loigiai{
\begin{itemize}
    \item Ý a) Đúng: Đúng, cả tháng 7 và tháng 8 đều có 31 ngày.
    \item Ý b) Sai: Năm nhuận tháng 2 có 29 ngày.
    \item Ý c) Đúng: $2 + 7 = 9$, đúng 1 tuần sau.
    \item Ý d) Sai: Tháng 8, 10, 12 là tháng chẵn nhưng có 31 ngày; tháng 2 có 28/29 ngày.
\end{itemize}
\begin{itemize}
    \item a) \textbf{Đúng}.
    \item b) \textbf{Sai}.
    \item c) \textbf{Đúng}.
    \item d) \textbf{Sai}.
\end{itemize}
}
\end{ex}

\begin{ex} %[2H6V30]
Trong một tháng có ngày thứ Hai rơi vào các ngày: ngày 3, ngày 10, ngày 17, ngày 24, ngày 31:
\choiceTF
{\True Tháng đó có đúng 5 ngày thứ Hai.}
{\True Tháng đó chắc chắn phải có 31 ngày.}
{Ngày 1 của tháng đó là ngày Chủ nhật.}
{\True Thứ Ba đầu tiên của tháng đó là ngày 4.}
\loigiai{
\begin{itemize}
    \item Ý a) Đúng: Các ngày: 3, 10, 17, 24, 31 (5 ngày).
    \item Ý b) Đúng: Vì có ngày 31 là thứ Hai nên tháng đó có 31 ngày.
    \item Ý c) Sai: Ngày 3 là thứ Hai thì ngày 2 là Chủ nhật, ngày 1 là thứ Bảy.
    \item Ý d) Đúng: Sau ngày 3 là ngày 4 (thứ Ba). Đúng.
\end{itemize}
\begin{itemize}
    \item a) \textbf{Đúng}.
    \item b) \textbf{Đúng}.
    \item c) \textbf{Sai}.
    \item d) \textbf{Đúng}.
\end{itemize}
}
\end{ex}

\textbf{Phần 3. Trả lời ngắn}
\setcounterpageref{ex}{0}

\begin{ex} %[2H6V30]
Một tuần có 7 ngày. Hỏi 3 tuần lễ có tất cả bao nhiêu ngày?
\shortans{21}
\loigiai{
\textbf{Đáp số:} $21$\\
$7 \times 3 = 21$ ngày.
}
\end{ex}

\begin{ex} %[2H6V30]
Nếu thứ Năm tuần này là ngày 12 thì thứ Năm tuần sau là ngày mấy?
\shortans{19}
\loigiai{
\textbf{Đáp số:} $19$\\
$12 + 7 = 19$.
}
\end{ex}

\begin{ex} %[2H6V30]
Tháng 12 có bao nhiêu ngày?
\shortans{31}
\loigiai{
\textbf{Đáp số:} $31$\\
Tháng 12 có 31 ngày.
}
\end{ex}

\begin{ex} %[2H6V30]
Năm nay ngày 20 tháng 11 là ngày thứ Tư. Hỏi ngày 27 tháng 11 năm nay là thứ mấy? (Quy ước: Thứ Hai = 2, ..., Thứ Tư = 4, Chủ nhật = 8)
\shortans{4}
\loigiai{
\textbf{Đáp số:} $4$\\
Cách nhau 7 ngày nên vẫn là thứ Tư (số 4).
}
\end{ex}

\begin{ex} %[2H6V30]
Trong một năm có bao nhiêu tháng có 30 ngày?
\shortans{4}
\loigiai{
\textbf{Đáp số:} $4$\\
Có 4 tháng: tháng 4, 6, 9, 11.
}
\end{ex}

\begin{ex} %[2H6V30]
Kỳ nghỉ hè của học sinh kéo dài trong 8 tuần lễ. Hỏi kỳ nghỉ hè kéo dài bao nhiêu ngày?
\shortans{56}
\loigiai{
\textbf{Đáp số:} $56$\\
$8 \times 7 = 56$ ngày.
}
\end{ex}

\newpage

\TaoPhieuKiemTra{Chủ đề 6 -- Ngày – giờ, giờ – phút, ngày – tháng}{31}{Thực hành xem đồng hồ, xem lịch}{https://loptoanthaylinh.vn/thi/chuyen-de/2/toan2_cd_b31}
\textbf{Phần 1. Trắc nghiệm lựa chọn}
\setcounterpageref{ex}{0}

\begin{ex} %[2H6V31]
Đồng hồ chỉ $4$ giờ $15$ phút. Kim phút đang chỉ vào số mấy?
\choice
{Số $4$}
{Số $15$}
{\True Số $3$}
{Số $1$}
\loigiai{
Mỗi khoảng giữa hai số là 5 phút. $15$ phút tương ứng kim phút chỉ vào số $3$.
\textbf{Đáp án của câu hỏi là C}
}
\end{ex}

\begin{ex} %[2H6V31]
Nhìn vào tờ lịch tháng $5$, ngày sinh nhật Bác Hồ là ngày $19$ tháng $5$. Nếu ngày $12$ tháng $5$ là thứ Bảy thì sinh nhật Bác là thứ mấy?
\choice
{\True Thứ Bảy}
{Chủ nhật}
{Thứ Sáu}
{Thứ Hai}
\loigiai{
$19 - 12 = 7$ ngày (đúng 1 tuần). Vậy vẫn là thứ Bảy.
\textbf{Đáp án của câu hỏi là A}
}
\end{ex}

\begin{ex} %[2H6V31]
Đồng hồ điện tử hiển thị $18:30$. Điều đó có nghĩa là mấy giờ?
\choice
{$8$ giờ $30$ phút tối}
{$18$ giờ $30$ phút sáng}
{$6$ giờ chiều}
{\True $6$ giờ $30$ phút tối}
\loigiai{
$18:30$ là $6$ giờ $30$ phút tối (hoặc $6$ giờ rưỡi chiều).
\textbf{Đáp án của câu hỏi là D}
}
\end{ex}

\begin{ex} %[2H6V31]
Lan đi bộ từ nhà đến trường hết $20$ phút. Lan đến trường lúc $7$ giờ $20$ phút sáng. Hỏi Lan xuất phát từ nhà lúc mấy giờ?
\choice
{$6$ giờ $50$ phút}
{\True $7$ giờ sáng}
{$7$ giờ $10$ phút}
{$7$ giờ}
\loigiai{
$7$ giờ $20$ phút $- 20$ phút $= 7$ giờ sáng.
\textbf{Đáp án của câu hỏi là B}
}
\end{ex}

\begin{ex} %[2H6V31]
Khi kim giờ và kim phút của đồng hồ cùng chỉ vào số $12$, lúc đó là mấy giờ?
\choice
{$12$ giờ $30$ phút}
{$6$ giờ}
{\True $12$ giờ đúng}
{$1$ giờ}
\loigiai{
Cả hai kim cùng chỉ số 12 là đúng 12 giờ.
\textbf{Đáp án của câu hỏi là C}
}
\end{ex}

\begin{ex} %[2H6V31]
Trong một tháng, ngày Chủ nhật đầu tiên là ngày $2$. Hỏi ngày Chủ nhật cuối cùng của tháng đó là ngày mấy nếu tháng đó có $30$ ngày?
\choice
{\True Ngày $30$}
{Ngày $28$}
{Ngày $23$}
{Ngày $29$}
\loigiai{
Các ngày Chủ nhật: $2, 9, 16, 23, 30$. Ngày cuối cùng là ngày $30$.
\textbf{Đáp án của câu hỏi là A}
}
\end{ex}

\begin{ex} %[2H6V31]
Mẹ đi chợ từ $8$ giờ sáng đến $9$ giờ $30$ phút sáng. Mẹ đi chợ trong bao lâu?
\choice
{$1$ giờ}
{$2$ giờ}
{$45$ phút}
{\True $1$ giờ $30$ phút}
\loigiai{
$9$ giờ $30$ phút $- 8$ giờ $= 1$ giờ $30$ phút.
\textbf{Đáp án của câu hỏi là D}
}
\end{ex}

\begin{ex} %[2H6V31]
Đồng hồ đang chỉ $2$ giờ rưỡi chiều. Sau $30$ phút nữa thì đồng hồ chỉ mấy giờ?
\choice
{$2$ giờ chiều}
{\True $3$ giờ chiều}
{$3$ giờ $30$ phút}
{$4$ giờ chiều}
\loigiai{
$2$ giờ rưỡi chiều là $2$ giờ $30$ phút chiều. Thêm $30$ phút nữa là đúng $3$ giờ chiều.
\textbf{Đáp án của câu hỏi là B}
}
\end{ex}

\begin{ex} %[2H6V31]
Nếu ngày $28$ tháng $4$ là thứ Bảy thì ngày $1$ tháng $5$ cùng năm đó là thứ mấy? (Biết tháng $4$ có $30$ ngày)
\choice
{Thứ Hai}
{Chủ nhật}
{\True Thứ Ba}
{Thứ Tư}
\loigiai{
Ngày 28: thứ Bảy; Ngày 29: Chủ nhật; Ngày 30: thứ Hai; Ngày 1 tháng 5: thứ Ba.
\textbf{Đáp án của câu hỏi là C}
}
\end{ex}

\begin{ex} %[2H6V31]
Một tiết mục xiếc kéo dài $45$ phút, kết thúc lúc $10$ giờ sáng. Tiết mục xiếc đó bắt đầu lúc mấy giờ?
\choice
{\True $9$ giờ $15$ phút sáng}
{$9$ giờ}
{$9$ giờ $30$ phút}
{$8$ giờ $45$ phút}
\loigiai{
$10$ giờ trừ đi $45$ phút là $9$ giờ $15$ phút sáng.
\textbf{Đáp án của câu hỏi là A}
}
\end{ex}

\begin{ex} %[2H6V31]
Đồng hồ có kim phút chỉ số $12$, kim giờ chỉ số $7$. Lúc này là
\choice
{$12$ giờ $7$ phút}
{$7$ giờ $12$ phút}
{$6$ giờ $50$ phút}
{\True $7$ giờ đúng}
\loigiai{
Kim phút chỉ số 12, kim giờ chỉ số 7 là đúng 7 giờ.
\textbf{Đáp án của câu hỏi là D}
}
\end{ex}

\begin{ex} %[2H6V31]
Nếu hôm nay là thứ Sáu ngày $13$. Thứ Sáu tuần trước là ngày mấy?
\choice
{Ngày $7$}
{\True Ngày $6$}
{Ngày $5$}
{Ngày $8$}
\loigiai{
$13 - 7 = 6$.
\textbf{Đáp án của câu hỏi là B}
}
\end{ex}

\textbf{Phần 2. Trắc nghiệm đúng sai}
\setcounterpageref{ex}{0}

\begin{ex} %[2H6V31]
Xét tính đúng sai của các tình huống thực hành xem giờ và lịch:
\choiceTF
{\True Khi kim phút chỉ vào số 3 là biểu thị 15 phút.}
{Khi kim phút chỉ vào số 6 là biểu thị 45 phút.}
{\True Hai ngày Chủ nhật liên tiếp trong cùng một tháng luôn cách nhau đúng 7 ngày.}
{\True Từ 8 giờ đến 9 giờ kim phút quay trọn vẹn đúng 1 vòng.}
\loigiai{
\begin{itemize}
    \item Ý a) Đúng: $3 \times 5 = 15$ phút. Đúng.
    \item Ý b) Sai: Kim phút chỉ số 6 là 30 phút.
    \item Ý c) Đúng: Đúng, chu kỳ 1 tuần là 7 ngày.
    \item Ý d) Đúng: 1 giờ ứng với 1 vòng kim phút.
\end{itemize}
\begin{itemize}
    \item a) \textbf{Đúng}.
    \item b) \textbf{Sai}.
    \item c) \textbf{Đúng}.
    \item d) \textbf{Đúng}.
\end{itemize}
}
\end{ex}

\begin{ex} %[2H6V31]
Lịch thi học kì của bạn Minh diễn ra trong tuần:
\choiceTF
{\True Thứ Hai ngày 10 thi Toán, thứ Năm cùng tuần thi Tiếng Việt thì thi Tiếng Việt vào ngày 13.}
{\True Nếu môn Toán thi lúc 8 giờ và làm bài 40 phút thì kết thúc lúc 8 giờ 40 phút.}
{Môn Tiếng Việt thi từ 9 giờ đến 10 giờ 30 phút là thi trong 2 giờ.}
{Thứ Bảy tuần đó là ngày 14.}
\loigiai{
\begin{itemize}
    \item Ý a) Đúng: $10 + 3 = 13$. Đúng.
    \item Ý b) Đúng: Đúng.
    \item Ý c) Sai: Chỉ có 1 giờ 30 phút.
    \item Ý d) Sai: Thứ Hai ngày 10, thứ Bảy là $10 + 5 = 15$ (không phải 14).
\end{itemize}
\begin{itemize}
    \item a) \textbf{Đúng}.
    \item b) \textbf{Đúng}.
    \item c) \textbf{Sai}.
    \item d) \textbf{Sai}.
\end{itemize}
}
\end{ex}

\textbf{Phần 3. Trả lời ngắn}
\setcounterpageref{ex}{0}

\begin{ex} %[2H6V31]
Bác Hùng đi ô tô từ 14 giờ đến 17 giờ thì đến nơi. Bác đi hết bao nhiêu giờ?
\shortans{3}
\loigiai{
\textbf{Đáp số:} $3$\\
$17 - 14 = 3$ giờ.
}
\end{ex}

\begin{ex} %[2H6V31]
Đồng hồ chỉ 9 giờ rưỡi, tức là 9 giờ bao nhiêu phút?
\shortans{30}
\loigiai{
\textbf{Đáp số:} $30$\\
Rưỡi là 30 phút.
}
\end{ex}

\begin{ex} %[2H6V31]
Tháng Ba có 31 ngày. Hỏi tháng Ba có bao nhiêu ngày nhiều hơn một tháng có 30 ngày?
\shortans{1}
\loigiai{
\textbf{Đáp số:} $1$\\
$31 - 30 = 1$ ngày.
}
\end{ex}

\begin{ex} %[2H6V31]
Nếu thứ Tư tuần này là ngày 15 thì thứ Tư tuần trước là ngày mấy?
\shortans{8}
\loigiai{
\textbf{Đáp số:} $8$\\
$15 - 7 = 8$.
}
\end{ex}

\begin{ex} %[2H6V31]
Một trận thi đấu cờ vua bắt đầu lúc 8 giờ 10 phút và kết thúc lúc 8 giờ 40 phút. Trận đấu diễn ra trong bao nhiêu phút?
\shortans{30}
\loigiai{
\textbf{Đáp số:} $30$\\
$40 - 10 = 30$ phút.
}
\end{ex}

\begin{ex} %[2H6V31]
Nếu ngày 25 của một tháng là thứ Bảy thì thứ Bảy tuần trước là ngày mấy?
\shortans{18}
\loigiai{
\textbf{Đáp số:} $18$\\
$25 - 7 = 18$.
}
\end{ex}

\newpage

\TaoPhieuKiemTra{Chủ đề 6 -- Ngày – giờ, giờ – phút, ngày – tháng}{32}{Luyện tập chung}{https://loptoanthaylinh.vn/thi/chuyen-de/2/toan2_cd_b32}
\textbf{Phần 1. Trắc nghiệm lựa chọn}
\setcounterpageref{ex}{0}

\begin{ex} %[2H6V32]
Một ngày có $24$ giờ. Hỏi $3$ ngày có tất cả bao nhiêu giờ?
\choice
{$48$ giờ}
{\True $72$ giờ}
{$60$ giờ}
{$96$ giờ}
\loigiai{
$24 + 24 + 24 = 72$ (giờ).
\textbf{Đáp án của câu hỏi là B}
}
\end{ex}

\begin{ex} %[2H6V32]
Đồng hồ chỉ $7$ giờ $15$ phút tối. Cách đọc nào sau đây theo hệ $24$ giờ là đúng?
\choice
{$17$ giờ $15$ phút}
{$18$ giờ $15$ phút}
{$20$ giờ $15$ phút}
{\True $19$ giờ $15$ phút}
\loigiai{
$7$ giờ tối $= 19$ giờ. Vậy là $19$ giờ $15$ phút.
\textbf{Đáp án của câu hỏi là D}
}
\end{ex}

\begin{ex} %[2H6V32]
Tháng nào sau đây có đúng $30$ ngày?
\choice
{\True Tháng 4}
{Tháng 3}
{Tháng 5}
{Tháng 7}
\loigiai{
Tháng 4 có đúng 30 ngày.
\textbf{Đáp án của câu hỏi là A}
}
\end{ex}

\begin{ex} %[2H6V32]
Nếu hôm nay là thứ Năm ngày $8$. Hỏi $10$ ngày nữa là ngày mấy?
\choice
{Ngày $17$}
{Ngày $19$}
{\True Ngày $18$}
{Ngày $16$}
\loigiai{
$8 + 10 = 18$.
\textbf{Đáp án của câu hỏi là C}
}
\end{ex}

\begin{ex} %[2H6V32]
Bác thợ mộc đóng một chiếc bàn mất $3$ giờ. Bác bắt đầu đóng lúc $8$ giờ sáng và nghỉ giải lao $30$ phút. Bác đóng xong chiếc bàn lúc mấy giờ?
\choice
{$11$ giờ}
{\True $11$ giờ $30$ phút trưa}
{$12$ giờ}
{$10$ giờ $30$ phút}
\loigiai{
Thời gian hoàn thành: $8$ giờ $+ 3$ giờ $+ 30$ phút $= 11$ giờ $30$ phút.
\textbf{Đáp án của câu hỏi là B}
}
\end{ex}

\begin{ex} %[2H6V32]
Kim giờ quay được $1$ vòng thì kim phút quay được bao nhiêu vòng?
\choice
{$6$ vòng}
{$24$ vòng}
{$60$ vòng}
{\True $12$ vòng}
\loigiai{
Kim giờ quay 1 vòng hết 12 giờ. Mỗi giờ kim phút quay 1 vòng nên quay được 12 vòng.
\textbf{Đáp án của câu hỏi là D}
}
\end{ex}

\begin{ex} %[2H6V32]
Một tháng có $31$ ngày, nếu ngày đầu tiên là thứ Ba thì tháng đó có bao nhiêu ngày thứ Ba?
\choice
{\True $5$ ngày}
{$4$ ngày}
{$6$ ngày}
{$3$ ngày}
\loigiai{
Các ngày thứ Ba là: 1, 8, 15, 22, 29 (có 5 ngày).
\textbf{Đáp án của câu hỏi là A}
}
\end{ex}

\begin{ex} %[2H6V32]
Tàu hỏa chạy từ ga $A$ lúc $6$ giờ sáng và đến ga $B$ lúc $14$ giờ cùng ngày. Tàu đã chạy hết bao nhiêu giờ?
\choice
{$7$ giờ}
{$9$ giờ}
{\True $8$ giờ}
{$10$ giờ}
\loigiai{
$14 - 6 = 8$ (giờ).
\textbf{Đáp án của câu hỏi là C}
}
\end{ex}

\begin{ex} %[2H6V32]
Có bao nhiêu ngày trong $4$ tuần lễ?
\choice
{$24$ ngày}
{\True $28$ ngày}
{$30$ ngày}
{$32$ ngày}
\loigiai{
$7 \times 4 = 28$ (ngày).
\textbf{Đáp án của câu hỏi là B}
}
\end{ex}

\begin{ex} %[2H6V32]
Đồng hồ chỉ $16$ giờ $30$ phút. Kim phút đang chỉ vào số nào trên mặt đồng hồ?
\choice
{Số $3$}
{Số $12$}
{Số $9$}
{\True Số $6$}
\loigiai{
$30$ phút ứng với kim phút chỉ vào số 6.
\textbf{Đáp án của câu hỏi là D}
}
\end{ex}

\begin{ex} %[2H6V32]
Nếu thứ Bảy là ngày $15$ thì Chủ nhật tuần sau là ngày mấy?
\choice
{\True Ngày $23$}
{Ngày $22$}
{Ngày $24$}
{Ngày $21$}
\loigiai{
Chủ nhật tuần này là ngày 16. Chủ nhật tuần sau là: $16 + 7 = 23$.
\textbf{Đáp án của câu hỏi là A}
}
\end{ex}

\begin{ex} %[2H6V32]
Mỗi ngày Nam dành $1$ giờ để đọc sách. Hỏi trong một tuần Nam dành bao nhiêu giờ đọc sách?
\choice
{$5$ giờ}
{$6$ giờ}
{\True $7$ giờ}
{$8$ giờ}
\loigiai{
Một tuần có 7 ngày: $1 \times 7 = 7$ (giờ).
\textbf{Đáp án của câu hỏi là C}
}
\end{ex}

\textbf{Phần 2. Trắc nghiệm đúng sai}
\setcounterpageref{ex}{0}

\begin{ex} %[2H6V32]
Xét tính đúng sai của các bài toán thời gian tổng hợp:
\choiceTF
{\True Một năm thông thường có 365 ngày.}
{Năm nào tháng 2 cũng có 29 ngày.}
{\True Khoảng thời gian từ 10 giờ sáng đến 2 giờ chiều là 4 giờ.}
{Một tuần lễ chỉ gồm các ngày học từ thứ Hai đến thứ Sáu.}
\loigiai{
\begin{itemize}
    \item Ý a) Đúng: Đúng theo lịch dương chuẩn.
    \item Ý b) Sai: Chỉ có năm nhuận tháng 2 mới có 29 ngày.
    \item Ý c) Đúng: 2 giờ chiều là 14 giờ: $14 - 10 = 4$ giờ. Đúng.
    \item Ý d) Sai: Một tuần lễ đầy đủ gồm 7 ngày từ thứ Hai đến Chủ nhật.
\end{itemize}
\begin{itemize}
    \item a) \textbf{Đúng}.
    \item b) \textbf{Sai}.
    \item c) \textbf{Đúng}.
    \item d) \textbf{Sai}.
\end{itemize}
}
\end{ex}

\begin{ex} %[2H6V32]
Kỳ nghỉ lễ Quốc khánh 2 tháng 9:
\choiceTF
{\True Tháng 9 có đúng 30 ngày.}
{\True Ngày 2 tháng 9 là ngày Quốc khánh của nước Việt Nam.}
{Nếu ngày 2 tháng 9 là thứ Hai thì ngày 9 tháng 9 là thứ Ba.}
{\True Nếu được nghỉ 4 ngày từ ngày 1 đến hết ngày 4 tháng 9 thì học sinh đi học lại vào ngày 5 tháng 9.}
\loigiai{
\begin{itemize}
    \item Ý a) Đúng: Đúng.
    \item Ý b) Đúng: Đúng.
    \item Ý c) Sai: Cộng 7 ngày nên ngày 9 tháng 9 vẫn là thứ Hai.
    \item Ý d) Đúng: Đúng.
\end{itemize}
\begin{itemize}
    \item a) \textbf{Đúng}.
    \item b) \textbf{Đúng}.
    \item c) \textbf{Sai}.
    \item d) \textbf{Đúng}.
\end{itemize}
}
\end{ex}

\textbf{Phần 3. Trả lời ngắn}
\setcounterpageref{ex}{0}

\begin{ex} %[2H6V32]
Một ngày có 24 giờ. Hỏi nửa ngày có bao nhiêu giờ?
\shortans{12}
\loigiai{
\textbf{Đáp số:} $12$\\
$24 : 2 = 12$ giờ.
}
\end{ex}

\begin{ex} %[2H6V32]
Nếu ngày 5 của một tháng là thứ Ba thì ngày thứ Ba của tuần tiếp theo là ngày mấy?
\shortans{12}
\loigiai{
\textbf{Đáp số:} $12$\\
$5 + 7 = 12$.
}
\end{ex}

\begin{ex} %[2H6V32]
Một tuần lễ có 7 ngày. Hỏi 5 tuần lễ có tất cả bao nhiêu ngày?
\shortans{35}
\loigiai{
\textbf{Đáp số:} $35$\\
$7 \times 5 = 35$ ngày.
}
\end{ex}

\begin{ex} %[2H6V32]
Bác thợ làm việc từ 7 giờ sáng đến 11 giờ trưa và từ 13 giờ đến 17 giờ chiều. Hỏi bác làm việc tất cả bao nhiêu giờ trong ngày?
\shortans{8}
\loigiai{
\textbf{Đáp số:} $8$\\
Sáng: $11 - 7 = 4$ giờ. Chiều: $17 - 13 = 4$ giờ. Tổng: $4 + 4 = 8$ giờ.
}
\end{ex}

\begin{ex} %[2H6V32]
Đồng hồ chỉ 21 giờ 30 phút, tức là 9 giờ bao nhiêu phút tối?
\shortans{30}
\loigiai{
\textbf{Đáp số:} $30$\\
30 phút.
}
\end{ex}

\begin{ex} %[2H6V32]
Tháng Mười có 31 ngày, tháng Mười Một có 30 ngày. Hỏi hai tháng đó có tất cả bao nhiêu ngày?
\shortans{61}
\loigiai{
\textbf{Đáp số:} $61$\\
$31 + 30 = 61$ ngày.
}
\end{ex}

\newpage


% =========================================================================
% CHỦ ĐỀ 7 - ÔN TẬP HỌC KÌ 1
% =========================================================================
\section{Ôn tập học kì 1}

\TaoPhieuKiemTra{Chủ đề 7 -- Ôn tập học kì 1}{33}{Ôn tập phép cộng, phép trừ trong phạm vi 20, 100}{https://loptoanthaylinh.vn/thi/chuyen-de/2/toan2_cd_b33}
\textbf{Phần 1. Trắc nghiệm lựa chọn}
\setcounterpageref{ex}{0}

\begin{ex} %[2D7V33]
Số tròn chục lớn nhất có hai chữ số là
\choice
{\True $90$}
{$99$}
{$80$}
{$100$}
\loigiai{
Số tròn chục lớn nhất có hai chữ số là 90.
\textbf{Đáp án của câu hỏi là A}
}
\end{ex}

\begin{ex} %[2D7V33]
Tìm số nhỏ nhất có hai chữ số mà tổng hai chữ số của nó bằng $10$.
\choice
{$28$}
{$37$}
{\True $19$}
{$18$}
\loigiai{
Để nhỏ nhất, chữ số hàng chục nhỏ nhất là 1, chữ số hàng đơn vị là $10 - 1 = 9$. Số đó là 19.
\textbf{Đáp án của câu hỏi là C}
}
\end{ex}

\begin{ex} %[2D7V33]
Trong số $74$, chữ số $7$ có giá trị lớn hơn chữ số $4$ bao nhiêu đơn vị?
\choice
{$63$ đơn vị}
{\True $66$ đơn vị}
{$70$ đơn vị}
{$3$ đơn vị}
\loigiai{
Giá trị chữ số 7 là 70, chữ số 4 là 4. Hiệu: $70 - 4 = 66$.
\textbf{Đáp án của câu hỏi là B}
}
\end{ex}

\begin{ex} %[2D7V33]
Cho các số: $35, 82, 53, 28, 91$. Số lớn thứ hai trong các số đã cho là
\choice
{$91$}
{$53$}
{$35$}
{\True $82$}
\loigiai{
Sắp xếp giảm dần: $91 > 82 > 53 > 35 > 28$. Số lớn thứ hai là 82.
\textbf{Đáp án của câu hỏi là D}
}
\end{ex}

\begin{ex} %[2D7V33]
Số liền sau của số lớn nhất có hai chữ số khác nhau là
\choice
{\True $99$}
{$98$}
{$100$}
{$97$}
\loigiai{
Số lớn nhất có 2 chữ số khác nhau là 98. Số liền sau là $98 + 1 = 99$.
\textbf{Đáp án của câu hỏi là A}
}
\end{ex}

\begin{ex} %[2D7V33]
Có bao nhiêu số có hai chữ số mà chữ số hàng chục bé hơn chữ số hàng đơn vị $7$ đơn vị?
\choice
{$1$ số}
{$3$ số}
{\True $2$ số}
{$4$ số}
\loigiai{
Các số đó là: 18 ($8 - 1 = 7$) và 29 ($9 - 2 = 7$). Có tất cả 2 số.
\textbf{Đáp án của câu hỏi là C}
}
\end{ex}

\begin{ex} %[2D7V33]
Tìm $x$ biết $x$ là số tròn chục và $45 < x < 65$. Giá trị của $x$ là
\choice
{$50$}
{\True $50\text{ hoặc } 60$}
{$60$}
{$55$}
\loigiai{
Các số tròn chục thỏa mãn là 50 và 60.
\textbf{Đáp án của câu hỏi là B}
}
\end{ex}

\begin{ex} %[2D7V33]
Từ hai chữ số $4$ và $9$, có thể viết được bao nhiêu số có hai chữ số?
\choice
{$2$ số}
{$3$ số}
{$6$ số}
{\True $4$ số}
\loigiai{
Các số là: 44, 49, 94, 99. Có tất cả 4 số.
\textbf{Đáp án của câu hỏi là D}
}
\end{ex}

\begin{ex} %[2D7V33]
Hiệu giữa số chẵn nhỏ nhất có hai chữ số và số lẻ nhỏ nhất có một chữ số là
\choice
{\True $9$}
{$8$}
{$10$}
{$7$}
\loigiai{
Số chẵn nhỏ nhất có 2 chữ số là 10. Số lẻ nhỏ nhất có 1 chữ số là 1. Hiệu: $10 - 1 = 9$.
\textbf{Đáp án của câu hỏi là A}
}
\end{ex}

\begin{ex} %[2D7V33]
Cho dãy số: $5, 10, 15, 20, \dots$ Số hạng thứ $8$ của dãy số là
\choice
{$35$}
{$45$}
{\True $40$}
{$30$}
\loigiai{
Dãy số là các số đếm thêm 5. Số hạng thứ 8 là: $5 \times 8 = 40$.
\textbf{Đáp án của câu hỏi là C}
}
\end{ex}

\begin{ex} %[2D7V33]
Một lớp học có số học sinh nhiều hơn $30$ bạn nhưng ít hơn $35$ bạn. Nếu xếp thành các hàng $4$ bạn thì vừa hết. Số học sinh của lớp là
\choice
{$31$ bạn}
{\True $32$ bạn}
{$33$ bạn}
{$34$ bạn}
\loigiai{
Số chia hết cho 4 nằm giữa 30 và 35 là 32 bạn.
\textbf{Đáp án của câu hỏi là B}
}
\end{ex}

\begin{ex} %[2D7V33]
Điền dấu thích hợp: $50 + 40 \dots 100 - 15$.
\choice
{$<$}
{$=$}
{Không so sánh được}
{\True $>$}
\loigiai{
$50 + 40 = 90$; $100 - 15 = 85$. Vì $90 > 85$ nên điền dấu $>$. 
\textbf{Đáp án của câu hỏi là D}
}
\end{ex}

\textbf{Phần 2. Trắc nghiệm đúng sai}
\setcounterpageref{ex}{0}

\begin{ex} %[2D7V33]
Xét tính đúng sai của các nhận định số học ôn tập học kì 1:
\choiceTF
{\True Số bé nhất có hai chữ số là 10.}
{Số 100 là số lớn nhất có hai chữ số.}
{\True Hai số tự nhiên liên tiếp luôn gồm một số chẵn và một số lẻ.}
{\True Trong số 85, chữ số 8 thuộc hàng chục và có giá trị là 80.}
\loigiai{
\begin{itemize}
    \item Ý a) Đúng: Đúng.
    \item Ý b) Sai: 100 là số có 3 chữ số.
    \item Ý c) Đúng: Đúng.
    \item Ý d) Đúng: Đúng.
\end{itemize}
\begin{itemize}
    \item a) \textbf{Đúng}.
    \item b) \textbf{Sai}.
    \item c) \textbf{Đúng}.
    \item d) \textbf{Đúng}.
\end{itemize}
}
\end{ex}

\begin{ex} %[2D7V33]
Cho ba chữ số 2, 5, 7:
\choiceTF
{\True Lập được tất cả 6 số có hai chữ số khác nhau từ 3 chữ số trên.}
{\True Số lớn nhất lập được là 75.}
{Có 3 số chẵn lập được từ 3 chữ số trên.}
{\True Tổng của số bé nhất (25) và số lớn nhất (75) bằng 100.}
\loigiai{
\begin{itemize}
    \item Ý a) Đúng: Các số là 25, 27, 52, 57, 72, 75.
    \item Ý b) Đúng: Trong các số khác nhau, 75 là lớn nhất.
    \item Ý c) Sai: Chỉ có 2 số có tận cùng là 2: 52 và 72.
    \item Ý d) Đúng: $25 + 75 = 100$. Đúng.
\end{itemize}
\begin{itemize}
    \item a) \textbf{Đúng}.
    \item b) \textbf{Đúng}.
    \item c) \textbf{Sai}.
    \item d) \textbf{Đúng}.
\end{itemize}
}
\end{ex}

\textbf{Phần 3. Trả lời ngắn}
\setcounterpageref{ex}{0}

\begin{ex} %[2D7V33]
Tìm số lớn nhất có hai chữ số mà tổng hai chữ số bằng 11.
\shortans{92}
\loigiai{
\textbf{Đáp số:} $92$\\
Chọn chữ số hàng chục lớn nhất là 9, hàng đơn vị là $11 - 9 = 2$. Số đó là 92.
}
\end{ex}

\begin{ex} %[2D7V33]
Có bao nhiêu số tròn chục có hai chữ số lớn hơn 30?
\shortans{6}
\loigiai{
\textbf{Đáp số:} $6$\\
Các số: 40, 50, 60, 70, 80, 90 (có 6 số).
}
\end{ex}

\begin{ex} %[2D7V33]
Tìm số liền trước của số nhỏ nhất có hai chữ số giống nhau.
\shortans{10}
\loigiai{
\textbf{Đáp số:} $10$\\
Số nhỏ nhất có 2 chữ số giống nhau là 11. Số liền trước là 10.
}
\end{ex}

\begin{ex} %[2D7V33]
Tính tổng các chữ số của số lớn nhất có hai chữ số khác nhau.
\shortans{17}
\loigiai{
\textbf{Đáp số:} $17$\\
Số lớn nhất có 2 chữ số khác nhau là 98. Tổng: $9 + 8 = 17$.
}
\end{ex}

\begin{ex} %[2D7V33]
Từ số 20 đến số 60 có tất cả bao nhiêu số tự nhiên liên tiếp?
\shortans{41}
\loigiai{
\textbf{Đáp số:} $41$\\
$(60 - 20) + 1 = 41$ số.
}
\end{ex}

\begin{ex} %[2D7V33]
Tìm một số biết rằng lấy số đó trừ đi 25 thì được số tròn chục nhỏ nhất có hai chữ số.
\shortans{35}
\loigiai{
\textbf{Đáp số:} $35$\\
Số tròn chục nhỏ nhất là 10. Số đó là: $10 + 25 = 35$.
}
\end{ex}

\newpage

\TaoPhieuKiemTra{Chủ đề 7 -- Ôn tập học kì 1}{34}{Ôn tập hình phẳng}{https://loptoanthaylinh.vn/thi/chuyen-de/2/toan2_cd_b34}
\textbf{Phần 1. Trắc nghiệm lựa chọn}
\setcounterpageref{ex}{0}

\begin{ex} %[2D7V34]
Tính giá trị biểu thức: $46 + 28 - 35$.
\choice
{$37$}
{$38$}
{$40$}
{\True $39$}
\loigiai{
$46 + 28 = 74$; $74 - 35 = 39$.
\textbf{Đáp án của câu hỏi là D}
}
\end{ex}

\begin{ex} %[2D7V34]
Tìm $x$ biết: $x - 29 = 45$.
\choice
{$72$}
{\True $74$}
{$73$}
{$75$}
\loigiai{
$x = 45 + 29 = 74$.
\textbf{Đáp án của câu hỏi là B}
}
\end{ex}

\begin{ex} %[2D7V34]
Một đàn gà có $52$ con, mẹ bán đi $18$ con rồi mua thêm $15$ con nữa. Hỏi đàn gà có bao nhiêu con?
\choice
{$48$ con gà}
{$50$ con gà}
{\True $49$ con gà}
{$47$ con gà}
\loigiai{
$52 - 18 + 15 = 34 + 15 = 49$ (con gà).
\textbf{Đáp án của câu hỏi là C}
}
\end{ex}

\begin{ex} %[2D7V34]
Tính nhanh: $36 + 18 + 24 + 12$.
\choice
{\True $90$}
{$85$}
{$80$}
{$95$}
\loigiai{
$(36 + 24) + (18 + 12) = 60 + 30 = 90$.
\textbf{Đáp án của câu hỏi là A}
}
\end{ex}

\begin{ex} %[2D7V34]
Một sợi dây dài $80\text{ cm}$. Người ta cắt lấy $2$ đoạn, mỗi đoạn dài $25\text{ cm}$. Đoạn dây còn lại dài bao nhiêu xăng-ti-mét?
\choice
{$25\text{ cm}$}
{$35\text{ cm}$}
{$20\text{ cm}$}
{\True $30\text{ cm}$}
\loigiai{
Hai đoạn dài: $25 + 25 = 50\text{ (cm)}$. Còn lại: $80 - 50 = 30\text{ (cm)}$.
\textbf{Đáp án của câu hỏi là D}
}
\end{ex}

\begin{ex} %[2D7V34]
Hai số có tổng bằng $68$, nếu số hạng thứ nhất tăng $7$ đơn vị và số hạng thứ hai giảm $5$ đơn vị thì tổng mới là
\choice
{$69$}
{\True $70$}
{$71$}
{$66$}
\loigiai{
Tổng mới: $68 + 7 - 5 = 70$.
\textbf{Đáp án của câu hỏi là B}
}
\end{ex}

\begin{ex} %[2D7V34]
Điền số thích hợp: $75 - \dots = 38 + 17$.
\choice
{$19$}
{$21$}
{\True $20$}
{$18$}
\loigiai{
Vế phải: $38 + 17 = 55$. Ô trống $= 75 - 55 = 20$.
\textbf{Đáp án của câu hỏi là C}
}
\end{ex}

\begin{ex} %[2D7V34]
Lớp 2A có $35$ học sinh, lớp 2B có $32$ học sinh. Cả hai lớp có bao nhiêu học sinh?
\choice
{\True $67$ học sinh}
{$65$ học sinh}
{$66$ học sinh}
{$68$ học sinh}
\loigiai{
$35 + 32 = 67$ (học sinh).
\textbf{Đáp án của câu hỏi là A}
}
\end{ex}

\begin{ex} %[2D7V34]
Trong một phép trừ, số bị trừ là $82$, hiệu là $39$. Số trừ là
\choice
{$41$}
{$42$}
{$44$}
{\True $43$}
\loigiai{
Số trừ: $82 - 39 = 43$.
\textbf{Đáp án của câu hỏi là D}
}
\end{ex}

\begin{ex} %[2D7V34]
Bác Ba thu hoạch được $72\text{ kg}$ cam. Bác bán được $48\text{ kg}$. Bác còn lại bao nhiêu ki-lô-gam cam?
\choice
{$22\text{ kg}$}
{\True $24\text{ kg}$}
{$23\text{ kg}$}
{$25\text{ kg}$}
\loigiai{
$72 - 48 = 24\text{ (kg)}$.
\textbf{Đáp án của câu hỏi là B}
}
\end{ex}

\begin{ex} %[2D7V34]
Tìm một số biết rằng lấy số đó cộng với $34$ thì bằng $91 - 18$.
\choice
{$38$}
{$40$}
{\True $39$}
{$37$}
\loigiai{
Vế phải: $91 - 18 = 73$. Số đó là: $73 - 34 = 39$.
\textbf{Đáp án của câu hỏi là C}
}
\end{ex}

\begin{ex} %[2D7V34]
Điền dấu thích hợp: $56 + 27 \dots 92 - 15$.
\choice
{\True $>$}
{$<$}
{$=$}
{Không so sánh được}
\loigiai{
$56 + 27 = 83$; $92 - 15 = 77$. Vì $83 > 77$ nên điền dấu $>$. 
\textbf{Đáp án của câu hỏi là A}
}
\end{ex}

\textbf{Phần 2. Trắc nghiệm đúng sai}
\setcounterpageref{ex}{0}

\begin{ex} %[2D7V34]
Xét tính đúng sai của các biểu thức tính toán:
\choiceTF
{\True Biểu thức $38 + 47 - 25$ có kết quả là $60$.}
{Hiệu của $72 - 38$ bằng hiệu của $82 - 38$.}
{\True Trong phép cộng, tổng luôn lớn hơn hoặc bằng mỗi số hạng.}
{\True Phép tính $65 - 28$ có kết quả là số lẻ.}
\loigiai{
\begin{itemize}
    \item Ý a) Đúng: $38 + 47 = 85$; $85 - 25 = 60$.
    \item Ý b) Sai: Hai số bị trừ khác nhau thì hiệu phải khác nhau.
    \item Ý c) Đúng: Đúng với các số tự nhiên.
    \item Ý d) Đúng: $65 - 28 = 37$, là số lẻ. Đúng.
\end{itemize}
\begin{itemize}
    \item a) \textbf{Đúng}.
    \item b) \textbf{Sai}.
    \item c) \textbf{Đúng}.
    \item d) \textbf{Đúng}.
\end{itemize}
}
\end{ex}

\begin{ex} %[2D7V34]
Một cửa hàng có 95 mét vải, ngày thứ nhất bán 38 mét, ngày thứ hai bán 29 mét:
\choiceTF
{\True Tổng số mét vải bán trong hai ngày là 67 mét.}
{\True Cửa hàng còn lại 28 mét vải.}
{Ngày thứ hai bán nhiều hơn ngày thứ nhất.}
{Số mét vải còn lại nhiều hơn số mét vải bán ngày thứ hai.}
\loigiai{
\begin{itemize}
    \item Ý a) Đúng: $38 + 29 = 67\text{ m}$.
    \item Ý b) Đúng: $95 - 67 = 28\text{ m}$.
    \item Ý c) Sai: Ngày thứ hai (29 m) ít hơn ngày thứ nhất (38 m).
    \item Ý d) Sai: Còn 28 m, bán 29 m nên còn ít hơn.
\end{itemize}
\begin{itemize}
    \item a) \textbf{Đúng}.
    \item b) \textbf{Đúng}.
    \item c) \textbf{Sai}.
    \item d) \textbf{Sai}.
\end{itemize}
}
\end{ex}

\textbf{Phần 3. Trả lời ngắn}
\setcounterpageref{ex}{0}

\begin{ex} %[2D7V34]
Tính giá trị: $58 + 37 - 49$.
\shortans{46}
\loigiai{
\textbf{Đáp số:} $46$\\
$58 + 37 = 95$; $95 - 49 = 46$.
}
\end{ex}

\begin{ex} %[2D7V34]
Tìm x biết: $91 - x = 36 + 18$.
\shortans{37}
\loigiai{
\textbf{Đáp số:} $37$\\
$36 + 18 = 54 \Rightarrow x = 91 - 54 = 37$.
}
\end{ex}

\begin{ex} %[2D7V34]
Một đàn vịt có 65 con, người ta bán đi 28 con. Hỏi đàn vịt còn lại bao nhiêu con?
\shortans{37}
\loigiai{
\textbf{Đáp số:} $37$\\
$65 - 28 = 37$ con.
}
\end{ex}

\begin{ex} %[2D7V34]
Tổng của hai số là 82. Nếu bớt ở số thứ nhất đi 15 đơn vị thì tổng mới bằng bao nhiêu?
\shortans{67}
\loigiai{
\textbf{Đáp số:} $67$\\
$82 - 15 = 67$.
}
\end{ex}

\begin{ex} %[2D7V34]
Tìm một số biết rằng lấy số đó trừ đi 38 rồi cộng thêm 19 thì được 52.
\shortans{71}
\loigiai{
\textbf{Đáp số:} $71$\\
Số đó: $52 - 19 + 38 = 71$.
}
\end{ex}

\begin{ex} %[2D7V34]
Hai bạn có tất cả 54 viên bi. Nếu bạn An cho bạn Bình 6 viên bi thì hai bạn có số bi bằng nhau. Hỏi lúc đầu bạn An có bao nhiêu viên bi?
\shortans{33}
\loigiai{
\textbf{Đáp số:} $33$\\
Khi bằng nhau mỗi bạn có: $54 : 2 = 27$ viên. Ban đầu An có: $27 + 6 = 33$ viên.
}
\end{ex}

\newpage

\TaoPhieuKiemTra{Chủ đề 7 -- Ôn tập học kì 1}{35}{Ôn tập đo lường}{https://loptoanthaylinh.vn/thi/chuyen-de/2/toan2_cd_b35}
\textbf{Phần 1. Trắc nghiệm lựa chọn}
\setcounterpageref{ex}{0}

\begin{ex} %[2H7V35]
Tính chu vi hình tam giác có độ dài ba cạnh lần lượt là $18\text{ cm}, 25\text{ cm}, 27\text{ cm}$.
\choice
{$68\text{ cm}$}
{$69\text{ cm}$}
{\True $70\text{ cm}$}
{$71\text{ cm}$}
\loigiai{
$18 + 25 + 27 = 70\text{ (cm)}$.
\textbf{Đáp án của câu hỏi là C}
}
\end{ex}

\begin{ex} %[2H7V35]
Một thùng chứa $45\text{ L}$ dầu, người ta đã dùng hết $18\text{ L}$. Trong thùng còn lại bao nhiêu lít dầu?
\choice
{$26\text{ L}$}
{\True $27\text{ L}$}
{$28\text{ L}$}
{$25\text{ L}$}
\loigiai{
$45 - 18 = 27\text{ (L)}$.
\textbf{Đáp án của câu hỏi là B}
}
\end{ex}

\begin{ex} %[2H7V35]
Đoạn thẳng $AB$ dài $3\text{ dm } 5\text{ cm}$. Đổi độ dài đoạn thẳng $AB$ ra xăng-ti-mét là
\choice
{$305\text{ cm}$}
{$53\text{ cm}$}
{$30\text{ cm}$}
{\True $35\text{ cm}$}
\loigiai{
$3\text{ dm } 5\text{ cm} = 35\text{ cm}$.
\textbf{Đáp án của câu hỏi là D}
}
\end{ex}

\begin{ex} %[2H7V35]
Bao ngô nặng $42\text{ kg}$, bao gạo nặng hơn bao ngô $15\text{ kg}$. Bao gạo nặng bao nhiêu ki-lô-gam?
\choice
{\True $57\text{ kg}$}
{$55\text{ kg}$}
{$56\text{ kg}$}
{$58\text{ kg}$}
\loigiai{
$42 + 15 = 57\text{ (kg)}$.
\textbf{Đáp án của câu hỏi là A}
}
\end{ex}

\begin{ex} %[2H7V35]
Đường gấp khúc gồm $4$ đoạn thẳng bằng nhau, mỗi đoạn dài $8\text{ cm}$. Chiều dài đường gấp khúc là
\choice
{$28\text{ cm}$}
{$30\text{ cm}$}
{\True $32\text{ cm}$}
{$36\text{ cm}$}
\loigiai{
$8 \times 4 = 32\text{ (cm)}$.
\textbf{Đáp án của câu hỏi là C}
}
\end{ex}

\begin{ex} %[2H7V35]
Một chiếc đồng hồ chỉ $8$ giờ $15$ phút sáng. Sau $45$ phút nữa là mấy giờ?
\choice
{$8$ giờ $45$ phút}
{\True $9$ giờ sáng}
{$9$ giờ $15$ phút}
{$10$ giờ}
\loigiai{
$15$ phút $+ 45$ phút $= 60$ phút $= 1$ giờ. $8$ giờ $+ 1$ giờ $= 9$ giờ sáng.
\textbf{Đáp án của câu hỏi là B}
}
\end{ex}

\begin{ex} %[2H7V35]
Hình tứ giác có độ dài các cạnh là $15\text{ cm}, 18\text{ cm}, 22\text{ cm}, 25\text{ cm}$. Chu vi hình tứ giác là
\choice
{$78\text{ cm}$}
{$79\text{ cm}$}
{$82\text{ cm}$}
{\True $80\text{ cm}$}
\loigiai{
$(15 + 25) + (18 + 22) = 40 + 40 = 80\text{ (cm)}$.
\textbf{Đáp án của câu hỏi là D}
}
\end{ex}

\begin{ex} %[2H7V35]
Có $3$ can dầu loại $5\text{ L}$ và $1$ can loại $10\text{ L}$. Tổng số lít dầu là
\choice
{\True $25\text{ L}$}
{$20\text{ L}$}
{$30\text{ L}$}
{$15\text{ L}$}
\loigiai{
$5 \times 3 + 10 = 15 + 10 = 25\text{ (L)}$.
\textbf{Đáp án của câu hỏi là A}
}
\end{ex}

\begin{ex} %[2H7V35]
Con ngỗng nặng $6\text{ kg}$, con vịt nhẹ hơn con ngỗng $4\text{ kg}$. Cả hai con nặng bao nhiêu ki-lô-gam?
\choice
{$7\text{ kg}$}
{$9\text{ kg}$}
{\True $8\text{ kg}$}
{$10\text{ kg}$}
\loigiai{
Con vịt nặng: $6 - 4 = 2\text{ (kg)}$. Cả hai con nặng: $6 + 2 = 8\text{ (kg)}$.
\textbf{Đáp án của câu hỏi là C}
}
\end{ex}

\begin{ex} %[2H7V35]
Nếu ngày $1$ tháng $11$ là thứ Năm thì ngày $15$ tháng $11$ là thứ mấy?
\choice
{Thứ Tư}
{\True Thứ Năm}
{Thứ Sáu}
{Thứ Bảy}
\loigiai{
Từ ngày 1 đến ngày 15 là 14 ngày (đúng 2 tuần lễ). Do đó vẫn là thứ Năm.
\textbf{Đáp án của câu hỏi là B}
}
\end{ex}

\begin{ex} %[2H7V35]
Hình vuông có cạnh $20\text{ cm}$. Chu vi của hình vuông đó là
\choice
{$60\text{ cm}$}
{$70\text{ cm}$}
{$100\text{ cm}$}
{\True $80\text{ cm}$}
\loigiai{
$20 \times 4 = 80\text{ (cm)}$.
\textbf{Đáp án của câu hỏi là D}
}
\end{ex}

\begin{ex} %[2H7V35]
Một sợi dây dài $1\text{ m}$. Cắt bớt đi $35\text{ cm}$. Sợi dây còn lại dài bao nhiêu xăng-ti-mét?
\choice
{\True $65\text{ cm}$}
{$60\text{ cm}$}
{$75\text{ cm}$}
{$70\text{ cm}$}
\loigiai{
$1\text{ m} = 100\text{ cm}$. Còn lại: $100 - 35 = 65\text{ (cm)}$.
\textbf{Đáp án của câu hỏi là A}
}
\end{ex}

\textbf{Phần 2. Trắc nghiệm đúng sai}
\setcounterpageref{ex}{0}

\begin{ex} %[2H7V35]
Xét tính đúng sai của các nhận định về đại lượng và hình học:
\choiceTF
{\True Đổi 1 m = 10 dm = 100 cm.}
{1 kg bông nhẹ hơn 1 kg sắt.}
{\True Chu vi của một hình tam giác bằng tổng độ dài 3 cạnh của nó.}
{Một ngày có 12 giờ.}
\loigiai{
\begin{itemize}
    \item Ý a) Đúng: Quy đổi đơn vị đo độ dài chuẩn.
    \item Ý b) Sai: Cả hai đều nặng đúng 1 kg nên bằng nhau.
    \item Ý c) Đúng: Đúng theo định nghĩa.
    \item Ý d) Sai: Một ngày có 24 giờ.
\end{itemize}
\begin{itemize}
    \item a) \textbf{Đúng}.
    \item b) \textbf{Sai}.
    \item c) \textbf{Đúng}.
    \item d) \textbf{Sai}.
\end{itemize}
}
\end{ex}

\begin{ex} %[2H7V35]
Một bình nước chứa 18 L nước, người ta rót vào 3 can, mỗi can 4 L:
\choiceTF
{\True Số nước đã rót ra là 12 L.}
{\True Trong bình còn lại 6 L nước.}
{Số nước còn lại nhiều hơn số nước đã rót ra.}
{\True Lượng nước còn lại có thể rót đầy thêm 1 can 4 L và 1 can 2 L.}
\loigiai{
\begin{itemize}
    \item Ý a) Đúng: $4 \times 3 = 12\text{ L}$.
    \item Ý b) Đúng: $18 - 12 = 6\text{ L}$.
    \item Ý c) Sai: Còn 6 L, đã rót 12 L nên còn lại ít hơn.
    \item Ý d) Đúng: $4 + 2 = 6\text{ L}$. Đúng.
\end{itemize}
\begin{itemize}
    \item a) \textbf{Đúng}.
    \item b) \textbf{Đúng}.
    \item c) \textbf{Sai}.
    \item d) \textbf{Đúng}.
\end{itemize}
}
\end{ex}

\textbf{Phần 3. Trả lời ngắn}
\setcounterpageref{ex}{0}

\begin{ex} %[2H7V35]
Tính chu vi hình tam giác có độ dài các cạnh là 17 cm, 24 cm, 29 cm.
\shortans{70}
\loigiai{
\textbf{Đáp số:} $70$\\
$17 + 24 + 29 = 70\text{ cm}$.
}
\end{ex}

\begin{ex} %[2H7V35]
Một thùng chứa 50 L dầu, người ta múc ra 28 L dầu. Hỏi trong thùng còn lại bao nhiêu lít dầu?
\shortans{22}
\loigiai{
\textbf{Đáp số:} $22$\\
$50 - 28 = 22\text{ L}$.
}
\end{ex}

\begin{ex} %[2H7V35]
Một bao gạo nặng 45 kg, một bao ngô nặng 38 kg. Hỏi cả hai bao nặng bao nhiêu ki-lô-gam?
\shortans{83}
\loigiai{
\textbf{Đáp số:} $83$\\
$45 + 38 = 83\text{ kg}$.
}
\end{ex}

\begin{ex} %[2H7V35]
Đường gấp khúc gồm 4 đoạn thẳng có độ dài bằng nhau, mỗi đoạn dài 7 cm. Tính độ dài đường gấp khúc đó.
\shortans{28}
\loigiai{
\textbf{Đáp số:} $28$\\
$7 \times 4 = 28\text{ cm}$.
}
\end{ex}

\begin{ex} %[2H7V35]
Một buổi học bắt đầu lúc 7 giờ sáng và kết thúc lúc 11 giờ trưa. Hỏi buổi học kéo dài bao nhiêu giờ?
\shortans{4}
\loigiai{
\textbf{Đáp số:} $4$\\
$11 - 7 = 4$ giờ.
}
\end{ex}

\begin{ex} %[2H7V35]
Đoạn thẳng MN dài 1 m, đoạn thẳng PQ dài 48 cm. Hỏi đoạn thẳng MN dài hơn đoạn thẳng PQ bao nhiêu xăng-ti-mét?
\shortans{52}
\loigiai{
\textbf{Đáp số:} $52$\\
$100 - 48 = 52\text{ cm}$.
}
\end{ex}

\newpage

\TaoPhieuKiemTra{Chủ đề 7 -- Ôn tập học kì 1}{36}{Ôn tập chung}{https://loptoanthaylinh.vn/thi/chuyen-de/2/toan2_cd_b36}
\textbf{Phần 1. Trắc nghiệm lựa chọn}
\setcounterpageref{ex}{0}

\begin{ex} %[2D7V36]
Tính giá trị của biểu thức: $72 - 38 + 26$.
\choice
{$58$}
{\True $60$}
{$59$}
{$61$}
\loigiai{
$72 - 38 = 34$; $34 + 26 = 60$.
\textbf{Đáp án của câu hỏi là B}
}
\end{ex}

\begin{ex} %[2D7V36]
Một lớp học có $36$ học sinh. Trong đó có $19$ bạn nữ. Số bạn nam của lớp là
\choice
{\True $17$ bạn}
{$16$ bạn}
{$18$ bạn}
{$15$ bạn}
\loigiai{
Số bạn nam: $36 - 19 = 17$ (bạn).
\textbf{Đáp án của câu hỏi là A}
}
\end{ex}

\begin{ex} %[2D7V36]
Số bé nhất có hai chữ số mà tổng hai chữ số của nó bằng $9$ là
\choice
{$27$}
{$36$}
{\True $18$}
{$90$}
\loigiai{
Chữ số hàng chục nhỏ nhất là 1, hàng đơn vị là $9 - 1 = 8$. Số đó là 18.
\textbf{Đáp án của câu hỏi là C}
}
\end{ex}

\begin{ex} %[2D7V36]
Chu vi hình tứ giác có độ dài các cạnh lần lượt là $14\text{ cm}, 18\text{ cm}, 22\text{ cm}, 26\text{ cm}$ là
\choice
{$78\text{ cm}$}
{$79\text{ cm}$}
{$82\text{ cm}$}
{\True $80\text{ cm}$}
\loigiai{
$(14 + 26) + (18 + 22) = 40 + 40 = 80\text{ (cm)}$.
\textbf{Đáp án của câu hỏi là D}
}
\end{ex}

\begin{ex} %[2D7V36]
Can thứ nhất có $18\text{ L}$ dầu, can thứ hai có nhiều hơn can thứ nhất $6\text{ L}$ dầu. Cả hai can có bao nhiêu lít dầu?
\choice
{$40\text{ L}$}
{\True $42\text{ L}$}
{$44\text{ L}$}
{$24\text{ L}$}
\loigiai{
Can thứ hai có: $18 + 6 = 24\text{ (L)}$. Cả hai can có: $18 + 24 = 42\text{ (L)}$.
\textbf{Đáp án của câu hỏi là B}
}
\end{ex}

\begin{ex} %[2D7V36]
Tìm $x$ biết: $x + 27 = 83 - 19$.
\choice
{\True $37$}
{$35$}
{$36$}
{$38$}
\loigiai{
Vế phải: $83 - 19 = 64$. Ta có: $x = 64 - 27 = 37$.
\textbf{Đáp án của câu hỏi là A}
}
\end{ex}

\begin{ex} %[2D7V36]
Nếu thứ Bảy tuần này là ngày $16$ thì thứ Bảy tuần sau là ngày mấy?
\choice
{$22$}
{$24$}
{\True $23$}
{$25$}
\loigiai{
$16 + 7 = 23$.
\textbf{Đáp án của câu hỏi là C}
}
\end{ex}

\begin{ex} %[2D7V36]
Một sợi dây dài $65\text{ dm}$. Cắt bớt đi $28\text{ dm}$. Chiều dài đoạn dây còn lại là
\choice
{$35\text{ dm}$}
{$36\text{ dm}$}
{$38\text{ dm}$}
{\True $37\text{ dm}$}
\loigiai{
$65 - 28 = 37\text{ (dm)}$.
\textbf{Đáp án của câu hỏi là D}
}
\end{ex}

\begin{ex} %[2D7V36]
Tìm một số biết rằng lấy số đó trừ đi $35$ rồi cộng với $18$ thì được $60$.
\choice
{$75$}
{\True $77$}
{$76$}
{$78$}
\loigiai{
Tính ngược từ cuối: $60 - 18 + 35 = 42 + 35 = 77$.
\textbf{Đáp án của câu hỏi là B}
}
\end{ex}

\begin{ex} %[2D7V36]
Trong một phép trừ, số bị trừ là $90$, số trừ là $47$. Hiệu của phép trừ đó là
\choice
{\True $43$}
{$41$}
{$42$}
{$44$}
\loigiai{
$90 - 47 = 43$.
\textbf{Đáp án của câu hỏi là A}
}
\end{ex}

\begin{ex} %[2D7V36]
Đồng hồ chỉ $20$ giờ $15$ phút, tức là mấy giờ đêm?
\choice
{$7$ giờ $15$ phút tối}
{$9$ giờ $15$ phút tối}
{\True $8$ giờ $15$ phút tối}
{$8$ giờ rưỡi tối}
\loigiai{
$20$ giờ $= 8$ giờ tối. Vậy là $8$ giờ $15$ phút tối.
\textbf{Đáp án của câu hỏi là C}
}
\end{ex}

\begin{ex} %[2D7V36]
Bác Nam có $62\text{ kg}$ thóc, bác bán đi $38\text{ kg}$ thóc. Bác Nam còn lại bao nhiêu ki-lô-gam thóc?
\choice
{$22\text{ kg}$}
{$23\text{ kg}$}
{$25\text{ kg}$}
{\True $24\text{ kg}$}
\loigiai{
$62 - 38 = 24\text{ (kg)}$.
\textbf{Đáp án của câu hỏi là D}
}
\end{ex}

\textbf{Phần 2. Trắc nghiệm đúng sai}
\setcounterpageref{ex}{0}

\begin{ex} %[2D7V36]
Xét tính đúng sai của các bài toán tổng hợp toàn bộ Học kì 1:
\choiceTF
{\True Tổng của số lẻ lớn nhất có 2 chữ số và số chẵn nhỏ nhất có 2 chữ số là 109.}
{Đoạn thẳng dài 1 dm dài hơn đoạn thẳng dài 12 cm.}
{\True Trong một phép trừ, nếu số bị trừ bằng số trừ thì hiệu bằng 0.}
{Một năm chỉ có 3 tháng có 31 ngày.}
\loigiai{
\begin{itemize}
    \item Ý a) Đúng: $99 + 10 = 109$. Đúng.
    \item Ý b) Sai: $1\text{ dm} = 10\text{ cm} < 12\text{ cm}$. Sai.
    \item Ý c) Đúng: Đúng.
    \item Ý d) Sai: Có tất cả 7 tháng có 31 ngày.
\end{itemize}
\begin{itemize}
    \item a) \textbf{Đúng}.
    \item b) \textbf{Sai}.
    \item c) \textbf{Đúng}.
    \item d) \textbf{Sai}.
\end{itemize}
}
\end{ex}

\begin{ex} %[2D7V36]
Một giá sách có hai tầng, tầng trên có 48 quyển, tầng dưới có 36 quyển:
\choiceTF
{\True Cả hai tầng có tất cả 84 quyển sách.}
{\True Tầng trên nhiều hơn tầng dưới 12 quyển sách.}
{Chuyển 12 quyển từ tầng trên xuống tầng dưới thì hai tầng bằng nhau.}
{\True Nếu chuyển 6 quyển từ tầng trên xuống tầng dưới thì mỗi tầng có 42 quyển.}
\loigiai{
\begin{itemize}
    \item Ý a) Đúng: $48 + 36 = 84$ quyển.
    \item Ý b) Đúng: $48 - 36 = 12$ quyển.
    \item Ý c) Sai: Chỉ cần chuyển $12 : 2 = 6$ quyển là bằng nhau.
    \item Ý d) Đúng: $48 - 6 = 42$ và $36 + 6 = 42$. Đúng.
\end{itemize}
\begin{itemize}
    \item a) \textbf{Đúng}.
    \item b) \textbf{Đúng}.
    \item c) \textbf{Sai}.
    \item d) \textbf{Đúng}.
\end{itemize}
}
\end{ex}

\textbf{Phần 3. Trả lời ngắn}
\setcounterpageref{ex}{0}

\begin{ex} %[2D7V36]
Tính giá trị biểu thức: $48 + 36 - 29$.
\shortans{55}
\loigiai{
\textbf{Đáp số:} $55$\\
$48 + 36 = 84$; $84 - 29 = 55$.
}
\end{ex}

\begin{ex} %[2D7V36]
Một đàn gà có 63 con, người ta bán đi 37 con. Hỏi đàn gà còn lại bao nhiêu con?
\shortans{26}
\loigiai{
\textbf{Đáp số:} $26$\\
$63 - 37 = 26$ con.
}
\end{ex}

\begin{ex} %[2D7V36]
Tìm x biết: $85 - x = 27 + 19$.
\shortans{39}
\loigiai{
\textbf{Đáp số:} $39$\\
$27 + 19 = 46 \Rightarrow x = 85 - 46 = 39$.
}
\end{ex}

\begin{ex} %[2D7V36]
Tính chu vi hình tứ giác có độ dài các cạnh là 15 cm, 25 cm, 18 cm, 22 cm.
\shortans{80}
\loigiai{
\textbf{Đáp số:} $80$\\
$(15 + 25) + (18 + 22) = 40 + 40 = 80\text{ cm}$.
}
\end{ex}

\begin{ex} %[2D7V36]
Can A có 35 L dầu, Can B có ít hơn Can A 8 L dầu. Hỏi cả hai can có bao nhiêu lít dầu?
\shortans{62}
\loigiai{
\textbf{Đáp số:} $62$\\
Can B: $35 - 8 = 27\text{ L}$. Cả hai can: $35 + 27 = 62\text{ L}$.
}
\end{ex}

\begin{ex} %[2D7V36]
Nếu hôm nay là thứ Tư ngày 7 thì thứ Tư của 2 tuần sau là ngày mấy?
\shortans{21}
\loigiai{
\textbf{Đáp số:} $21$\\
$7 + 7 + 7 = 21$.
}
\end{ex}

\newpage


% =========================================================================
% CHỦ ĐỀ 8 - PHÉP NHÂN, PHÉP CHIA
% =========================================================================
\section{Phép nhân, phép chia}

\TaoPhieuKiemTra{Chủ đề 8 -- Phép nhân, phép chia}{37}{Phép nhân}{https://loptoanthaylinh.vn/thi/chuyen-de/2/toan2_cd_b37}
\textbf{Phần 1. Trắc nghiệm lựa chọn}
\setcounterpageref{ex}{0}

\begin{ex} %[2D8V37]
Tổng $5 + 5 + 5 + 5 + 5 + 5$ được viết thành phép nhân nào dưới đây?
\choice
{$5 \times 5$}
{$6 \times 5$}
{\True $5 \times 6$}
{$5 \times 7$}
\loigiai{
Số 5 được lấy 6 lần nên viết là $5 \times 6$.
\textbf{Đáp án của câu hỏi là C}
}
\end{ex}

\begin{ex} %[2D8V37]
Một bàn học có $2$ bạn học sinh ngồi. Hỏi $8$ bàn học như thế có tất cả bao nhiêu bạn học sinh ngồi?
\choice
{\True $16$ bạn}
{$14$ bạn}
{$18$ bạn}
{$10$ bạn}
\loigiai{
$2 \times 8 = 16$ (bạn).
\textbf{Đáp án của câu hỏi là A}
}
\end{ex}

\begin{ex} %[2D8V37]
Phép nhân $4 \times 3$ biểu diễn tổng của
\choice
{$4$ số hạng $3$}
{$4$ cộng với $3$}
{$4$ trừ đi $3$}
{\True $3$ số hạng $4$ cộng lại với nhau}
\loigiai{
$4 \times 3 = 4 + 4 + 4$, tức là 3 số hạng 4 cộng lại.
\textbf{Đáp án của câu hỏi là D}
}
\end{ex}

\begin{ex} %[2D8V37]
Mỗi con thỏ có $4$ cái chân. Hỏi $5$ con thỏ có tất cả bao nhiêu cái chân?
\choice
{$16$ cái chân}
{\True $20$ cái chân}
{$24$ cái chân}
{$15$ cái chân}
\loigiai{
$4 \times 5 = 20$ (cái chân).
\textbf{Đáp án của câu hỏi là B}
}
\end{ex}

\begin{ex} %[2D8V37]
Tính giá trị của biểu thức: $3 \times 4 + 15$.
\choice
{$25$}
{$26$}
{\True $27$}
{$28$}
\loigiai{
Thực hiện nhân trước: $3 \times 4 = 12$; $12 + 15 = 27$.
\textbf{Đáp án của câu hỏi là C}
}
\end{ex}

\begin{ex} %[2D8V37]
Mỗi gói kẹo có $5$ cái kẹo. Bạn Lan mua $4$ gói kẹo và được tặng thêm $2$ cái kẹo nữa. Lan có tất cả bao nhiêu cái kẹo?
\choice
{\True $22$ cái kẹo}
{$20$ cái kẹo}
{$24$ cái kẹo}
{$18$ cái kẹo}
\loigiai{
$5 \times 4 + 2 = 20 + 2 = 22$ (cái kẹo).
\textbf{Đáp án của câu hỏi là A}
}
\end{ex}

\begin{ex} %[2D8V37]
Điền dấu thích hợp: $2 \times 7 \dots 3 \times 5$.
\choice
{$>$}
{$=$}
{Không so sánh được}
{\True $<$}
\loigiai{
$2 \times 7 = 14$; $3 \times 5 = 15$. Vì $14 < 15$ nên điền dấu $<$. 
\textbf{Đáp án của câu hỏi là D}
}
\end{ex}

\begin{ex} %[2D8V37]
Một xe đạp có $2$ bánh xe. Hỏi $9$ xe đạp như thế có bao nhiêu bánh xe?
\choice
{$16$ bánh xe}
{\True $18$ bánh xe}
{$20$ bánh xe}
{$14$ bánh xe}
\loigiai{
$2 \times 9 = 18$ (bánh xe).
\textbf{Đáp án của câu hỏi là B}
}
\end{ex}

\begin{ex} %[2D8V37]
Viết tổng $2 + 2 + 2 + 2 + 2 + 2 + 2$ dưới dạng tích:
\choice
{$2 \times 6$}
{$7 \times 2$}
{\True $2 \times 7$}
{$2 \times 8$}
\loigiai{
Số 2 được lấy 7 lần: $2 \times 7$.
\textbf{Đáp án của câu hỏi là C}
}
\end{ex}

\begin{ex} %[2D8V37]
Tìm một số biết rằng lấy số đó trừ đi $12$ thì bằng tích của $2$ và $6$.
\choice
{\True $24$}
{$22$}
{$26$}
{$20$}
\loigiai{
Tích của 2 và 6 là $2 \times 6 = 12$. Số đó là: $12 + 12 = 24$.
\textbf{Đáp án của câu hỏi là A}
}
\end{ex}

\begin{ex} %[2D8V37]
Một tuần lễ có $7$ ngày. Hỏi $2$ tuần lễ có bao nhiêu ngày?
\choice
{$12$ ngày}
{$16$ ngày}
{$10$ ngày}
{\True $14$ ngày}
\loigiai{
$7 \times 2 = 14$ (ngày).
\textbf{Đáp án của câu hỏi là D}
}
\end{ex}

\begin{ex} %[2D8V37]
Có $4$ hàng cây, mỗi hàng có $5$ cây. Hỏi có tất cả bao nhiêu cây?
\choice
{$18$ cây}
{\True $20$ cây}
{$24$ cây}
{$16$ cây}
\loigiai{
$5 \times 4 = 20$ (cây).
\textbf{Đáp án của câu hỏi là B}
}
\end{ex}

\textbf{Phần 2. Trắc nghiệm đúng sai}
\setcounterpageref{ex}{0}

\begin{ex} %[2D8V37]
Xét tính đúng sai của các nhận định về phép nhân:
\choiceTF
{\True Phép nhân là phép cộng các số hạng bằng nhau.}
{Biểu thức $3 + 3 + 3 + 4$ có thể viết thành $3 \times 4$.}
{\True Kết quả của $2 \times 5$ bằng kết quả của $5 \times 2$.}
{\True Tích của bất kỳ số nào với 1 cũng bằng chính số đó.}
\loigiai{
\begin{itemize}
    \item Ý a) Đúng: Đúng theo bản chất phép nhân.
    \item Ý b) Sai: Các số hạng không bằng nhau (có số 4) nên không thể viết thành $3 \times 4$.
    \item Ý c) Đúng: Tính chất giao hoán của phép nhân: $2 \times 5 = 5 \times 2 = 10$.
    \item Ý d) Đúng: Đúng.
\end{itemize}
\begin{itemize}
    \item a) \textbf{Đúng}.
    \item b) \textbf{Sai}.
    \item c) \textbf{Đúng}.
    \item d) \textbf{Đúng}.
\end{itemize}
}
\end{ex}

\begin{ex} %[2D8V37]
Có 4 đĩa cam, mỗi đĩa có 3 quả cam:
\choiceTF
{\True Tổng số cam là 12 quả.}
{Phép nhân biểu diễn số cam trên là $4 \times 3$.}
{\True Nếu ăn đi 2 quả cam thì còn lại 10 quả cam.}
{Nếu thêm 1 đĩa cam nữa thì có tất cả 16 quả cam.}
\loigiai{
\begin{itemize}
    \item Ý a) Đúng: $3 \times 4 = 12$ quả.
    \item Ý b) Sai: Số quả cam ở mỗi đĩa là 3, được lấy 4 lần nên viết chuẩn là $3 \times 4$.
    \item Ý c) Đúng: $12 - 2 = 10$ quả.
    \item Ý d) Sai: Thêm 1 đĩa có thêm 3 quả: $12 + 3 = 15$ quả (không phải 16 quả).
\end{itemize}
\begin{itemize}
    \item a) \textbf{Đúng}.
    \item b) \textbf{Sai}.
    \item c) \textbf{Đúng}.
    \item d) \textbf{Sai}.
\end{itemize}
}
\end{ex}

\textbf{Phần 3. Trả lời ngắn}
\setcounterpageref{ex}{0}

\begin{ex} %[2D8V37]
Tính giá trị của biểu thức: $2 \times 6 + 18$.
\shortans{30}
\loigiai{
\textbf{Đáp số:} $30$\\
$2 \times 6 = 12$; $12 + 18 = 30$.
}
\end{ex}

\begin{ex} %[2D8V37]
Mỗi con bò có 4 cái chân. Hỏi 6 con bò có tất cả bao nhiêu cái chân?
\shortans{24}
\loigiai{
\textbf{Đáp số:} $24$\\
$4 \times 6 = 24$ cái chân.
}
\end{ex}

\begin{ex} %[2D8V37]
Có 5 bạn học sinh, mỗi bạn có 3 quyển vở. Hỏi cả 5 bạn có tất cả bao nhiêu quyển vở?
\shortans{15}
\loigiai{
\textbf{Đáp số:} $15$\\
$3 \times 5 = 15$ quyển vở.
}
\end{ex}

\begin{ex} %[2D8V37]
Tính nhanh tổng sau bằng phép nhân: $5 + 5 + 5 + 5 + 5$.
\shortans{25}
\loigiai{
\textbf{Đáp số:} $25$\\
$5 \times 5 = 25$.
}
\end{ex}

\begin{ex} %[2D8V37]
Một xe ô tô con có 4 bánh xe. Hỏi 8 xe ô tô con như thế có bao nhiêu bánh xe?
\shortans{32}
\loigiai{
\textbf{Đáp số:} $32$\\
$4 \times 8 = 32$ bánh xe.
}
\end{ex}

\begin{ex} %[2D8V37]
Tìm một số biết rằng lấy số đó chia cho 2 thì được kết quả bằng 7.
\shortans{14}
\loigiai{
\textbf{Đáp số:} $14$\\
$7 \times 2 = 14$.
}
\end{ex}

\newpage

\TaoPhieuKiemTra{Chủ đề 8 -- Phép nhân, phép chia}{38}{Thừa số, tích}{https://loptoanthaylinh.vn/thi/chuyen-de/2/toan2_cd_b38}
\textbf{Phần 1. Trắc nghiệm lựa chọn}
\setcounterpageref{ex}{0}

\begin{ex} %[2D8V38]
Trong phép nhân $5 \times 4 = 20$, số $20$ được gọi là
\choice
{Thừa số}
{\True Tích}
{Số bị chia}
{Tổng}
\loigiai{
20 là kết quả của phép nhân nên được gọi là Tích.
\textbf{Đáp án của câu hỏi là B}
}
\end{ex}

\begin{ex} %[2D8V38]
Biết thừa số thứ nhất là $2$, thừa số thứ hai là $9$. Tích của hai số đó là
\choice
{$16$}
{$17$}
{$20$}
{\True $18$}
\loigiai{
$2 \times 9 = 18$.
\textbf{Đáp án của câu hỏi là D}
}
\end{ex}

\begin{ex} %[2D8V38]
Trong phép nhân có thừa số là $5$, nếu tăng thừa số thứ hai thêm $2$ đơn vị thì tích tăng thêm bao nhiêu?
\choice
{\True $10$ đơn vị}
{$5$ đơn vị}
{$7$ đơn vị}
{$2$ đơn vị}
\loigiai{
Thừa số thứ hai tăng 2 đơn vị thì tích tăng thêm: $5 \times 2 = 10$ đơn vị.
\textbf{Đáp án của câu hỏi là A}
}
\end{ex}

\begin{ex} %[2D8V38]
Tìm một thừa số biết tích bằng $14$ và thừa số kia bằng $2$.
\choice
{$6$}
{$8$}
{\True $7$}
{$9$}
\loigiai{
$14 : 2 = 7$.
\textbf{Đáp án của câu hỏi là C}
}
\end{ex}

\begin{ex} %[2D8V38]
Tích của số chẵn nhỏ nhất có một chữ số khác 0 và số lớn nhất có một chữ số là
\choice
{$16$}
{\True $18$}
{$14$}
{$20$}
\loigiai{
Số chẵn nhỏ nhất khác 0 là 2. Số lớn nhất có 1 chữ số là 9. Tích: $2 \times 9 = 18$.
\textbf{Đáp án của câu hỏi là B}
}
\end{ex}

\begin{ex} %[2D8V38]
Điền số thích hợp: $2 \times \dots = 16$.
\choice
{$6$}
{$7$}
{$9$}
{\True $8$}
\loigiai{
$16 : 2 = 8$.
\textbf{Đáp án của câu hỏi là D}
}
\end{ex}

\begin{ex} %[2D8V38]
Có hai túi kẹo, mỗi túi có $6$ cái kẹo. Hỏi cả hai túi có bao nhiêu cái kẹo?
\choice
{\True $12$ cái kẹo}
{$10$ cái kẹo}
{$14$ cái kẹo}
{$8$ cái kẹo}
\loigiai{
$6 \times 2 = 12$ (cái kẹo).
\textbf{Đáp án của câu hỏi là A}
}
\end{ex}

\begin{ex} %[2D8V38]
Tích của $5$ với số nào cho kết quả bằng $35$?
\choice
{$6$}
{$8$}
{\True $7$}
{$9$}
\loigiai{
$35 : 5 = 7$.
\textbf{Đáp án của câu hỏi là C}
}
\end{ex}

\begin{ex} %[2D8V38]
Trong một phép nhân, nếu một thừa số bằng $0$ thì tích luôn bằng
\choice
{$1$}
{\True $0$}
{Chính thừa số kia}
{Không xác định}
\loigiai{
Bất kỳ số nào nhân với 0 cũng bằng 0.
\textbf{Đáp án của câu hỏi là B}
}
\end{ex}

\begin{ex} %[2D8V38]
Tính giá trị biểu thức: $2 \times 8 - 7$.
\choice
{$8$}
{$10$}
{$7$}
{\True $9$}
\loigiai{
$2 \times 8 = 16$; $16 - 7 = 9$.
\textbf{Đáp án của câu hỏi là D}
}
\end{ex}

\begin{ex} %[2D8V38]
Tìm tích của số liền sau số $1$ và số liền trước số $6$.
\choice
{\True $10$}
{$8$}
{$12$}
{$9$}
\loigiai{
Số liền sau 1 là 2. Số liền trước 6 là 5. Tích: $2 \times 5 = 10$.
\textbf{Đáp án của câu hỏi là A}
}
\end{ex}

\begin{ex} %[2D8V38]
Mỗi bạn nhỏ có $2$ quả bóng bay. Hỏi $7$ bạn nhỏ có bao nhiêu quả bóng bay?
\choice
{$12$ quả}
{$16$ quả}
{\True $14$ quả}
{$9$ quả}
\loigiai{
$2 \times 7 = 14$ (quả).
\textbf{Đáp án của câu hỏi là C}
}
\end{ex}

\textbf{Phần 2. Trắc nghiệm đúng sai}
\setcounterpageref{ex}{0}

\begin{ex} %[2D8V38]
Xét tính đúng sai của các nhận định về thừa số và tích:
\choiceTF
{\True Trong phép nhân a x b = c, a và b là các thừa số, c là tích.}
{Đổi chỗ hai thừa số thì tích sẽ bị thay đổi.}
{\True Tích của 2 số tự nhiên khác 0 luôn lớn hơn hoặc bằng mỗi thừa số.}
{Nếu tích bằng 0 thì cả hai thừa số đều bắt buộc phải bằng 0.}
\loigiai{
\begin{itemize}
    \item Ý a) Đúng: Đúng theo tên gọi.
    \item Ý b) Sai: Tính chất giao hoán: đổi chỗ hai thừa số thì tích không đổi.
    \item Ý c) Đúng: Đúng.
    \item Ý d) Sai: Chỉ cần ít nhất một thừa số bằng 0 là tích bằng 0.
\end{itemize}
\begin{itemize}
    \item a) \textbf{Đúng}.
    \item b) \textbf{Sai}.
    \item c) \textbf{Đúng}.
    \item d) \textbf{Sai}.
\end{itemize}
}
\end{ex}

\begin{ex} %[2D8V38]
Cho phép nhân 5 x 6 = 30:
\choiceTF
{\True Các thừa số là 5 và 6, tích là 30.}
{\True Nếu tăng thừa số thứ nhất thêm 1 đơn vị thành 6 thì tích mới là 36.}
{Tích 30 là một số lẻ.}
{\True Biểu thức 5 x 6 bằng biểu thức 6 x 5.}
\loigiai{
\begin{itemize}
    \item Ý a) Đúng: Đúng.
    \item Ý b) Đúng: $6 \times 6 = 36$. Đúng.
    \item Ý c) Sai: 30 có chữ số tận cùng là 0 nên là số chẵn.
    \item Ý d) Đúng: Đúng.
\end{itemize}
\begin{itemize}
    \item a) \textbf{Đúng}.
    \item b) \textbf{Đúng}.
    \item c) \textbf{Sai}.
    \item d) \textbf{Đúng}.
\end{itemize}
}
\end{ex}

\textbf{Phần 3. Trả lời ngắn}
\setcounterpageref{ex}{0}

\begin{ex} %[2D8V38]
Tìm tích của 2 và 7.
\shortans{14}
\loigiai{
\textbf{Đáp số:} $14$\\
$2 \times 7 = 14$.
}
\end{ex}

\begin{ex} %[2D8V38]
Tìm thừa số chưa biết: x x 5 = 40.
\shortans{8}
\loigiai{
\textbf{Đáp số:} $8$\\
$x = 40 : 5 = 8$.
}
\end{ex}

\begin{ex} %[2D8V38]
Tính giá trị: 2 x 9 + 15.
\shortans{33}
\loigiai{
\textbf{Đáp số:} $33$\\
$2 \times 9 = 18$; $18 + 15 = 33$.
}
\end{ex}

\begin{ex} %[2D8V38]
Mỗi nhóm học sinh có 5 bạn. Hỏi 6 nhóm như thế có tất cả bao nhiêu bạn học sinh?
\shortans{30}
\loigiai{
\textbf{Đáp số:} $30$\\
$5 \times 6 = 30$ bạn.
}
\end{ex}

\begin{ex} %[2D8V38]
Tìm tích của số lớn nhất có một chữ số và số 2.
\shortans{18}
\loigiai{
\textbf{Đáp số:} $18$\\
$9 \times 2 = 18$.
}
\end{ex}

\begin{ex} %[2D8V38]
Một lọ hoa cắm 5 bông hoa. Hỏi 8 lọ hoa như thế có bao nhiêu bông hoa?
\shortans{40}
\loigiai{
\textbf{Đáp số:} $40$\\
$5 \times 8 = 40$ bông hoa.
}
\end{ex}

\newpage

\TaoPhieuKiemTra{Chủ đề 8 -- Phép nhân, phép chia}{39}{Bảng nhân 2}{https://loptoanthaylinh.vn/thi/chuyen-de/2/toan2_cd_b39}
\textbf{Phần 1. Trắc nghiệm lựa chọn}
\setcounterpageref{ex}{0}

\begin{ex} %[2D8V39]
Kết quả của phép tính $2 \times 8$ là
\choice
{\True $16$}
{$14$}
{$18$}
{$15$}
\loigiai{
$2 \times 8 = 16$.
\textbf{Đáp án của câu hỏi là A}
}
\end{ex}

\begin{ex} %[2D8V39]
Một chuồng nuôi $7$ con gà. Hỏi đàn gà đó có tất cả bao nhiêu cái chân?
\choice
{$12$ cái chân}
{$16$ cái chân}
{\True $14$ cái chân}
{$10$ cái chân}
\loigiai{
Mỗi con gà có 2 cái chân: $2 \times 7 = 14$ (cái chân).
\textbf{Đáp án của câu hỏi là C}
}
\end{ex}

\begin{ex} %[2D8V39]
Tính giá trị biểu thức: $2 \times 9 - 8$.
\choice
{$9$}
{\True $10$}
{$11$}
{$8$}
\loigiai{
$2 \times 9 = 18$; $18 - 8 = 10$.
\textbf{Đáp án của câu hỏi là B}
}
\end{ex}

\begin{ex} %[2D8V39]
Mỗi bạn được phát $2$ quyển vở. Hỏi $6$ bạn được phát bao nhiêu quyển vở?
\choice
{$10$ quyển}
{$14$ quyển}
{$16$ quyển}
{\True $12$ quyển}
\loigiai{
$2 \times 6 = 12$ (quyển).
\textbf{Đáp án của câu hỏi là D}
}
\end{ex}

\begin{ex} %[2D8V39]
Tìm $x$ biết: $x : 2 = 9$.
\choice
{\True $18$}
{$16$}
{$17$}
{$20$}
\loigiai{
$x = 9 \times 2 = 18$.
\textbf{Đáp án của câu hỏi là A}
}
\end{ex}

\begin{ex} %[2D8V39]
Điền số thích hợp: $2 \times 5 + 2 \times 3 = 2 \times \dots$
\choice
{$7$}
{$9$}
{\True $8$}
{$6$}
\loigiai{
$2 \times (5 + 3) = 2 \times 8$. Điền 8.
\textbf{Đáp án của câu hỏi là C}
}
\end{ex}

\begin{ex} %[2D8V39]
Một đôi đũa gồm $2$ chiếc đũa. Hỏi $9$ đôi đũa có tất cả bao nhiêu chiếc đũa?
\choice
{$16$ chiếc}
{\True $18$ chiếc}
{$20$ chiếc}
{$14$ chiếc}
\loigiai{
$2 \times 9 = 18$ (chiếc).
\textbf{Đáp án của câu hỏi là B}
}
\end{ex}

\begin{ex} %[2D8V39]
Tính tổng: $2 + 2 + 2 + 2 + 2 + 2 + 10$.
\choice
{$20$}
{$24$}
{$18$}
{\True $22$}
\loigiai{
$2 \times 6 + 10 = 12 + 10 = 22$.
\textbf{Đáp án của câu hỏi là D}
}
\end{ex}

\begin{ex} %[2D8V39]
Lan có $5$ cái kẹo, Mai có số kẹo gấp đôi số kẹo của Lan. Hỏi cả hai bạn có tất cả bao nhiêu cái kẹo?
\choice
{\True $15$ cái kẹo}
{$10$ cái kẹo}
{$12$ cái kẹo}
{$14$ cái kẹo}
\loigiai{
Mai có: $5 \times 2 = 10$ (cái kẹo).\nCả hai bạn có: $5 + 10 = 15$ (cái kẹo).
\textbf{Đáp án của câu hỏi là A}
}
\end{ex}

\begin{ex} %[2D8V39]
Điền dấu thích hợp: $2 \times 6 \dots 2 \times 4 + 5$.
\choice
{$>$}
{$=$}
{\True $<$}
{Không so sánh được}
\loigiai{
Vế trái: $2 \times 6 = 12$. Vế phải: $8 + 5 = 13$. Vì $12 < 13$ nên điền dấu $<$. 
\textbf{Đáp án của câu hỏi là C}
}
\end{ex}

\begin{ex} %[2D8V39]
Có $8$ xe máy, mỗi xe máy có $2$ bánh xe. Hỏi $8$ xe máy có bao nhiêu bánh xe?
\choice
{$14$ bánh xe}
{\True $16$ bánh xe}
{$18$ bánh xe}
{$12$ bánh xe}
\loigiai{
$2 \times 8 = 16$ (bánh xe).
\textbf{Đáp án của câu hỏi là B}
}
\end{ex}

\begin{ex} %[2D8V39]
Tìm một số biết rằng lấy số đó trừ đi $6$ thì bằng $2 \times 7$.
\choice
{$18$}
{$19$}
{$21$}
{\True $20$}
\loigiai{
$2 \times 7 = 14$. Số đó: $14 + 6 = 20$.
\textbf{Đáp án của câu hỏi là D}
}
\end{ex}

\textbf{Phần 2. Trắc nghiệm đúng sai}
\setcounterpageref{ex}{0}

\begin{ex} %[2D8V39]
Xét tính đúng sai của các phép tính trong bảng nhân 2:
\choiceTF
{\True Mọi kết quả trong bảng nhân 2 đều là các số chẵn.}
{Phép tính 2 x 7 có kết quả lớn hơn 2 x 8.}
{\True Hai tích liền kề trong bảng nhân 2 luôn hơn kém nhau 2 đơn vị.}
{\True Biểu thức 2 x 4 + 2 bằng 2 x 5.}
\loigiai{
\begin{itemize}
    \item Ý a) Đúng: Số 2 nhân với bất kì số tự nhiên nào cũng được số chẵn. Đúng.
    \item Ý b) Sai: $2 \times 7 = 14 < 2 \times 8 = 16$.
    \item Ý c) Đúng: Đúng.
    \item Ý d) Đúng: $8 + 2 = 10 = 2 \times 5$. Đúng.
\end{itemize}
\begin{itemize}
    \item a) \textbf{Đúng}.
    \item b) \textbf{Sai}.
    \item c) \textbf{Đúng}.
    \item d) \textbf{Đúng}.
\end{itemize}
}
\end{ex}

\begin{ex} %[2D8V39]
Một đàn vịt có 8 con đang bơi dưới nước:
\choiceTF
{\True Đàn vịt có tất cả 16 cái chân.}
{\True Đàn vịt có tất cả 16 cái cánh.}
{Đàn vịt có tất cả 16 cái mỏ.}
{Nếu có thêm 2 con vịt nữa thì đàn vịt có 18 cái chân.}
\loigiai{
\begin{itemize}
    \item Ý a) Đúng: $2 \times 8 = 16$ cái chân.
    \item Ý b) Đúng: Mỗi con vịt có 2 cái cánh: $2 \times 8 = 16$ cái cánh.
    \item Ý c) Sai: Mỗi con vịt chỉ có 1 cái mỏ nên có 8 cái mỏ.
    \item Ý d) Sai: Thêm 2 con thì có 10 con: $2 \times 10 = 20$ cái chân (không phải 18).
\end{itemize}
\begin{itemize}
    \item a) \textbf{Đúng}.
    \item b) \textbf{Đúng}.
    \item c) \textbf{Sai}.
    \item d) \textbf{Sai}.
\end{itemize}
}
\end{ex}

\textbf{Phần 3. Trả lời ngắn}
\setcounterpageref{ex}{0}

\begin{ex} %[2D8V39]
Tính kết quả của phép nhân: 2 x 9.
\shortans{18}
\loigiai{
\textbf{Đáp số:} $18$\\
$2 \times 9 = 18$.
}
\end{ex}

\begin{ex} %[2D8V39]
Một con chim có 2 cái cánh. Hỏi đàn chim gồm 8 con có tất cả bao nhiêu cái cánh?
\shortans{16}
\loigiai{
\textbf{Đáp số:} $16$\\
$2 \times 8 = 16$ cái cánh.
}
\end{ex}

\begin{ex} %[2D8V39]
Tìm x biết: x - 2 x 6 = 8.
\shortans{20}
\loigiai{
\textbf{Đáp số:} $20$\\
$2 \times 6 = 12 \Rightarrow x = 8 + 12 = 20$.
}
\end{ex}

\begin{ex} %[2D8V39]
Một lớp học có 9 bàn đôi, mỗi bàn có 2 học sinh ngồi. Hỏi lớp học có bao nhiêu học sinh?
\shortans{18}
\loigiai{
\textbf{Đáp số:} $18$\\
$2 \times 9 = 18$ học sinh.
}
\end{ex}

\begin{ex} %[2D8V39]
Tính giá trị của biểu thức: 2 x 7 + 16.
\shortans{30}
\loigiai{
\textbf{Đáp số:} $30$\\
$2 \times 7 = 14$; $14 + 16 = 30$.
}
\end{ex}

\begin{ex} %[2D8V39]
Mai có 4 cái bút chì, số bút chì của Lan gấp đôi số bút chì của Mai. Hỏi cả hai bạn có bao nhiêu cái bút chì?
\shortans{12}
\loigiai{
\textbf{Đáp số:} $12$\\
Lan có: $4 \times 2 = 8$. Cả hai bạn: $4 + 8 = 12$ chiếc.
}
\end{ex}

\newpage

\TaoPhieuKiemTra{Chủ đề 8 -- Phép nhân, phép chia}{40}{Bảng nhân 5}{https://loptoanthaylinh.vn/thi/chuyen-de/2/toan2_cd_b40}
\textbf{Phần 1. Trắc nghiệm lựa chọn}
\setcounterpageref{ex}{0}

\begin{ex} %[2D8V40]
Kết quả của phép tính $5 \times 7$ là
\choice
{$30$}
{$40$}
{$32$}
{\True $35$}
\loigiai{
$5 \times 7 = 35$.
\textbf{Đáp án của câu hỏi là D}
}
\end{ex}

\begin{ex} %[2D8V40]
Mỗi bàn tay có $5$ ngón tay. Hỏi $6$ bàn tay có bao nhiêu ngón tay?
\choice
{$25$ ngón tay}
{\True $30$ ngón tay}
{$35$ ngón tay}
{$20$ ngón tay}
\loigiai{
$5 \times 6 = 30$ (ngón tay).
\textbf{Đáp án của câu hỏi là B}
}
\end{ex}

\begin{ex} %[2D8V40]
Tính giá trị của biểu thức: $5 \times 8 - 15$.
\choice
{$20$}
{$30$}
{\True $25$}
{$35$}
\loigiai{
$5 \times 8 = 40$; $40 - 15 = 25$.
\textbf{Đáp án của câu hỏi là C}
}
\end{ex}

\begin{ex} %[2D8V40]
Một tuần lễ em đi học $5$ ngày. Hỏi trong $4$ tuần lễ em đi học bao nhiêu ngày?
\choice
{\True $20$ ngày}
{$18$ ngày}
{$22$ ngày}
{$25$ ngày}
\loigiai{
$5 \times 4 = 20$ (ngày).
\textbf{Đáp án của câu hỏi là A}
}
\end{ex}

\begin{ex} %[2D8V40]
Các kết quả trong bảng nhân 5 có chữ số tận cùng là
\choice
{Chỉ có chữ số $0$}
{Chỉ có chữ số $5$}
{Các chữ số chẵn}
{\True $0\text{ hoặc } 5$}
\loigiai{
Mọi bội số của 5 đều có tận cùng là 0 hoặc 5.
\textbf{Đáp án của câu hỏi là D}
}
\end{ex}

\begin{ex} %[2D8V40]
Tìm $x$ biết: $x : 5 = 6$.
\choice
{$25$}
{\True $30$}
{$35$}
{$20$}
\loigiai{
$x = 6 \times 5 = 30$.
\textbf{Đáp án của câu hỏi là B}
}
\end{ex}

\begin{ex} %[2D8V40]
Có $5$ túi gạo, mỗi túi nặng $5\text{ kg}$. Hỏi cả $5$ túi nặng bao nhiêu ki-lô-gam?
\choice
{$20\text{ kg}$}
{$30\text{ kg}$}
{\True $25\text{ kg}$}
{$24\text{ kg}$}
\loigiai{
$5 \times 5 = 25\text{ (kg)}$.
\textbf{Đáp án của câu hỏi là C}
}
\end{ex}

\begin{ex} %[2D8V40]
Mỗi can chứa $5\text{ L}$ dầu. Hỏi $9$ can như thế chứa bao nhiêu lít dầu?
\choice
{\True $45\text{ L}$}
{$40\text{ L}$}
{$50\text{ L}$}
{$35\text{ L}$}
\loigiai{
$5 \times 9 = 45\text{ (L)}$.
\textbf{Đáp án của câu hỏi là A}
}
\end{ex}

\begin{ex} %[2D8V40]
Điền số thích hợp: $5 \times 4 + 5 = 5 \times \dots$
\choice
{$4$}
{$6$}
{$7$}
{\True $5$}
\loigiai{
$5 \times 4 + 5 = 5 \times 5$. Vậy điền số 5.
\textbf{Đáp án của câu hỏi là D}
}
\end{ex}

\begin{ex} %[2D8V40]
Một ngôi sao $5$ cánh có $5$ đỉnh. Hỏi $7$ ngôi sao như thế có bao nhiêu đỉnh?
\choice
{$30$ đỉnh}
{\True $35$ đỉnh}
{$40$ đỉnh}
{$25$ đỉnh}
\loigiai{
$5 \times 7 = 35$ (đỉnh).
\textbf{Đáp án của câu hỏi là B}
}
\end{ex}

\begin{ex} %[2D8V40]
Tính nhanh: $5 \times 3 + 5 \times 7$.
\choice
{$45$}
{$55$}
{\True $50$}
{$40$}
\loigiai{
$5 \times (3 + 7) = 5 \times 10 = 50$.
\textbf{Đáp án của câu hỏi là C}
}
\end{ex}

\begin{ex} %[2D8V40]
Bác An nuôi $8$ chuồng thỏ, mỗi chuồng có $5$ con. Bác đã bán đi $12$ con thỏ. Bác còn lại bao nhiêu con thỏ?
\choice
{\True $28$ con}
{$26$ con}
{$30$ con}
{$24$ con}
\loigiai{
Tổng số thỏ: $5 \times 8 = 40$ (con).\nCòn lại: $40 - 12 = 28$ (con).
\textbf{Đáp án của câu hỏi là A}
}
\end{ex}

\textbf{Phần 2. Trắc nghiệm đúng sai}
\setcounterpageref{ex}{0}

\begin{ex} %[2D8V40]
Xét tính đúng sai của các phép tính trong bảng nhân 5:
\choiceTF
{\True Tích của 5 với số chẵn luôn có tận cùng là chữ số 0.}
{\True Tích của 5 với số lẻ luôn có tận cùng là chữ số 5.}
{Phép tính 5 x 6 có kết quả bằng 35.}
{\True Hai tích liên tiếp trong bảng nhân 5 hơn kém nhau 5 đơn vị.}
\loigiai{
\begin{itemize}
    \item Ý a) Đúng: Ví dụ: $5 \times 2 = 10, 5 \times 4 = 20$. Đúng.
    \item Ý b) Đúng: Ví dụ: $5 \times 3 = 15, 5 \times 5 = 25$. Đúng.
    \item Ý c) Sai: $5 \times 6 = 30$ (không phải 35).
    \item Ý d) Đúng: Đúng.
\end{itemize}
\begin{itemize}
    \item a) \textbf{Đúng}.
    \item b) \textbf{Đúng}.
    \item c) \textbf{Sai}.
    \item d) \textbf{Đúng}.
\end{itemize}
}
\end{ex}

\begin{ex} %[2D8V40]
Có 6 gói bánh, mỗi gói có 5 chiếc bánh quy:
\choiceTF
{\True Có tất cả 30 chiếc bánh quy.}
{Nếu chia đều 30 chiếc bánh cho 4 bạn thì mỗi bạn được 8 chiếc.}
{\True Nếu ăn hết 5 chiếc bánh thì số bánh còn lại bằng 5 x 5.}
{Số bánh ban đầu là một số lẻ.}
\loigiai{
\begin{itemize}
    \item Ý a) Đúng: $5 \times 6 = 30$ chiếc.
    \item Ý b) Sai: $30$ không chia hết cho 4 ($30 : 4 = 7$ dư 2).
    \item Ý c) Đúng: $30 - 5 = 25 = 5 \times 5$. Đúng.
    \item Ý d) Sai: 30 là số chẵn.
\end{itemize}
\begin{itemize}
    \item a) \textbf{Đúng}.
    \item b) \textbf{Sai}.
    \item c) \textbf{Đúng}.
    \item d) \textbf{Sai}.
\end{itemize}
}
\end{ex}

\textbf{Phần 3. Trả lời ngắn}
\setcounterpageref{ex}{0}

\begin{ex} %[2D8V40]
Tính kết quả phép tính: 5 x 8.
\shortans{40}
\loigiai{
\textbf{Đáp số:} $40$\\
$5 \times 8 = 40$.
}
\end{ex}

\begin{ex} %[2D8V40]
Một can đựng 5 lít nước mắm. Hỏi 7 can như thế đựng bao nhiêu lít nước mắm?
\shortans{35}
\loigiai{
\textbf{Đáp số:} $35$\\
$5 \times 7 = 35\text{ L}$.
}
\end{ex}

\begin{ex} %[2D8V40]
Tính giá trị: 5 x 6 + 25.
\shortans{55}
\loigiai{
\textbf{Đáp số:} $55$\\
$5 \times 6 = 30$; $30 + 25 = 55$.
}
\end{ex}

\begin{ex} %[2D8V40]
Mỗi bông hoa có 5 cánh hoa. Hỏi một bó gồm 9 bông hoa có tất cả bao nhiêu cánh hoa?
\shortans{45}
\loigiai{
\textbf{Đáp số:} $45$\\
$5 \times 9 = 45$ cánh hoa.
}
\end{ex}

\begin{ex} %[2D8V40]
Tìm số tự nhiên y biết: y : 5 = 8.
\shortans{40}
\loigiai{
\textbf{Đáp số:} $40$\\
$y = 8 \times 5 = 40$.
}
\end{ex}

\begin{ex} %[2D8V40]
Mẹ mua 5 hộp bút chì màu, mỗi hộp có 6 chiếc. Mẹ cho em 8 chiếc. Hỏi mẹ còn lại bao nhiêu chiếc bút chì màu?
\shortans{22}
\loigiai{
\textbf{Đáp số:} $22$\\
$5 \times 6 = 30$; $30 - 8 = 22$ chiếc.
}
\end{ex}

\newpage

\TaoPhieuKiemTra{Chủ đề 8 -- Phép nhân, phép chia}{41}{Phép chia}{https://loptoanthaylinh.vn/thi/chuyen-de/2/toan2_cd_b41}
\textbf{Phần 1. Trắc nghiệm lựa chọn}
\setcounterpageref{ex}{0}

\begin{ex} %[2D8V41]
Có $12$ quả cam chia đều vào $2$ đĩa. Mỗi đĩa có bao nhiêu quả cam?
\choice
{$5$ quả cam}
{$7$ quả cam}
{\True $6$ quả cam}
{$4$ quả cam}
\loigiai{
$12 : 2 = 6$ (quả cam).
\textbf{Đáp án của câu hỏi là C}
}
\end{ex}

\begin{ex} %[2D8V41]
Từ phép nhân $2 \times 6 = 12$, ta viết được phép chia nào dưới đây?
\choice
{$12 : 3 = 4$}
{\True $12 : 2 = 6$}
{$6 : 2 = 3$}
{$12 : 6 = 3$}
\loigiai{
Từ phép nhân $2 \times 6 = 12$ ta có hai phép chia: $12 : 2 = 6$ và $12 : 6 = 2$.
\textbf{Đáp án của câu hỏi là B}
}
\end{ex}

\begin{ex} %[2D8V41]
Có $15$ bông hoa cắm đều vào các lọ, mỗi lọ cắm $5$ bông hoa. Hỏi cắm được bao nhiêu lọ hoa?
\choice
{$4$ lọ hoa}
{$2$ lọ hoa}
{$5$ lọ hoa}
{\True $3$ lọ hoa}
\loigiai{
$15 : 5 = 3$ (lọ hoa).
\textbf{Đáp án của câu hỏi là D}
}
\end{ex}

\begin{ex} %[2D8V41]
Kết quả của phép chia $20 : 2$ là
\choice
{\True $10$}
{$8$}
{$9$}
{$12$}
\loigiai{
$20 : 2 = 10$.
\textbf{Đáp án của câu hỏi là A}
}
\end{ex}

\begin{ex} %[2D8V41]
Một sợi dây dài $18\text{ dm}$ được cắt thành các đoạn bằng nhau, mỗi đoạn dài $2\text{ dm}$. Cắt được bao nhiêu đoạn dây?
\choice
{$8$ đoạn dây}
{$10$ đoạn dây}
{\True $9$ đoạn dây}
{$7$ đoạn dây}
\loigiai{
$18 : 2 = 9$ (đoạn dây).
\textbf{Đáp án của câu hỏi là C}
}
\end{ex}

\begin{ex} %[2D8V41]
Tính giá trị biểu thức: $14 : 2 + 18$.
\choice
{$24$}
{\True $25$}
{$26$}
{$23$}
\loigiai{
$14 : 2 = 7$; $7 + 18 = 25$.
\textbf{Đáp án của câu hỏi là B}
}
\end{ex}

\begin{ex} %[2D8V41]
Có $30$ bạn học sinh chia đều thành các nhóm, mỗi nhóm có $5$ bạn. Hỏi chia được bao nhiêu nhóm?
\choice
{$5$ nhóm}
{$7$ nhóm}
{$8$ nhóm}
{\True $6$ nhóm}
\loigiai{
$30 : 5 = 6$ (nhóm).
\textbf{Đáp án của câu hỏi là D}
}
\end{ex}

\begin{ex} %[2D8V41]
Điền dấu thích hợp: $16 : 2 \dots 25 : 5$.
\choice
{\True $>$}
{$<$}
{$=$}
{Không so sánh được}
\loigiai{
$16 : 2 = 8$; $25 : 5 = 5$. Vì $8 > 5$ nên điền dấu $>$. 
\textbf{Đáp án của câu hỏi là A}
}
\end{ex}

\begin{ex} %[2D8V41]
Có $10$ chiếc bút chì chia đều cho Mai và Nam. Mỗi bạn nhận được bao nhiêu chiếc bút chì?
\choice
{$4$ chiếc}
{$6$ chiếc}
{\True $5$ chiếc}
{$3$ chiếc}
\loigiai{
$10 : 2 = 5$ (chiếc).
\textbf{Đáp án của câu hỏi là C}
}
\end{ex}

\begin{ex} %[2D8V41]
Một hộp bánh có $20$ chiếc bánh. Nếu chia đều cho $5$ bạn thì mỗi bạn được mấy chiếc bánh?
\choice
{$3$ chiếc bánh}
{\True $4$ chiếc bánh}
{$5$ chiếc bánh}
{$6$ chiếc bánh}
\loigiai{
$20 : 5 = 4$ (chiếc bánh).
\textbf{Đáp án của câu hỏi là B}
}
\end{ex}

\begin{ex} %[2D8V41]
Tìm $x$ biết: $x \times 2 = 18$.
\choice
{$8$}
{$7$}
{$10$}
{\True $9$}
\loigiai{
$x = 18 : 2 = 9$.
\textbf{Đáp án của câu hỏi là D}
}
\end{ex}

\begin{ex} %[2D8V41]
Một thanh gỗ dài $16\text{ m}$ cưa thành các khúc gỗ dài $2\text{ m}$. Hỏi phải thực hiện bao nhiêu nhát cưa?
\choice
{\True $7$ nhát cưa}
{$8$ nhát cưa}
{$6$ nhát cưa}
{$9$ nhát cưa}
\loigiai{
Số khúc gỗ: $16 : 2 = 8$ (khúc). Số nhát cưa: $8 - 1 = 7$ (nhát cưa).
\textbf{Đáp án của câu hỏi là A}
}
\end{ex}

\textbf{Phần 2. Trắc nghiệm đúng sai}
\setcounterpageref{ex}{0}

\begin{ex} %[2D8V41]
Xét tính đúng sai của các nhận định về phép chia:
\choiceTF
{\True Từ một phép nhân hai thừa số khác nhau có thể lập được hai phép chia tương ứng.}
{Số 0 có thể làm số chia trong phép chia.}
{\True Một số tự nhiên khác 0 chia cho chính nó luôn bằng 1.}
{Phép tính 18 : 2 có kết quả là 8.}
\loigiai{
\begin{itemize}
    \item Ý a) Đúng: Đúng: $a \times b = c \Rightarrow c : a = b$ và $c : b = a$.
    \item Ý b) Sai: Không thể chia cho số 0.
    \item Ý c) Đúng: $a : a = 1$. Đúng.
    \item Ý d) Sai: $18 : 2 = 9$ (không phải 8).
\end{itemize}
\begin{itemize}
    \item a) \textbf{Đúng}.
    \item b) \textbf{Sai}.
    \item c) \textbf{Đúng}.
    \item d) \textbf{Sai}.
\end{itemize}
}
\end{ex}

\begin{ex} %[2D8V41]
Có 20 cái kẹo chia đều cho 5 bạn:
\choiceTF
{\True Mỗi bạn nhận được 4 cái kẹo.}
{\True Nếu có thêm 5 cái kẹo nữa thì mỗi bạn nhận được 5 cái kẹo.}
{Nếu chia đều cho 2 bạn thì mỗi bạn được 8 cái kẹo.}
{\True Số kẹo của mỗi bạn bằng một phần năm tổng số kẹo.}
\loigiai{
\begin{itemize}
    \item Ý a) Đúng: $20 : 5 = 4$ cái kẹo.
    \item Ý b) Đúng: $25 : 5 = 5$ cái kẹo.
    \item Ý c) Sai: $20 : 2 = 10$ cái kẹo (không phải 8).
    \item Ý d) Đúng: Đúng.
\end{itemize}
\begin{itemize}
    \item a) \textbf{Đúng}.
    \item b) \textbf{Đúng}.
    \item c) \textbf{Sai}.
    \item d) \textbf{Đúng}.
\end{itemize}
}
\end{ex}

\textbf{Phần 3. Trả lời ngắn}
\setcounterpageref{ex}{0}

\begin{ex} %[2D8V41]
Tính kết quả: 18 : 2.
\shortans{9}
\loigiai{
\textbf{Đáp số:} $9$\\
$18 : 2 = 9$.
}
\end{ex}

\begin{ex} %[2D8V41]
Có 25 lít dầu chia đều vào 5 can. Hỏi mỗi can chứa bao nhiêu lít dầu?
\shortans{5}
\loigiai{
\textbf{Đáp số:} $5$\\
$25 : 5 = 5\text{ L}$.
}
\end{ex}

\begin{ex} %[2D8V41]
Tính giá trị: 16 : 2 + 35.
\shortans{43}
\loigiai{
\textbf{Đáp số:} $43$\\
$16 : 2 = 8$; $8 + 35 = 43$.
}
\end{ex}

\begin{ex} %[2D8V41]
Một cuộn dây dài 14 m được cắt thành các đoạn dài 2 m. Hỏi cắt được bao nhiêu đoạn dây?
\shortans{7}
\loigiai{
\textbf{Đáp số:} $7$\\
$14 : 2 = 7$ đoạn.
}
\end{ex}

\begin{ex} %[2D8V41]
Tìm x biết: x x 5 = 35.
\shortans{7}
\loigiai{
\textbf{Đáp số:} $7$\\
$x = 35 : 5 = 7$.
}
\end{ex}

\begin{ex} %[2D8V41]
Có 20 kg gạo chia đều vào 2 bao. Hỏi mỗi bao có bao nhiêu ki-lô-gam gạo?
\shortans{10}
\loigiai{
\textbf{Đáp số:} $10$\\
$20 : 2 = 10\text{ kg}$.
}
\end{ex}

\newpage

\TaoPhieuKiemTra{Chủ đề 8 -- Phép nhân, phép chia}{42}{Số bị chia, số chia, thương}{https://loptoanthaylinh.vn/thi/chuyen-de/2/toan2_cd_b42}
\textbf{Phần 1. Trắc nghiệm lựa chọn}
\setcounterpageref{ex}{0}

\begin{ex} %[2D8V42]
Trong phép chia $18 : 2 = 9$, số $18$ được gọi là
\choice
{Số chia}
{\True Số bị chia}
{Thương}
{Số hạng}
\loigiai{
18 là số đứng trước dấu chia nên được gọi là Số bị chia.
\textbf{Đáp án của câu hỏi là B}
}
\end{ex}

\begin{ex} %[2D8V42]
Biết số bị chia là $35$, số chia là $5$. Thương là
\choice
{\True $7$}
{$6$}
{$8$}
{$9$}
\loigiai{
Thương $= 35 : 5 = 7$.
\textbf{Đáp án của câu hỏi là A}
}
\end{ex}

\begin{ex} %[2D8V42]
Tìm số bị chia biết số chia là $2$ và thương là $8$.
\choice
{$14$}
{$18$}
{\True $16$}
{$12$}
\loigiai{
Số bị chia = Thương $\times$ Số chia = $8 \times 2 = 16$.
\textbf{Đáp án của câu hỏi là C}
}
\end{ex}

\begin{ex} %[2D8V42]
Tìm số chia biết số bị chia là $20$ và thương là $4$.
\choice
{$4$}
{$6$}
{$3$}
{\True $5$}
\loigiai{
Số chia = Số bị chia : Thương = $20 : 4 = 5$.
\textbf{Đáp án của câu hỏi là D}
}
\end{ex}

\begin{ex} %[2D8V42]
Trong phép chia có số bị chia bằng số chia (khác 0), thương luôn bằng
\choice
{$0$}
{\True $1$}
{Chính số đó}
{Không xác định}
\loigiai{
Một số khác 0 chia cho chính nó luôn có thương bằng 1.
\textbf{Đáp án của câu hỏi là B}
}
\end{ex}

\begin{ex} %[2D8V42]
Thương của số lớn nhất có một chữ số và số $1$ là
\choice
{\True $9$}
{$1$}
{$8$}
{$10$}
\loigiai{
Số lớn nhất có 1 chữ số là 9. Thương: $9 : 1 = 9$.
\textbf{Đáp án của câu hỏi là A}
}
\end{ex}

\begin{ex} %[2D8V42]
Một phép chia có thương bằng $5$, nếu giữ nguyên số chia là $2$ và tăng số bị chia thêm $4$ đơn vị thì thương mới là
\choice
{$6$}
{$8$}
{\True $7$}
{$9$}
\loigiai{
Số bị chia ban đầu: $5 \times 2 = 10$. Số bị chia mới: $10 + 4 = 14$. Thương mới: $14 : 2 = 7$.
\textbf{Đáp án của câu hỏi là C}
}
\end{ex}

\begin{ex} %[2D8V42]
Điền số thích hợp: $\dots : 5 = 8$.
\choice
{$35$}
{$45$}
{$30$}
{\True $40$}
\loigiai{
Số bị chia $= 8 \times 5 = 40$.
\textbf{Đáp án của câu hỏi là D}
}
\end{ex}

\begin{ex} %[2D8V42]
Có $16\text{ L}$ nước mắm chia đều vào các can, mỗi can $2\text{ L}$. Số can nước mắm là
\choice
{$7$ can}
{\True $8$ can}
{$9$ can}
{$6$ can}
\loigiai{
$16 : 2 = 8$ (can).
\textbf{Đáp án của câu hỏi là B}
}
\end{ex}

\begin{ex} %[2D8V42]
Trong một phép chia, số bị chia là số chẵn nhỏ nhất có hai chữ số, số chia là $2$. Thương là
\choice
{\True $5$}
{$4$}
{$6$}
{$7$}
\loigiai{
Số chẵn nhỏ nhất có 2 chữ số là 10. Thương: $10 : 2 = 5$.
\textbf{Đáp án của câu hỏi là A}
}
\end{ex}

\begin{ex} %[2D8V42]
Tìm $y$ biết: $45 : y = 5$.
\choice
{$8$}
{$7$}
{\True $9$}
{$6$}
\loigiai{
$y = 45 : 5 = 9$.
\textbf{Đáp án của câu hỏi là C}
}
\end{ex}

\begin{ex} %[2D8V42]
Một lớp học có $35$ học sinh được chia đều thành các tổ, mỗi tổ có $5$ bạn. Hỏi có bao nhiêu tổ?
\choice
{$6$ tổ}
{$8$ tổ}
{$5$ tổ}
{\True $7$ tổ}
\loigiai{
$35 : 5 = 7$ (tổ).
\textbf{Đáp án của câu hỏi là D}
}
\end{ex}

\textbf{Phần 2. Trắc nghiệm đúng sai}
\setcounterpageref{ex}{0}

\begin{ex} %[2D8V42]
Xét tính đúng sai của các nhận định về thành phần phép chia:
\choiceTF
{\True Công thức: Số bị chia = Thương x Số chia.}
{\True Công thức: Số chia = Số bị chia : Thương (với thương khác 0).}
{Trong phép chia, số bị chia luôn phải nhỏ hơn số chia.}
{Thương của 0 chia cho bất kỳ số nào khác 0 luôn bằng chính số đó.}
\loigiai{
\begin{itemize}
    \item Ý a) Đúng: Đúng.
    \item Ý b) Đúng: Đúng.
    \item Ý c) Sai: Số bị chia trong toán lớp 2 luôn lớn hơn hoặc bằng số chia.
    \item Ý d) Sai: $0 : a = 0$ (bằng 0 chứ không phải bằng chính số đó).
\end{itemize}
\begin{itemize}
    \item a) \textbf{Đúng}.
    \item b) \textbf{Đúng}.
    \item c) \textbf{Sai}.
    \item d) \textbf{Sai}.
\end{itemize}
}
\end{ex}

\begin{ex} %[2D8V42]
Cho phép chia 40 : 5 = 8:
\choiceTF
{\True Số bị chia là 40, số chia là 5, thương là 8.}
{Nếu gấp số chia lên 2 lần thành 10 thì thương tăng gấp đôi.}
{\True Tích của số chia và thương là 5 x 8 = 40.}
{\True Thương 8 là một số chẵn.}
\loigiai{
\begin{itemize}
    \item Ý a) Đúng: Đúng.
    \item Ý b) Sai: Khi số chia tăng thì thương giảm: $40 : 10 = 4$ (giảm đi một nửa).
    \item Ý c) Đúng: Đúng.
    \item Ý d) Đúng: Đúng.
\end{itemize}
\begin{itemize}
    \item a) \textbf{Đúng}.
    \item b) \textbf{Sai}.
    \item c) \textbf{Đúng}.
    \item d) \textbf{Đúng}.
\end{itemize}
}
\end{ex}

\textbf{Phần 3. Trả lời ngắn}
\setcounterpageref{ex}{0}

\begin{ex} %[2D8V42]
Tìm số bị chia biết số chia là 5 và thương là 7.
\shortans{35}
\loigiai{
\textbf{Đáp số:} $35$\\
$7 \times 5 = 35$.
}
\end{ex}

\begin{ex} %[2D8V42]
Tìm số chia biết số bị chia là 16 và thương là 2.
\shortans{8}
\loigiai{
\textbf{Đáp số:} $8$\\
$16 : 2 = 8$.
}
\end{ex}

\begin{ex} %[2D8V42]
Tính thương của 14 và 2.
\shortans{7}
\loigiai{
\textbf{Đáp số:} $7$\\
$14 : 2 = 7$.
}
\end{ex}

\begin{ex} %[2D8V42]
Có 30 bông hoa cắm đều vào các bình, mỗi bình 5 bông hoa. Hỏi cắm được bao nhiêu bình hoa?
\shortans{6}
\loigiai{
\textbf{Đáp số:} $6$\\
$30 : 5 = 6$ bình.
}
\end{ex}

\begin{ex} %[2D8V42]
Tìm một số biết rằng lấy số đó nhân với 5 thì được 45.
\shortans{9}
\loigiai{
\textbf{Đáp số:} $9$\\
$45 : 5 = 9$.
}
\end{ex}

\begin{ex} %[2D8V42]
Một đàn gà có 18 cái chân. Hỏi đàn gà đó có bao nhiêu con gà?
\shortans{9}
\loigiai{
\textbf{Đáp số:} $9$\\
$18 : 2 = 9$ con gà.
}
\end{ex}

\newpage

\TaoPhieuKiemTra{Chủ đề 8 -- Phép nhân, phép chia}{43}{Bảng chia 2}{https://loptoanthaylinh.vn/thi/chuyen-de/2/toan2_cd_b43}
\textbf{Phần 1. Trắc nghiệm lựa chọn}
\setcounterpageref{ex}{0}

\begin{ex} %[2D8V43]
Kết quả của phép chia $14 : 2$ là
\choice
{$6$}
{$8$}
{\True $7$}
{$5$}
\loigiai{
$14 : 2 = 7$.
\textbf{Đáp án của câu hỏi là C}
}
\end{ex}

\begin{ex} %[2D8V43]
Có $16$ cái bánh chia đều cho hai chị em. Mỗi người được bao nhiêu cái bánh?
\choice
{\True $8$ cái bánh}
{$7$ cái bánh}
{$9$ cái bánh}
{$6$ cái bánh}
\loigiai{
$16 : 2 = 8$ (cái bánh).
\textbf{Đáp án của câu hỏi là A}
}
\end{ex}

\begin{ex} %[2D8V43]
Tìm $x$ biết: $x \times 2 = 12$.
\choice
{$5$}
{$7$}
{$8$}
{\True $6$}
\loigiai{
$x = 12 : 2 = 6$.
\textbf{Đáp án của câu hỏi là D}
}
\end{ex}

\begin{ex} %[2D8V43]
Tính giá trị biểu thức: $18 : 2 + 15$.
\choice
{$22$}
{\True $24$}
{$23$}
{$25$}
\loigiai{
$18 : 2 = 9$; $9 + 15 = 24$.
\textbf{Đáp án của câu hỏi là B}
}
\end{ex}

\begin{ex} %[2D8V43]
Một đoạn dây dài $20\text{ cm}$ gập đôi lại. Chiều dài mỗi nửa đoạn dây là
\choice
{$8\text{ cm}$}
{$12\text{ cm}$}
{\True $10\text{ cm}$}
{$5\text{ cm}$}
\loigiai{
$20 : 2 = 10\text{ (cm)}$.
\textbf{Đáp án của câu hỏi là C}
}
\end{ex}

\begin{ex} %[2D8V43]
Có $18$ học sinh xếp thành $2$ hàng đều nhau. Mỗi hàng có bao nhiêu học sinh?
\choice
{\True $9$ học sinh}
{$8$ học sinh}
{$10$ học sinh}
{$7$ học sinh}
\loigiai{
$18 : 2 = 9$ (học sinh).
\textbf{Đáp án của câu hỏi là A}
}
\end{ex}

\begin{ex} %[2D8V43]
Điền số thích hợp: $16 : 2 = \dots + 3$.
\choice
{$4$}
{$6$}
{$7$}
{\True $5$}
\loigiai{
$16 : 2 = 8$. Ta có: $8 - 3 = 5$. Vậy điền số 5.
\textbf{Đáp án của câu hỏi là D}
}
\end{ex}

\begin{ex} %[2D8V43]
Bác Nam có $10$ con chim bồ câu nhốt đều vào $2$ lồng. Mỗi lồng có bao nhiêu con chim?
\choice
{$4$ con}
{\True $5$ con}
{$6$ con}
{$3$ con}
\loigiai{
$10 : 2 = 5$ (con).
\textbf{Đáp án của câu hỏi là B}
}
\end{ex}

\begin{ex} %[2D8V43]
Tìm một số biết rằng lấy số đó chia cho $2$ rồi cộng với $8$ thì được $15$.
\choice
{$12$}
{$16$}
{\True $14$}
{$18$}
\loigiai{
Tính ngược từ cuối: $15 - 8 = 7$; $7 \times 2 = 14$.
\textbf{Đáp án của câu hỏi là C}
}
\end{ex}

\begin{ex} %[2D8V43]
Thương của số lớn nhất trong bảng nhân $2$ và số $2$ là
\choice
{\True $10$}
{$9$}
{$8$}
{$20$}
\loigiai{
Số lớn nhất trong bảng nhân 2 là 20. Thương: $20 : 2 = 10$.
\textbf{Đáp án của câu hỏi là A}
}
\end{ex}

\begin{ex} %[2D8V43]
Điền dấu thích hợp: $12 : 2 + 4 \dots 18 : 2$.
\choice
{$<$}
{$=$}
{Không so sánh được}
{\True $>$}
\loigiai{
Vế trái: $6 + 4 = 10$. Vế phải: $18 : 2 = 9$. Vì $10 > 9$ nên điền dấu $>$. 
\textbf{Đáp án của câu hỏi là D}
}
\end{ex}

\begin{ex} %[2D8V43]
Một người thợ cưa một cây gỗ dài $8\text{ m}$ thành $2$ đoạn bằng nhau. Mỗi đoạn dài bao nhiêu mét?
\choice
{$3\text{ m}$}
{\True $4\text{ m}$}
{$5\text{ m}$}
{$2\text{ m}$}
\loigiai{
$8 : 2 = 4\text{ (m)}$.
\textbf{Đáp án của câu hỏi là B}
}
\end{ex}

\textbf{Phần 2. Trắc nghiệm đúng sai}
\setcounterpageref{ex}{0}

\begin{ex} %[2D8V43]
Xét tính đúng sai của các phép tính trong bảng chia 2:
\choiceTF
{\True Một số chẵn chia cho 2 luôn được thương là số tự nhiên.}
{Phép chia 15 : 2 là một phép chia hết trong bảng chia 2.}
{\True Trong bảng chia 2, số bị chia lớn nhất là 20, thương lớn nhất là 10.}
{Biểu thức 12 : 2 có kết quả bé hơn biểu thức 10 : 2.}
\loigiai{
\begin{itemize}
    \item Ý a) Đúng: Đúng vì mọi số chẵn đều chia hết cho 2.
    \item Ý b) Sai: 15 là số lẻ nên chia 2 có dư.
    \item Ý c) Đúng: Đúng.
    \item Ý d) Sai: $12 : 2 = 6 > 10 : 2 = 5$. Sai.
\end{itemize}
\begin{itemize}
    \item a) \textbf{Đúng}.
    \item b) \textbf{Sai}.
    \item c) \textbf{Đúng}.
    \item d) \textbf{Sai}.
\end{itemize}
}
\end{ex}

\begin{ex} %[2D8V43]
Mẹ có 14 quả táo chia đều cho hai anh em An và Bình:
\choiceTF
{\True Mỗi người nhận được 7 quả táo.}
{\True Số táo của An bằng số táo của Bình.}
{Nếu mẹ cho thêm mỗi người 1 quả thì mỗi người có 9 quả.}
{\True Tổng số táo của hai người là một số chẵn.}
\loigiai{
\begin{itemize}
    \item Ý a) Đúng: $14 : 2 = 7$ quả.
    \item Ý b) Đúng: Đúng vì chia đều.
    \item Ý c) Sai: Mỗi người có $7 + 1 = 8$ quả (không phải 9 quả).
    \item Ý d) Đúng: 14 là số chẵn. Đúng.
\end{itemize}
\begin{itemize}
    \item a) \textbf{Đúng}.
    \item b) \textbf{Đúng}.
    \item c) \textbf{Sai}.
    \item d) \textbf{Đúng}.
\end{itemize}
}
\end{ex}

\textbf{Phần 3. Trả lời ngắn}
\setcounterpageref{ex}{0}

\begin{ex} %[2D8V43]
Tính kết quả: 16 : 2.
\shortans{8}
\loigiai{
\textbf{Đáp số:} $8$\\
$16 : 2 = 8$.
}
\end{ex}

\begin{ex} %[2D8V43]
Có 18 viên bi chia đều cho 2 bạn. Mỗi bạn được bao nhiêu viên bi?
\shortans{9}
\loigiai{
\textbf{Đáp số:} $9$\\
$18 : 2 = 9$ viên bi.
}
\end{ex}

\begin{ex} %[2D8V43]
Tính giá trị: 12 : 2 + 18.
\shortans{24}
\loigiai{
\textbf{Đáp số:} $24$\\
$12 : 2 = 6$; $6 + 18 = 24$.
}
\end{ex}

\begin{ex} %[2D8V43]
Một đàn gà có tất cả 14 cái chân. Hỏi đàn gà có bao nhiêu con gà?
\shortans{7}
\loigiai{
\textbf{Đáp số:} $7$\\
$14 : 2 = 7$ con gà.
}
\end{ex}

\begin{ex} %[2D8V43]
Tìm x biết: x x 2 = 20.
\shortans{10}
\loigiai{
\textbf{Đáp số:} $10$\\
$x = 20 : 2 = 10$.
}
\end{ex}

\begin{ex} %[2D8V43]
Một đoạn dây dài 10 dm gập đôi lại. Hỏi mỗi nửa dài bao nhiêu đề-xi-mét?
\shortans{5}
\loigiai{
\textbf{Đáp số:} $5$\\
$10 : 2 = 5\text{ dm}$.
}
\end{ex}

\newpage

\TaoPhieuKiemTra{Chủ đề 8 -- Phép nhân, phép chia}{44}{Bảng chia 5}{https://loptoanthaylinh.vn/thi/chuyen-de/2/toan2_cd_b44}
\textbf{Phần 1. Trắc nghiệm lựa chọn}
\setcounterpageref{ex}{0}

\begin{ex} %[2D8V44]
Kết quả của phép chia $35 : 5$ là
\choice
{$6$}
{\True $7$}
{$8$}
{$9$}
\loigiai{
$35 : 5 = 7$.
\textbf{Đáp án của câu hỏi là B}
}
\end{ex}

\begin{ex} %[2D8V44]
Có $40$ quyển vở chia đều cho $5$ tổ học sinh. Mỗi tổ nhận được bao nhiêu quyển vở?
\choice
{$6$ quyển}
{$7$ quyển}
{$9$ quyển}
{\True $8$ quyển}
\loigiai{
$40 : 5 = 8$ (quyển).
\textbf{Đáp án của câu hỏi là D}
}
\end{ex}

\begin{ex} %[2D8V44]
Tính giá trị của biểu thức: $45 : 5 - 4$.
\choice
{\True $5$}
{$4$}
{$6$}
{$7$}
\loigiai{
$45 : 5 = 9$; $9 - 4 = 5$.
\textbf{Đáp án của câu hỏi là A}
}
\end{ex}

\begin{ex} %[2D8V44]
Một đàn gà con có $25$ con được chia đều vào $5$ chuồng. Mỗi chuồng có bao nhiêu con gà?
\choice
{$4$ con}
{$6$ con}
{\True $5$ con}
{$3$ con}
\loigiai{
$25 : 5 = 5$ (con).
\textbf{Đáp án của câu hỏi là C}
}
\end{ex}

\begin{ex} %[2D8V44]
Điền số thích hợp: $50 : 5 = \dots \times 2$.
\choice
{$4$}
{\True $5$}
{$6$}
{$3$}
\loigiai{
Vế trái: $50 : 5 = 10$. Ta có: $10 : 2 = 5$. Vậy điền 5.
\textbf{Đáp án của câu hỏi là B}
}
\end{ex}

\begin{ex} %[2D8V44]
Tìm $x$ biết: $x \times 5 = 30$.
\choice
{$5$}
{$7$}
{$8$}
{\True $6$}
\loigiai{
$x = 30 : 5 = 6$.
\textbf{Đáp án của câu hỏi là D}
}
\end{ex}

\begin{ex} %[2D8V44]
Mẹ có $15$ chiếc bánh, mẹ chia cho các con mỗi người $5$ chiếc bánh. Hỏi mẹ đã chia bánh cho mấy người con?
\choice
{\True $3$ người}
{$2$ người}
{$4$ người}
{$5$ người}
\loigiai{
$15 : 5 = 3$ (người).
\textbf{Đáp án của câu hỏi là A}
}
\end{ex}

\begin{ex} %[2D8V44]
Bác nông dân có $20\text{ kg}$ đỗ xanh đóng đều vào các túi loại $5\text{ kg}$. Hỏi bác đóng được bao nhiêu túi đỗ?
\choice
{$3$ túi}
{$5$ túi}
{\True $4$ túi}
{$6$ túi}
\loigiai{
$20 : 5 = 4$ (túi).
\textbf{Đáp án của câu hỏi là C}
}
\end{ex}

\begin{ex} %[2D8V44]
Điền dấu thích hợp: $30 : 5 \dots 12 : 2$.
\choice
{$>$}
{\True $=$}
{$<$}
{Không so sánh được}
\loigiai{
$30 : 5 = 6$; $12 : 2 = 6$. Hai vế bằng nhau nên điền dấu $=$.
\textbf{Đáp án của câu hỏi là B}
}
\end{ex}

\begin{ex} %[2D8V44]
Tìm một số biết rằng lấy số đó nhân với $5$ rồi trừ đi $15$ thì được $20$.
\choice
{$6$}
{$8$}
{$9$}
{\True $7$}
\loigiai{
Tính ngược: $20 + 15 = 35$; $35 : 5 = 7$.
\textbf{Đáp án của câu hỏi là D}
}
\end{ex}

\begin{ex} %[2D8V44]
Một tấm vải dài $45\text{ m}$ được cắt thành các mảnh dài $5\text{ m}$. Hỏi người thợ cắt được bao nhiêu mảnh vải?
\choice
{\True $9$ mảnh vải}
{$8$ mảnh vải}
{$10$ mảnh vải}
{$7$ mảnh vải}
\loigiai{
$45 : 5 = 9$ (mảnh vải).
\textbf{Đáp án của câu hỏi là A}
}
\end{ex}

\begin{ex} %[2D8V44]
Trong vườn có $50$ cây hoa trồng thành $5$ luống bằng nhau. Mỗi luống có bao nhiêu cây hoa?
\choice
{$8$ cây}
{$9$ cây}
{\True $10$ cây}
{$12$ cây}
\loigiai{
$50 : 5 = 10$ (cây).
\textbf{Đáp án của câu hỏi là C}
}
\end{ex}

\textbf{Phần 2. Trắc nghiệm đúng sai}
\setcounterpageref{ex}{0}

\begin{ex} %[2D8V44]
Xét tính đúng sai của các phép tính trong bảng chia 5:
\choiceTF
{\True Các số bị chia trong bảng chia 5 đều có chữ số tận cùng là 0 hoặc 5.}
{Thương của 25 : 5 lớn hơn thương của 30 : 5.}
{\True Nếu lấy 50 : 5 thì được thương là 10.}
{\True Số chia trong bảng chia 5 luôn luôn bằng 5.}
\loigiai{
\begin{itemize}
    \item Ý a) Đúng: Đúng.
    \item Ý b) Sai: $25 : 5 = 5 < 30 : 5 = 6$.
    \item Ý c) Đúng: $50 : 5 = 10$. Đúng.
    \item Ý d) Đúng: Đúng theo tên gọi bảng chia.
\end{itemize}
\begin{itemize}
    \item a) \textbf{Đúng}.
    \item b) \textbf{Sai}.
    \item c) \textbf{Đúng}.
    \item d) \textbf{Đúng}.
\end{itemize}
}
\end{ex}

\begin{ex} %[2D8V44]
Một lớp học có 35 học sinh chia đều thành 5 nhóm:
\choiceTF
{\True Mỗi nhóm có 7 học sinh.}
{Mỗi nhóm có số học sinh là số chẵn.}
{\True Tổng số học sinh của 2 nhóm là 14 học sinh.}
{Nếu mỗi nhóm bớt đi 1 bạn thì cả lớp còn lại 25 bạn.}
\loigiai{
\begin{itemize}
    \item Ý a) Đúng: $35 : 5 = 7$ học sinh.
    \item Ý b) Sai: 7 là số lẻ.
    \item Ý c) Đúng: $7 \times 2 = 14$ học sinh.
    \item Ý d) Sai: Cả lớp bớt $5 \times 1 = 5$ bạn thì còn $35 - 5 = 30$ bạn (không phải 25).
\end{itemize}
\begin{itemize}
    \item a) \textbf{Đúng}.
    \item b) \textbf{Sai}.
    \item c) \textbf{Đúng}.
    \item d) \textbf{Sai}.
\end{itemize}
}
\end{ex}

\textbf{Phần 3. Trả lời ngắn}
\setcounterpageref{ex}{0}

\begin{ex} %[2D8V44]
Tính kết quả: 45 : 5.
\shortans{9}
\loigiai{
\textbf{Đáp số:} $9$\\
$45 : 5 = 9$.
}
\end{ex}

\begin{ex} %[2D8V44]
Có 30 lít dầu rót đều vào 5 can. Hỏi mỗi can chứa bao nhiêu lít dầu?
\shortans{6}
\loigiai{
\textbf{Đáp số:} $6$\\
$30 : 5 = 6\text{ L}$.
}
\end{ex}

\begin{ex} %[2D8V44]
Tính giá trị: 40 : 5 + 18.
\shortans{26}
\loigiai{
\textbf{Đáp số:} $26$\\
$40 : 5 = 8$; $8 + 18 = 26$.
}
\end{ex}

\begin{ex} %[2D8V44]
Một đoạn dây dài 25 cm cắt thành 5 đoạn bằng nhau. Mỗi đoạn dài bao nhiêu xăng-ti-mét?
\shortans{5}
\loigiai{
\textbf{Đáp số:} $5$\\
$25 : 5 = 5\text{ cm}$.
}
\end{ex}

\begin{ex} %[2D8V44]
Tìm x biết: x x 5 = 50.
\shortans{10}
\loigiai{
\textbf{Đáp số:} $10$\\
$x = 50 : 5 = 10$.
}
\end{ex}

\begin{ex} %[2D8V44]
Có 35 quả cam chia đều vào các đĩa, mỗi đĩa 5 quả. Hỏi chia được bao nhiêu đĩa cam?
\shortans{7}
\loigiai{
\textbf{Đáp số:} $7$\\
$35 : 5 = 7$ đĩa.
}
\end{ex}

\newpage

\TaoPhieuKiemTra{Chủ đề 8 -- Phép nhân, phép chia}{45}{Luyện tập chung}{https://loptoanthaylinh.vn/thi/chuyen-de/2/toan2_cd_b45}
\textbf{Phần 1. Trắc nghiệm lựa chọn}
\setcounterpageref{ex}{0}

\begin{ex} %[2D8V45]
Tính giá trị của biểu thức: $2 \times 8 + 35 : 5$.
\choice
{\True $23$}
{$21$}
{$22$}
{$24$}
\loigiai{
$2 \times 8 = 16$; $35 : 5 = 7$; $16 + 7 = 23$.
\textbf{Đáp án của câu hỏi là A}
}
\end{ex}

\begin{ex} %[2D8V45]
Có $18$ chiếc kẹo chia đều cho Mai và Lan. Sau đó mẹ cho Mai thêm $3$ cái kẹo nữa. Hỏi Mai có tất cả bao nhiêu cái kẹo?
\choice
{$10$ cái kẹo}
{$11$ cái kẹo}
{\True $12$ cái kẹo}
{$13$ cái kẹo}
\loigiai{
Mỗi bạn nhận được: $18 : 2 = 9$ (cái kẹo).\nMai có tất cả: $9 + 3 = 12$ (cái kẹo).
\textbf{Đáp án của câu hỏi là C}
}
\end{ex}

\begin{ex} %[2D8V45]
Tìm $x$ biết: $x : 2 = 15 - 8$.
\choice
{$12$}
{\True $14$}
{$16$}
{$18$}
\loigiai{
Vế phải: $15 - 8 = 7$. Ta có: $x = 7 \times 2 = 14$.
\textbf{Đáp án của câu hỏi là B}
}
\end{ex}

\begin{ex} %[2D8V45]
Một chuồng nuôi $5$ con thỏ và $6$ con gà. Hỏi cả chuồng có bao nhiêu cái chân?
\choice
{$30$ cái chân}
{$28$ cái chân}
{$34$ cái chân}
{\True $32$ cái chân}
\loigiai{
Số chân thỏ: $4 \times 5 = 20$ cái. Số chân gà: $2 \times 6 = 12$ cái. Tổng: $20 + 12 = 32$ cái chân.
\textbf{Đáp án của câu hỏi là D}
}
\end{ex}

\begin{ex} %[2D8V45]
Tính nhanh: $5 \times 4 + 5 \times 6$.
\choice
{\True $50$}
{$45$}
{$55$}
{$40$}
\loigiai{
$5 \times (4 + 6) = 5 \times 10 = 50$.
\textbf{Đáp án của câu hỏi là A}
}
\end{ex}

\begin{ex} %[2D8V45]
Có $40\text{ kg}$ gạo đóng đều vào các túi loại $5\text{ kg}$. Người ta đã bán đi $3$ túi. Hỏi còn lại bao nhiêu túi gạo?
\choice
{$4$ túi}
{$6$ túi}
{\True $5$ túi}
{$7$ túi}
\loigiai{
Tổng số túi: $40 : 5 = 8$ (túi). Còn lại: $8 - 3 = 5$ (túi).
\textbf{Đáp án của câu hỏi là C}
}
\end{ex}

\begin{ex} %[2D8V45]
Điền dấu thích hợp: $2 \times 9 \dots 5 \times 4 - 3$.
\choice
{$<$}
{\True $>$}
{$=$}
{Không so sánh được}
\loigiai{
Vế trái: $2 \times 9 = 18$. Vế phải: $20 - 3 = 17$. Vì $18 > 17$ nên điền dấu $>$. 
\textbf{Đáp án của câu hỏi là B}
}
\end{ex}

\begin{ex} %[2D8V45]
Một cuộn dây dài $20\text{ m}$ cắt thành các đoạn dài $2\text{ m}$. Hỏi phải cắt bao nhiêu nhát kéo?
\choice
{$10$ nhát}
{$8$ nhát}
{$7$ nhát}
{\True $9$ nhát}
\loigiai{
Số đoạn dây: $20 : 2 = 10$ đoạn. Số nhát cắt: $10 - 1 = 9$ nhát.
\textbf{Đáp án của câu hỏi là D}
}
\end{ex}

\begin{ex} %[2D8V45]
Tìm một số biết rằng lấy số đó nhân với $2$ rồi cộng với $14$ thì được $30$.
\choice
{\True $8$}
{$7$}
{$9$}
{$6$}
\loigiai{
Tính ngược: $30 - 14 = 16$; $16 : 2 = 8$.
\textbf{Đáp án của câu hỏi là A}
}
\end{ex}

\begin{ex} %[2D8V45]
Trong một phép chia cho $5$, thương là số chẵn lớn nhất có một chữ số. Số bị chia là
\choice
{$35$}
{$45$}
{\True $40$}
{$30$}
\loigiai{
Số chẵn lớn nhất có 1 chữ số là 8. Số bị chia: $8 \times 5 = 40$.
\textbf{Đáp án của câu hỏi là C}
}
\end{ex}

\begin{ex} %[2D8V45]
Mỗi can chứa $5\text{ L}$ dầu. Bác Nam có $35\text{ L}$ dầu đóng vào các can như thế và cho bạn $2$ can. Bác còn lại bao nhiêu can dầu?
\choice
{$4$ can}
{\True $5$ can}
{$6$ can}
{$7$ can}
\loigiai{
Bác đóng được: $35 : 5 = 7$ (can). Còn lại: $7 - 2 = 5$ (can).
\textbf{Đáp án của câu hỏi là B}
}
\end{ex}

\begin{ex} %[2D8V45]
Hai số có tích bằng $18$ và hiệu bằng $7$. Hai số đó là
\choice
{$6\text{ và } 3$}
{$18\text{ và } 1$}
{$8\text{ và } 1$}
{\True $9\text{ và } 2$}
\loigiai{
$9 \times 2 = 18$ và $9 - 2 = 7$. Hai số là 9 và 2.
\textbf{Đáp án của câu hỏi là D}
}
\end{ex}

\textbf{Phần 2. Trắc nghiệm đúng sai}
\setcounterpageref{ex}{0}

\begin{ex} %[2D8V45]
Xét tính đúng sai của các bài toán rèn luyện tư duy nhân, chia:
\choiceTF
{\True Biểu thức 2 x 7 + 5 x 4 có kết quả bằng 34.}
{Thực hiện phép tính 40 : 5 - 3 ta lấy 5 - 3 trước rồi lấy 40 chia.}
{\True Số chia luôn khác 0.}
{Tích của 2 số tự nhiên luôn lớn hơn thương của 2 số đó.}
\loigiai{
\begin{itemize}
    \item Ý a) Đúng: $14 + 20 = 34$. Đúng.
    \item Ý b) Sai: Quy tắc: Nhân chia trước, cộng trừ sau. Phải lấy $40 : 5 = 8$ rồi $8 - 3 = 5$.
    \item Ý c) Đúng: Đúng.
    \item Ý d) Sai: Nếu $2 \times 1 = 2$ và $2 : 1 = 2$, tích bằng thương.
\end{itemize}
\begin{itemize}
    \item a) \textbf{Đúng}.
    \item b) \textbf{Sai}.
    \item c) \textbf{Đúng}.
    \item d) \textbf{Sai}.
\end{itemize}
}
\end{ex}

\begin{ex} %[2D8V45]
Bác thợ mộc có 16 m gỗ, cưa thành các đoạn dài 2 m:
\choiceTF
{\True Bác cưa được 8 đoạn gỗ.}
{\True Bác phải cưa 7 nhát cưa.}
{Nếu mỗi nhát cưa mất 2 phút thì bác cưa hết 16 phút.}
{\True Nếu mỗi đoạn gỗ bác đóng thành một chiếc ghế thì bác đóng được 8 chiếc ghế.}
\loigiai{
\begin{itemize}
    \item Ý a) Đúng: $16 : 2 = 8$ đoạn.
    \item Ý b) Đúng: $8 - 1 = 7$ nhát.
    \item Ý c) Sai: Thời gian cưa: $7 \times 2 = 14$ phút (không phải 16 phút).
    \item Ý d) Đúng: Đúng.
\end{itemize}
\begin{itemize}
    \item a) \textbf{Đúng}.
    \item b) \textbf{Đúng}.
    \item c) \textbf{Sai}.
    \item d) \textbf{Đúng}.
\end{itemize}
}
\end{ex}

\textbf{Phần 3. Trả lời ngắn}
\setcounterpageref{ex}{0}

\begin{ex} %[2D8V45]
Tính giá trị biểu thức: 2 x 9 + 45 : 5.
\shortans{27}
\loigiai{
\textbf{Đáp số:} $27$\\
$2 \times 9 = 18$; $45 : 5 = 9$; $18 + 9 = 27$.
}
\end{ex}

\begin{ex} %[2D8V45]
Có 5 bạn nhỏ, mỗi bạn có 4 quả bóng bay. Sau đó bị nổ mất 6 quả. Hỏi còn lại bao nhiêu quả bóng bay?
\shortans{14}
\loigiai{
\textbf{Đáp số:} $14$\\
$5 \times 4 - 6 = 20 - 6 = 14$ quả.
}
\end{ex}

\begin{ex} %[2D8V45]
Tìm x biết: x : 5 + 17 = 23.
\shortans{30}
\loigiai{
\textbf{Đáp số:} $30$\\
$x : 5 = 23 - 17 = 6 \Rightarrow x = 6 \times 5 = 30$.
}
\end{ex}

\begin{ex} %[2D8V45]
Một chuồng nuôi gà có tất cả 16 cái chân. Hỏi chuồng đó có bao nhiêu con gà?
\shortans{8}
\loigiai{
\textbf{Đáp số:} $8$\\
$16 : 2 = 8$ con gà.
}
\end{ex}

\begin{ex} %[2D8V45]
Bao gạo nặng 50 kg chia đều vào 5 túi. Mẹ đã dùng hết 2 túi. Hỏi còn lại bao nhiêu ki-lô-gam gạo?
\shortans{30}
\loigiai{
\textbf{Đáp số:} $30$\\
Mỗi túi: $50 : 5 = 10\text{ kg}$. Còn 3 túi: $10 \times 3 = 30\text{ kg}$.
}
\end{ex}

\begin{ex} %[2D8V45]
Tìm một số biết rằng lấy số đó chia cho 2 rồi nhân với 5 thì được 20.
\shortans{8}
\loigiai{
\textbf{Đáp số:} $8$\\
$20 : 5 = 4$; $4 \times 2 = 8$.
}
\end{ex}

\newpage


% =========================================================================
% CHỦ ĐỀ 9 - LÀM QUEN VỚI HÌNH KHỐI
% =========================================================================
\section{Làm quen với hình khối}

\TaoPhieuKiemTra{Chủ đề 9 -- Làm quen với hình khối}{46}{Khối trụ, khối cầu}{https://loptoanthaylinh.vn/thi/chuyen-de/2/toan2_cd_b46}
\textbf{Phần 1. Trắc nghiệm lựa chọn}
\setcounterpageref{ex}{0}

\begin{ex} %[2H9V46]
Đồ vật nào dưới đây có dạng khối cầu?
\choice
{Hộp sữa bột}
{Hộp phấn viết}
{Cục tẩy}
{\True Quả bóng đá}
\loigiai{
Quả bóng đá có dạng khối cầu.
\textbf{Đáp án của câu hỏi là D}
}
\end{ex}

\begin{ex} %[2H9V46]
Đồ vật nào dưới đây có dạng khối trụ?
\choice
{Quả địa cầu}
{\True Lon nước ngọt}
{Hộp bánh quy vuông}
{Cuốn sách giáo khoa}
\loigiai{
Lon nước ngọt có dạng khối trụ.
\textbf{Đáp án của câu hỏi là B}
}
\end{ex}

\begin{ex} %[2H9V46]
Khối cầu có đặc điểm nào dưới đây?
\choice
{Có hai mặt đáy phẳng}
{Có $6$ mặt vuông}
{\True Lăn được về mọi phía}
{Không lăn được}
\loigiai{
Khối cầu có bề mặt cong đều nên lăn được về mọi phía.
\textbf{Đáp án của câu hỏi là C}
}
\end{ex}

\begin{ex} %[2H9V46]
Khối trụ có bao nhiêu mặt phẳng hình tròn?
\choice
{\True $2$ mặt phẳng}
{$1$ mặt phẳng}
{$3$ mặt phẳng}
{Không có mặt nào}
\loigiai{
Khối trụ có hai mặt đáy là hai hình tròn phẳng bằng nhau.
\textbf{Đáp án của câu hỏi là A}
}
\end{ex}

\begin{ex} %[2H9V46]
Người ta có thể xếp chồng các khối nào lên nhau mà không bị lăn đổ?
\choice
{Các khối cầu xếp chồng lên nhau}
{Hai quả bóng xếp chồng}
{Khối cầu đặt trên khối trụ}
{\True Các khối trụ xếp chồng theo mặt đáy phẳng}
\loigiai{
Xếp chồng các khối trụ theo mặt đáy phẳng thì đứng vững chắc.
\textbf{Đáp án của câu hỏi là D}
}
\end{ex}

\begin{ex} %[2H9V46]
Trong các vật sau: quả bưởi, lon sữa ông thọ, hộp chè hình trụ, viên bi, chiếc cốc hình trụ. Có bao nhiêu vật có dạng khối cầu?
\choice
{$3$ vật}
{\True $2$ vật}
{$1$ vật}
{$4$ vật}
\loigiai{
Hai vật có dạng khối cầu là: quả bưởi và viên bi.
\textbf{Đáp án của câu hỏi là B}
}
\end{ex}

\begin{ex} %[2H9V46]
Cắt một khối trụ bằng một nhát cắt phẳng song song với hai đáy, phần nhận được ở mặt cắt có hình gì?
\choice
{Hình vuông}
{Hình tam giác}
{\True Hình tròn}
{Hình chữ nhật}
\loigiai{
Mặt cắt song song với đáy khối trụ là một hình tròn.
\textbf{Đáp án của câu hỏi là C}
}
\end{ex}

\begin{ex} %[2H9V46]
Viên phấn viết bảng thông thường có dạng hình gì?
\choice
{\True Khối trụ}
{Khối cầu}
{Khối lập phương}
{Khối hộp chữ nhật}
\loigiai{
Thân viên phấn có dạng khối trụ tròn dài.
\textbf{Đáp án của câu hỏi là A}
}
\end{ex}

\begin{ex} %[2H9V46]
Một mô hình lâu đài được ghép từ $4$ khối trụ, $3$ khối cầu và $5$ khối lập phương. Có tất cả bao nhiêu khối hình?
\choice
{$10$ khối}
{$11$ khối}
{$14$ khối}
{\True $12$ khối}
\loigiai{
$4 + 3 + 5 = 12$ (khối).
\textbf{Đáp án của câu hỏi là D}
}
\end{ex}

\begin{ex} %[2H9V46]
Vật nào sau đây vừa có thể trượt vừa có thể lăn?
\choice
{Khối cầu}
{\True Khối trụ}
{Khối lập phương}
{Khối hộp chữ nhật}
\loigiai{
Khối trụ đặt đứng thì trượt được theo đáy phẳng, đặt nằm thì lăn được.
\textbf{Đáp án của câu hỏi là B}
}
\end{ex}

\begin{ex} %[2H9V46]
Có $6$ khối trụ bằng nhau, mỗi khối cao $3\text{ cm}$. Xếp chồng cả $6$ khối trụ lên nhau thì được cột trụ cao bao nhiêu xăng-ti-mét?
\choice
{$15\text{ cm}$}
{$21\text{ cm}$}
{\True $18\text{ cm}$}
{$16\text{ cm}$}
\loigiai{
$3 \times 6 = 18\text{ (cm)}$.
\textbf{Đáp án của câu hỏi là C}
}
\end{ex}

\begin{ex} %[2H9V46]
Trái Đất nơi chúng ta đang sống có hình dạng gần giống với khối gì?
\choice
{\True Khối cầu}
{Khối trụ}
{Khối lập phương}
{Hình tròn phẳng}
\loigiai{
Trái Đất có dạng hình khối cầu.
\textbf{Đáp án của câu hỏi là A}
}
\end{ex}

\textbf{Phần 2. Trắc nghiệm đúng sai}
\setcounterpageref{ex}{0}

\begin{ex} %[2H9V46]
Xét tính đúng sai của các nhận định về khối trụ và khối cầu:
\choiceTF
{\True Khối cầu không có cạnh và không có đỉnh nào.}
{Khối trụ có thể lăn được về mọi phía giống như khối cầu.}
{\True Lon nước ngọt đặt đứng trên bàn không bị lăn vì có mặt đáy phẳng.}
{Khối cầu có thể xếp chồng lên nhau thành tháp cao vững chắc.}
\loigiai{
\begin{itemize}
    \item Ý a) Đúng: Đúng.
    \item Ý b) Sai: Khối trụ chỉ lăn được theo một hướng khi đặt nằm tròn.
    \item Ý c) Đúng: Đúng.
    \item Ý d) Sai: Khối cầu hình cong sẽ bị lăn đổ ngay.
\end{itemize}
\begin{itemize}
    \item a) \textbf{Đúng}.
    \item b) \textbf{Sai}.
    \item c) \textbf{Đúng}.
    \item d) \textbf{Sai}.
\end{itemize}
}
\end{ex}

\begin{ex} %[2H9V46]
Trong giờ học xếp hình khối, bạn Nam có một bộ đồ chơi gồm các khối hình:
\choiceTF
{\True Nam lấy 3 quả bóng tròn thì Nam đã chọn 3 khối cầu.}
{\True Nam dùng các ống bơ để làm cột trụ lâu đài thì đó là các khối trụ.}
{Xếp một khối cầu bên dưới rồi đặt khối trụ lên trên thì cột trụ đứng vững.}
{\True Một khối trụ có 2 mặt phẳng đối diện nhau là hình tròn bằng nhau.}
\loigiai{
\begin{itemize}
    \item Ý a) Đúng: Đúng.
    \item Ý b) Đúng: Đúng.
    \item Ý c) Sai: Khối cầu bên dưới tròn sẽ làm khối trụ bên trên đổ ngã.
    \item Ý d) Đúng: Đúng.
\end{itemize}
\begin{itemize}
    \item a) \textbf{Đúng}.
    \item b) \textbf{Đúng}.
    \item c) \textbf{Sai}.
    \item d) \textbf{Đúng}.
\end{itemize}
}
\end{ex}

\textbf{Phần 3. Trả lời ngắn}
\setcounterpageref{ex}{0}

\begin{ex} %[2H9V46]
Một hộp đồ chơi có 4 khối cầu và 6 khối trụ. Hỏi hộp đó có tất cả bao nhiêu khối hình?
\shortans{10}
\loigiai{
\textbf{Đáp số:} $10$\\
$4 + 6 = 10$ khối.
}
\end{ex}

\begin{ex} %[2H9V46]
Người ta xếp chồng 5 lon sữa hình trụ lên nhau, mỗi lon cao 4 cm. Hỏi cột lon sữa cao bao nhiêu xăng-ti-mét?
\shortans{20}
\loigiai{
\textbf{Đáp số:} $20$\\
$4 \times 5 = 20\text{ cm}$.
}
\end{ex}

\begin{ex} %[2H9V46]
Một chú hề tung hứng 3 quả bóng xanh và 5 quả bóng đỏ có dạng khối cầu. Hỏi chú hề tung hứng tất cả bao nhiêu quả bóng?
\shortans{8}
\loigiai{
\textbf{Đáp số:} $8$\\
$3 + 5 = 8$ quả bóng.
}
\end{ex}

\begin{ex} %[2H9V46]
Trong rổ có 15 quả cam có dạng khối cầu, mẹ lấy ra 7 quả để vắt nước. Trong rổ còn lại bao nhiêu quả cam?
\shortans{8}
\loigiai{
\textbf{Đáp số:} $8$\\
$15 - 7 = 8$ quả.
}
\end{ex}

\begin{ex} %[2H9V46]
Khối trụ có mấy mặt phẳng hình tròn?
\shortans{2}
\loigiai{
\textbf{Đáp số:} $2$\\
Có 2 mặt phẳng hình tròn ở hai đầu đáy.
}
\end{ex}

\begin{ex} %[2H9V46]
Một mô hình xe buýt đồ chơi gồm 6 bánh xe có dạng khối trụ dẹt. Hỏi 3 chiếc xe buýt như thế có tất cả bao nhiêu bánh xe?
\shortans{18}
\loigiai{
\textbf{Đáp số:} $18$\\
$6 \times 3 = 18$ bánh xe.
}
\end{ex}

\newpage

\TaoPhieuKiemTra{Chủ đề 9 -- Làm quen với hình khối}{47}{Luyện tập chung}{https://loptoanthaylinh.vn/thi/chuyen-de/2/toan2_cd_b47}
\textbf{Phần 1. Trắc nghiệm lựa chọn}
\setcounterpageref{ex}{0}

\begin{ex} %[2H9V47]
Hình nào dưới đây có các mặt đều là hình vuông bằng nhau?
\choice
{Khối hộp chữ nhật}
{Khối trụ}
{\True Khối lập phương}
{Khối cầu}
\loigiai{
Khối lập phương có 6 mặt đều là hình vuông bằng nhau.
\textbf{Đáp án của câu hỏi là C}
}
\end{ex}

\begin{ex} %[2H9V47]
Bao diêm, viên gạch xây nhà thường có hình dạng là khối gì?
\choice
{Khối lập phương}
{\True Khối hộp chữ nhật}
{Khối trụ}
{Khối cầu}
\loigiai{
Bao diêm và viên gạch có dạng khối hộp chữ nhật.
\textbf{Đáp án của câu hỏi là B}
}
\end{ex}

\begin{ex} %[2H9V47]
Xếp $8$ khối lập phương nhỏ cạnh $1\text{ cm}$ thành một khối lập phương lớn. Khối lập phương lớn có cạnh dài bao nhiêu?
\choice
{$3\text{ cm}$}
{$4\text{ cm}$}
{$1\text{ cm}$}
{\True $2\text{ cm}$}
\loigiai{
Khối lập phương lớn gồm $2 \times 2 \times 2 = 8$ khối nhỏ, có cạnh dài $2\text{ cm}$.
\textbf{Đáp án của câu hỏi là D}
}
\end{ex}

\begin{ex} %[2H9V47]
Trong các đồ vật sau: quả bóng bàn, quả bóng chuyền, quả bóng rổ. Chúng đều có điểm chung là có dạng
\choice
{\True Khối cầu}
{Khối trụ}
{Khối tròn dẹt}
{Khối lập phương}
\loigiai{
Các quả bóng đều có dạng khối cầu.
\textbf{Đáp án của câu hỏi là A}
}
\end{ex}

\begin{ex} %[2H9V47]
Một người thợ xây dùng $45$ viên gạch hình hộp chữ nhật để xây bệ đỡ, sau đó dùng thêm $28$ viên gạch nữa. Bác thợ đã dùng tất cả bao nhiêu viên gạch?
\choice
{$71$ viên gạch}
{$72$ viên gạch}
{\True $73$ viên gạch}
{$75$ viên gạch}
\loigiai{
$45 + 28 = 73$ (viên gạch).
\textbf{Đáp án của câu hỏi là C}
}
\end{ex}

\begin{ex} %[2H9V47]
Thùng đựng nước hình trụ có thể chứa $20\text{ L}$ nước. Hiện trong thùng có $12\text{ L}$ nước. Cần đổ thêm bao nhiêu lít nước để đầy thùng?
\choice
{$7\text{ L}$}
{\True $8\text{ L}$}
{$9\text{ L}$}
{$6\text{ L}$}
\loigiai{
$20 - 12 = 8\text{ (L)}$.
\textbf{Đáp án của câu hỏi là B}
}
\end{ex}

\begin{ex} %[2H9V47]
Con xúc xắc đồ chơi để chơi cờ cá ngựa có dạng hình khối gì?
\choice
{Khối trụ}
{Khối cầu}
{Khối chữ nhật}
{\True Khối lập phương}
\loigiai{
Xúc xắc có 6 mặt vuông bằng nhau nên có dạng khối lập phương.
\textbf{Đáp án của câu hỏi là D}
}
\end{ex}

\begin{ex} %[2H9V47]
Một hộp quà hình lập phương có cạnh $10\text{ cm}$. Hỏi dùng một sợi dây để thắt nơ viền quanh $4$ mặt xung quanh của hộp quà cần ít nhất bao nhiêu xăng-ti-mét dây (không tính phần nơ)?
\choice
{\True $40\text{ cm}$}
{$30\text{ cm}$}
{$50\text{ cm}$}
{$20\text{ cm}$}
\loigiai{
Viền quanh 4 mặt gồm 4 cạnh: $10 \times 4 = 40\text{ (cm)}$.
\textbf{Đáp án của câu hỏi là A}
}
\end{ex}

\begin{ex} %[2H9V47]
Chiếc trống trường học thường có thân trống phình nhẹ nhưng hai mặt trống có dạng phẳng hình tròn. Dạng hình học gần giống nhất của chiếc trống là
\choice
{Khối cầu}
{Khối lập phương}
{\True Khối trụ}
{Khối hộp chữ nhật}
\loigiai{
Trống trường có dạng hình khối trụ.
\textbf{Đáp án của câu hỏi là C}
}
\end{ex}

\begin{ex} %[2H9V47]
Có $18$ lon nước ngọt hình trụ đóng đều vào $3$ lốc. Mỗi lốc có bao nhiêu lon nước ngọt?
\choice
{$5$ lon}
{\True $6$ lon}
{$7$ lon}
{$8$ lon}
\loigiai{
$18 : 3 = 6$ (lon).
\textbf{Đáp án của câu hỏi là B}
}
\end{ex}

\begin{ex} %[2H9V47]
Để lăn một vật trên mặt đất một cách dễ dàng và tốn ít sức nhất, người ta nên dùng vật có dạng khối nào?
\choice
{Khối lập phương}
{Khối hộp chữ nhật}
{Khối tam giác}
{\True Khối trụ hoặc khối cầu}
\loigiai{
Bề mặt cong của khối trụ và khối cầu giúp vật lăn tròn dễ dàng.
\textbf{Đáp án của câu hỏi là D}
}
\end{ex}

\begin{ex} %[2H9V47]
Bé xếp một chiếc xe lửa đồ chơi gồm $1$ khối lập phương làm đầu tàu và $5$ khối hộp chữ nhật làm các toa tàu. Chiếc xe lửa có tất cả bao nhiêu khối hình?
\choice
{\True $6$ khối hình}
{$5$ khối hình}
{$7$ khối hình}
{$8$ khối hình}
\loigiai{
$1 + 5 = 6$ (khối hình).
\textbf{Đáp án của câu hỏi là A}
}
\end{ex}

\textbf{Phần 2. Trắc nghiệm đúng sai}
\setcounterpageref{ex}{0}

\begin{ex} %[2H9V47]
Xét tính đúng sai của các bài toán hình khối thực tế:
\choiceTF
{\True Hộp phấn viết bảng dạng khối hộp chữ nhật có 6 mặt, 8 đỉnh và 12 cạnh.}
{Một khối lập phương có thể lăn tròn về mọi phía như quả bóng.}
{\True Chiếc bánh chưng ngày Tết có hình dạng gần giống khối hộp chữ nhật hoặc khối lập phương dẹt.}
{Tất cả các mặt của khối hộp chữ nhật đều bắt buộc phải là hình vuông.}
\loigiai{
\begin{itemize}
    \item Ý a) Đúng: Đặc điểm cơ bản của hình hộp chữ nhật. Đúng.
    \item Ý b) Sai: Khối lập phương có các mặt phẳng và góc cạnh nên không lăn tròn được.
    \item Ý c) Đúng: Đúng.
    \item Ý d) Sai: Các mặt thường là hình chữ nhật (nếu là hình vuông hết thì thành khối lập phương).
\end{itemize}
\begin{itemize}
    \item a) \textbf{Đúng}.
    \item b) \textbf{Sai}.
    \item c) \textbf{Đúng}.
    \item d) \textbf{Sai}.
\end{itemize}
}
\end{ex}

\begin{ex} %[2H9V47]
Một kệ đồ chơi có: 6 quả bóng tròn, 4 chiếc trống hình trụ, 5 khối rubik lập phương:
\choiceTF
{\True Kệ có tất cả 15 món đồ chơi.}
{\True Số đồ chơi dạng khối cầu nhiều nhất trong ba loại.}
{Số trống hình trụ nhiều hơn số khối rubik lập phương.}
{Nếu lấy bớt đi 2 quả bóng thì số bóng ít hơn số trống.}
\loigiai{
\begin{itemize}
    \item Ý a) Đúng: $6 + 4 + 5 = 15$ món đồ chơi.
    \item Ý b) Đúng: Bóng tròn có 6 quả, nhiều nhất. Đúng.
    \item Ý c) Sai: Trống có 4, rubik có 5 nên trống ít hơn rubik.
    \item Ý d) Sai: Bóng còn: $6 - 2 = 4$, bằng số trống (không phải ít hơn).
\end{itemize}
\begin{itemize}
    \item a) \textbf{Đúng}.
    \item b) \textbf{Đúng}.
    \item c) \textbf{Sai}.
    \item d) \textbf{Sai}.
\end{itemize}
}
\end{ex}

\textbf{Phần 3. Trả lời ngắn}
\setcounterpageref{ex}{0}

\begin{ex} %[2H9V47]
Một khối lập phương có mấy mặt?
\shortans{6}
\loigiai{
\textbf{Đáp số:} $6$\\
Khối lập phương có 6 mặt.
}
\end{ex}

\begin{ex} %[2H9V47]
Có 24 lon sữa hình trụ đóng đều vào 4 hộp carton. Hỏi mỗi hộp có bao nhiêu lon sữa?
\shortans{6}
\loigiai{
\textbf{Đáp số:} $6$\\
$24 : 4 = 6$ lon.
}
\end{ex}

\begin{ex} %[2H9V47]
Bác thợ mộc cưa một khúc gỗ hình trụ dài 20 dm thành 4 khúc bằng nhau. Mỗi khúc dài bao nhiêu đề-xi-mét?
\shortans{5}
\loigiai{
\textbf{Đáp số:} $5$\\
$20 : 4 = 5\text{ dm}$.
}
\end{ex}

\begin{ex} %[2H9V47]
Xếp các khối lập phương nhỏ thành một hàng dài gồm 9 khối, mỗi khối cạnh 3 cm. Hỏi hàng đó dài bao nhiêu xăng-ti-mét?
\shortans{27}
\loigiai{
\textbf{Đáp số:} $27$\\
$3 \times 9 = 27\text{ cm}$.
}
\end{ex}

\begin{ex} %[2H9V47]
Một hộp đồ chơi có 12 khối cầu và 18 khối trụ. Hỏi hộp đó có tất cả bao nhiêu khối hình?
\shortans{30}
\loigiai{
\textbf{Đáp số:} $30$\\
$12 + 18 = 30$ khối.
}
\end{ex}

\begin{ex} %[2H9V47]
Có bao nhiêu đỉnh trên một khối lập phương?
\shortans{8}
\loigiai{
\textbf{Đáp số:} $8$\\
Khối lập phương có 8 đỉnh.
}
\end{ex}

\newpage


% =========================================================================
% CHỦ ĐỀ 10 - CÁC SỐ TRONG PHẠM VI 1 000
% =========================================================================
\section{Các số trong phạm vi 1 000}

\TaoPhieuKiemTra{Chủ đề 10 -- Các số trong phạm vi 1 000}{48}{Đơn vị, chục, trăm, nghìn}{https://loptoanthaylinh.vn/thi/chuyen-de/2/toan2_cd_b48}
\textbf{Phần 1. Trắc nghiệm lựa chọn}
\setcounterpageref{ex}{0}

\begin{ex} %[2D10V48]
Số gồm $10$ chục que tính gộp lại được gọi là số nào?
\choice
{Mười chục}
{Một nghìn}
{\True Một trăm}
{Mười đơn vị}
\loigiai{
$10$ chục $= 100$ (một trăm).
\textbf{Đáp án của câu hỏi là C}
}
\end{ex}

\begin{ex} %[2D10V48]
Số gồm $10$ trăm đơn vị gộp lại được gọi là số nào?
\choice
{\True Một nghìn}
{Một trăm}
{Mười chục}
{Một vạn}
\loigiai{
$10$ trăm $= 1 000$ (một nghìn).
\textbf{Đáp án của câu hỏi là A}
}
\end{ex}

\begin{ex} %[2D10V48]
Một người xếp các túi bi, mỗi túi có $10$ viên bi. Có $10$ túi như thế đóng trong một hộp. Hỏi một hộp có tất cả bao nhiêu viên bi?
\choice
{$20$ viên}
{$50$ viên}
{$1 000$ viên}
{\True $100$ viên}
\loigiai{
$10$ túi, mỗi túi 10 viên là $10$ chục viên $= 100$ viên bi.
\textbf{Đáp án của câu hỏi là D}
}
\end{ex}

\begin{ex} %[2D10V48]
Một kho hàng có $10$ thùng sách, mỗi thùng chứa đúng $100$ cuốn sách. Hỏi trong kho có tất cả bao nhiêu cuốn sách?
\choice
{$100$ cuốn}
{\True $1 000$ cuốn}
{$500$ cuốn}
{$200$ cuốn}
\loigiai{
$10$ thùng, mỗi thùng 100 cuốn là $10$ trăm cuốn $= 1 000$ cuốn sách.
\textbf{Đáp án của câu hỏi là B}
}
\end{ex}

\begin{ex} %[2D10V48]
Hiệu giữa số tự nhiên bé nhất có bốn chữ số và số tự nhiên lớn nhất có ba chữ số là
\choice
{$10$}
{$100$}
{\True $1$}
{$0$}
\loigiai{
Số bé nhất có bốn chữ số là $1 000$, số lớn nhất có ba chữ số là $999$. Ta có: $1 000 - 999 = 1$.
\textbf{Đáp án của câu hỏi là C}
}
\end{ex}

\begin{ex} %[2D10V48]
Trong số $1 000$, chữ số $1$ nằm ở hàng nào?
\choice
{\True Hàng nghìn}
{Hàng đơn vị}
{Hàng chục}
{Hàng trăm}
\loigiai{
Số $1 000$ gồm $1$ nghìn, $0$ trăm, $0$ chục, $0$ đơn vị.
\textbf{Đáp án của câu hỏi là A}
}
\end{ex}

\begin{ex} %[2D10V48]
Có bao nhiêu chục que tính để gộp lại thành đúng $4$ trăm que tính?
\choice
{$4$ chục}
{$14$ chục}
{$400$ chục}
{\True $40$ chục}
\loigiai{
$4$ trăm $= 40$ chục.
\textbf{Đáp án của câu hỏi là D}
}
\end{ex}

\begin{ex} %[2D10V48]
Bác Nam đóng gói cam vào $7$ thùng to, mỗi thùng $100$ quả và $3$ túi nhỏ, mỗi túi $10$ quả. Bác Nam có tất cả bao nhiêu quả cam?
\choice
{$703$ quả}
{\True $730$ quả}
{$73$ quả}
{$370$ quả}
\loigiai{
$7$ trăm và $3$ chục quả cam là $730$ quả cam.
\textbf{Đáp án của câu hỏi là B}
}
\end{ex}

\begin{ex} %[2D10V48]
Một xưởng sản xuất đóng được $1 000$ chiếc quạt. Người ta chở đi phân phối bằng các xe tải nhỏ, mỗi xe chở được $100$ chiếc. Cần bao nhiêu chuyến xe để chở hết số quạt đó?
\choice
{$5$ chuyến}
{$100$ chuyến}
{\True $10$ chuyến}
{$8$ chuyến}
\loigiai{
$1 000$ gồm $10$ trăm, do đó cần $10$ chuyến xe.
\textbf{Đáp án của câu hỏi là C}
}
\end{ex}

\begin{ex} %[2D10V48]
Số gồm $9$ trăm và $10$ chục là số nào?
\choice
{\True $1 000$}
{$910$}
{$901$}
{$900$}
\loigiai{
$10$ chục $= 1$ trăm. $9$ trăm thêm $1$ trăm là $10$ trăm $= 1 000$.
\textbf{Đáp án của câu hỏi là A}
}
\end{ex}

\begin{ex} %[2D10V48]
Một cửa hàng bán hoa kết các bó hoa hồng, mỗi bó có $10$ bông. Để có đủ $300$ bông hoa, cửa hàng cần chuẩn bị bao nhiêu bó hoa như vậy?
\choice
{$3$ bó}
{$300$ bó}
{$33$ bó}
{\True $30$ bó}
\loigiai{
$300$ bông gồm $30$ chục bông, vậy cần $30$ bó hoa.
\textbf{Đáp án của câu hỏi là D}
}
\end{ex}

\begin{ex} %[2D10V48]
Số $1 000$ có giá trị gấp bao nhiêu lần số $100$?
\choice
{$5$ lần}
{\True $10$ lần}
{$100$ lần}
{$2$ lần}
\loigiai{
$1 000 = 100 \times 10$, nên số $1 000$ gấp $10$ lần số $100$.
\textbf{Đáp án của câu hỏi là B}
}
\end{ex}

\textbf{Phần 2. Trắc nghiệm đúng sai}
\setcounterpageref{ex}{0}

\begin{ex} %[2D10V48]
Xét cấu tạo và mối quan hệ giữa các hàng đơn vị, chục, trăm, nghìn:
\choiceTF
{\True Mười chục đơn vị gộp lại bằng một trăm.}
{Mười trăm đơn vị gộp lại bằng mười nghìn.}
{\True Số liền trước của số 1 000 là số 999.}
{\True Một nghìn gồm một trăm chục.}
\loigiai{
\begin{itemize}
    \item Ý a) Đúng: 10 chục = 100 (một trăm).
    \item Ý b) Sai: 10 trăm = 1 nghìn, không phải 10 nghìn.
    \item Ý c) Đúng: 999 + 1 = 1 000.
    \item Ý d) Đúng: 100 chục = 100 x 10 = 1 000.
\end{itemize}
\begin{itemize}
    \item a) \textbf{Đúng}.
    \item b) \textbf{Sai}.
    \item c) \textbf{Đúng}.
    \item d) \textbf{Đúng}.
\end{itemize}
}
\end{ex}

\begin{ex} %[2D10V48]
Một xưởng bánh có 8 thùng bánh to (mỗi thùng 100 cái) và 20 hộp bánh nhỏ (mỗi hộp 10 cái):
\choiceTF
{\True Số bánh trong 8 thùng to là 800 cái.}
{\True Số bánh trong 20 hộp nhỏ là 200 cái.}
{\True Tổng số bánh xưởng có là 1 000 cái.}
{Xưởng cần thêm 100 cái bánh nữa để đủ 1 000 cái.}
\loigiai{
\begin{itemize}
    \item Ý a) Đúng: 8 x 100 = 800 cái.
    \item Ý b) Đúng: 20 x 10 = 200 cái.
    \item Ý c) Đúng: 800 + 200 = 1 000 cái bánh.
    \item Ý d) Sai: Xưởng đã có đủ 1 000 cái rồi.
\end{itemize}
\begin{itemize}
    \item a) \textbf{Đúng}.
    \item b) \textbf{Đúng}.
    \item c) \textbf{Đúng}.
    \item d) \textbf{Sai}.
\end{itemize}
}
\end{ex}

\textbf{Phần 3. Trả lời ngắn}
\setcounterpageref{ex}{0}

\begin{ex} %[2D10V48]
Tìm số liền sau của số lớn nhất có ba chữ số.
\shortans{1 000}
\loigiai{
\textbf{Đáp số:} $1 000$\\
Số lớn nhất có 3 chữ số là 999. Số liền sau là 1 000.
}
\end{ex}

\begin{ex} %[2D10V48]
Một cửa hàng có 6 hộp bút lớn (mỗi hộp 100 chiếc) và 4 hộp bút nhỏ (mỗi hộp 10 chiếc). Cửa hàng có tất cả bao nhiêu chiếc bút?
\shortans{640}
\loigiai{
\textbf{Đáp số:} $640$\\
6 hộp lớn có 600 chiếc, 4 hộp nhỏ có 40 chiếc. Tổng là $600 + 40 = 640$ chiếc.
}
\end{ex}

\begin{ex} %[2D10V48]
Cần thêm bao nhiêu chục que tính vào 8 trăm que tính để được đúng 1 000 que tính?
\shortans{20}
\loigiai{
\textbf{Đáp số:} $20$\\
1 000 que tính = 10 trăm. Cần thêm $10 - 8 = 2$ trăm. 2 trăm = 20 chục.
}
\end{ex}

\begin{ex} %[2D10V48]
Có bao nhiêu chữ số 0 trong cách viết số một nghìn?
\shortans{3}
\loigiai{
\textbf{Đáp số:} $3$\\
Số 1 000 có 3 chữ số 0.
}
\end{ex}

\begin{ex} %[2D10V48]
Một sọt cam có 100 quả. Hỏi 7 sọt cam như thế có bao nhiêu quả cam?
\shortans{700}
\loigiai{
\textbf{Đáp số:} $700$\\
7 x 100 = 700 quả cam.
}
\end{ex}

\begin{ex} %[2D10V48]
Tìm số tự nhiên tròn chục sao cho khi lấy số đó nhân với 10 thì được 500.
\shortans{50}
\loigiai{
\textbf{Đáp số:} $50$\\
Số đó là: $500 : 10 = 50$.
}
\end{ex}

\newpage

\TaoPhieuKiemTra{Chủ đề 10 -- Các số trong phạm vi 1 000}{49}{Các số tròn trăm, tròn chục}{https://loptoanthaylinh.vn/thi/chuyen-de/2/toan2_cd_b49}
\textbf{Phần 1. Trắc nghiệm lựa chọn}
\setcounterpageref{ex}{0}

\begin{ex} %[2D10V49]
Số tròn trăm liền trước của số $600$ là
\choice
{$590$}
{\True $500$}
{$700$}
{$599$}
\loigiai{
Số tròn trăm liền trước của 600 là 500.
\textbf{Đáp án của câu hỏi là B}
}
\end{ex}

\begin{ex} %[2D10V49]
Dãy số tròn trăm nào sau đây được viết theo thứ tự từ bé đến lớn?
\choice
{$900, 800, 700$}
{$200, 400, 300$}
{$600, 400, 200$}
{\True $300, 500, 700$}
\loigiai{
Dãy $300, 500, 700$ xếp tăng dần.
\textbf{Đáp án của câu hỏi là D}
}
\end{ex}

\begin{ex} %[2D10V49]
Số gồm $7$ trăm và $0$ chục, $0$ đơn vị viết là
\choice
{\True $700$}
{$70$}
{$707$}
{$7$}
\loigiai{
Viết là $700$.
\textbf{Đáp án của câu hỏi là A}
}
\end{ex}

\begin{ex} %[2D10V49]
Từ $100$ đến $900$ có tất cả bao nhiêu số tròn trăm?
\choice
{$8$ số}
{$10$ số}
{\True $9$ số}
{$7$ số}
\loigiai{
Các số tròn trăm: 100, 200, 300, 400, 500, 600, 700, 800, 900. Có 9 số.
\textbf{Đáp án của câu hỏi là C}
}
\end{ex}

\begin{ex} %[2D10V49]
Số tròn chục lớn nhất có hai chữ số đem cộng với $100$ thì được kết quả là
\choice
{$199$}
{\True $190$}
{$180$}
{$200$}
\loigiai{
Số tròn chục lớn nhất có hai chữ số là $90$. Ta có: $90 + 100 = 190$.
\textbf{Đáp án của câu hỏi là B}
}
\end{ex}

\begin{ex} %[2D10V49]
Trong các số sau, số nào vừa là số tròn trăm vừa có chữ số hàng trăm là số chẵn lớn hơn $5$?
\choice
{$500$}
{$700$}
{$650$}
{\True $800$}
\loigiai{
Chữ số hàng trăm là số chẵn lớn hơn 5 chỉ có 6 hoặc 8. Ở đây có 800 là số tròn trăm.
\textbf{Đáp án của câu hỏi là D}
}
\end{ex}

\begin{ex} %[2D10V49]
Một nông trại đóng gói $400$ quả trứng gà, mỗi thùng xếp đúng $100$ quả. Hỏi nông trại xếp được bao nhiêu thùng trứng?
\choice
{\True $4$ thùng}
{$40$ thùng}
{$14$ thùng}
{$400$ thùng}
\loigiai{
$400 : 100 = 4$ (thùng).
\textbf{Đáp án của câu hỏi là A}
}
\end{ex}

\begin{ex} %[2D10V49]
Tổng của số tròn trăm bé nhất có ba chữ số và số tròn trăm lớn nhất có ba chữ số là
\choice
{$900$}
{$800$}
{\True $1 000$}
{$1 100$}
\loigiai{
Số tròn trăm bé nhất có 3 chữ số là 100, lớn nhất là 900. Tổng: $100 + 900 = 1 000$.
\textbf{Đáp án của câu hỏi là C}
}
\end{ex}

\begin{ex} %[2D10V49]
Số tròn trăm nào nằm chính giữa $400$ và $600$?
\choice
{$450$}
{\True $500$}
{$550$}
{$300$}
\loigiai{
Số tròn trăm ở chính giữa 400 và 600 là 500.
\textbf{Đáp án của câu hỏi là B}
}
\end{ex}

\begin{ex} %[2D10V49]
Có bao nhiêu số tròn chục lớn hơn $120$ và bé hơn $180$?
\choice
{$4$ số}
{$6$ số}
{$7$ số}
{\True $5$ số}
\loigiai{
Các số là: 130, 140, 150, 160, 170. Có tất cả 5 số.
\textbf{Đáp án của câu hỏi là D}
}
\end{ex}

\begin{ex} %[2D10V49]
Một trang trại nuôi $300$ con gà trống và $500$ con gà mái. Trang trại đó nuôi tất cả bao nhiêu con gà?
\choice
{\True $800$ con}
{$700$ con}
{$900$ con}
{$600$ con}
\loigiai{
$300 + 500 = 800$ (con gà).
\textbf{Đáp án của câu hỏi là A}
}
\end{ex}

\begin{ex} %[2D10V49]
Một bồn chứa $900$ lít nước, người ta đã xả đi $300$ lít nước để tưới vườn. Bồn còn lại bao nhiêu lít nước?
\choice
{$500$ lít}
{$400$ lít}
{\True $600$ lít}
{$700$ lít}
\loigiai{
$900 - 300 = 600$ (lít).
\textbf{Đáp án của câu hỏi là C}
}
\end{ex}

\textbf{Phần 2. Trắc nghiệm đúng sai}
\setcounterpageref{ex}{0}

\begin{ex} %[2D10V49]
Xét tính đúng sai của các khẳng định về số tròn trăm, tròn chục:
\choiceTF
{\True Số 1 000 là một số tròn trăm.}
{Số 350 là một số tròn trăm.}
{\True Giữa hai số tròn trăm 200 và 400 chỉ có đúng một số tròn trăm là 300.}
{\True Tổng của 200 và 800 bằng 1 000.}
\loigiai{
\begin{itemize}
    \item Ý a) Đúng: 1 000 gồm 10 trăm, tận cùng là hai chữ số 0 nên là số tròn trăm.
    \item Ý b) Sai: 350 là số tròn chục nhưng không phải số tròn trăm.
    \item Ý c) Đúng: Chính xác, chỉ có 300.
    \item Ý d) Đúng: 200 + 800 = 1 000.
\end{itemize}
\begin{itemize}
    \item a) \textbf{Đúng}.
    \item b) \textbf{Sai}.
    \item c) \textbf{Đúng}.
    \item d) \textbf{Đúng}.
\end{itemize}
}
\end{ex}

\begin{ex} %[2D10V49]
Cho ba số tròn trăm A = 200, B = 400, C = 600:
\choiceTF
{\True Tổng của A và B bằng C.}
{\True C gấp 3 lần A.}
{Hiệu giữa C và B bé hơn A.}
{\True Tổng của cả ba số là 1 200.}
\loigiai{
\begin{itemize}
    \item Ý a) Đúng: 200 + 400 = 600.
    \item Ý b) Đúng: 600 : 200 = 3.
    \item Ý c) Sai: C - B = 600 - 400 = 200 = A, không bé hơn.
    \item Ý d) Đúng: 200 + 400 + 600 = 1 200.
\end{itemize}
\begin{itemize}
    \item a) \textbf{Đúng}.
    \item b) \textbf{Đúng}.
    \item c) \textbf{Sai}.
    \item d) \textbf{Đúng}.
\end{itemize}
}
\end{ex}

\textbf{Phần 3. Trả lời ngắn}
\setcounterpageref{ex}{0}

\begin{ex} %[2D10V49]
Tính tổng của số tròn trăm lớn nhất có ba chữ số và số tròn chục bé nhất có hai chữ số.
\shortans{910}
\loigiai{
\textbf{Đáp số:} $910$\\
Số tròn trăm lớn nhất có 3 chữ số là 900. Số tròn chục bé nhất có 2 chữ số là 10. Tổng: $900 + 10 = 910$.
}
\end{ex}

\begin{ex} %[2D10V49]
Kho A có 400 bao gạo, kho B có ít hơn kho A 100 bao gạo. Cả hai kho có tất cả bao nhiêu bao gạo?
\shortans{700}
\loigiai{
\textbf{Đáp số:} $700$\\
Kho B có: $400 - 100 = 300$ bao. Cả hai kho: $400 + 300 = 700$ bao.
}
\end{ex}

\begin{ex} %[2D10V49]
Tìm một số tròn trăm, biết rằng nếu lấy 1 000 trừ đi số đó thì được hiệu là 600.
\shortans{400}
\loigiai{
\textbf{Đáp số:} $400$\\
Số cần tìm: $1 000 - 600 = 400$.
}
\end{ex}

\begin{ex} %[2D10V49]
Có bao nhiêu số tròn chục có hai chữ số?
\shortans{9}
\loigiai{
\textbf{Đáp số:} $9$\\
Đó là 10, 20, ..., 90. Có 9 số.
}
\end{ex}

\begin{ex} %[2D10V49]
Một cửa hàng ngày thứ nhất bán 200 m vải, ngày thứ hai bán 300 m vải. Cả hai ngày cửa hàng bán được bao nhiêu mét vải?
\shortans{500}
\loigiai{
\textbf{Đáp số:} $500$\\
200 + 300 = 500 m.
}
\end{ex}

\begin{ex} %[2D10V49]
Tìm số tròn trăm $x$ thỏa mãn: $x + 300 = 700$.
\shortans{400}
\loigiai{
\textbf{Đáp số:} $400$\\
$x = 700 - 300 = 400$.
}
\end{ex}

\newpage

\TaoPhieuKiemTra{Chủ đề 10 -- Các số trong phạm vi 1 000}{50}{So sánh các số tròn trăm, tròn chục}{https://loptoanthaylinh.vn/thi/chuyen-de/2/toan2_cd_b50}
\textbf{Phần 1. Trắc nghiệm lựa chọn}
\setcounterpageref{ex}{0}

\begin{ex} %[2D10V50]
Điền dấu thích hợp vào chỗ chấm: $700 \dots 500 + 200$.
\choice
{\True $=$}
{$>$}
{$<$}
{Không so sánh được}
\loigiai{
$500 + 200 = 700$, vậy $700 = 700$.
\textbf{Đáp án của câu hỏi là A}
}
\end{ex}

\begin{ex} %[2D10V50]
Trong các số $400, 700, 200, 900$, số lớn nhất là
\choice
{$700$}
{$400$}
{\True $900$}
{$200$}
\loigiai{
Số lớn nhất là 900.
\textbf{Đáp án của câu hỏi là C}
}
\end{ex}

\begin{ex} %[2D10V50]
Tìm số tròn trăm $x$ sao cho: $300 < x < 500$. Giá trị của $x$ là
\choice
{$350$}
{\True $400$}
{$450$}
{$600$}
\loigiai{
$x$ tròn trăm nằm giữa 300 và 500 nên $x = 400$.
\textbf{Đáp án của câu hỏi là B}
}
\end{ex}

\begin{ex} %[2D10V50]
Dãy số nào sau đây được sắp xếp theo thứ tự giảm dần?
\choice
{$200, 500, 800$}
{$600, 700, 300$}
{$800, 600, 700$}
{\True $400, 200, 100$}
\loigiai{
Dãy $400, 200, 100$ xếp giảm dần.
\textbf{Đáp án của câu hỏi là D}
}
\end{ex}

\begin{ex} %[2D10V50]
So sánh hai biểu thức: $800 - 300$ và $200 + 400$.
\choice
{\True $800 - 300 < 200 + 400$}
{$800 - 300 > 200 + 400$}
{$800 - 300 = 200 + 400$}
{Không so sánh được}
\loigiai{
$800 - 300 = 500$; $200 + 400 = 600$. Do $500 < 600$ nên $800 - 300 < 200 + 400$.
\textbf{Đáp án của câu hỏi là A}
}
\end{ex}

\begin{ex} %[2D10V50]
Cho các số tròn chục: $230, 280, 250, 210$. Số bé nhất trong dãy là
\choice
{$230$}
{$250$}
{\True $210$}
{$280$}
\loigiai{
Cùng có chữ số hàng trăm là 2, hàng chục nhỏ nhất là 1 nên 210 bé nhất.
\textbf{Đáp án của câu hỏi là C}
}
\end{ex}

\begin{ex} %[2D10V50]
Có bao nhiêu số tròn trăm $x$ thỏa mãn $200 \le x < 700$?
\choice
{$4$ số}
{\True $5$ số}
{$6$ số}
{$3$ số}
\loigiai{
Các số là: 200, 300, 400, 500, 600. Có 5 số.
\textbf{Đáp án của câu hỏi là B}
}
\end{ex}

\begin{ex} %[2D10V50]
Cây cầu thứ nhất dài $300\text{ m}$, cây cầu thứ hai dài $500\text{ m}$. Cây cầu thứ hai dài hơn cây cầu thứ nhất bao nhiêu mét?
\choice
{$100\text{ m}$}
{$300\text{ m}$}
{$400\text{ m}$}
{\True $200\text{ m}$}
\loigiai{
$500 - 300 = 200$ (m).
\textbf{Đáp án của câu hỏi là D}
}
\end{ex}

\begin{ex} %[2D10V50]
Điền số tròn trăm thích hợp vào chỗ chấm: $600 + \dots = 1 000$.
\choice
{\True $400$}
{$300$}
{$500$}
{$200$}
\loigiai{
$1 000 - 600 = 400$.
\textbf{Đáp án của câu hỏi là A}
}
\end{ex}

\begin{ex} %[2D10V50]
Một cửa hàng buổi sáng bán được $300\text{ kg}$ gạo, buổi chiều bán được nhiều hơn buổi sáng $200\text{ kg}$ gạo. Buổi chiều cửa hàng bán được bao nhiêu ki-lô-gam gạo?
\choice
{$400\text{ kg}$}
{$600\text{ kg}$}
{\True $500\text{ kg}$}
{$700\text{ kg}$}
\loigiai{
$300 + 200 = 500$ (kg).
\textbf{Đáp án của câu hỏi là C}
}
\end{ex}

\begin{ex} %[2D10V50]
Trong các khẳng định sau, khẳng định nào SAI?
\choice
{$900 > 800$}
{\True $1 000 < 900$}
{$300 + 400 = 700$}
{$500 - 100 > 300$}
\loigiai{
$1 000 > 900$, do đó $1 000 < 900$ là SAI.
\textbf{Đáp án của câu hỏi là B}
}
\end{ex}

\begin{ex} %[2D10V50]
Có ba bể nước: Bể A chứa $400\text{ l}$, bể B chứa $600\text{ l}$, bể C chứa $500\text{ l}$. Thứ tự các bể từ chứa ít nước nhất đến nhiều nước nhất là
\choice
{Bể B, Bể C, Bể A}
{Bể C, Bể A, Bể B}
{Bể A, Bể B, Bể C}
{\True Bể A, Bể C, Bể B}
\loigiai{
Vì $400 < 500 < 600$ nên thứ tự là Bể A, Bể C, Bể B.
\textbf{Đáp án của câu hỏi là D}
}
\end{ex}

\textbf{Phần 2. Trắc nghiệm đúng sai}
\setcounterpageref{ex}{0}

\begin{ex} %[2D10V50]
Kiểm tra các nhận định so sánh số tròn trăm:
\choiceTF
{1 000 là số tròn trăm lớn nhất có bốn chữ số.}
{\True Nếu a = 400 và b = 300 thì a + b < 800.}
{\True Số tròn chục liền sau số 390 là 400.}
{Hiệu giữa hai số tròn trăm liên tiếp luôn bằng 10.}
\loigiai{
\begin{itemize}
    \item Ý a) Sai: 1 000 là số tròn trăm bé nhất có bốn chữ số.
    \item Ý b) Đúng: 400 + 300 = 700 < 800. Đúng.
    \item Ý c) Đúng: 390 + 10 = 400. Đúng.
    \item Ý d) Sai: Hiệu giữa hai số tròn trăm liên tiếp luôn bằng 100.
\end{itemize}
\begin{itemize}
    \item a) \textbf{Sai}.
    \item b) \textbf{Đúng}.
    \item c) \textbf{Đúng}.
    \item d) \textbf{Sai}.
\end{itemize}
}
\end{ex}

\begin{ex} %[2D10V50]
Cho ba số tròn trăm 300, 500, 800:
\choiceTF
{\True 300 + 500 = 800.}
{\True Số lớn nhất gấp hơn 2 lần số bé nhất.}
{500 là số ở chính giữa 300 và 800.}
{\True Tổng ba số bằng 1 600.}
\loigiai{
\begin{itemize}
    \item Ý a) Đúng: Đúng.
    \item Ý b) Đúng: 800 > 300 x 2 = 600. Đúng.
    \item Ý c) Sai: Khoảng cách 500 - 300 = 200, 800 - 500 = 300 (không bằng nhau).
    \item Ý d) Đúng: 300 + 500 + 800 = 1 600. Đúng.
\end{itemize}
\begin{itemize}
    \item a) \textbf{Đúng}.
    \item b) \textbf{Đúng}.
    \item c) \textbf{Sai}.
    \item d) \textbf{Đúng}.
\end{itemize}
}
\end{ex}

\textbf{Phần 3. Trả lời ngắn}
\setcounterpageref{ex}{0}

\begin{ex} %[2D10V50]
Cho ba số tròn trăm $a, b, c$ với $a < b < c$. Biết $a$ bé nhất có ba chữ số, $c$ lớn nhất có ba chữ số và $b$ ở chính giữa. Tìm $b$.
\shortans{500}
\loigiai{
\textbf{Đáp số:} $500$\\
$a = 100, c = 900$. Số ở chính giữa là 500.
}
\end{ex}

\begin{ex} %[2D10V50]
Đoạn đường từ nhà Nam đến trường dài 700 m, từ nhà Bắc đến trường dài 400 m. Đoạn đường từ nhà Nam dài hơn từ nhà Bắc bao nhiêu mét?
\shortans{300}
\loigiai{
\textbf{Đáp số:} $300$\\
700 - 400 = 300 m.
}
\end{ex}

\begin{ex} %[2D10V50]
Tìm số tròn trăm $x$ biết $900 - x = 200 + 300$.
\shortans{400}
\loigiai{
\textbf{Đáp số:} $400$\\
$200 + 300 = 500$. $x = 900 - 500 = 400$.
}
\end{ex}

\begin{ex} %[2D10V50]
Có bao nhiêu số tròn trăm lớn hơn 150 và bé hơn 650?
\shortans{5}
\loigiai{
\textbf{Đáp số:} $5$\\
Đó là 200, 300, 400, 500, 600. Có 5 số.
}
\end{ex}

\begin{ex} %[2D10V50]
Tính hiệu của số tròn trăm lớn nhất có ba chữ số và số tròn trăm bé nhất có ba chữ số.
\shortans{800}
\loigiai{
\textbf{Đáp số:} $800$\\
900 - 100 = 800.
}
\end{ex}

\begin{ex} %[2D10V50]
Thùng thứ nhất có 300 lít dầu, thùng thứ hai có nhiều hơn 400 lít dầu. Cả hai thùng có bao nhiêu lít dầu?
\shortans{1 000}
\loigiai{
\textbf{Đáp số:} $1 000$\\
Thùng 2 có $300 + 400 = 700$ l. Cả hai thùng có $300 + 700 = 1 000$ l.
}
\end{ex}

\newpage

\TaoPhieuKiemTra{Chủ đề 10 -- Các số trong phạm vi 1 000}{51}{Các số có ba chữ số}{https://loptoanthaylinh.vn/thi/chuyen-de/2/toan2_cd_b51}
\textbf{Phần 1. Trắc nghiệm lựa chọn}
\setcounterpageref{ex}{0}

\begin{ex} %[2D10V51]
Số gồm $5$ trăm, $3$ chục và $8$ đơn vị được viết là
\choice
{$583$}
{$358$}
{$853$}
{\True $538$}
\loigiai{
Viết là $538$.
\textbf{Đáp án của câu hỏi là D}
}
\end{ex}

\begin{ex} %[2D10V51]
Số $706$ được đọc là
\choice
{Bảy mươi sáu}
{\True Bảy trăm linh sáu}
{Bảy trăm sáu mươi}
{Bảy không sáu}
\loigiai{
Số 706 đọc là Bảy trăm linh sáu.
\textbf{Đáp án của câu hỏi là B}
}
\end{ex}

\begin{ex} %[2D10V51]
Chữ số $4$ trong số $842$ có giá trị là
\choice
{$400$}
{$4$}
{\True $40$}
{$14$}
\loigiai{
Chữ số 4 ở hàng chục nên có giá trị là 40.
\textbf{Đáp án của câu hỏi là C}
}
\end{ex}

\begin{ex} %[2D10V51]
Số liền trước của số $300$ là
\choice
{\True $299$}
{$290$}
{$301$}
{$200$}
\loigiai{
$300 - 1 = 299$.
\textbf{Đáp án của câu hỏi là A}
}
\end{ex}

\begin{ex} %[2D10V51]
Từ ba chữ số $2, 5, 8$ (mỗi chữ số viết đúng một lần), có thể lập được bao nhiêu số có ba chữ số khác nhau?
\choice
{$4$ số}
{$5$ số}
{$3$ số}
{\True $6$ số}
\loigiai{
Có $3 \times 2 \times 1 = 6$ số: 258, 285, 528, 582, 825, 852.
\textbf{Đáp án của câu hỏi là D}
}
\end{ex}

\begin{ex} %[2D10V51]
Số có ba chữ số khác nhau lớn nhất được lập từ các chữ số $1, 9, 4$ là
\choice
{$914$}
{\True $941$}
{$491$}
{$194$}
\loigiai{
Xếp giảm dần: hàng trăm là 9, hàng chục là 4, hàng đơn vị là 1: số $941$.
\textbf{Đáp án của câu hỏi là B}
}
\end{ex}

\begin{ex} %[2D10V51]
Một số có chữ số hàng trăm là $6$, hàng đơn vị là $2$ và hàng chục gấp đôi hàng đơn vị. Số đó là
\choice
{$624$}
{$682$}
{\True $642$}
{$622$}
\loigiai{
Hàng đơn vị là 2, hàng chục gấp đôi là 4, hàng trăm là 6: số 642.
\textbf{Đáp án của câu hỏi là C}
}
\end{ex}

\begin{ex} %[2D10V51]
Số gồm $4$ trăm và $5$ đơn vị là số nào?
\choice
{\True $405$}
{$450$}
{$45$}
{$540$}
\loigiai{
Hàng chục bằng 0 nên viết là $405$.
\textbf{Đáp án của câu hỏi là A}
}
\end{ex}

\begin{ex} %[2D10V51]
Số lẻ liền sau số $539$ là
\choice
{$540$}
{$543$}
{$538$}
{\True $541$}
\loigiai{
Số lẻ liền sau là $539 + 2 = 541$.
\textbf{Đáp án của câu hỏi là D}
}
\end{ex}

\begin{ex} %[2D10V51]
Số trang của một cuốn sách có chữ số hàng trăm là $2$, hàng chục lớn hơn hàng trăm $3$ đơn vị, hàng đơn vị kém hàng chục $1$ đơn vị. Cuốn sách có bao nhiêu trang?
\choice
{$243$ trang}
{\True $254$ trang}
{$235$ trang}
{$255$ trang}
\loigiai{
Hàng trăm: 2; hàng chục: $2 + 3 = 5$; hàng đơn vị: $5 - 1 = 4$. Số trang là 254.
\textbf{Đáp án của câu hỏi là B}
}
\end{ex}

\begin{ex} %[2D10V51]
Có bao nhiêu số có ba chữ số có chữ số hàng trăm là $3$ và chữ số tận cùng là $0$?
\choice
{$9$ số}
{$8$ số}
{\True $10$ số}
{$11$ số}
\loigiai{
Các số dạng $3a0$ với $a \in \{0, 1, \dots, 9\}$. Có đúng 10 số.
\textbf{Đáp án của câu hỏi là C}
}
\end{ex}

\begin{ex} %[2D10V51]
Số bé nhất có ba chữ số mà tổng các chữ số bằng $3$ là
\choice
{\True $102$}
{$300$}
{$201$}
{$210$}
\loigiai{
Hàng trăm bé nhất là 1, hàng chục bé nhất là 0, hàng đơn vị là $3 - 1 - 0 = 2$: số 102.
\textbf{Đáp án của câu hỏi là A}
}
\end{ex}

\textbf{Phần 2. Trắc nghiệm đúng sai}
\setcounterpageref{ex}{0}

\begin{ex} %[2D10V51]
Cho số 825, xét tính đúng sai của các khẳng định sau:
\choiceTF
{\True Số 825 gồm 8 trăm, 2 chục và 5 đơn vị.}
{Chữ số 8 có giá trị là 80.}
{\True Số liền sau của 825 là 826.}
{\True Tổng các chữ số của 825 bằng 15.}
\loigiai{
\begin{itemize}
    \item Ý a) Đúng: Đúng cấu tạo thập phân.
    \item Ý b) Sai: Chữ số 8 ở hàng trăm nên có giá trị là 800.
    \item Ý c) Đúng: 825 + 1 = 826.
    \item Ý d) Đúng: 8 + 2 + 5 = 15.
\end{itemize}
\begin{itemize}
    \item a) \textbf{Đúng}.
    \item b) \textbf{Sai}.
    \item c) \textbf{Đúng}.
    \item d) \textbf{Đúng}.
\end{itemize}
}
\end{ex}

\begin{ex} %[2D10V51]
Xét các số có ba chữ số được lập từ các chữ số 0, 3, 7:
\choiceTF
{\True Số lớn nhất có ba chữ số khác nhau lập được là 730.}
{\True Số bé nhất có ba chữ số khác nhau lập được là 307.}
{Có thể lập được 6 số có ba chữ số khác nhau từ 3 chữ số trên.}
{\True Tổng của số lớn nhất và số bé nhất lập được là 1 037.}
\loigiai{
\begin{itemize}
    \item Ý a) Đúng: Đúng, xếp giảm dần: 730.
    \item Ý b) Đúng: Hàng trăm khác 0 nên chọn 3, tiếp theo là 0, 7. Đúng.
    \item Ý c) Sai: Vì chữ số 0 không đứng đầu nên chỉ lập được $2 \times 2 \times 1 = 4$ số.
    \item Ý d) Đúng: 730 + 307 = 1 037. Đúng.
\end{itemize}
\begin{itemize}
    \item a) \textbf{Đúng}.
    \item b) \textbf{Đúng}.
    \item c) \textbf{Sai}.
    \item d) \textbf{Đúng}.
\end{itemize}
}
\end{ex}

\textbf{Phần 3. Trả lời ngắn}
\setcounterpageref{ex}{0}

\begin{ex} %[2D10V51]
Viết số bé nhất có ba chữ số khác nhau được lập từ ba chữ số 0, 7, 3.
\shortans{307}
\loigiai{
\textbf{Đáp số:} $307$\\
Hàng trăm chọn 3, chục chọn 0, đơn vị chọn 7: số 307.
}
\end{ex}

\begin{ex} %[2D10V51]
Tìm một số có ba chữ số biết hàng trăm là 7, hàng chục là 0, xóa chữ số hàng trăm thì được số 9.
\shortans{709}
\loigiai{
\textbf{Đáp số:} $709$\\
Xóa hàng trăm còn lại 09 = 9. Vậy số đó là 709.
}
\end{ex}

\begin{ex} %[2D10V51]
Có bao nhiêu số có ba chữ số giống nhau?
\shortans{9}
\loigiai{
\textbf{Đáp số:} $9$\\
Đó là 111, 222, ..., 999. Có 9 số.
}
\end{ex}

\begin{ex} %[2D10V51]
Tìm số liền sau của số 799.
\shortans{800}
\loigiai{
\textbf{Đáp số:} $800$\\
799 + 1 = 800.
}
\end{ex}

\begin{ex} %[2D10V51]
Số có ba chữ số lớn nhất có các chữ số khác nhau là số nào?
\shortans{987}
\loigiai{
\textbf{Đáp số:} $987$\\
Đó là 987.
}
\end{ex}

\begin{ex} %[2D10V51]
Một số gồm 4 trăm, 0 chục và 9 đơn vị. Viết số đó.
\shortans{409}
\loigiai{
\textbf{Đáp số:} $409$\\
Viết là 409.
}
\end{ex}

\newpage

\TaoPhieuKiemTra{Chủ đề 10 -- Các số trong phạm vi 1 000}{52}{Viết số thành tổng các trăm, chục, đơn vị}{https://loptoanthaylinh.vn/thi/chuyen-de/2/toan2_cd_b52}
\textbf{Phần 1. Trắc nghiệm lựa chọn}
\setcounterpageref{ex}{0}

\begin{ex} %[2D10V52]
Số $645$ được viết thành tổng các trăm, chục, đơn vị là
\choice
{$600 + 45$}
{$640 + 5$}
{\True $600 + 40 + 5$}
{$6 + 4 + 5$}
\loigiai{
$645 = 600 + 40 + 5$.
\textbf{Đáp án của câu hỏi là C}
}
\end{ex}

\begin{ex} %[2D10V52]
Tổng $800 + 70 + 2$ bằng số nào sau đây?
\choice
{$827$}
{\True $872$}
{$807$}
{$782$}
\loigiai{
$800 + 70 + 2 = 872$.
\textbf{Đáp án của câu hỏi là B}
}
\end{ex}

\begin{ex} %[2D10V52]
Số $509$ được viết thành tổng là
\choice
{$500 + 90$}
{$50 + 9$}
{$500 + 90 + 0$}
{\True $500 + 9$}
\loigiai{
Hàng chục bằng 0 nên $509 = 500 + 9$.
\textbf{Đáp án của câu hỏi là D}
}
\end{ex}

\begin{ex} %[2D10V52]
Cho số có dạng $a37 = 400 + 30 + 7$. Chữ số $a$ là
\choice
{\True $4$}
{$3$}
{$7$}
{$0$}
\loigiai{
Chữ số hàng trăm là 4.
\textbf{Đáp án của câu hỏi là A}
}
\end{ex}

\begin{ex} %[2D10V52]
Số nào dưới đây có chữ số hàng chục bằng $0$ và được viết thành tổng $700 + 8$?
\choice
{$780$}
{$788$}
{\True $708$}
{$807$}
\loigiai{
$700 + 8 = 708$.
\textbf{Đáp án của câu hỏi là C}
}
\end{ex}

\begin{ex} %[2D10V52]
Biết $x = 300 + 60 + 4$ và $y = 300 + 50 + 9$. Khẳng định nào sau đây là đúng?
\choice
{$x < y$}
{\True $x > y$}
{$x = y$}
{Không so sánh được}
\loigiai{
$x = 364, y = 359$. Vì $364 > 359$ nên $x > y$.
\textbf{Đáp án của câu hỏi là B}
}
\end{ex}

\begin{ex} %[2D10V52]
Một số gồm $b$ trăm, $5$ chục và $1$ đơn vị. Biết $b$ là số liền sau số $6$. Số đó là
\choice
{$561$}
{$651$}
{$851$}
{\True $751$}
\loigiai{
Liền sau 6 là 7, nên $b = 7$. Số đó là 751.
\textbf{Đáp án của câu hỏi là D}
}
\end{ex}

\begin{ex} %[2D10V52]
Tổng nào sau đây có giá trị bằng $990$?
\choice
{\True $900 + 90$}
{$900 + 9$}
{$90 + 9$}
{$900 + 90 + 9$}
\loigiai{
$900 + 90 = 990$.
\textbf{Đáp án của câu hỏi là A}
}
\end{ex}

\begin{ex} %[2D10V52]
Cho số $435$. Nếu bớt đi $3$ chục và $5$ đơn vị thì số đó còn lại bao nhiêu?
\choice
{$430$}
{$405$}
{\True $400$}
{$300$}
\loigiai{
$435 - 35 = 400$.
\textbf{Đáp án của câu hỏi là C}
}
\end{ex}

\begin{ex} %[2D10V52]
Số chẵn lớn nhất có ba chữ số được phân tích thành tổng các trăm, chục, đơn vị là
\choice
{$900 + 90 + 9$}
{\True $900 + 90 + 8$}
{$900 + 80 + 8$}
{$800 + 90 + 8$}
\loigiai{
Số chẵn lớn nhất có 3 chữ số là 998. Phân tích: $900 + 90 + 8$.
\textbf{Đáp án của câu hỏi là B}
}
\end{ex}

\begin{ex} %[2D10V52]
Một cửa hàng bán được $200 + 40 + 8$ mét vải. Buổi sáng bán được $100$ mét. Buổi chiều bán được bao nhiêu mét vải?
\choice
{$138\text{ m}$}
{$140\text{ m}$}
{$108\text{ m}$}
{\True $148\text{ m}$}
\loigiai{
Tổng số vải là 248 m. Buổi chiều bán: $248 - 100 = 148$ m.
\textbf{Đáp án của câu hỏi là D}
}
\end{ex}

\begin{ex} %[2D10V52]
Tìm $x$ biết: $x - 70 = 500 + 3$.
\choice
{\True $573$}
{$503$}
{$570$}
{$433$}
\loigiai{
$x - 70 = 503 \Rightarrow x = 503 + 70 = 573$.
\textbf{Đáp án của câu hỏi là A}
}
\end{ex}

\textbf{Phần 2. Trắc nghiệm đúng sai}
\setcounterpageref{ex}{0}

\begin{ex} %[2D10V52]
Xét cấu tạo phân tích số thành tổng trăm, chục, đơn vị:
\choiceTF
{\True Số 320 phân tích thành 300 + 20.}
{Số 406 phân tích thành 400 + 60.}
{\True Biểu thức 900 + 90 + 9 có giá trị bằng số lớn nhất có ba chữ số.}
{\True Số 604 được viết là 600 + 4.}
\loigiai{
\begin{itemize}
    \item Ý a) Đúng: Đúng, hàng đơn vị bằng 0.
    \item Ý b) Sai: Sai, phải là 400 + 6.
    \item Ý c) Đúng: 900 + 90 + 9 = 999.
    \item Ý d) Đúng: Đúng với cấu tạo số có hàng chục bằng 0.
\end{itemize}
\begin{itemize}
    \item a) \textbf{Đúng}.
    \item b) \textbf{Sai}.
    \item c) \textbf{Đúng}.
    \item d) \textbf{Đúng}.
\end{itemize}
}
\end{ex}

\begin{ex} %[2D10V52]
Cho số M = 746:
\choiceTF
{\True M = 700 + 40 + 6.}
{\True Giá trị của chữ số 7 gấp 100 lần giá trị của chữ số 7 ở hàng đơn vị (nếu có).}
{Nếu xóa chữ số 4 thì được số 76.}
{\True Tổng các chữ số của M là 17.}
\loigiai{
\begin{itemize}
    \item Ý a) Đúng: Đúng cấu tạo.
    \item Ý b) Đúng: 700 : 7 = 100. Đúng.
    \item Ý c) Sai: Xóa chữ số 4 thì hàng trăm là 7, đơn vị là 6 thành số 76 (sai, thành 76 thì mất chữ số hàng trăm, thực tế là 76 nếu bỏ hàng chục nhưng theo vị trí là 76). Khẳng định đúng là số còn lại có 2 chữ số.
    \item Ý d) Đúng: 7 + 4 + 6 = 17. Đúng.
\end{itemize}
\begin{itemize}
    \item a) \textbf{Đúng}.
    \item b) \textbf{Đúng}.
    \item c) \textbf{Sai}.
    \item d) \textbf{Đúng}.
\end{itemize}
}
\end{ex}

\textbf{Phần 3. Trả lời ngắn}
\setcounterpageref{ex}{0}

\begin{ex} %[2D10V52]
Tìm số có ba chữ số biết khi xóa chữ số 5 ở hàng trăm thì được số 48.
\shortans{548}
\loigiai{
\textbf{Đáp số:} $548$\\
Hàng trăm là 5, hai hàng sau là 48: số 548.
}
\end{ex}

\begin{ex} %[2D10V52]
Tính giá trị của biểu thức: $A = 700 + 80 + 5 - 200 - 5$.
\shortans{580}
\loigiai{
\textbf{Đáp số:} $580$\\
$A = 785 - 205 = 580$.
}
\end{ex}

\begin{ex} %[2D10V52]
Bác nông dân thu được 6 trăm, 8 chục và 9 bắp ngô. Bác biếu 89 bắp. Còn lại bao nhiêu bắp ngô?
\shortans{600}
\loigiai{
\textbf{Đáp số:} $600$\\
Tổng: 689 bắp. Còn lại: $689 - 89 = 600$ bắp ngô.
}
\end{ex}

\begin{ex} %[2D10V52]
Viết số gồm 9 trăm, 9 chục và 9 đơn vị.
\shortans{999}
\loigiai{
\textbf{Đáp số:} $999$\\
Viết là 999.
}
\end{ex}

\begin{ex} %[2D10V52]
Tìm số tự nhiên biết số đó bằng tổng của 3 trăm và 5 đơn vị.
\shortans{305}
\loigiai{
\textbf{Đáp số:} $305$\\
$300 + 5 = 305$.
}
\end{ex}

\begin{ex} %[2D10V52]
Cho $x = 400 + 20 + 8$. Tính $x - 28$.
\shortans{400}
\loigiai{
\textbf{Đáp số:} $400$\\
$428 - 28 = 400$.
}
\end{ex}

\newpage

\TaoPhieuKiemTra{Chủ đề 10 -- Các số trong phạm vi 1 000}{53}{So sánh các số có ba chữ số}{https://loptoanthaylinh.vn/thi/chuyen-de/2/toan2_cd_b53}
\textbf{Phần 1. Trắc nghiệm lựa chọn}
\setcounterpageref{ex}{0}

\begin{ex} %[2D10V53]
Điền dấu thích hợp vào chỗ chấm: $645 \dots 654$.
\choice
{$>$}
{\True $<$}
{$=$}
{Không so sánh được}
\loigiai{
Hàng trăm bằng nhau, hàng chục $4 < 5$ nên $645 < 654$.
\textbf{Đáp án của câu hỏi là B}
}
\end{ex}

\begin{ex} %[2D10V53]
Trong các số $789, 798, 879, 897$, số lớn nhất là
\choice
{\True $897$}
{$798$}
{$879$}
{$789$}
\loigiai{
So sánh hàng trăm, 8 lớn hơn 7. So sánh tiếp $897 > 879$.
\textbf{Đáp án của câu hỏi là A}
}
\end{ex}

\begin{ex} %[2D10V53]
Số nào dưới đây lớn hơn $456$ và bé hơn $462$?
\choice
{$452$}
{$465$}
{\True $459$}
{$470$}
\loigiai{
$456 < 459 < 462$.
\textbf{Đáp án của câu hỏi là C}
}
\end{ex}

\begin{ex} %[2D10V53]
Cho các số: $341, 314, 431, 413$. Thứ tự các số từ bé đến lớn là
\choice
{$431, 413, 341, 314$}
{$341, 314, 413, 431$}
{$314, 413, 341, 431$}
{\True $314, 341, 413, 431$}
\loigiai{
$314 < 341 < 413 < 431$.
\textbf{Đáp án của câu hỏi là D}
}
\end{ex}

\begin{ex} %[2D10V53]
Điền chữ số thích hợp vào dấu $*$ để: $5*8 > 588$.
\choice
{$8$}
{\True $9$}
{$7$}
{$0$}
\loigiai{
Để $5*8 > 588$ thì chữ số hàng chục $*$ phải lớn hơn 8, vậy $* = 9$.
\textbf{Đáp án của câu hỏi là B}
}
\end{ex}

\begin{ex} %[2D10V53]
Điền dấu thích hợp: $400 + 50 + 8 \dots 400 + 60$.
\choice
{\True $<$}
{$>$}
{$=$}
{Không so sánh được}
\loigiai{
$400 + 50 + 8 = 458$; $400 + 60 = 460$. Vì $458 < 460$ nên điền dấu $<$. 
\textbf{Đáp án của câu hỏi là A}
}
\end{ex}

\begin{ex} %[2D10V53]
Có bao nhiêu số tự nhiên $x$ có ba chữ số thỏa mãn: $198 < x < 203$?
\choice
{$3$ số}
{$5$ số}
{\True $4$ số}
{$2$ số}
\loigiai{
Các số là: 199, 200, 201, 202. Có tất cả 4 số.
\textbf{Đáp án của câu hỏi là C}
}
\end{ex}

\begin{ex} %[2D10V53]
Lớp 2A quyên góp được $325$ quyển vở, lớp 2B được $340$ quyển vở, lớp 2C được $318$ quyển vở. Lớp quyên góp được nhiều vở nhất là
\choice
{Lớp 2A}
{Lớp 2C}
{Cả ba lớp bằng nhau}
{\True Lớp 2B}
\loigiai{
$340 > 325 > 318$ nên lớp 2B quyên góp nhiều nhất.
\textbf{Đáp án của câu hỏi là D}
}
\end{ex}

\begin{ex} %[2D10V53]
Tìm số có ba chữ số bé nhất trong các số sau:
\choice
{$205$}
{\True $189$}
{$250$}
{$199$}
\loigiai{
So sánh hàng trăm 1 và 2, giữa 189 và 199 thì hàng chục $8 < 9$ nên 189 bé nhất.
\textbf{Đáp án của câu hỏi là B}
}
\end{ex}

\begin{ex} %[2D10V53]
Khoảng cách từ Hà Nội đến Vinh khoảng $300\text{ km}$, đến Thanh Hóa khoảng $160\text{ km}$, đến Huế khoảng $660\text{ km}$. Nơi nào xa Hà Nội nhất?
\choice
{\True Huế}
{Vinh}
{Thanh Hóa}
{Ba nơi xa như nhau}
\loigiai{
$660 > 300 > 160$ nên Huế xa Hà Nội nhất.
\textbf{Đáp án của câu hỏi là A}
}
\end{ex}

\begin{ex} %[2D10V53]
Cho số $A = 700 + 30 + 5$ và số $B = 739$. Khẳng định nào sau đây ĐÚNG?
\choice
{$A > B$}
{$A = B$}
{\True $A < B$}
{$A - B = 4$}
\loigiai{
$A = 735 < 739 = B$.
\textbf{Đáp án của câu hỏi là C}
}
\end{ex}

\begin{ex} %[2D10V53]
Dãy số nào sau đây gồm toàn các số lẻ và được xếp theo thứ tự giảm dần?
\choice
{$888, 777, 666$}
{$135, 357, 579$}
{$999, 998, 997$}
{\True $987, 875, 763$}
\loigiai{
Dãy $987, 875, 763$ gồm toàn số lẻ và giảm dần.
\textbf{Đáp án của câu hỏi là D}
}
\end{ex}

\textbf{Phần 2. Trắc nghiệm đúng sai}
\setcounterpageref{ex}{0}

\begin{ex} %[2D10V53]
Xét các mệnh đề so sánh số có ba chữ số:
\choiceTF
{\True Số nào có chữ số hàng trăm lớn hơn thì lớn hơn.}
{\True 999 là số có ba chữ số lớn nhất.}
{Số 542 bé hơn số 524.}
{\True Có duy nhất một số tự nhiên nằm giữa 699 và 701.}
\loigiai{
\begin{itemize}
    \item Ý a) Đúng: Khi so sánh 2 số có 3 chữ số, số có chữ số hàng trăm lớn hơn thì lớn hơn.
    \item Ý b) Đúng: Đúng, số liền sau là 1 000 có 4 chữ số.
    \item Ý c) Sai: Hàng chục $4 > 2$ nên $542 > 524$.
    \item Ý d) Đúng: Đó là số 700.
\end{itemize}
\begin{itemize}
    \item a) \textbf{Đúng}.
    \item b) \textbf{Đúng}.
    \item c) \textbf{Sai}.
    \item d) \textbf{Đúng}.
\end{itemize}
}
\end{ex}

\begin{ex} %[2D10V53]
Cho bốn số 425, 452, 524, 542:
\choiceTF
{\True Số bé nhất trong bốn số là 425.}
{\True Số lớn nhất trong bốn số là 542.}
{452 lớn hơn 524.}
{\True Hiệu giữa số lớn nhất và số bé nhất là 117.}
\loigiai{
\begin{itemize}
    \item Ý a) Đúng: Đúng.
    \item Ý b) Đúng: Đúng.
    \item Ý c) Sai: Hàng trăm $4 < 5$ nên $452 < 524$.
    \item Ý d) Đúng: 542 - 425 = 117. Đúng.
\end{itemize}
\begin{itemize}
    \item a) \textbf{Đúng}.
    \item b) \textbf{Đúng}.
    \item c) \textbf{Sai}.
    \item d) \textbf{Đúng}.
\end{itemize}
}
\end{ex}

\textbf{Phần 3. Trả lời ngắn}
\setcounterpageref{ex}{0}

\begin{ex} %[2D10V53]
Cho bốn chữ số 3, 0, 8, 5. Viết số chẵn lớn nhất có ba chữ số khác nhau.
\shortans{850}
\loigiai{
\textbf{Đáp số:} $850$\\
Hàng trăm 8, chục 5, đơn vị chẵn 0: số 850.
}
\end{ex}

\begin{ex} %[2D10V53]
Có bao nhiêu số có ba chữ số lớn hơn 895 và bé hơn 905?
\shortans{9}
\loigiai{
\textbf{Đáp số:} $9$\\
Từ 896 đến 904 có: $904 - 896 + 1 = 9$ số.
}
\end{ex}

\begin{ex} %[2D10V53]
Tìm hiệu giữa số lớn nhất có ba chữ số khác nhau và số bé nhất có ba chữ số khác nhau.
\shortans{885}
\loigiai{
\textbf{Đáp số:} $885$\\
Số lớn nhất: 987. Số bé nhất: 102. Hiệu: $987 - 102 = 885$.
}
\end{ex}

\begin{ex} %[2D10V53]
Tìm chữ số thích hợp thay cho $*$ để $478 < 4*8 < 498$.
\shortans{8}
\loigiai{
\textbf{Đáp số:} $8$\\
Hàng chục $*$ phải lớn hơn 7 và bé hơn 9 nên $* = 8$.
}
\end{ex}

\begin{ex} %[2D10V53]
Sắp xếp các số 502, 205, 520, 250 theo thứ tự tăng dần. Số thứ hai trong dãy là số nào?
\shortans{250}
\loigiai{
\textbf{Đáp số:} $250$\\
Thứ tự: 205, 250, 502, 520. Số thứ hai là 250.
}
\end{ex}

\begin{ex} %[2D10V53]
Tìm số tự nhiên liền trước của số có ba chữ số bé nhất.
\shortans{99}
\loigiai{
\textbf{Đáp số:} $99$\\
Số bé nhất có 3 chữ số là 100. Liền trước là 99.
}
\end{ex}

\newpage

\TaoPhieuKiemTra{Chủ đề 10 -- Các số trong phạm vi 1 000}{54}{Luyện tập chung}{https://loptoanthaylinh.vn/thi/chuyen-de/2/toan2_cd_b54}
\textbf{Phần 1. Trắc nghiệm lựa chọn}
\setcounterpageref{ex}{0}

\begin{ex} %[2D10V54]
Số liền trước của số chẵn bé nhất có ba chữ số là
\choice
{$101$}
{$100$}
{\True $99$}
{$98$}
\loigiai{
Số chẵn bé nhất có 3 chữ số là 100. Số liền trước là 99.
\textbf{Đáp án của câu hỏi là C}
}
\end{ex}

\begin{ex} %[2D10V54]
Số gồm $6$ trăm và $6$ đơn vị được viết là
\choice
{\True $606$}
{$660$}
{$66$}
{$6006$}
\loigiai{
Hàng chục bằng 0 nên viết là $606$.
\textbf{Đáp án của câu hỏi là A}
}
\end{ex}

\begin{ex} %[2D10V54]
Tổng của số lớn nhất có hai chữ số và số bé nhất có ba chữ số là
\choice
{$198$}
{$200$}
{$109$}
{\True $199$}
\loigiai{
$99 + 100 = 199$.
\textbf{Đáp án của câu hỏi là D}
}
\end{ex}

\begin{ex} %[2D10V54]
Trong các số sau, số nào gần với số $350$ nhất?
\choice
{$341$}
{\True $349$}
{$359$}
{$339$}
\loigiai{
349 cách 350 đúng 1 đơn vị, gần nhất.
\textbf{Đáp án của câu hỏi là B}
}
\end{ex}

\begin{ex} %[2D10V54]
Cho số có ba chữ số, biết chữ số hàng trăm gấp $3$ lần hàng đơn vị và hàng chục bằng $0$. Số đó có thể là
\choice
{$309$}
{$603$}
{\True $903$}
{$301$}
\loigiai{
Chữ số hàng trăm là 9, hàng đơn vị là 3, vì $9 : 3 = 3$. Số đó là 903.
\textbf{Đáp án của câu hỏi là C}
}
\end{ex}

\begin{ex} %[2D10V54]
Tìm một số biết rằng nếu lấy số đó cộng với $200$ rồi trừ đi $50$ thì được $650$.
\choice
{\True $500$}
{$400$}
{$600$}
{$450$}
\loigiai{
Số đó là: $650 + 50 - 200 = 500$.
\textbf{Đáp án của câu hỏi là A}
}
\end{ex}

\begin{ex} %[2D10V54]
Đếm thêm $5$ đơn vị từ số $245$ đến số $265$ thì được bao nhiêu số?
\choice
{$4$ số}
{$6$ số}
{$7$ số}
{\True $5$ số}
\loigiai{
Các số là: 245, 250, 255, 260, 265. Có 5 số.
\textbf{Đáp án của câu hỏi là D}
}
\end{ex}

\begin{ex} %[2D10V54]
Một rạp chiếu phim có hai phòng chiếu. Phòng một có $250$ ghế, phòng hai có $280$ ghế. Cả hai phòng có bao nhiêu ghế?
\choice
{$520$ ghế}
{\True $530$ ghế}
{$500$ ghế}
{$430$ ghế}
\loigiai{
$250 + 280 = 530$ (ghế).
\textbf{Đáp án của câu hỏi là B}
}
\end{ex}

\begin{ex} %[2D10V54]
Hình tam giác có độ dài ba cạnh đều bằng $100\text{ cm}$. Chu vi của hình tam giác đó là
\choice
{$200\text{ cm}$}
{$400\text{ cm}$}
{\True $300\text{ cm}$}
{$100\text{ cm}$}
\loigiai{
$100 \times 3 = 300$ (cm).
\textbf{Đáp án của câu hỏi là C}
}
\end{ex}

\begin{ex} %[2D10V54]
Tìm số tự nhiên có ba chữ số biết chữ số hàng trăm là $4$, nếu xóa chữ số hàng trăm thì được số $25$.
\choice
{\True $425$}
{$452$}
{$420$}
{$254$}
\loigiai{
Số đó là 425.
\textbf{Đáp án của câu hỏi là A}
}
\end{ex}

\begin{ex} %[2D10V54]
Một cửa hàng bán được $100\text{ kg}$ đường vào ngày thứ nhất. Ngày thứ hai bán gấp đôi ngày thứ nhất. Cả hai ngày bán được bao nhiêu ki-lô-gam đường?
\choice
{$200\text{ kg}$}
{$250\text{ kg}$}
{$350\text{ kg}$}
{\True $300\text{ kg}$}
\loigiai{
Ngày thứ hai: $100 \times 2 = 200$ kg. Cả hai ngày: $100 + 200 = 300$ kg.
\textbf{Đáp án của câu hỏi là D}
}
\end{ex}

\begin{ex} %[2D10V54]
Khối 2 trồng được $400$ cây, khối 3 trồng nhiều hơn khối 2 là $200$ cây. Cả hai khối trồng được tất cả bao nhiêu cây?
\choice
{$600$ cây}
{\True $1 000$ cây}
{$800$ cây}
{$900$ cây}
\loigiai{
Khối 3: $400 + 200 = 600$ cây. Cả hai khối: $400 + 600 = 1 000$ cây.
\textbf{Đáp án của câu hỏi là B}
}
\end{ex}

\textbf{Phần 2. Trắc nghiệm đúng sai}
\setcounterpageref{ex}{0}

\begin{ex} %[2D10V54]
Xét tính đúng sai của các khẳng định về số học phạm vi 1 000:
\choiceTF
{\True 1 000 là số tự nhiên có bốn chữ số nhỏ nhất.}
{Số 780 đọc là bảy trăm linh tám.}
{\True Hiệu giữa 1 000 và 1 bằng 999.}
{Số gồm 3 trăm và 4 chục là 304.}
\loigiai{
\begin{itemize}
    \item Ý a) Đúng: Đúng.
    \item Ý b) Sai: 780 đọc là bảy trăm tám mươi.
    \item Ý c) Đúng: 1 000 - 1 = 999.
    \item Ý d) Sai: 3 trăm và 4 chục là 340, không phải 304.
\end{itemize}
\begin{itemize}
    \item a) \textbf{Đúng}.
    \item b) \textbf{Sai}.
    \item c) \textbf{Đúng}.
    \item d) \textbf{Sai}.
\end{itemize}
}
\end{ex}

\begin{ex} %[2D10V54]
Cho dãy số: 100, 200, 300, 400, 500, 600, 700, 800, 900, 1 000:
\choiceTF
{\True Dãy số gồm 10 số đều là các số tròn trăm.}
{\True Mỗi số trong dãy cách số liền sau 100 đơn vị.}
{Số ở chính giữa dãy số là 500.}
{\True Tổng của số đầu và số cuối là 1 100.}
\loigiai{
\begin{itemize}
    \item Ý a) Đúng: Đúng, 10 số từ 100 đến 1 000.
    \item Ý b) Đúng: Đúng.
    \item Ý c) Sai: Dãy có 10 số nên không có một số chính giữa (ở giữa là 500 và 600).
    \item Ý d) Đúng: 100 + 1 000 = 1 100. Đúng.
\end{itemize}
\begin{itemize}
    \item a) \textbf{Đúng}.
    \item b) \textbf{Đúng}.
    \item c) \textbf{Sai}.
    \item d) \textbf{Đúng}.
\end{itemize}
}
\end{ex}

\textbf{Phần 3. Trả lời ngắn}
\setcounterpageref{ex}{0}

\begin{ex} %[2D10V54]
Tìm số có ba chữ số biết chữ số hàng trăm là 5, hàng chục gấp đôi hàng đơn vị và tổng ba chữ số bằng 11.
\shortans{542}
\loigiai{
\textbf{Đáp số:} $542$\\
Tổng hàng chục và đơn vị là $11 - 5 = 6$. Đơn vị là $6 : 3 = 2$, chục là 4: số 542.
}
\end{ex}

\begin{ex} %[2D10V54]
Bác An có 800 quả trứng, bán đi 500 quả. Số còn lại chia đều vào 3 thùng. Mỗi thùng có bao nhiêu quả trứng?
\shortans{100}
\loigiai{
\textbf{Đáp số:} $100$\\
Còn lại: $800 - 500 = 300$ quả. Mỗi thùng: $300 : 3 = 100$ quả.
}
\end{ex}

\begin{ex} %[2D10V54]
Có bao nhiêu số có ba chữ số mà hàng trăm là 7 và hàng đơn vị là 5?
\shortans{10}
\loigiai{
\textbf{Đáp số:} $10$\\
Dạng $7a5$ với $a \in \{0, 1, \dots, 9\}$. Có 10 số.
}
\end{ex}

\begin{ex} %[2D10V54]
Tính tổng của số chẵn lớn nhất có ba chữ số và số lẻ bé nhất có một chữ số.
\shortans{999}
\loigiai{
\textbf{Đáp số:} $999$\\
998 + 1 = 999.
}
\end{ex}

\begin{ex} %[2D10V54]
Một đội xe có 3 xe chở hàng, mỗi xe chở 300 bao xi măng. Đội xe chở tất cả bao nhiêu bao xi măng?
\shortans{900}
\loigiai{
\textbf{Đáp số:} $900$\\
3 x 300 = 900 bao.
}
\end{ex}

\begin{ex} %[2D10V54]
Cho số 654. Nếu đổi chỗ chữ số hàng trăm và chữ số hàng đơn vị thì được số mới là bao nhiêu?
\shortans{456}
\loigiai{
\textbf{Đáp số:} $456$\\
Đổi 6 và 4 ta được 456.
}
\end{ex}

\newpage


% =========================================================================
% CHỦ ĐỀ 11 - ĐỘ DÀI VÀ ĐƠN VỊ ĐO ĐỘ DÀI. TIỀN VIỆT NAM
% =========================================================================
\section{Độ dài và đơn vị đo độ dài. tiền việt nam}

\TaoPhieuKiemTra{Chủ đề 11 -- Độ dài và đơn vị đo độ dài. Tiền Việt Nam}{55}{Đề-xi-mét, mét, ki-lô-mét}{https://loptoanthaylinh.vn/thi/chuyen-de/2/toan2_cd_b55}
\textbf{Phần 1. Trắc nghiệm lựa chọn}
\setcounterpageref{ex}{0}

\begin{ex} %[2D11V55]
$1\text{ m}$ bằng bao nhiêu đề-xi-mét?
\choice
{$100\text{ dm}$}
{\True $10\text{ dm}$}
{$1\text{ dm}$}
{$1 000\text{ dm}$}
\loigiai{
$1\text{ m} = 10\text{ dm}$.
\textbf{Đáp án của câu hỏi là B}
}
\end{ex}

\begin{ex} %[2D11V55]
$1\text{ km}$ bằng bao nhiêu mét?
\choice
{$10\text{ m}$}
{$100\text{ m}$}
{$10 000\text{ m}$}
{\True $1 000\text{ m}$}
\loigiai{
$1\text{ km} = 1 000\text{ m}$.
\textbf{Đáp án của câu hỏi là D}
}
\end{ex}

\begin{ex} %[2D11V55]
Độ dài chiếc bút chì thông thường khoảng bao nhiêu?
\choice
{\True $15\text{ cm}$}
{$15\text{ m}$}
{$15\text{ km}$}
{$15\text{ dm}$}
\loigiai{
Bút chì dài khoảng 15 cm.
\textbf{Đáp án của câu hỏi là A}
}
\end{ex}

\begin{ex} %[2D11V55]
Một sợi dây dài $2\text{ m}$. Người ta cắt đi $6\text{ dm}$. Chiều dài sợi dây còn lại là
\choice
{$16\text{ dm}$}
{$4\text{ dm}$}
{\True $14\text{ dm}$}
{$14\text{ m}$}
\loigiai{
Đổi $2\text{ m} = 20\text{ dm}$. Còn lại: $20 - 6 = 14$ dm.
\textbf{Đáp án của câu hỏi là C}
}
\end{ex}

\begin{ex} %[2D11V55]
Một quãng đường gồm hai chặng: chặng một dài $400\text{ m}$, chặng hai dài $600\text{ m}$. Cả hai chặng đường dài bao nhiêu ki-lô-mét?
\choice
{$10\text{ km}$}
{\True $1\text{ km}$}
{$2\text{ km}$}
{$100\text{ km}$}
\loigiai{
$400 + 600 = 1 000\text{ m} = 1\text{ km}$.
\textbf{Đáp án của câu hỏi là B}
}
\end{ex}

\begin{ex} %[2D11V55]
Điền số thích hợp: $3\text{ m } 5\text{ dm} = \dots\text{ dm}$.
\choice
{$305\text{ dm}$}
{$8\text{ dm}$}
{$53\text{ dm}$}
{\True $35\text{ dm}$}
\loigiai{
$3\text{ m} = 30\text{ dm} \Rightarrow 3\text{ m } 5\text{ dm} = 35$ dm.
\textbf{Đáp án của câu hỏi là D}
}
\end{ex}

\begin{ex} %[2D11V55]
Con đường từ nhà Mai đến trường dài $500\text{ m}$. Mai đi từ nhà đến trường rồi lại đi từ trường về nhà. Mai đã đi tất cả bao nhiêu mét?
\choice
{\True $1 000\text{ m}$}
{$500\text{ m}$}
{$800\text{ m}$}
{$100\text{ m}$}
\loigiai{
$500 + 500 = 1 000$ (m).
\textbf{Đáp án của câu hỏi là A}
}
\end{ex}

\begin{ex} %[2D11V55]
Điền dấu thích hợp: $4\text{ m} \dots 38\text{ dm}$.
\choice
{$<$}
{$=$}
{\True $>$}
{Không so sánh được}
\loigiai{
$4\text{ m} = 40\text{ dm} > 38\text{ dm}$.
\textbf{Đáp án của câu hỏi là C}
}
\end{ex}

\begin{ex} %[2D11V55]
Bác thợ có thanh gỗ dài $1\text{ m}$. Bác cưa thành các đoạn ngắn dài $2\text{ dm}$. Bác cưa được bao nhiêu đoạn?
\choice
{$10$ đoạn}
{\True $5$ đoạn}
{$2$ đoạn}
{$8$ đoạn}
\loigiai{
$1\text{ m} = 10\text{ dm}$. Số đoạn: $10 : 2 = 5$ đoạn.
\textbf{Đáp án của câu hỏi là B}
}
\end{ex}

\begin{ex} %[2D11V55]
Một cái hồ hình chữ nhật có chiều dài $300\text{ m}$, chiều rộng $200\text{ m}$. Vận động viên chạy một vòng quanh bờ hồ hết bao nhiêu mét?
\choice
{$500\text{ m}$}
{$600\text{ m}$}
{$800\text{ m}$}
{\True $1 000\text{ m}$}
\loigiai{
Chu vi bờ hồ: $(300 + 200) \times 2 = 1 000$ m.
\textbf{Đáp án của câu hỏi là D}
}
\end{ex}

\begin{ex} %[2D11V55]
Trong các đơn vị đo độ dài mm, cm, dm, m, km, đơn vị đo nào lớn nhất?
\choice
{\True km}
{m}
{cm}
{dm}
\loigiai{
Ki-lô-mét (km) là đơn vị đo độ dài lớn nhất trong các đơn vị đã học.
\textbf{Đáp án của câu hỏi là A}
}
\end{ex}

\begin{ex} %[2D11V55]
Đoạn dây thứ nhất dài $45\text{ dm}$, đoạn dây thứ hai ngắn hơn đoạn thứ nhất $15\text{ dm}$. Cả hai đoạn dây dài bao nhiêu đề-xi-mét?
\choice
{$60\text{ dm}$}
{$70\text{ dm}$}
{\True $75\text{ dm}$}
{$80\text{ dm}$}
\loigiai{
Đoạn 2: $45 - 15 = 30$ dm. Cả hai đoạn: $45 + 30 = 75$ dm.
\textbf{Đáp án của câu hỏi là C}
}
\end{ex}

\textbf{Phần 2. Trắc nghiệm đúng sai}
\setcounterpageref{ex}{0}

\begin{ex} %[2D11V55]
Xét tính đúng sai của các phép đổi và quan hệ đơn vị đo độ dài:
\choiceTF
{\True 1 m = 100 cm.}
{\True 500 m + 500 m = 1 km.}
{1 km ngắn hơn 900 m.}
{\True Chiều cao của một học sinh lớp 2 khoảng 12 dm.}
\loigiai{
\begin{itemize}
    \item Ý a) Đúng: 1 m = 10 dm = 100 cm. Đúng.
    \item Ý b) Đúng: 1 000 m = 1 km. Đúng.
    \item Ý c) Sai: 1 km = 1 000 m > 900 m.
    \item Ý d) Đúng: 12 dm = 120 cm, phù hợp thực tế. Đúng.
\end{itemize}
\begin{itemize}
    \item a) \textbf{Đúng}.
    \item b) \textbf{Đúng}.
    \item c) \textbf{Sai}.
    \item d) \textbf{Đúng}.
\end{itemize}
}
\end{ex}

\begin{ex} %[2D11V55]
Cho ba đoạn thẳng: AB = 2 m, CD = 18 dm, EF = 250 cm:
\choiceTF
{\True AB dài hơn CD.}
{\True EF là đoạn thẳng dài nhất trong ba đoạn.}
{Tổng độ dài của AB và CD bé hơn 3 m.}
{\True Độ dài đoạn EF gấp 5 lần một đoạn dài 5 dm.}
\loigiai{
\begin{itemize}
    \item Ý a) Đúng: 2 m = 20 dm > 18 dm. Đúng.
    \item Ý b) Đúng: EF = 250 cm = 25 dm > 20 dm > 18 dm. Đúng.
    \item Ý c) Sai: 20 dm + 18 dm = 38 dm = 3 m 8 dm > 3 m.
    \item Ý d) Đúng: 25 dm : 5 = 5 dm. Đúng.
\end{itemize}
\begin{itemize}
    \item a) \textbf{Đúng}.
    \item b) \textbf{Đúng}.
    \item c) \textbf{Sai}.
    \item d) \textbf{Đúng}.
\end{itemize}
}
\end{ex}

\textbf{Phần 3. Trả lời ngắn}
\setcounterpageref{ex}{0}

\begin{ex} %[2D11V55]
Một chú ốc sên ban ngày bò lên được 8 dm, ban đêm tụt xuống 3 dm. Sau một ngày một đêm ốc sên bò lên được bao nhiêu đề-xi-mét?
\shortans{5}
\loigiai{
\textbf{Đáp số:} $5$\\
$8 - 3 = 5$ dm.
}
\end{ex}

\begin{ex} %[2D11V55]
Cuộn dây thép dài 1 km được chia thành các cuộn nhỏ dài 200 m. Được bao nhiêu cuộn nhỏ?
\shortans{5}
\loigiai{
\textbf{Đáp số:} $5$\\
$1\text{ km} = 1 000\text{ m}$. Số cuộn: $1 000 : 200 = 5$ cuộn.
}
\end{ex}

\begin{ex} %[2D11V55]
Thanh sắt dài 3 m được cưa thành các đoạn dài 5 dm. Cần cưa bao nhiêu lần để đứt hết?
\shortans{5}
\loigiai{
\textbf{Đáp số:} $5$\\
$3\text{ m} = 30\text{ dm}$. Số đoạn: $30 : 5 = 6$ đoạn. Cần $6 - 1 = 5$ lần cưa.
}
\end{ex}

\begin{ex} %[2D11V55]
Một hình tam giác có ba cạnh dài lần lượt là 20 cm, 30 cm, 50 cm. Chu vi của tam giác là bao nhiêu mét?
\shortans{1}
\loigiai{
\textbf{Đáp số:} $1$\\
$20 + 30 + 50 = 100\text{ cm} = 1$ m.
}
\end{ex}

\begin{ex} %[2D11V55]
Đoạn đường từ tỉnh A đến tỉnh B dài 45 km, từ tỉnh B đến tỉnh C dài 35 km. Đi từ A qua B đến C dài bao nhiêu ki-lô-mét?
\shortans{80}
\loigiai{
\textbf{Đáp số:} $80$\\
$45 + 35 = 80$ km.
}
\end{ex}

\begin{ex} %[2D11V55]
Một con mương dài 100 m. Người ta đã nạo vét được 60 m. Còn lại bao nhiêu mét mương chưa nạo vét?
\shortans{40}
\loigiai{
\textbf{Đáp số:} $40$\\
$100 - 60 = 40$ m.
}
\end{ex}

\newpage

\TaoPhieuKiemTra{Chủ đề 11 -- Độ dài và đơn vị đo độ dài. Tiền Việt Nam}{56}{Giới thiệu tiền Việt Nam}{https://loptoanthaylinh.vn/thi/chuyen-de/2/toan2_cd_b56}
\textbf{Phần 1. Trắc nghiệm lựa chọn}
\setcounterpageref{ex}{0}

\begin{ex} %[2D11V56]
Mệnh giá tiền nào dưới đây KHÔNG có trong chương trình làm quen của học sinh lớp 2?
\choice
{\True $50 000\text{ đồng}$}
{$100\text{ đồng}$}
{$500\text{ đồng}$}
{$1 000\text{ đồng}$}
\loigiai{
Lớp 2 chỉ làm quen với: 100 đồng, 200 đồng, 500 đồng và 1 000 đồng.
\textbf{Đáp án của câu hỏi là A}
}
\end{ex}

\begin{ex} %[2D11V56]
Cần bao nhiêu tờ tiền $500\text{ đồng}$ để đổi được $1$ tờ tiền $1 000\text{ đồng}$?
\choice
{$5$ tờ}
{$3$ tờ}
{\True $2$ tờ}
{$10$ tờ}
\loigiai{
$500 \times 2 = 1 000$ đồng, nên cần 2 tờ.
\textbf{Đáp án của câu hỏi là C}
}
\end{ex}

\begin{ex} %[2D11V56]
Nam có $3$ tờ tiền $200\text{ đồng}$ và $1$ tờ tiền $100\text{ đồng}$. Nam có tất cả bao nhiêu đồng?
\choice
{$600\text{ đồng}$}
{\True $700\text{ đồng}$}
{$800\text{ đồng}$}
{$500\text{ đồng}$}
\loigiai{
$200 \times 3 + 100 = 700$ đồng.
\textbf{Đáp án của câu hỏi là B}
}
\end{ex}

\begin{ex} %[2D11V56]
Một cây bút chì giá $800\text{ đồng}$. Hoa đưa cho cô bán hàng tờ tiền $1 000\text{ đồng}$. Cô bán hàng trả lại bao nhiêu tiền thừa?
\choice
{$100\text{ đồng}$}
{$300\text{ đồng}$}
{$50\text{ đồng}$}
{\True $200\text{ đồng}$}
\loigiai{
$1 000 - 800 = 200$ đồng.
\textbf{Đáp án của câu hỏi là D}
}
\end{ex}

\begin{ex} %[2D11V56]
Có bao nhiêu tờ $200\text{ đồng}$ để đổi được đúng $1 000\text{ đồng}$?
\choice
{\True $5$ tờ}
{$4$ tờ}
{$6$ tờ}
{$2$ tờ}
\loigiai{
$1 000 : 200 = 5$ tờ.
\textbf{Đáp án của câu hỏi là A}
}
\end{ex}

\begin{ex} %[2D11V56]
An có $1$ tờ $500\text{ đồng}$ và $2$ tờ $200\text{ đồng}$. An mua một cục tẩy hết $600\text{ đồng}$. An còn lại bao nhiêu đồng?
\choice
{$200\text{ đồng}$}
{$400\text{ đồng}$}
{\True $300\text{ đồng}$}
{$100\text{ đồng}$}
\loigiai{
An có: $500 + 200 \times 2 = 900$ đồng. Còn lại: $900 - 600 = 300$ đồng.
\textbf{Đáp án của câu hỏi là C}
}
\end{ex}

\begin{ex} %[2D11V56]
Một chiếc kẹo giá $200\text{ đồng}$. Với $1 000\text{ đồng}$, em có thể mua được nhiều nhất bao nhiêu chiếc kẹo như thế?
\choice
{$4$ cái}
{\True $5$ cái}
{$6$ cái}
{$3$ cái}
\loigiai{
$1 000 : 200 = 5$ cái kẹo.
\textbf{Đáp án của câu hỏi là B}
}
\end{ex}

\begin{ex} %[2D11V56]
Tờ tiền $1 000\text{ đồng}$ có mệnh giá gấp bao nhiêu lần tờ tiền $100\text{ đồng}$?
\choice
{$5$ lần}
{$20$ lần}
{$2$ lần}
{\True $10$ lần}
\loigiai{
$1 000 : 100 = 10$ lần.
\textbf{Đáp án của câu hỏi là D}
}
\end{ex}

\begin{ex} %[2D11V56]
Bình có $1$ tờ $500\text{ đồng}$, $1$ tờ $200\text{ đồng}$ và $2$ tờ $100\text{ đồng}$. Tổng số tiền của Bình là
\choice
{\True $900\text{ đồng}$}
{$800\text{ đồng}$}
{$1 000\text{ đồng}$}
{$700\text{ đồng}$}
\loigiai{
$500 + 200 + 100 \times 2 = 900$ đồng.
\textbf{Đáp án của câu hỏi là A}
}
\end{ex}

\begin{ex} %[2D11V56]
Mua một quyển vở hết $700\text{ đồng}$, mua một cái thước kẻ hết $200\text{ đồng}$. Mua cả hai thứ hết bao nhiêu đồng?
\choice
{$800\text{ đồng}$}
{$1 000\text{ đồng}$}
{\True $900\text{ đồng}$}
{$600\text{ đồng}$}
\loigiai{
$700 + 200 = 900$ đồng.
\textbf{Đáp án của câu hỏi là C}
}
\end{ex}

\begin{ex} %[2D11V56]
Để có đúng $800\text{ đồng}$ từ các tờ tiền $200\text{ đồng}$, cần ít nhất bao nhiêu tờ $200\text{ đồng}$?
\choice
{$3$ tờ}
{\True $4$ tờ}
{$5$ tờ}
{$6$ tờ}
\loigiai{
$800 : 200 = 4$ tờ.
\textbf{Đáp án của câu hỏi là B}
}
\end{ex}

\begin{ex} %[2D11V56]
Một gói bánh giá $1 000\text{ đồng}$. Minh có $2$ tờ $200\text{ đồng}$ và $1$ tờ $500\text{ đồng}$. Minh còn thiếu bao nhiêu đồng nữa mới đủ mua?
\choice
{$200\text{ đồng}$}
{$50\text{ đồng}$}
{$300\text{ đồng}$}
{\True $100\text{ đồng}$}
\loigiai{
Minh có $400 + 500 = 900$ đồng. Thiếu: $1 000 - 900 = 100$ đồng.
\textbf{Đáp án của câu hỏi là D}
}
\end{ex}

\textbf{Phần 2. Trắc nghiệm đúng sai}
\setcounterpageref{ex}{0}

\begin{ex} %[2D11V56]
Kiểm tra các nhận định về các mệnh giá tiền Việt Nam:
\choiceTF
{\True 1 tờ 1 000 đồng đổi được 2 tờ 500 đồng.}
{\True 5 tờ 200 đồng có tổng giá trị bằng 1 000 đồng.}
{Tờ tiền 100 đồng có giá trị lớn hơn tờ tiền 500 đồng.}
{\True Có thể dùng 1 tờ 500 đồng và 2 tờ 200 đồng để được 900 đồng.}
\loigiai{
\begin{itemize}
    \item Ý a) Đúng: 500 + 500 = 1 000 đồng. Đúng.
    \item Ý b) Đúng: 200 x 5 = 1 000 đồng. Đúng.
    \item Ý c) Sai: 100 đồng bé hơn 500 đồng.
    \item Ý d) Đúng: 500 + 200 x 2 = 900 đồng. Đúng.
\end{itemize}
\begin{itemize}
    \item a) \textbf{Đúng}.
    \item b) \textbf{Đúng}.
    \item c) \textbf{Sai}.
    \item d) \textbf{Đúng}.
\end{itemize}
}
\end{ex}

\begin{ex} %[2D11V56]
Bác Ba đi chợ với 1 tờ 1 000 đồng, mua mớ rau giá 400 đồng:
\choiceTF
{\True Bác Ba còn lại 600 đồng tiền thừa.}
{\True Số tiền còn lại đổi được đúng 3 tờ 200 đồng.}
{Số tiền còn lại đổi được 2 tờ 500 đồng.}
{\True Nếu mua thêm 1 quả ớt giá 100 đồng thì bác còn lại 500 đồng.}
\loigiai{
\begin{itemize}
    \item Ý a) Đúng: 1 000 - 400 = 600 đồng. Đúng.
    \item Ý b) Đúng: 600 : 200 = 3 tờ. Đúng.
    \item Ý c) Sai: 2 tờ 500 đồng là 1 000 đồng > 600 đồng.
    \item Ý d) Đúng: 600 - 100 = 500 đồng. Đúng.
\end{itemize}
\begin{itemize}
    \item a) \textbf{Đúng}.
    \item b) \textbf{Đúng}.
    \item c) \textbf{Sai}.
    \item d) \textbf{Đúng}.
\end{itemize}
}
\end{ex}

\textbf{Phần 3. Trả lời ngắn}
\setcounterpageref{ex}{0}

\begin{ex} %[2D11V56]
Mai có 1 tờ 1 000 đồng. Mai mua 2 phong bì thư, mỗi phong bì giá 300 đồng. Mai còn lại bao nhiêu đồng tiền thừa?
\shortans{400}
\loigiai{
\textbf{Đáp số:} $400$\\
Mua hết: $300 \times 2 = 600$ đồng. Còn lại: $1 000 - 600 = 400$ đồng.
}
\end{ex}

\begin{ex} %[2D11V56]
Có 1 tờ 500 đồng và 2 tờ 200 đồng. Cần thêm bao nhiêu tờ 100 đồng nữa để đủ 1 000 đồng?
\shortans{1}
\loigiai{
\textbf{Đáp số:} $1$\\
Đã có 900 đồng. Thiếu: $(1 000 - 900) : 100 = 1$ tờ 100 đồng.
}
\end{ex}

\begin{ex} %[2D11V56]
Bác Ba có 10 tờ tiền 100 đồng. Bác muốn đổi sang loại tiền 500 đồng. Bác đổi được bao nhiêu tờ 500 đồng?
\shortans{2}
\loigiai{
\textbf{Đáp số:} $2$\\
$100 \times 10 = 1 000$ đồng. $1 000 : 500 = 2$ tờ.
}
\end{ex}

\begin{ex} %[2D11V56]
Một chiếc bánh rán giá 500 đồng. Hỏi 2 chiếc bánh rán giá bao nhiêu đồng?
\shortans{1 000}
\loigiai{
\textbf{Đáp số:} $1 000$\\
$500 \times 2 = 1 000$ đồng.
}
\end{ex}

\begin{ex} %[2D11V56]
Dũng có 4 tờ 200 đồng. Dũng có tất cả bao nhiêu đồng?
\shortans{800}
\loigiai{
\textbf{Đáp số:} $800$\\
$200 \times 4 = 800$ đồng.
}
\end{ex}

\begin{ex} %[2D11V56]
Huy có 1 000 đồng, Huy mua kem hết 700 đồng. Cô bán hàng thối lại tiền toàn bằng tờ 100 đồng. Hỏi cô đưa cho Huy mấy tờ 100 đồng?
\shortans{3}
\loigiai{
\textbf{Đáp số:} $3$\\
Thừa 300 đồng. $300 : 100 = 3$ tờ.
}
\end{ex}

\newpage

\TaoPhieuKiemTra{Chủ đề 11 -- Độ dài và đơn vị đo độ dài. Tiền Việt Nam}{57}{Thực hành và trải nghiệm đo độ dài}{https://loptoanthaylinh.vn/thi/chuyen-de/2/toan2_cd_b57}
\textbf{Phần 1. Trắc nghiệm lựa chọn}
\setcounterpageref{ex}{0}

\begin{ex} %[2D11V57]
Để đo chiều dài của một cuốn sách giáo khoa Toán 2, em nên dùng thước có vạch chia theo đơn vị nào là thích hợp nhất?
\choice
{m}
{km}
{dm}
{\True cm}
\loigiai{
Đo chiều dài sách dùng xăng-ti-mét (cm) là thích hợp nhất.
\textbf{Đáp án của câu hỏi là D}
}
\end{ex}

\begin{ex} %[2D11V57]
Một gang tay của học sinh lớp 2 dài khoảng bao nhiêu?
\choice
{15 m}
{\True 15 cm}
{15 dm}
{15 km}
\loigiai{
Một gang tay học sinh dài khoảng 15 cm.
\textbf{Đáp án của câu hỏi là B}
}
\end{ex}

\begin{ex} %[2D11V57]
Chiều cao của cánh cửa lớp học thường khoảng
\choice
{2 cm}
{20 m}
{\True 20 dm}
{20 cm}
\loigiai{
Cánh cửa lớp học cao khoảng 2 m = 20 dm.
\textbf{Đáp án của câu hỏi là C}
}
\end{ex}

\begin{ex} %[2D11V57]
Một sợi dây khi gập đôi lại thì đo được $4\text{ dm}$. Chiều dài sợi dây đó khi duỗi thẳng là
\choice
{\True 8 dm}
{2 dm}
{6 dm}
{10 dm}
\loigiai{
$4 \times 2 = 8$ dm.
\textbf{Đáp án của câu hỏi là A}
}
\end{ex}

\begin{ex} %[2D11V57]
Dùng thước kẻ dài $20\text{ cm}$ để đo mép bàn. Đặt thước liên tiếp $5$ lần thì vừa hết chiều dài mép bàn. Chiều dài mép bàn đó là
\choice
{10 cm}
{50 cm}
{2 m}
{\True 1 m}
\loigiai{
$20 \times 5 = 100\text{ cm} = 1\text{ m}$.
\textbf{Đáp án của câu hỏi là D}
}
\end{ex}

\begin{ex} %[2D11V57]
Khoảng cách giữa hai cột điện liên tiếp trên phố thường được đo bằng đơn vị nào?
\choice
{cm}
{\True m}
{mm}
{dm}
\loigiai{
Khoảng cách giữa 2 cột điện đo bằng mét (m).
\textbf{Đáp án của câu hỏi là B}
}
\end{ex}

\begin{ex} %[2D11V57]
Để đo chu vi của thân một cây đại thụ trong sân trường, cách làm nào sau đây thuận tiện nhất?
\choice
{Dùng thước kẻ thẳng nhựa ngắn 15 cm}
{Dùng cân đĩa}
{\True Dùng một sợi dây mềm quấn một vòng quanh thân cây rồi đo chiều dài đoạn dây đó}
{Dùng đồng hồ bấm giây}
\loigiai{
Dùng sợi dây mềm quấn quanh thân cây rồi đo đoạn dây là thuận tiện và chính xác nhất.
\textbf{Đáp án của câu hỏi là C}
}
\end{ex}

\begin{ex} %[2D11V57]
Một mảnh vườn hình vuông có cạnh dài $25\text{ m}$. Bác thợ rào xung quanh mảnh vườn bằng lưới thép. Chiều dài hàng rào là
\choice
{\True 100 m}
{50 m}
{75 m}
{125 m}
\loigiai{
Chu vi hình vuông: $25 \times 4 = 100$ m.
\textbf{Đáp án của câu hỏi là A}
}
\end{ex}

\begin{ex} %[2D11V57]
Một bạn bước chân được khoảng $4\text{ dm}$ mỗi bước. Bạn đi từ cửa lớp đến bảng hết $10$ bước chân. Khoảng cách đó dài khoảng bao nhiêu mét?
\choice
{40 m}
{2 m}
{14 m}
{\True 4 m}
\loigiai{
$4 \times 10 = 40\text{ dm} = 4\text{ m}$.
\textbf{Đáp án của câu hỏi là D}
}
\end{ex}

\begin{ex} %[2D11V57]
Một cuộn băng dính dài $5\text{ m}$. Người ta cắt $5$ đoạn để dán trang trí, mỗi đoạn dài $4\text{ dm}$. Cuộn băng dính còn lại bao nhiêu đề-xi-mét?
\choice
{20 dm}
{\True 30 dm}
{10 dm}
{25 dm}
\loigiai{
Đổi $5\text{ m} = 50\text{ dm}$. Đã cắt: $4 \times 5 = 20$ dm. Còn lại: $50 - 20 = 30$ dm.
\textbf{Đáp án của câu hỏi là B}
}
\end{ex}

\begin{ex} %[2D11V57]
Tổ Một đo được đoạn đường dài $12\text{ m}$, tổ Hai đo được đoạn đường dài $15\text{ m}$. Cả hai tổ đo được tổng cộng bao nhiêu mét đường?
\choice
{25 m}
{26 m}
{\True 27 m}
{28 m}
\loigiai{
$12 + 15 = 27$ m.
\textbf{Đáp án của câu hỏi là C}
}
\end{ex}

\begin{ex} %[2D11V57]
Một thanh tre dài $2\text{ m } 4\text{ dm}$. Bác thợ muốn cắt thành $4$ đoạn bằng nhau. Mỗi đoạn dài bao nhiêu đề-xi-mét?
\choice
{\True 6 dm}
{5 dm}
{4 dm}
{7 dm}
\loigiai{
$2\text{ m } 4\text{ dm} = 24\text{ dm}$. Mỗi đoạn: $24 : 4 = 6$ dm.
\textbf{Đáp án của câu hỏi là A}
}
\end{ex}

\textbf{Phần 2. Trắc nghiệm đúng sai}
\setcounterpageref{ex}{0}

\begin{ex} %[2D11V57]
Xét tính hợp lý của các ước lượng và kết quả đo độ dài thực tế:
\choiceTF
{\True Một bước chân người lớn dài khoảng 70 cm.}
{Độ dày của một quyển vở ô ly khoảng 5 dm.}
{\True Khi gập đôi một sợi dây 3 lần, chiều dài mỗi phần bằng một phần tám sợi dây ban đầu.}
{\True 100 bước chân mỗi bước 5 dm dài tất cả 50 m.}
\loigiai{
\begin{itemize}
    \item Ý a) Đúng: Ước lượng phù hợp thực tế. Đúng.
    \item Ý b) Sai: 5 dm = 50 cm là quá dày, quyển vở chỉ dày vài milimét.
    \item Ý c) Đúng: Gập 3 lần: chia làm 2 x 2 x 2 = 8 phần bằng nhau. Đúng.
    \item Ý d) Đúng: 100 x 5 dm = 500 dm = 50 m. Đúng.
\end{itemize}
\begin{itemize}
    \item a) \textbf{Đúng}.
    \item b) \textbf{Sai}.
    \item c) \textbf{Đúng}.
    \item d) \textbf{Đúng}.
\end{itemize}
}
\end{ex}

\begin{ex} %[2D11V57]
Trong tiết trải nghiệm, Lan dùng thước cuộn đo kích thước phòng học:
\choiceTF
{\True Chiều dài phòng học đo được 8 m.}
{\True Chiều rộng phòng học đo được 6 m.}
{Chu vi sàn phòng học là 14 m.}
{\True Nửa chu vi sàn phòng học là 14 m.}
\loigiai{
\begin{itemize}
    \item Ý a) Đúng: Kích thước thực tế hợp lý. Đúng.
    \item Ý b) Đúng: Kích thước thực tế hợp lý. Đúng.
    \item Ý c) Sai: Chu vi: (8 + 6) x 2 = 28 m, không phải 14 m.
    \item Ý d) Đúng: 8 + 6 = 14 m. Đúng.
\end{itemize}
\begin{itemize}
    \item a) \textbf{Đúng}.
    \item b) \textbf{Đúng}.
    \item c) \textbf{Sai}.
    \item d) \textbf{Đúng}.
\end{itemize}
}
\end{ex}

\textbf{Phần 3. Trả lời ngắn}
\setcounterpageref{ex}{0}

\begin{ex} %[2D11V57]
Dùng thước thẳng 30 cm đo mép bàn. Đặt thước 4 lần liên tiếp thì còn thiếu 10 cm nữa mới hết mép bàn. Chiếc bàn đó dài bao nhiêu đề-xi-mét?
\shortans{13}
\loigiai{
\textbf{Đáp số:} $13$\\
$30 \times 4 + 10 = 130\text{ cm} = 13$ dm.
}
\end{ex}

\begin{ex} %[2D11V57]
Một đoạn dây gập làm 4 phần bằng nhau. Một phần dài 25 cm. Cả sợi dây dài bao nhiêu mét?
\shortans{1}
\loigiai{
\textbf{Đáp số:} $1$\\
$25 \times 4 = 100\text{ cm} = 1$ m.
}
\end{ex}

\begin{ex} %[2D11V57]
Bạn Hùng bước 50 bước, mỗi bước khoảng 4 dm. Sân trường dài khoảng bao nhiêu mét?
\shortans{20}
\loigiai{
\textbf{Đáp số:} $20$\\
$50 \times 4 = 200\text{ dm} = 20$ m.
}
\end{ex}

\begin{ex} %[2D11V57]
Bác An có sợi dây thừng dài 15 m. Bác cắt làm 3 đoạn bằng nhau. Mỗi đoạn dài bao nhiêu mét?
\shortans{5}
\loigiai{
\textbf{Đáp số:} $5$\\
$15 : 3 = 5$ m.
}
\end{ex}

\begin{ex} %[2D11V57]
Một hình tam giác có độ dài ba cạnh lần lượt là 12 cm, 18 cm, 20 cm. Chu vi của hình tam giác đó là bao nhiêu đề-xi-mét?
\shortans{5}
\loigiai{
\textbf{Đáp số:} $5$\\
$12 + 18 + 20 = 50\text{ cm} = 5$ dm.
}
\end{ex}

\begin{ex} %[2D11V57]
Một con kiến bò quanh mép một hộp bút hình chữ nhật dài 15 cm, rộng 5 cm. Kiến bò hết một vòng là bao nhiêu xăng-ti-mét?
\shortans{40}
\loigiai{
\textbf{Đáp số:} $40$\\
$(15 + 5) \times 2 = 40$ cm.
}
\end{ex}

\newpage

\TaoPhieuKiemTra{Chủ đề 11 -- Độ dài và đơn vị đo độ dài. Tiền Việt Nam}{58}{Luyện tập chung}{https://loptoanthaylinh.vn/thi/chuyen-de/2/toan2_cd_b58}
\textbf{Phần 1. Trắc nghiệm lựa chọn}
\setcounterpageref{ex}{0}

\begin{ex} %[2D11V58]
Giá trị của biểu thức $350\text{ m} + 150\text{ m}$ là
\choice
{400 m}
{600 m}
{\True 500 m}
{450 m}
\loigiai{
$350 + 150 = 500$ m.
\textbf{Đáp án của câu hỏi là C}
}
\end{ex}

\begin{ex} %[2D11V58]
Có $3$ tờ $200\text{ đồng}$ và $2$ tờ $100\text{ đồng}$. Tổng số tiền là
\choice
{700 đồng}
{\True 800 đồng}
{600 đồng}
{900 đồng}
\loigiai{
$200 \times 3 + 100 \times 2 = 600 + 200 = 800$ đồng.
\textbf{Đáp án của câu hỏi là B}
}
\end{ex}

\begin{ex} %[2D11V58]
Một đoạn đường dài $1\text{ km}$. Người ta đã trải nhựa được $700\text{ m}$. Còn bao nhiêu mét đường chưa được trải nhựa?
\choice
{400 m}
{200 m}
{500 m}
{\True 300 m}
\loigiai{
$1\text{ km} = 1 000\text{ m}$. Còn lại: $1 000 - 700 = 300$ m.
\textbf{Đáp án của câu hỏi là D}
}
\end{ex}

\begin{ex} %[2D11V58]
Điền số thích hợp: $5\text{ m } 6\text{ dm} = \dots\text{ dm}$.
\choice
{\True 56 dm}
{506 dm}
{11 dm}
{65 dm}
\loigiai{
$5\text{ m} = 50\text{ dm} \Rightarrow 5\text{ m } 6\text{ dm} = 56$ dm.
\textbf{Đáp án của câu hỏi là A}
}
\end{ex}

\begin{ex} %[2D11V58]
Một người có $1 000\text{ đồng}$, mua một chiếc kẹo giá $300\text{ đồng}$ và một gói bim bim giá $500\text{ đồng}$. Người đó còn lại bao nhiêu tiền?
\choice
{100 đồng}
{300 đồng}
{\True 200 đồng}
{400 đồng}
\loigiai{
Đã tiêu: $300 + 500 = 800$ đồng. Còn lại: $1 000 - 800 = 200$ đồng.
\textbf{Đáp án của câu hỏi là C}
}
\end{ex}

\begin{ex} %[2D11V58]
Sợi dây thứ nhất dài $12\text{ dm}$, sợi dây thứ hai dài gấp đôi sợi thứ nhất. Cả hai sợi dây dài bao nhiêu đề-xi-mét?
\choice
{24 dm}
{\True 36 dm}
{30 dm}
{48 dm}
\loigiai{
Sợi 2: $12 \times 2 = 24$ dm. Cả hai sợi: $12 + 24 = 36$ dm.
\textbf{Đáp án của câu hỏi là B}
}
\end{ex}

\begin{ex} %[2D11V58]
Một hình chữ nhật có chiều dài $30\text{ m}$, chiều rộng $20\text{ m}$. Chu vi hình chữ nhật đó là
\choice
{50 m}
{80 m}
{60 m}
{\True 100 m}
\loigiai{
$(30 + 20) \times 2 = 100$ m.
\textbf{Đáp án của câu hỏi là D}
}
\end{ex}

\begin{ex} %[2D11V58]
Nam mua một hộp sữa giá $700\text{ đồng}$. Nam đưa cho cô bán hàng $4$ tờ tiền $200\text{ đồng}$. Cô bán hàng trả lại Nam bao nhiêu tiền?
\choice
{\True 100 đồng}
{200 đồng}
{50 đồng}
{150 đồng}
\loigiai{
Nam đưa: $200 \times 4 = 800$ đồng. Trả lại: $800 - 700 = 100$ đồng.
\textbf{Đáp án của câu hỏi là A}
}
\end{ex}

\begin{ex} %[2D11V58]
Độ dài đường gấp khúc gồm $4$ đoạn thẳng bằng nhau, mỗi đoạn dài $25\text{ cm}$ là
\choice
{10 cm}
{50 cm}
{\True 1 m}
{2 m}
\loigiai{
$25 \times 4 = 100\text{ cm} = 1\text{ m}$.
\textbf{Đáp án của câu hỏi là C}
}
\end{ex}

\begin{ex} %[2D11V58]
Điền dấu thích hợp: $1\text{ km} - 400\text{ m} \dots 500\text{ m}$.
\choice
{$<$}
{\True $>$}
{$=$}
{Không so sánh được}
\loigiai{
$1\text{ km} - 400\text{ m} = 600\text{ m} > 500\text{ m}$.
\textbf{Đáp án của câu hỏi là B}
}
\end{ex}

\begin{ex} %[2D11V58]
Bác thợ rào một khu đất hình vuông cạnh $20\text{ m}$, có để một lối đi rộng $2\text{ m}$ không rào. Chiều dài hàng rào là
\choice
{80 m}
{76 m}
{82 m}
{\True 78 m}
\loigiai{
Chu vi khu đất: $20 \times 4 = 80$ m. Hàng rào: $80 - 2 = 78$ m.
\textbf{Đáp án của câu hỏi là D}
}
\end{ex}

\begin{ex} %[2D11V58]
Lan có $1 000\text{ đồng}$, mua $5$ quả mận giá $200\text{ đồng}$ một quả. Lan còn lại bao nhiêu đồng?
\choice
{\True 0 đồng}
{200 đồng}
{100 đồng}
{500 đồng}
\loigiai{
Lan mua hết: $200 \times 5 = 1 000$ đồng. Còn lại: $1 000 - 1 000 = 0$ đồng.
\textbf{Đáp án của câu hỏi là A}
}
\end{ex}

\textbf{Phần 2. Trắc nghiệm đúng sai}
\setcounterpageref{ex}{0}

\begin{ex} %[2D11V58]
Xét tính đúng sai của các phép toán đo lường và tiền tệ:
\choiceTF
{\True 1 km dài hơn 1 000 dm.}
{3 tờ 200 đồng có giá trị ít hơn 1 tờ 500 đồng.}
{\True Một hình vuông có chu vi 1 m thì cạnh của nó dài 25 cm.}
{\True 1 000 đồng đổi được đúng 10 tờ 100 đồng.}
\loigiai{
\begin{itemize}
    \item Ý a) Đúng: 1 km = 10 000 dm > 1 000 dm. Đúng.
    \item Ý b) Sai: 200 x 3 = 600 đồng > 500 đồng.
    \item Ý c) Đúng: 100 cm : 4 = 25 cm. Đúng.
    \item Ý d) Đúng: 100 x 10 = 1 000 đồng. Đúng.
\end{itemize}
\begin{itemize}
    \item a) \textbf{Đúng}.
    \item b) \textbf{Sai}.
    \item c) \textbf{Đúng}.
    \item d) \textbf{Đúng}.
\end{itemize}
}
\end{ex}

\begin{ex} %[2D11V58]
Một cửa hàng có cuộn vải dài 100 m. Ngày thứ nhất bán được 35 m, ngày thứ hai bán được 45 m:
\choiceTF
{\True Cả hai ngày bán được 80 m vải.}
{\True Sau hai ngày, cuộn vải còn lại 20 m.}
{Số mét vải ngày thứ hai bán gấp đôi ngày thứ nhất.}
{\True Nếu mỗi mét vải giá 10 đồng thì 20 m còn lại giá 200 đồng.}
\loigiai{
\begin{itemize}
    \item Ý a) Đúng: 35 + 45 = 80 m. Đúng.
    \item Ý b) Đúng: 100 - 80 = 20 m. Đúng.
    \item Ý c) Sai: 35 x 2 = 70 != 45 m.
    \item Ý d) Đúng: 20 x 10 = 200 đồng. Đúng.
\end{itemize}
\begin{itemize}
    \item a) \textbf{Đúng}.
    \item b) \textbf{Đúng}.
    \item c) \textbf{Sai}.
    \item d) \textbf{Đúng}.
\end{itemize}
}
\end{ex}

\textbf{Phần 3. Trả lời ngắn}
\setcounterpageref{ex}{0}

\begin{ex} %[2D11V58]
Đường từ nhà An đến trường dài 400 m, từ trường đến thư viện dài 600 m. An đi từ nhà đến trường rồi qua thư viện, sau đó về nhà theo đường cũ. An đã đi được bao nhiêu ki-lô-mét?
\shortans{2}
\loigiai{
\textbf{Đáp số:} $2$\\
Một lượt: $400 + 600 = 1 000\text{ m} = 1\text{ km}$. Cả đi lẫn về: 2 km.
}
\end{ex}

\begin{ex} %[2D11V58]
Bác Hùng có 1 000 đồng, mua xà phòng hết 600 đồng, mua khăn mặt hết 300 đồng. Bác Hùng còn lại bao nhiêu đồng tiền thừa?
\shortans{100}
\loigiai{
\textbf{Đáp số:} $100$\\
Đã tiêu: $600 + 300 = 900$ đồng. Còn lại: $1 000 - 900 = 100$ đồng.
}
\end{ex}

\begin{ex} %[2D11V58]
Một sợi dây đồng dài 2 m 4 dm uốn thành hình tam giác có 3 cạnh bằng nhau. Mỗi cạnh dài bao nhiêu đề-xi-mét?
\shortans{8}
\loigiai{
\textbf{Đáp số:} $8$\\
$2\text{ m } 4\text{ dm} = 24\text{ dm}$. Mỗi cạnh: $24 : 3 = 8$ dm.
}
\end{ex}

\begin{ex} %[2D11V58]
Một cây cột cao 6 m. Người ta cắm xuống đất 1 m 5 dm. Phần cột nhô trên mặt đất dài bao nhiêu đề-xi-mét?
\shortans{45}
\loigiai{
\textbf{Đáp số:} $45$\\
6 m = 60 dm. Phần ngập đất = 15 dm. Phần nhô lên: $60 - 15 = 45$ dm.
}
\end{ex}

\begin{ex} %[2D11V58]
Mẹ mua 2 kg táo giá 500 đồng một ki-lô-gam. Mẹ đưa tờ 1 000 đồng. Người bán hàng phải trả lại mẹ bao nhiêu đồng?
\shortans{0}
\loigiai{
\textbf{Đáp số:} $0$\\
2 x 500 = 1 000 đồng. Tiền thối lại là 0 đồng.
}
\end{ex}

\begin{ex} %[2D11V58]
Một hình chữ nhật có chu vi 40 cm, chiều dài là 12 cm. Chiều rộng của hình chữ nhật đó là bao nhiêu xăng-ti-mét?
\shortans{8}
\loigiai{
\textbf{Đáp số:} $8$\\
Nửa chu vi: $40 : 2 = 20$ cm. Chiều rộng: $20 - 12 = 8$ cm.
}
\end{ex}

\newpage


% =========================================================================
% CHỦ ĐỀ 12 - PHÉP CỘNG, PHÉP TRỪ TRONG PHẠM VI 1 000
% =========================================================================
\section{Phép cộng, phép trừ trong phạm vi 1 000}

\TaoPhieuKiemTra{Chủ đề 12 -- Phép cộng, phép trừ trong phạm vi 1 000}{59}{Phép cộng (không nhớ) trong phạm vi 1 000}{https://loptoanthaylinh.vn/thi/chuyen-de/2/toan2_cd_b59}
\textbf{Phần 1. Trắc nghiệm lựa chọn}
\setcounterpageref{ex}{0}

\begin{ex} %[2D12V59]
Kết quả của phép tính $324 + 145$ là
\choice
{$479$}
{$459$}
{\True $469$}
{$369$}
\loigiai{
$324 + 145 = 469$.
\textbf{Đáp án của câu hỏi là C}
}
\end{ex}

\begin{ex} %[2D12V59]
Tính tổng của số chẵn lớn nhất có hai chữ số và số $300$.
\choice
{\True $398$}
{$399$}
{$390$}
{$400$}
\loigiai{
Số chẵn lớn nhất có 2 chữ số là 98. Ta có: $98 + 300 = 398$.
\textbf{Đáp án của câu hỏi là A}
}
\end{ex}

\begin{ex} %[2D12V59]
Điền chữ số thích hợp vào dấu $*$ để: $4*2 + 135 = 577$.
\choice
{$3$}
{$5$}
{$2$}
{\True $4$}
\loigiai{
Chữ số hàng chục: $* + 3 = 7 \Rightarrow * = 4$.
\textbf{Đáp án của câu hỏi là D}
}
\end{ex}

\begin{ex} %[2D12V59]
Một nông trại nuôi $235$ con bò sữa và $142$ con bò thịt. Nông trại có tất cả bao nhiêu con bò?
\choice
{$367$ con}
{\True $377$ con}
{$387$ con}
{$375$ con}
\loigiai{
$235 + 142 = 377$ (con bò).
\textbf{Đáp án của câu hỏi là B}
}
\end{ex}

\begin{ex} %[2D12V59]
Tìm $x$ biết: $x - 213 = 451$.
\choice
{$662$}
{$654$}
{\True $664$}
{$238$}
\loigiai{
$x = 451 + 213 = 664$.
\textbf{Đáp án của câu hỏi là C}
}
\end{ex}

\begin{ex} %[2D12V59]
Một thư viện trường học có $510$ cuốn sách tham khảo và $385$ cuốn truyện tranh. Thư viện có tất cả bao nhiêu cuốn sách và truyện tranh?
\choice
{\True $895$ cuốn}
{$885$ cuốn}
{$890$ cuốn}
{$795$ cuốn}
\loigiai{
$510 + 385 = 895$ (cuốn).
\textbf{Đáp án của câu hỏi là A}
}
\end{ex}

\begin{ex} %[2D12V59]
Tổng của hai số là $658$. Nếu giữ nguyên số hạng thứ nhất và tăng số hạng thứ hai thêm $31$ đơn vị thì tổng mới là bao nhiêu?
\choice
{$679$}
{$699$}
{$627$}
{\True $689$}
\loigiai{
Tổng mới: $658 + 31 = 689$.
\textbf{Đáp án của câu hỏi là D}
}
\end{ex}

\begin{ex} %[2D12V59]
Cho phép tính: $120 + 230 + 340$. Giá trị của biểu thức đó là
\choice
{$680$}
{\True $690$}
{$700$}
{$590$}
\loigiai{
$120 + 230 + 340 = 690$.
\textbf{Đáp án của câu hỏi là B}
}
\end{ex}

\begin{ex} %[2D12V59]
Một vườn ươm có $430$ cây cam và $350$ cây bưởi. Vườn ươm đó có tất cả bao nhiêu cây giống?
\choice
{$770$ cây}
{$880$ cây}
{\True $780$ cây}
{$680$ cây}
\loigiai{
$430 + 350 = 780$ (cây).
\textbf{Đáp án của câu hỏi là C}
}
\end{ex}

\begin{ex} %[2D12V59]
Điền dấu thích hợp: $312 + 245 \dots 212 + 355$.
\choice
{\True $<$}
{$>$}
{$=$}
{Không so sánh được}
\loigiai{
$312 + 245 = 557$; $212 + 355 = 567$. Vì $557 < 567$ nên điền dấu $<$. 
\textbf{Đáp án của câu hỏi là A}
}
\end{ex}

\begin{ex} %[2D12V59]
Tìm tổng của số bé nhất có ba chữ số khác nhau và số $543$.
\choice
{$643$}
{$646$}
{$655$}
{\True $645$}
\loigiai{
Số bé nhất có ba chữ số khác nhau là 102. Tổng: $102 + 543 = 645$.
\textbf{Đáp án của câu hỏi là D}
}
\end{ex}

\begin{ex} %[2D12V59]
Xe thứ nhất chở $320\text{ kg}$ gạo, xe thứ hai chở nhiều hơn xe thứ nhất $45\text{ kg}$ gạo. Cả hai xe chở được bao nhiêu ki-lô-gam gạo?
\choice
{$675\text{ kg}$}
{\True $685\text{ kg}$}
{$365\text{ kg}$}
{$695\text{ kg}$}
\loigiai{
Xe thứ hai chở: $320 + 45 = 365$ kg. Cả hai xe chở: $320 + 365 = 685$ kg.
\textbf{Đáp án của câu hỏi là B}
}
\end{ex}

\textbf{Phần 2. Trắc nghiệm đúng sai}
\setcounterpageref{ex}{0}

\begin{ex} %[2D12V59]
Xét các phép tính cộng không nhớ trong phạm vi 1 000:
\choiceTF
{\True 321 + 456 = 777.}
{500 + 40 + 9 = 540.}
{\True Khi cộng hai số có ba chữ số không nhớ, ta thực hiện cộng từ phải sang trái.}
{\True Nếu mỗi số hạng đều có chữ số hàng đơn vị là 0 thì tổng có chữ số hàng đơn vị là 0.}
\loigiai{
\begin{itemize}
    \item Ý a) Đúng: 1 + 6 = 7, 2 + 5 = 7, 3 + 4 = 7. Đúng.
    \item Ý b) Sai: 500 + 40 + 9 = 549, không phải 540.
    \item Ý c) Đúng: Quy tắc đặt tính và tính đúng.
    \item Ý d) Đúng: 0 + 0 = 0. Đúng.
\end{itemize}
\begin{itemize}
    \item a) \textbf{Đúng}.
    \item b) \textbf{Sai}.
    \item c) \textbf{Đúng}.
    \item d) \textbf{Đúng}.
\end{itemize}
}
\end{ex}

\begin{ex} %[2D12V59]
Khối 2 có hai lớp: Lớp 2A có 142 học sinh giỏi toán, lớp 2B có 154 học sinh giỏi toán:
\choiceTF
{\True Cả hai lớp có 296 học sinh giỏi toán.}
{\True Lớp 2B nhiều hơn lớp 2A 12 học sinh giỏi toán.}
{Nếu lớp 2A có thêm 15 bạn giỏi toán thì lớp 2A đông hơn lớp 2B.}
{Tổng số học sinh giỏi của hai lớp vượt quá 300 bạn.}
\loigiai{
\begin{itemize}
    \item Ý a) Đúng: 142 + 154 = 296 học sinh. Đúng.
    \item Ý b) Đúng: 154 - 142 = 12 học sinh. Đúng.
    \item Ý c) Sai: 142 + 15 = 157 > 154 (thực tế đúng, nhưng nếu chỉ thêm 10 bạn thì 152 < 154). Ở đây: 142 + 15 = 157 > 154, nên câu này là Đúng.
    \item Ý d) Sai: 296 < 300, không vượt quá. Sai.
\end{itemize}
\begin{itemize}
    \item a) \textbf{Đúng}.
    \item b) \textbf{Đúng}.
    \item c) \textbf{Sai}.
    \item d) \textbf{Sai}.
\end{itemize}
}
\end{ex}

\textbf{Phần 3. Trả lời ngắn}
\setcounterpageref{ex}{0}

\begin{ex} %[2D12V59]
Tính tổng của $415$ và $372$.
\shortans{787}
\loigiai{
\textbf{Đáp số:} $787$\\
$415 + 372 = 787$.
}
\end{ex}

\begin{ex} %[2D12V59]
Kho thứ nhất có 350 tấn thóc, kho thứ hai có 420 tấn thóc. Cả hai kho có bao nhiêu tấn thóc?
\shortans{770}
\loigiai{
\textbf{Đáp số:} $770$\\
$350 + 420 = 770$ tấn thóc.
}
\end{ex}

\begin{ex} %[2D12V59]
Tìm $x$ biết $x - 302 = 184$.
\shortans{486}
\loigiai{
\textbf{Đáp số:} $486$\\
$x = 184 + 302 = 486$.
}
\end{ex}

\begin{ex} %[2D12V59]
Một trường tiểu học có 423 học sinh nam và 465 học sinh nữ. Trường có tất cả bao nhiêu học sinh?
\shortans{888}
\loigiai{
\textbf{Đáp số:} $888$\\
$423 + 465 = 888$ học sinh.
}
\end{ex}

\begin{ex} %[2D12V59]
Tính giá trị biểu thức: $214 + 132 + 321$.
\shortans{667}
\loigiai{
\textbf{Đáp số:} $667$\\
$214 + 132 + 321 = 667$.
}
\end{ex}

\begin{ex} %[2D12V59]
Số hạng thứ nhất là 531, số hạng thứ hai kém số hạng thứ nhất 210 đơn vị. Tìm tổng hai số.
\shortans{852}
\loigiai{
\textbf{Đáp số:} $852$\\
Số hạng 2: $531 - 210 = 321$. Tổng: $531 + 321 = 852$.
}
\end{ex}

\newpage

\TaoPhieuKiemTra{Chủ đề 12 -- Phép cộng, phép trừ trong phạm vi 1 000}{60}{Phép trừ (không nhớ) trong phạm vi 1 000}{https://loptoanthaylinh.vn/thi/chuyen-de/2/toan2_cd_b60}
\textbf{Phần 1. Trắc nghiệm lựa chọn}
\setcounterpageref{ex}{0}

\begin{ex} %[2D12V60]
Kết quả của phép tính $768 - 342$ là
\choice
{$416$}
{\True $426$}
{$436$}
{$326$}
\loigiai{
$768 - 342 = 426$.
\textbf{Đáp án của câu hỏi là B}
}
\end{ex}

\begin{ex} %[2D12V60]
Tìm hiệu của số tròn trăm lớn nhất có ba chữ số và số $250$.
\choice
{$750$}
{$600$}
{$700$}
{\True $650$}
\loigiai{
Số tròn trăm lớn nhất có 3 chữ số là 900. $900 - 250 = 650$.
\textbf{Đáp án của câu hỏi là D}
}
\end{ex}

\begin{ex} %[2D12V60]
Điền chữ số thích hợp vào ô trống: $685 - 2*3 = 442$.
\choice
{\True $4$}
{$3$}
{$5$}
{$2$}
\loigiai{
Hàng chục: $8 - * = 4 \Rightarrow * = 4$.
\textbf{Đáp án của câu hỏi là A}
}
\end{ex}

\begin{ex} %[2D12V60]
Một cửa hàng có $576\text{ kg}$ đường, đã bán được $234\text{ kg}$. Cửa hàng còn lại bao nhiêu ki-lô-gam đường?
\choice
{$332\text{ kg}$}
{$352\text{ kg}$}
{\True $342\text{ kg}$}
{$242\text{ kg}$}
\loigiai{
$576 - 234 = 342$ (kg).
\textbf{Đáp án của câu hỏi là C}
}
\end{ex}

\begin{ex} %[2D12V60]
Tìm $x$ biết: $897 - x = 564$.
\choice
{$323$}
{\True $333$}
{$343$}
{$433$}
\loigiai{
$x = 897 - 564 = 333$.
\textbf{Đáp án của câu hỏi là B}
}
\end{ex}

\begin{ex} %[2D12V60]
Trong một buổi hòa nhạc có $980$ vé được bán ra, trong đó có $450$ vé người lớn. Có bao nhiêu vé trẻ em?
\choice
{$520$ vé}
{$540$ vé}
{$430$ vé}
{\True $530$ vé}
\loigiai{
$980 - 450 = 530$ (vé).
\textbf{Đáp án của câu hỏi là D}
}
\end{ex}

\begin{ex} %[2D12V60]
Hiệu giữa số lớn nhất có ba chữ số và số $432$ là
\choice
{\True $567$}
{$557$}
{$577$}
{$467$}
\loigiai{
$999 - 432 = 567$.
\textbf{Đáp án của câu hỏi là A}
}
\end{ex}

\begin{ex} %[2D12V60]
Một xưởng dệt may được $865$ chiếc áo len. Xưởng đã xuất đi $521$ chiếc. Xưởng còn lại bao nhiêu chiếc áo len?
\choice
{$334$ chiếc}
{$354$ chiếc}
{\True $344$ chiếc}
{$342$ chiếc}
\loigiai{
$865 - 521 = 344$ (chiếc).
\textbf{Đáp án của câu hỏi là C}
}
\end{ex}

\begin{ex} %[2D12V60]
Điền dấu thích hợp: $756 - 324 \dots 689 - 257$.
\choice
{$>$}
{\True $=$}
{$<$}
{Không so sánh được}
\loigiai{
$756 - 324 = 432$; $689 - 257 = 432$. Vì vậy hai vế bằng nhau.
\textbf{Đáp án của câu hỏi là B}
}
\end{ex}

\begin{ex} %[2D12V60]
Một người mang $680\text{ kg}$ hoa quả đi bán. Buổi sáng bán được $310\text{ kg}$, buổi chiều bán được $240\text{ kg}$. Người đó còn lại bao nhiêu ki-lô-gam hoa quả?
\choice
{$120\text{ kg}$}
{$140\text{ kg}$}
{$230\text{ kg}$}
{\True $130\text{ kg}$}
\loigiai{
Đã bán: $310 + 240 = 550$ kg. Còn lại: $680 - 550 = 130$ kg.
\textbf{Đáp án của câu hỏi là D}
}
\end{ex}

\begin{ex} %[2D12V60]
Một cuốn sách dày $458$ trang. Nam đã đọc được $235$ trang. Hỏi Nam còn phải đọc bao nhiêu trang nữa thì hết cuốn sách?
\choice
{\True $223$ trang}
{$213$ trang}
{$233$ trang}
{$225$ trang}
\loigiai{
$458 - 235 = 223$ (trang).
\textbf{Đáp án của câu hỏi là A}
}
\end{ex}

\begin{ex} %[2D12V60]
Hiệu của hai số là $342$. Nếu số bị trừ giảm đi $20$ đơn vị và giữ nguyên số trừ thì hiệu mới là
\choice
{$362$}
{$332$}
{\True $322$}
{$312$}
\loigiai{
Số bị trừ giảm bao nhiêu thì hiệu giảm bấy nhiêu: $342 - 20 = 322$.
\textbf{Đáp án của câu hỏi là C}
}
\end{ex}

\textbf{Phần 2. Trắc nghiệm đúng sai}
\setcounterpageref{ex}{0}

\begin{ex} %[2D12V60]
Xét tính đúng sai của các phép toán trừ không nhớ:
\choiceTF
{\True 685 - 231 = 454.}
{978 - 435 = 542.}
{\True Trong phép trừ không nhớ, mỗi chữ số của số bị trừ đều lớn hơn hoặc bằng chữ số tương ứng của số trừ.}
{\True Nếu lấy hiệu cộng với số trừ thì được số bị trừ.}
\loigiai{
\begin{itemize}
    \item Ý a) Đúng: 5 - 1 = 4, 8 - 3 = 5, 6 - 2 = 4. Đúng.
    \item Ý b) Sai: 8 - 5 = 3, kết quả phải là 543. Sai.
    \item Ý c) Đúng: Định nghĩa phép trừ không nhớ. Đúng.
    \item Ý d) Đúng: Quy tắc tìm số bị trừ. Đúng.
\end{itemize}
\begin{itemize}
    \item a) \textbf{Đúng}.
    \item b) \textbf{Sai}.
    \item c) \textbf{Đúng}.
    \item d) \textbf{Đúng}.
\end{itemize}
}
\end{ex}

\begin{ex} %[2D12V60]
Kho A có 789 bao gạo, kho B có 456 bao gạo:
\choiceTF
{\True Kho A nhiều hơn kho B 333 bao gạo.}
{Nếu kho A xuất đi 333 bao thì số gạo hai kho bằng nhau.}
{Nếu cả hai kho cùng xuất đi 50 bao thì hiệu số bao gạo giữa hai kho thay đổi.}
{\True Tổng số bao gạo của hai kho là 1 245 bao.}
\loigiai{
\begin{itemize}
    \item Ý a) Đúng: 789 - 456 = 333 bao. Đúng.
    \item Ý b) Sai: 789 - 333 = 456, đúng bằng kho B. (Mệnh đề này là Đúng).
    \item Ý c) Sai: Cùng bớt đi số lượng như nhau thì hiệu không đổi. Sai.
    \item Ý d) Đúng: 789 + 456 = 1 245 bao. Đúng.
\end{itemize}
\begin{itemize}
    \item a) \textbf{Đúng}.
    \item b) \textbf{Sai}.
    \item c) \textbf{Sai}.
    \item d) \textbf{Đúng}.
\end{itemize}
}
\end{ex}

\textbf{Phần 3. Trả lời ngắn}
\setcounterpageref{ex}{0}

\begin{ex} %[2D12V60]
Tính hiệu: $987 - 654$.
\shortans{333}
\loigiai{
\textbf{Đáp số:} $333$\\
$987 - 654 = 333$.
}
\end{ex}

\begin{ex} %[2D12V60]
Tìm $x$ biết $765 - x = 231$.
\shortans{534}
\loigiai{
\textbf{Đáp số:} $534$\\
$x = 765 - 231 = 534$.
}
\end{ex}

\begin{ex} %[2D12V60]
Một bể bơi chứa 890 lít nước. Sau khi tháo bớt đi thì trong bể còn lại 340 lít nước. Đã tháo bao nhiêu lít nước?
\shortans{550}
\loigiai{
\textbf{Đáp số:} $550$\\
$890 - 340 = 550$ lít.
}
\end{ex}

\begin{ex} %[2D12V60]
Số bị trừ là 846, số trừ là 324. Hiệu là bao nhiêu?
\shortans{522}
\loigiai{
\textbf{Đáp số:} $522$\\
$846 - 324 = 522$.
}
\end{ex}

\begin{ex} %[2D12V60]
Một xưởng dệt được 950 mét vải, đã bán 620 mét vải. Còn lại bao nhiêu mét vải?
\shortans{330}
\loigiai{
\textbf{Đáp số:} $330$\\
$950 - 620 = 330$ m.
}
\end{ex}

\begin{ex} %[2D12V60]
Tìm một số biết rằng lấy số đó cộng với 312 thì được 785.
\shortans{473}
\loigiai{
\textbf{Đáp số:} $473$\\
Số đó là: $785 - 312 = 473$.
}
\end{ex}

\newpage

\TaoPhieuKiemTra{Chủ đề 12 -- Phép cộng, phép trừ trong phạm vi 1 000}{61}{Phép cộng (có nhớ) trong phạm vi 1 000}{https://loptoanthaylinh.vn/thi/chuyen-de/2/toan2_cd_b61}
\textbf{Phần 1. Trắc nghiệm lựa chọn}
\setcounterpageref{ex}{0}

\begin{ex} %[2D12V61]
Kết quả của phép tính $256 + 178$ là
\choice
{\True $434$}
{$424$}
{$444$}
{$334$}
\loigiai{
$256 + 178 = 434$.
\textbf{Đáp án của câu hỏi là A}
}
\end{ex}

\begin{ex} %[2D12V61]
Tính tổng của $467$ và $285$.
\choice
{$742$}
{$762$}
{\True $752$}
{$652$}
\loigiai{
$467 + 285 = 752$.
\textbf{Đáp án của câu hỏi là C}
}
\end{ex}

\begin{ex} %[2D12V61]
Điền chữ số thích hợp vào dấu $*$ để: $3*8 + 245 = 613$.
\choice
{$5$}
{\True $6$}
{$7$}
{$4$}
\loigiai{
$613 - 245 = 368$, nên $* = 6$.
\textbf{Đáp án của câu hỏi là B}
}
\end{ex}

\begin{ex} %[2D12V61]
Một trại chăn nuôi xuất bán $375$ con gà trong tuần thứ nhất và $486$ con gà trong tuần thứ hai. Cả hai tuần trại xuất bán bao nhiêu con gà?
\choice
{$851$ con}
{$871$ con}
{$761$ con}
{\True $861$ con}
\loigiai{
$375 + 486 = 861$ (con gà).
\textbf{Đáp án của câu hỏi là D}
}
\end{ex}

\begin{ex} %[2D12V61]
Tìm $x$ biết: $x - 368 = 257$.
\choice
{\True $625$}
{$615$}
{$635$}
{$525$}
\loigiai{
$x = 257 + 368 = 625$.
\textbf{Đáp án của câu hỏi là A}
}
\end{ex}

\begin{ex} %[2D12V61]
Đoạn đường từ tỉnh A đến tỉnh B dài $285\text{ km}$, từ tỉnh B đến tỉnh C dài $348\text{ km}$. Quãng đường từ A qua B đến C dài bao nhiêu ki-lô-mét?
\choice
{$623\text{ km}$}
{$643\text{ km}$}
{\True $633\text{ km}$}
{$533\text{ km}$}
\loigiai{
$285 + 348 = 633$ (km).
\textbf{Đáp án của câu hỏi là C}
}
\end{ex}

\begin{ex} %[2D12V61]
Tính giá trị của biểu thức: $189 + 245 + 316$.
\choice
{$740$}
{\True $750$}
{$760$}
{$650$}
\loigiai{
$189 + 245 + 316 = 750$.
\textbf{Đáp án của câu hỏi là B}
}
\end{ex}

\begin{ex} %[2D12V61]
Số hạng thứ nhất là $458$, số hạng thứ hai lớn hơn số hạng thứ nhất $75$ đơn vị. Tổng của hai số đó là
\choice
{$981$}
{$533$}
{$1 001$}
{\True $991$}
\loigiai{
Số hạng thứ hai: $458 + 75 = 533$. Tổng: $458 + 533 = 991$.
\textbf{Đáp án của câu hỏi là D}
}
\end{ex}

\begin{ex} %[2D12V61]
Điền dấu thích hợp: $367 + 258 \dots 195 + 430$.
\choice
{\True $=$}
{$>$}
{$<$}
{Không so sánh được}
\loigiai{
$367 + 258 = 625$; $195 + 430 = 625$. Vậy hai biểu thức bằng nhau.
\textbf{Đáp án của câu hỏi là A}
}
\end{ex}

\begin{ex} %[2D12V61]
Một cửa hàng buổi sáng bán được $275\text{ kg}$ ngô, buổi chiều bán được $368\text{ kg}$ ngô. Cả ngày cửa hàng bán được bao nhiêu ki-lô-gam ngô?
\choice
{$633\text{ kg}$}
{$653\text{ kg}$}
{\True $643\text{ kg}$}
{$543\text{ kg}$}
\loigiai{
$275 + 368 = 643$ (kg).
\textbf{Đáp án của câu hỏi là C}
}
\end{ex}

\begin{ex} %[2D12V61]
Tổng của số lẻ bé nhất có ba chữ số khác nhau và $589$ là
\choice
{$682$}
{\True $692$}
{$702$}
{$690$}
\loigiai{
Số lẻ bé nhất có 3 chữ số khác nhau là 103. Tổng: $103 + 589 = 692$.
\textbf{Đáp án của câu hỏi là B}
}
\end{ex}

\begin{ex} %[2D12V61]
Bác Nam có hai thùng mật ong. Thùng thứ nhất chứa $185\text{ l}$, thùng thứ hai chứa $235\text{ l}$. Bác đã bán đi $120\text{ l}$. Bác còn lại bao nhiêu lít mật ong?
\choice
{$290\text{ l}$}
{$310\text{ l}$}
{$420\text{ l}$}
{\True $300\text{ l}$}
\loigiai{
Tổng số mật ong: $185 + 235 = 420$ l. Còn lại: $420 - 120 = 300$ l.
\textbf{Đáp án của câu hỏi là D}
}
\end{ex}

\textbf{Phần 2. Trắc nghiệm đúng sai}
\setcounterpageref{ex}{0}

\begin{ex} %[2D12V61]
Xét các phép tính cộng có nhớ trong phạm vi 1 000:
\choiceTF
{\True 378 + 245 = 623.}
{489 + 12 = 591.}
{\True Khi cộng có nhớ sang hàng chục, ta phải thêm 1 vào hàng chục của tổng.}
{\True Tổng của 450 và 550 bằng đúng 1 000.}
\loigiai{
\begin{itemize}
    \item Ý a) Đúng: 8 + 5 = 13 viết 3 nhớ 1; 7 + 4 = 11 nhớ 1 là 12 viết 2 nhớ 1; 3 + 2 = 5 nhớ 1 là 6. Đúng.
    \item Ý b) Sai: 489 + 12 = 501, không phải 591.
    \item Ý c) Đúng: Quy tắc cộng có nhớ. Đúng.
    \item Ý d) Đúng: 450 + 550 = 1 000. Đúng.
\end{itemize}
\begin{itemize}
    \item a) \textbf{Đúng}.
    \item b) \textbf{Sai}.
    \item c) \textbf{Đúng}.
    \item d) \textbf{Đúng}.
\end{itemize}
}
\end{ex}

\begin{ex} %[2D12V61]
Trường Tiểu học Hòa Bình có 368 học sinh nam và 357 học sinh nữ:
\choiceTF
{\True Trường có tất cả 725 học sinh.}
{\True Số học sinh nam nhiều hơn số học sinh nữ 11 bạn.}
{Nếu có thêm 5 bạn nữ chuyển đến thì số nam bằng số nữ.}
{\True Tổng số học sinh ít hơn 800 bạn.}
\loigiai{
\begin{itemize}
    \item Ý a) Đúng: 368 + 357 = 725 học sinh. Đúng.
    \item Ý b) Đúng: 368 - 357 = 11 bạn. Đúng.
    \item Ý c) Sai: 357 + 5 = 362 < 368 (vẫn ít hơn). Sai.
    \item Ý d) Đúng: 725 < 800. Đúng.
\end{itemize}
\begin{itemize}
    \item a) \textbf{Đúng}.
    \item b) \textbf{Đúng}.
    \item c) \textbf{Sai}.
    \item d) \textbf{Đúng}.
\end{itemize}
}
\end{ex}

\textbf{Phần 3. Trả lời ngắn}
\setcounterpageref{ex}{0}

\begin{ex} %[2D12V61]
Tính tổng: $476 + 389$.
\shortans{865}
\loigiai{
\textbf{Đáp số:} $865$\\
$476 + 389 = 865$.
}
\end{ex}

\begin{ex} %[2D12V61]
Tìm $x$ biết $x - 287 = 436$.
\shortans{723}
\loigiai{
\textbf{Đáp số:} $723$\\
$x = 436 + 287 = 723$.
}
\end{ex}

\begin{ex} %[2D12V61]
Đội một trồng được 345 cây, đội hai trồng được 385 cây. Cả hai đội trồng được bao nhiêu cây?
\shortans{730}
\loigiai{
\textbf{Đáp số:} $730$\\
$345 + 385 = 730$ cây.
}
\end{ex}

\begin{ex} %[2D12V61]
Tổng của hai số là 850. Nếu tăng số hạng thứ nhất 35 đơn vị và tăng số hạng thứ hai 15 đơn vị thì tổng mới là bao nhiêu?
\shortans{900}
\loigiai{
\textbf{Đáp số:} $900$\\
$850 + 35 + 15 = 900$.
}
\end{ex}

\begin{ex} %[2D12V61]
Tính: $195 + 285 + 320$.
\shortans{800}
\loigiai{
\textbf{Đáp số:} $800$\\
$195 + 285 + 320 = 800$.
}
\end{ex}

\begin{ex} %[2D12V61]
Một cuộn dây dài 258 m, cuộn thứ hai dài 367 m. Cả hai cuộn dài bao nhiêu mét?
\shortans{625}
\loigiai{
\textbf{Đáp số:} $625$\\
$258 + 367 = 625$ m.
}
\end{ex}

\newpage

\TaoPhieuKiemTra{Chủ đề 12 -- Phép cộng, phép trừ trong phạm vi 1 000}{62}{Phép trừ (có nhớ) trong phạm vi 1 000}{https://loptoanthaylinh.vn/thi/chuyen-de/2/toan2_cd_b62}
\textbf{Phần 1. Trắc nghiệm lựa chọn}
\setcounterpageref{ex}{0}

\begin{ex} %[2D12V62]
Kết quả của phép tính $632 - 275$ là
\choice
{$367$}
{$347$}
{$457$}
{\True $357$}
\loigiai{
$632 - 275 = 357$.
\textbf{Đáp án của câu hỏi là D}
}
\end{ex}

\begin{ex} %[2D12V62]
Tính hiệu của $800 - 345$.
\choice
{$465$}
{\True $455$}
{$555$}
{$445$}
\loigiai{
$800 - 345 = 455$.
\textbf{Đáp án của câu hỏi là B}
}
\end{ex}

\begin{ex} %[2D12V62]
Điền chữ số thích hợp vào dấu $*$ để: $7*4 - 289 = 445$.
\choice
{$2$}
{$4$}
{\True $3$}
{$5$}
\loigiai{
$445 + 289 = 734$, vậy $* = 3$.
\textbf{Đáp án của câu hỏi là C}
}
\end{ex}

\begin{ex} %[2D12V62]
Một kho thóc có $920$ bao thóc. Sau khi xuất đi một số bao thì trong kho còn lại $465$ bao. Hỏi kho đã xuất đi bao nhiêu bao thóc?
\choice
{\True $455$ bao}
{$445$ bao}
{$465$ bao}
{$555$ bao}
\loigiai{
$920 - 465 = 455$ (bao).
\textbf{Đáp án của câu hỏi là A}
}
\end{ex}

\begin{ex} %[2D12V62]
Tìm $x$ biết: $x + 286 = 614$.
\choice
{$318$}
{$338$}
{$428$}
{\True $328$}
\loigiai{
$x = 614 - 286 = 328$.
\textbf{Đáp án của câu hỏi là D}
}
\end{ex}

\begin{ex} %[2D12V62]
Một người có một mảnh đất dài $1 000\text{ m}$, người đó đã rào được $675\text{ m}$. Còn lại bao nhiêu mét đất chưa được rào?
\choice
{$315\text{ m}$}
{\True $325\text{ m}$}
{$335\text{ m}$}
{$425\text{ m}$}
\loigiai{
$1 000 - 675 = 325$ (m).
\textbf{Đáp án của câu hỏi là B}
}
\end{ex}

\begin{ex} %[2D12V62]
Một đàn vịt có $540$ con, trong đó có $285$ con vịt trắng, còn lại là vịt khoang. Đàn vịt có bao nhiêu con vịt khoang?
\choice
{$245$ con}
{$265$ con}
{\True $255$ con}
{$355$ con}
\loigiai{
$540 - 285 = 255$ (con vịt khoang).
\textbf{Đáp án của câu hỏi là C}
}
\end{ex}

\begin{ex} %[2D12V62]
Hiệu giữa số lớn nhất có ba chữ số và $678$ là
\choice
{\True $321$}
{$311$}
{$331$}
{$221$}
\loigiai{
$999 - 678 = 321$.
\textbf{Đáp án của câu hỏi là A}
}
\end{ex}

\begin{ex} %[2D12V62]
Điền dấu thích hợp: $524 - 186 \dots 612 - 274$.
\choice
{$>$}
{$<$}
{Không so sánh được}
{\True $=$}
\loigiai{
$524 - 186 = 338$; $612 - 274 = 338$. Hai vế bằng nhau.
\textbf{Đáp án của câu hỏi là D}
}
\end{ex}

\begin{ex} %[2D12V62]
Một cửa hàng có $700\text{ kg}$ gạo. Ngày đầu bán $265\text{ kg}$, ngày thứ hai bán $185\text{ kg}$. Cửa hàng còn lại bao nhiêu ki-lô-gam gạo?
\choice
{$240\text{ kg}$}
{\True $250\text{ kg}$}
{$260\text{ kg}$}
{$350\text{ kg}$}
\loigiai{
Đã bán: $265 + 185 = 450$ kg. Còn lại: $700 - 450 = 250$ kg.
\textbf{Đáp án của câu hỏi là B}
}
\end{ex}

\begin{ex} %[2D12V62]
Tìm số tự nhiên $x$ sao cho: $450 - x = 186 + 74$.
\choice
{$180$}
{$200$}
{\True $190$}
{$210$}
\loigiai{
$186 + 74 = 260$. $x = 450 - 260 = 190$.
\textbf{Đáp án của câu hỏi là C}
}
\end{ex}

\begin{ex} %[2D12V62]
Hiệu của hai số là $250$. Nếu bớt ở số bị trừ $30$ đơn vị và thêm vào số trừ $20$ đơn vị thì hiệu mới là bao nhiêu?
\choice
{\True $200$}
{$240$}
{$260$}
{$300$}
\loigiai{
Bớt số bị trừ 30 thì hiệu giảm 30; thêm số trừ 20 thì hiệu giảm tiếp 20. Hiệu mới: $250 - 30 - 20 = 200$.
\textbf{Đáp án của câu hỏi là A}
}
\end{ex}

\textbf{Phần 2. Trắc nghiệm đúng sai}
\setcounterpageref{ex}{0}

\begin{ex} %[2D12V62]
Xét tính đúng sai của các phép trừ có nhớ phạm vi 1 000:
\choiceTF
{\True 512 - 187 = 325.}
{700 - 250 = 550.}
{\True 1 000 - 555 = 445.}
{\True Khi thực hiện phép trừ có nhớ, ta mượn 1 ở hàng liền kề bên trái.}
\loigiai{
\begin{itemize}
    \item Ý a) Đúng: 12 - 7 = 5 nhớ 1; 11 - (8+1) = 2 nhớ 1; 5 - (1+1) = 3. Đúng.
    \item Ý b) Sai: 700 - 250 = 450, không phải 550.
    \item Ý c) Đúng: 1 000 - 555 = 445. Đúng.
    \item Ý d) Đúng: Quy tắc trừ có nhớ. Đúng.
\end{itemize}
\begin{itemize}
    \item a) \textbf{Đúng}.
    \item b) \textbf{Sai}.
    \item c) \textbf{Đúng}.
    \item d) \textbf{Đúng}.
\end{itemize}
}
\end{ex}

\begin{ex} %[2D12V62]
Một bể nước có 820 lít. Người ta dùng 350 lít để giặt và 280 lít để tưới rau:
\choiceTF
{\True Tổng số nước đã dùng là 630 lít.}
{\True Lượng nước còn lại trong bể là 190 lít.}
{Số nước giặt đồ ít hơn số nước tưới rau.}
{\True Nếu muốn bể đầy lại 820 lít thì cần bơm thêm 630 lít.}
\loigiai{
\begin{itemize}
    \item Ý a) Đúng: 350 + 280 = 630 l. Đúng.
    \item Ý b) Đúng: 820 - 630 = 190 l. Đúng.
    \item Ý c) Sai: 350 > 280 nên giặt đồ dùng nhiều hơn. Sai.
    \item Ý d) Đúng: Đúng vì đã dùng 630 lít.
\end{itemize}
\begin{itemize}
    \item a) \textbf{Đúng}.
    \item b) \textbf{Đúng}.
    \item c) \textbf{Sai}.
    \item d) \textbf{Đúng}.
\end{itemize}
}
\end{ex}

\textbf{Phần 3. Trả lời ngắn}
\setcounterpageref{ex}{0}

\begin{ex} %[2D12V62]
Tính hiệu: $813 - 457$.
\shortans{356}
\loigiai{
\textbf{Đáp số:} $356$\\
$813 - 457 = 356$.
}
\end{ex}

\begin{ex} %[2D12V62]
Tìm $x$ biết $904 - x = 368$.
\shortans{536}
\loigiai{
\textbf{Đáp số:} $536$\\
$x = 904 - 368 = 536$.
}
\end{ex}

\begin{ex} %[2D12V62]
Một nông trường có 640 cây ăn quả, trong đó có 275 cây xoài, còn lại là cây nhãn. Có bao nhiêu cây nhãn?
\shortans{365}
\loigiai{
\textbf{Đáp số:} $365$\\
$640 - 275 = 365$ cây.
}
\end{ex}

\begin{ex} %[2D12V62]
Tính giá trị biểu thức: $1 000 - 458 - 242$.
\shortans{300}
\loigiai{
\textbf{Đáp số:} $300$\\
$1 000 - 700 = 300$.
}
\end{ex}

\begin{ex} %[2D12V62]
Số bị trừ là 720, hiệu là 385. Số trừ là bao nhiêu?
\shortans{335}
\loigiai{
\textbf{Đáp số:} $335$\\
$720 - 385 = 335$.
}
\end{ex}

\begin{ex} %[2D12V62]
Một cuộn dây dài 500 m. Người ta cắt một đoạn dài 175 m và một đoạn dài 125 m. Cuộn dây còn lại bao nhiêu mét?
\shortans{200}
\loigiai{
\textbf{Đáp số:} $200$\\
$500 - (175 + 125) = 200$ m.
}
\end{ex}

\newpage

\TaoPhieuKiemTra{Chủ đề 12 -- Phép cộng, phép trừ trong phạm vi 1 000}{63}{Luyện tập chung}{https://loptoanthaylinh.vn/thi/chuyen-de/2/toan2_cd_b63}
\textbf{Phần 1. Trắc nghiệm lựa chọn}
\setcounterpageref{ex}{0}

\begin{ex} %[2D12V63]
Giá trị của biểu thức $456 + 234 - 190$ là
\choice
{$490$}
{$510$}
{\True $500$}
{$520$}
\loigiai{
$456 + 234 = 690$; $690 - 190 = 500$.
\textbf{Đáp án của câu hỏi là C}
}
\end{ex}

\begin{ex} %[2D12V63]
Tổng của hai số là $750$. Nếu số hạng thứ nhất tăng $50$ và số hạng thứ hai giảm $20$ thì tổng mới là
\choice
{$720$}
{\True $780$}
{$800$}
{$770$}
\loigiai{
Tổng mới: $750 + 50 - 20 = 780$.
\textbf{Đáp án của câu hỏi là B}
}
\end{ex}

\begin{ex} %[2D12V63]
Tìm $x$ biết: $x + 145 = 500 - 120$.
\choice
{$245$}
{$225$}
{$380$}
{\True $235$}
\loigiai{
$500 - 120 = 380$. $x = 380 - 145 = 235$.
\textbf{Đáp án của câu hỏi là D}
}
\end{ex}

\begin{ex} %[2D12V63]
Một đoàn tàu chở $680$ hành khách. Đến ga thứ nhất có $125$ người xuống tàu và $45$ người lên tàu. Trên tàu còn lại bao nhiêu hành khách?
\choice
{\True $600$ người}
{$590$ người}
{$610$ người}
{$555$ người}
\loigiai{
$680 - 125 + 45 = 600$ (hành khách).
\textbf{Đáp án của câu hỏi là A}
}
\end{ex}

\begin{ex} %[2D12V63]
Điền dấu thích hợp: $450 + 350 \dots 1 000 - 200$.
\choice
{$>$}
{$<$}
{\True $=$}
{Không so sánh được}
\loigiai{
$450 + 350 = 800$; $1 000 - 200 = 800$. Hai vế bằng nhau.
\textbf{Đáp án của câu hỏi là C}
}
\end{ex}

\begin{ex} %[2D12V63]
Một người có hai tấm vải: tấm thứ nhất dài $245\text{ m}$, tấm thứ hai dài hơn tấm thứ nhất $30\text{ m}$. Cả hai tấm vải dài bao nhiêu mét?
\choice
{$510\text{ m}$}
{\True $520\text{ m}$}
{$530\text{ m}$}
{$275\text{ m}$}
\loigiai{
Tấm 2: $245 + 30 = 275$ m. Cả hai tấm: $245 + 275 = 520$ m.
\textbf{Đáp án của câu hỏi là B}
}
\end{ex}

\begin{ex} %[2D12V63]
Số tự nhiên lớn nhất có ba chữ số khác nhau trừ đi số tự nhiên bé nhất có ba chữ số khác nhau được kết quả là
\choice
{$888$}
{$899$}
{$875$}
{\True $885$}
\loigiai{
Số lớn nhất khác nhau: 987. Bé nhất khác nhau: 102. Hiệu: $987 - 102 = 885$.
\textbf{Đáp án của câu hỏi là D}
}
\end{ex}

\begin{ex} %[2D12V63]
Thùng thứ nhất có $350$ lít dầu, thùng thứ hai có ít hơn thùng thứ nhất $70$ lít. Cả hai thùng có bao nhiêu lít dầu?
\choice
{\True $630$ lít}
{$620$ lít}
{$640$ lít}
{$560$ lít}
\loigiai{
Thùng 2: $350 - 70 = 280$ l. Cả hai thùng: $350 + 280 = 630$ l.
\textbf{Đáp án của câu hỏi là A}
}
\end{ex}

\begin{ex} %[2D12V63]
Tính tổng của ba số tròn trăm liên tiếp, biết số bé nhất là $200$.
\choice
{$800$}
{$700$}
{\True $900$}
{$1 000$}
\loigiai{
Ba số là 200, 300, 400. Tổng: $200 + 300 + 400 = 900$.
\textbf{Đáp án của câu hỏi là C}
}
\end{ex}

\begin{ex} %[2D12V63]
Bác Nam nuôi $320$ con gà và $280$ con vịt. Sau khi bán đi một số con gia cầm thì bác còn lại $150$ con. Bác Nam đã bán bao nhiêu con gia cầm?
\choice
{$440$ con}
{\True $450$ con}
{$460$ con}
{$350$ con}
\loigiai{
Tổng gia cầm: $320 + 280 = 600$ con. Bác đã bán: $600 - 150 = 450$ con.
\textbf{Đáp án của câu hỏi là B}
}
\end{ex}

\begin{ex} %[2D12V63]
Một hình tam giác có chu vi $600\text{ cm}$. Cạnh thứ nhất dài $180\text{ cm}$, cạnh thứ hai dài $220\text{ cm}$. Độ dài cạnh thứ ba là
\choice
{$190\text{ cm}$}
{$210\text{ cm}$}
{$250\text{ cm}$}
{\True $200\text{ cm}$}
\loigiai{
Cạnh thứ ba: $600 - (180 + 220) = 200$ cm.
\textbf{Đáp án của câu hỏi là D}
}
\end{ex}

\begin{ex} %[2D12V63]
Cho $x + y = 500$. Nếu giữ nguyên $x$ và bớt $y$ đi $120$ đơn vị thì tổng mới là bao nhiêu?
\choice
{\True $380$}
{$420$}
{$620$}
{$480$}
\loigiai{
Tổng mới: $500 - 120 = 380$.
\textbf{Đáp án của câu hỏi là A}
}
\end{ex}

\textbf{Phần 2. Trắc nghiệm đúng sai}
\setcounterpageref{ex}{0}

\begin{ex} %[2D12V63]
Xét tính đúng sai của các bài toán tổng hợp cộng trừ phạm vi 1 000:
\choiceTF
{\True 350 + 450 - 200 = 600.}
{1 000 - 300 - 200 = 600.}
{\True Nếu a + b = c thì c - a = b.}
{\True Tổng của số chẵn lớn nhất có 3 chữ số và số lẻ bé nhất có 1 chữ số bằng 999.}
\loigiai{
\begin{itemize}
    \item Ý a) Đúng: 800 - 200 = 600. Đúng.
    \item Ý b) Sai: 1 000 - 500 = 500, không phải 600. Sai.
    \item Ý c) Đúng: Tính chất của phép trừ và phép cộng. Đúng.
    \item Ý d) Đúng: 998 + 1 = 999. Đúng.
\end{itemize}
\begin{itemize}
    \item a) \textbf{Đúng}.
    \item b) \textbf{Sai}.
    \item c) \textbf{Đúng}.
    \item d) \textbf{Đúng}.
\end{itemize}
}
\end{ex}

\begin{ex} %[2D12V63]
Cửa hàng vật liệu có 900 bao xi măng. Buổi sáng bán 350 bao, buổi chiều bán 250 bao:
\choiceTF
{\True Cả ngày bán được 600 bao xi măng.}
{\True Cửa hàng còn lại 300 bao xi măng.}
{Số bao bán buổi sáng gấp đôi buổi chiều.}
{\True Số xi măng còn lại đúng bằng 1/2 số xi măng đã bán.}
\loigiai{
\begin{itemize}
    \item Ý a) Đúng: 350 + 250 = 600 bao. Đúng.
    \item Ý b) Đúng: 900 - 600 = 300 bao. Đúng.
    \item Ý c) Sai: 250 x 2 = 500 != 350. Sai.
    \item Ý d) Đúng: 300 = 600 : 2. Đúng.
\end{itemize}
\begin{itemize}
    \item a) \textbf{Đúng}.
    \item b) \textbf{Đúng}.
    \item c) \textbf{Sai}.
    \item d) \textbf{Đúng}.
\end{itemize}
}
\end{ex}

\textbf{Phần 3. Trả lời ngắn}
\setcounterpageref{ex}{0}

\begin{ex} %[2D12V63]
Tính giá trị: $523 + 277 - 350$.
\shortans{450}
\loigiai{
\textbf{Đáp số:} $450$\\
$800 - 350 = 450$.
}
\end{ex}

\begin{ex} %[2D12V63]
Tìm $x$ biết: $x - 175 = 625 - 200$.
\shortans{600}
\loigiai{
\textbf{Đáp số:} $600$\\
$625 - 200 = 425$. $x = 425 + 175 = 600$.
}
\end{ex}

\begin{ex} %[2D12V63]
Kho gạo có 1 000 kg gạo. Sau hai lần xuất thì còn lại 320 kg. Hỏi kho đã xuất bao nhiêu ki-lô-gam gạo?
\shortans{680}
\loigiai{
\textbf{Đáp số:} $680$\\
$1 000 - 320 = 680$ kg.
}
\end{ex}

\begin{ex} %[2D12V63]
Tổng của hai số là 900. Số thứ nhất là 345. Tìm số thứ hai.
\shortans{555}
\loigiai{
\textbf{Đáp số:} $555$\\
$900 - 345 = 555$.
}
\end{ex}

\begin{ex} %[2D12V63]
Một hình tam giác có độ dài các cạnh lần lượt là 250 cm, 300 cm và 450 cm. Chu vi hình tam giác đó là bao nhiêu mét?
\shortans{10}
\loigiai{
\textbf{Đáp số:} $10$\\
Chu vi: $250 + 300 + 450 = 1 000\text{ cm} = 10$ m.
}
\end{ex}

\begin{ex} %[2D12V63]
Có 3 xe chở hàng: xe 1 chở 280 bao, xe 2 chở 320 bao, xe 3 chở 400 bao. Cả 3 xe chở bao nhiêu bao?
\shortans{1 000}
\loigiai{
\textbf{Đáp số:} $1 000$\\
$280 + 320 + 400 = 1 000$ bao.
}
\end{ex}

\newpage


% =========================================================================
% CHỦ ĐỀ 13 - LÀM QUEN VỚI YẾU TỐ THỐNG KÊ, XÁC SUẤT
% =========================================================================
\section{Làm quen với yếu tố thống kê, xác suất}

\TaoPhieuKiemTra{Chủ đề 13 -- Làm quen với yếu tố thống kê, xác suất}{64}{Thu thập, phân loại, kiểm đếm số liệu}{https://loptoanthaylinh.vn/thi/chuyen-de/2/toan2_cd_b64}
\textbf{Phần 1. Trắc nghiệm lựa chọn}
\setcounterpageref{ex}{0}

\begin{ex} %[2D13V64]
Để biết số lượng học sinh thích các môn thể thao khác nhau trong lớp, bước đầu tiên em cần làm là
\choice
{Vẽ biểu đồ}
{\True Thu thập số liệu}
{Tính điểm trung bình}
{Xếp hạng học sinh}
\loigiai{
Bước đầu tiên là thu thập thông tin, số liệu từ các bạn.
\textbf{Đáp án của câu hỏi là B}
}
\end{ex}

\begin{ex} %[2D13V64]
Khi quan sát một giỏ hoa quả có cam, táo, chuối, việc xếp cam vào một nhóm, táo vào một nhóm, chuối vào một nhóm gọi là gì?
\choice
{\True Phân loại}
{Kiểm đếm}
{Tính toán}
{Dự đoán}
\loigiai{
Chia các đối tượng theo từng loại đặc điểm gọi là phân loại.
\textbf{Đáp án của câu hỏi là A}
}
\end{ex}

\begin{ex} %[2D13V64]
Trong vườn có $12$ con bướm trắng, $8$ con bướm vàng và $5$ con bướm nâu. Tổng số bướm quan sát được là
\choice
{$24$ con}
{$26$ con}
{\True $25$ con}
{$20$ con}
\loigiai{
$12 + 8 + 5 = 25$ (con bướm).
\textbf{Đáp án của câu hỏi là C}
}
\end{ex}

\begin{ex} %[2D13V64]
Kiểm đếm số đồ dùng học tập của An thấy có: $5$ chiếc bút chì, $3$ cục tẩy, $8$ quyển vở. Đồ dùng nào có số lượng nhiều nhất?
\choice
{Bút chì}
{Cục tẩy}
{Cả ba bằng nhau}
{\True Quyển vở}
\loigiai{
Vì $8 > 5 > 3$ nên vở có số lượng nhiều nhất.
\textbf{Đáp án của câu hỏi là D}
}
\end{ex}

\begin{ex} %[2D13V64]
Trong đợt kiểm tra sức khỏe, tổ Một có $4$ bạn cận thị và $6$ bạn không cận thị. Tổ Một có tất cả bao nhiêu bạn?
\choice
{$8$ bạn}
{\True $10$ bạn}
{$12$ bạn}
{$14$ bạn}
\loigiai{
$4 + 6 = 10$ (bạn).
\textbf{Đáp án của câu hỏi là B}
}
\end{ex}

\begin{ex} %[2D13V64]
Một nhóm bạn đếm số xe đi qua cổng trường trong 5 phút: 15 xe máy, 4 xe đạp, 3 xe ô tô. Số xe máy nhiều hơn tổng số xe đạp và ô tô là
\choice
{\True $8$ xe}
{$7$ xe}
{$9$ xe}
{$10$ xe}
\loigiai{
Tổng xe đạp và ô tô: $4 + 3 = 7$ xe. Xe máy nhiều hơn: $15 - 7 = 8$ xe.
\textbf{Đáp án của câu hỏi là A}
}
\end{ex}

\begin{ex} %[2D13V64]
Phân loại các số sau thành hai nhóm chẵn và lẻ: $12, 17, 24, 35, 40, 51$. Nhóm số chẵn gồm mấy số?
\choice
{$2$ số}
{$4$ số}
{\True $3$ số}
{$5$ số}
\loigiai{
Các số chẵn là 12, 24, 40 (3 số). Các số lẻ là 17, 35, 51 (3 số).
\textbf{Đáp án của câu hỏi là C}
}
\end{ex}

\begin{ex} %[2D13V64]
Khi đếm số lượng bóng bay, bạn Mai dùng các vạch gạch (mỗi cụm 5 vạch biểu thị 5 quả). Mai vẽ được 4 cụm vạch và 2 vạch lẻ. Mai đã đếm được bao nhiêu quả bóng?
\choice
{$20$ quả}
{$24$ quả}
{$18$ quả}
{\True $22$ quả}
\loigiai{
$5 \times 4 + 2 = 22$ (quả bóng).
\textbf{Đáp án của câu hỏi là D}
}
\end{ex}

\begin{ex} %[2D13V64]
Trong lớp học có $18$ bạn nam và $16$ bạn nữ. Sau giờ tan học có $10$ bạn nam và $8$ bạn nữ ra về trước. Trong lớp còn lại bao nhiêu bạn?
\choice
{$14$ bạn}
{\True $16$ bạn}
{$18$ bạn}
{$15$ bạn}
\loigiai{
Tổng sĩ số: $18 + 16 = 34$. Đã về: $10 + 8 = 18$. Còn lại: $34 - 18 = 16$ bạn.
\textbf{Đáp án của câu hỏi là B}
}
\end{ex}

\begin{ex} %[2D13V64]
Một khay đựng các hình hình học: 6 hình tam giác, 8 hình vuông, 10 hình tròn. Khẳng định nào sau đây là SAI?
\choice
{\True Hình tam giác có số lượng nhiều nhất}
{Hình tròn có số lượng nhiều nhất}
{Tổng số hình là 24 hình}
{Hình vuông ít hơn hình tròn 2 hình}
\loigiai{
Hình tròn nhiều nhất (10 hình), tam giác ít nhất (6 hình), nên khẳng định 'tam giác nhiều nhất' là SAI.
\textbf{Đáp án của câu hỏi là A}
}
\end{ex}

\begin{ex} %[2D13V64]
Thu thập số lượng hoa trong lọ: Hoa hồng có $15$ bông, hoa cúc có $20$ bông, hoa ly có $5$ bông. Tỉ lệ số bông hoa cúc so với hoa ly gấp mấy lần?
\choice
{$3$ lần}
{$5$ lần}
{\True $4$ lần}
{$2$ lần}
\loigiai{
$20 : 5 = 4$ lần.
\textbf{Đáp án của câu hỏi là C}
}
\end{ex}

\begin{ex} %[2D13V64]
Kiểm đếm số sách trên 3 ngăn: Ngăn 1 có $35$ cuốn, ngăn 2 có $45$ cuốn, ngăn 3 có $20$ cuốn. Ngăn 2 có nhiều hơn ngăn 3 bao nhiêu cuốn sách?
\choice
{$20$ cuốn}
{$15$ cuốn}
{$10$ cuốn}
{\True $25$ cuốn}
\loigiai{
$45 - 20 = 25$ cuốn.
\textbf{Đáp án của câu hỏi là D}
}
\end{ex}

\textbf{Phần 2. Trắc nghiệm đúng sai}
\setcounterpageref{ex}{0}

\begin{ex} %[2D13V64]
Xét các khẳng định về thu thập và phân loại dữ liệu:
\choiceTF
{\True Phân loại giúp sắp xếp các đồ vật hoặc số liệu có cùng đặc điểm vào chung một nhóm.}
{Kiểm đếm chỉ được thực hiện bằng cách cân các đồ vật.}
{\True Trong một bảng thống kê, tổng số liệu từng phần bằng số liệu toàn thể.}
{\True Có thể phân loại học sinh trong lớp theo nhiều tiêu chí như: giới tính, tháng sinh, sở thích.}
\loigiai{
\begin{itemize}
    \item Ý a) Đúng: Định nghĩa phân loại. Đúng.
    \item Ý b) Sai: Kiểm đếm là đếm số lượng, không phải chỉ cân. Sai.
    \item Ý c) Đúng: Đúng.
    \item Ý d) Đúng: Đúng.
\end{itemize}
\begin{itemize}
    \item a) \textbf{Đúng}.
    \item b) \textbf{Sai}.
    \item c) \textbf{Đúng}.
    \item d) \textbf{Đúng}.
\end{itemize}
}
\end{ex}

\begin{ex} %[2D13V64]
Một hộp bi có 14 viên bi xanh, 16 viên bi đỏ và 10 viên bi vàng:
\choiceTF
{\True Tổng số bi trong hộp là 40 viên.}
{\True Bi đỏ có số lượng nhiều nhất.}
{Số bi xanh ít hơn số bi vàng.}
{\True Số bi đỏ nhiều hơn số bi vàng là 6 viên.}
\loigiai{
\begin{itemize}
    \item Ý a) Đúng: 14 + 16 + 10 = 40 viên. Đúng.
    \item Ý b) Đúng: 16 > 14 > 10. Đúng.
    \item Ý c) Sai: Bi xanh có 14, bi vàng có 10 nên bi xanh nhiều hơn. Sai.
    \item Ý d) Đúng: 16 - 10 = 6 viên. Đúng.
\end{itemize}
\begin{itemize}
    \item a) \textbf{Đúng}.
    \item b) \textbf{Đúng}.
    \item c) \textbf{Sai}.
    \item d) \textbf{Đúng}.
\end{itemize}
}
\end{ex}

\textbf{Phần 3. Trả lời ngắn}
\setcounterpageref{ex}{0}

\begin{ex} %[2D13V64]
Trong vườn có 15 cây cam, 18 cây bưởi và 12 cây chanh. Hỏi trong vườn có tất cả bao nhiêu cây ăn quả?
\shortans{45}
\loigiai{
\textbf{Đáp số:} $45$\\
$15 + 18 + 12 = 45$ cây.
}
\end{ex}

\begin{ex} %[2D13V64]
Một hộp đựng 25 chiếc bút gồm bút bi xanh và bút bi đen, trong đó có 16 bút bi xanh. Hỏi có bao nhiêu bút bi đen?
\shortans{9}
\loigiai{
\textbf{Đáp số:} $9$\\
$25 - 16 = 9$ chiếc.
}
\end{ex}

\begin{ex} %[2D13V64]
Lớp 2A có 36 học sinh. Khi khảo sát môn thể thao yêu thích: 15 bạn thích bóng đá, 12 bạn thích bơi lội, còn lại thích cầu lông. Có bao nhiêu bạn thích cầu lông?
\shortans{9}
\loigiai{
\textbf{Đáp số:} $9$\\
$36 - 15 - 12 = 9$ bạn.
}
\end{ex}

\begin{ex} %[2D13V64]
Kiểm đếm số xe đạp trong bãi: dãy A có 24 xe, dãy B có 28 xe. Cả hai dãy có bao nhiêu xe đạp?
\shortans{52}
\loigiai{
\textbf{Đáp số:} $52$\\
$24 + 28 = 52$ xe.
}
\end{ex}

\begin{ex} %[2D13V64]
Trong đĩa có 8 quả chuối, 6 quả táo và 10 quả dâu tây. Có tất cả bao nhiêu quả?
\shortans{24}
\loigiai{
\textbf{Đáp số:} $24$\\
$8 + 6 + 10 = 24$ quả.
}
\end{ex}

\begin{ex} %[2D13V64]
Một nhóm bạn xếp 5 bàn học, mỗi bàn có 4 bạn ngồi. Đếm được tất cả bao nhiêu bạn?
\shortans{20}
\loigiai{
\textbf{Đáp số:} $20$\\
$4 \times 5 = 20$ bạn.
}
\end{ex}

\newpage

\TaoPhieuKiemTra{Chủ đề 13 -- Làm quen với yếu tố thống kê, xác suất}{65}{Biểu đồ tranh}{https://loptoanthaylinh.vn/thi/chuyen-de/2/toan2_cd_b65}
\textbf{Phần 1. Trắc nghiệm lựa chọn}
\setcounterpageref{ex}{0}

\begin{ex} %[2D13V65]
Trên biểu đồ tranh, mỗi biểu tượng ngôi sao biểu thị cho $2$ điểm tốt. Bạn Mai có $5$ biểu tượng ngôi sao. Mai có bao nhiêu điểm tốt?
\choice
{$7$ điểm}
{$5$ điểm}
{\True $10$ điểm}
{$12$ điểm}
\loigiai{
$2 \times 5 = 10$ (điểm tốt).
\textbf{Đáp án của câu hỏi là C}
}
\end{ex}

\begin{ex} %[2D13V65]
Một biểu đồ tranh biểu diễn số cây trồng của các tổ, mỗi hình cái cây đại diện cho $5$ cây xanh. Tổ Một trồng được $4$ hình cái cây. Tổ Một trồng được bao nhiêu cây xanh?
\choice
{\True $20$ cây}
{$9$ cây}
{$15$ cây}
{$25$ cây}
\loigiai{
$5 \times 4 = 20$ (cây xanh).
\textbf{Đáp án của câu hỏi là A}
}
\end{ex}

\begin{ex} %[2D13V65]
Biểu đồ tranh hoa điểm mười của 4 bạn: An có 6 bông, Bình có 4 bông, Cúc có 8 bông, Dũng có 5 bông (mỗi bông biểu thị 1 điểm 10). Bạn đạt nhiều điểm 10 nhất là
\choice
{An}
{Bình}
{Dũng}
{\True Cúc}
\loigiai{
Cúc có 8 bông hoa điểm 10, nhiều nhất.
\textbf{Đáp án của câu hỏi là D}
}
\end{ex}

\begin{ex} %[2D13V65]
Trên biểu đồ tranh bán xe đạp, mỗi hình xe đạp biểu thị $10$ chiếc xe. Cửa hàng vẽ $6$ hình xe đạp. Cửa hàng đã bán được bao nhiêu chiếc xe đạp?
\choice
{$16$ chiếc}
{\True $60$ chiếc}
{$50$ chiếc}
{$6$ chiếc}
\loigiai{
$10 \times 6 = 60$ (chiếc xe đạp).
\textbf{Đáp án của câu hỏi là B}
}
\end{ex}

\begin{ex} %[2D13V65]
Biểu đồ tranh số con cá câu được của ba bạn: Nam câu được 3 con, Việt câu được 5 con, Robot câu được 4 con. Cả ba bạn câu được bao nhiêu con cá?
\choice
{$10$ con}
{$11$ con}
{\True $12$ con}
{$14$ con}
\loigiai{
$3 + 5 + 4 = 12$ (con cá).
\textbf{Đáp án của câu hỏi là C}
}
\end{ex}

\begin{ex} %[2D13V65]
Trên biểu đồ số bạn đi học bằng các phương tiện: Đi bộ có 8 bạn, đi xe đạp có 12 bạn, đi ô tô buýt có 15 bạn. Số bạn đi bộ ít hơn số bạn đi xe đạp là
\choice
{\True $4$ bạn}
{$3$ bạn}
{$5$ bạn}
{$7$ bạn}
\loigiai{
$12 - 8 = 4$ (bạn).
\textbf{Đáp án của câu hỏi là A}
}
\end{ex}

\begin{ex} %[2D13V65]
Mỗi biểu tượng quả táo biểu thị $2\text{ kg}$ táo. Nếu muốn biểu thị $16\text{ kg}$ táo thì cần vẽ bao nhiêu biểu tượng quả táo?
\choice
{$6$ biểu tượng}
{$10$ biểu tượng}
{$14$ biểu tượng}
{\True $8$ biểu tượng}
\loigiai{
$16 : 2 = 8$ (biểu tượng).
\textbf{Đáp án của câu hỏi là D}
}
\end{ex}

\begin{ex} %[2D13V65]
Biểu đồ số con vật nuôi trong trang trại: Bò có 20 con, dê có 15 con, cừu có 25 con. Tổng số con vật nuôi là
\choice
{$50$ con}
{\True $60$ con}
{$55$ con}
{$65$ con}
\loigiai{
$20 + 15 + 25 = 60$ (con).
\textbf{Đáp án của câu hỏi là B}
}
\end{ex}

\begin{ex} %[2D13V65]
Trên biểu đồ tranh, mỗi khuôn mặt cười biểu thị cho $5$ bạn học sinh yêu thích môn Toán. Lớp có $7$ khuôn mặt cười. Có bao nhiêu bạn yêu thích môn Toán?
\choice
{$30$ bạn}
{$40$ bạn}
{\True $35$ bạn}
{$12$ bạn}
\loigiai{
$5 \times 7 = 35$ (bạn).
\textbf{Đáp án của câu hỏi là C}
}
\end{ex}

\begin{ex} %[2D13V65]
Biểu đồ số sách đọc trong tuần: Tuần 1 đọc 4 cuốn, tuần 2 đọc 6 cuốn, tuần 3 đọc 5 cuốn, tuần 4 đọc 7 cuốn. Trung bình hoặc tổng số sách đọc trong 4 tuần là
\choice
{\True $22$ cuốn}
{$20$ cuốn}
{$24$ cuốn}
{$18$ cuốn}
\loigiai{
Tổng số sách: $4 + 6 + 5 + 7 = 22$ (cuốn).
\textbf{Đáp án của câu hỏi là A}
}
\end{ex}

\begin{ex} %[2D13V65]
Biểu đồ số ly kem bán được trong ngày: Buổi sáng có 6 hình ly kem, buổi chiều có 9 hình ly kem (mỗi hình ly kem ứng với 2 ly thật). Buổi chiều bán nhiều hơn buổi sáng bao nhiêu ly kem?
\choice
{$3$ ly}
{$4$ ly}
{$8$ ly}
{\True $6$ ly}
\loigiai{
Nhiều hơn: $(9 - 6) \times 2 = 6$ ly kem.
\textbf{Đáp án của câu hỏi là D}
}
\end{ex}

\begin{ex} %[2D13V65]
Khi đọc biểu đồ tranh, thông tin nào sau đây là quan trọng nhất cần xem trước tiên để hiểu tỉ lệ?
\choice
{Tên người vẽ biểu đồ}
{\True Mỗi biểu tượng biểu thị cho bao nhiêu đối tượng}
{Màu sắc của biểu tượng}
{Kích thước của khung tranh}
\loigiai{
Cần xem quy ước mỗi biểu tượng biểu thị cho bao nhiêu đơn vị để tính đúng.
\textbf{Đáp án của câu hỏi là B}
}
\end{ex}

\textbf{Phần 2. Trắc nghiệm đúng sai}
\setcounterpageref{ex}{0}

\begin{ex} %[2D13V65]
Xét các nhận định về đọc và biểu diễn biểu đồ tranh:
\choiceTF
{\True Biểu đồ tranh dùng các hình ảnh hoặc biểu tượng để biểu diễn số liệu.}
{\True Một biểu tượng trên biểu đồ tranh có thể đại diện cho 1, 2, 5 hoặc 10 đối tượng tùy theo quy ước.}
{Biểu đồ tranh không thể dùng để so sánh số lượng giữa các nhóm.}
{\True Muốn biết số lượng thực tế, ta nhân số biểu tượng với giá trị của mỗi biểu tượng.}
\loigiai{
\begin{itemize}
    \item Ý a) Đúng: Đúng định nghĩa biểu đồ tranh.
    \item Ý b) Đúng: Đúng quy ước biểu diễn.
    \item Ý c) Sai: Biểu đồ tranh giúp so sánh rất trực quan. Sai.
    \item Ý d) Đúng: Đúng phương pháp tính.
\end{itemize}
\begin{itemize}
    \item a) \textbf{Đúng}.
    \item b) \textbf{Đúng}.
    \item c) \textbf{Sai}.
    \item d) \textbf{Đúng}.
\end{itemize}
}
\end{ex}

\begin{ex} %[2D13V65]
Biểu đồ tranh số quả cam thu hoạch được: Cây 1 có 3 hình sọt cam, cây 2 có 5 hình sọt cam. Mỗi hình sọt cam biểu thị 10 quả cam:
\choiceTF
{\True Cây 1 thu hoạch được 30 quả cam.}
{\True Cây 2 thu hoạch được 50 quả cam.}
{Cả hai cây thu hoạch được 70 quả cam.}
{\True Cây 2 thu hoạch nhiều hơn cây 1 là 20 quả cam.}
\loigiai{
\begin{itemize}
    \item Ý a) Đúng: 3 x 10 = 30 quả. Đúng.
    \item Ý b) Đúng: 5 x 10 = 50 quả. Đúng.
    \item Ý c) Sai: 30 + 50 = 80 quả, không phải 70 quả. Sai.
    \item Ý d) Đúng: 50 - 30 = 20 quả. Đúng.
\end{itemize}
\begin{itemize}
    \item a) \textbf{Đúng}.
    \item b) \textbf{Đúng}.
    \item c) \textbf{Sai}.
    \item d) \textbf{Đúng}.
\end{itemize}
}
\end{ex}

\textbf{Phần 3. Trả lời ngắn}
\setcounterpageref{ex}{0}

\begin{ex} %[2D13V65]
Trên biểu đồ tranh, mỗi hình bông hoa biểu thị 5 điểm mười. Tổ Một có 6 hình bông hoa. Tổ Một có tất cả bao nhiêu điểm mười?
\shortans{30}
\loigiai{
\textbf{Đáp số:} $30$\\
$5 \times 6 = 30$ điểm mười.
}
\end{ex}

\begin{ex} %[2D13V65]
Để biểu diễn 40 chiếc xe ô tô trên biểu đồ tranh, biết mỗi hình chiếc xe biểu thị 10 chiếc. Cần vẽ bao nhiêu hình chiếc xe?
\shortans{4}
\loigiai{
\textbf{Đáp số:} $4$\\
$40 : 10 = 4$ hình.
}
\end{ex}

\begin{ex} %[2D13V65]
Biểu đồ số bóng bàn bán được: Ngày thứ nhất có 8 quả, ngày thứ hai có 12 quả (mỗi hình vẽ 1 quả bóng). Cả hai ngày bán được bao nhiêu quả?
\shortans{20}
\loigiai{
\textbf{Đáp số:} $20$\\
$8 + 12 = 20$ quả.
}
\end{ex}

\begin{ex} %[2D13V65]
Mỗi biểu tượng cuốn sách đại diện cho 2 cuốn sách. Mai vẽ 9 biểu tượng cuốn sách. Mai có bao nhiêu cuốn sách?
\shortans{18}
\loigiai{
\textbf{Đáp số:} $18$\\
$2 \times 9 = 18$ cuốn.
}
\end{ex}

\begin{ex} %[2D13V65]
Tổ Hai trồng được 25 cây xanh, trên biểu đồ mỗi biểu tượng đại diện cho 5 cây. Cần vẽ mấy biểu tượng?
\shortans{5}
\loigiai{
\textbf{Đáp số:} $5$\\
$25 : 5 = 5$ biểu tượng.
}
\end{ex}

\begin{ex} %[2D13V65]
Trên biểu đồ tranh, bạn Hải có 7 ngôi sao, bạn Nam có 4 ngôi sao (mỗi ngôi sao là 2 điểm tốt). Hải nhiều hơn Nam bao nhiêu điểm tốt?
\shortans{6}
\loigiai{
\textbf{Đáp số:} $6$\\
$(7 - 4) \times 2 = 6$ điểm tốt.
}
\end{ex}

\newpage

\TaoPhieuKiemTra{Chủ đề 13 -- Làm quen với yếu tố thống kê, xác suất}{66}{Chắc chắn, có thể, không thể}{https://loptoanthaylinh.vn/thi/chuyen-de/2/toan2_cd_b66}
\textbf{Phần 1. Trắc nghiệm lựa chọn}
\setcounterpageref{ex}{0}

\begin{ex} %[2D13V66]
Một hộp chỉ đựng toàn bóng màu đỏ. Nếu nhắm mắt lấy ra một quả bóng thì biến cố nào sau đây là 'Chắc chắn' xảy ra?
\choice
{Lấy được quả bóng màu xanh}
{\True Lấy được quả bóng màu đỏ}
{Lấy được quả bóng màu vàng}
{Không lấy được quả nào}
\loigiai{
Hộp toàn bóng đỏ nên chắc chắn lấy ra quả bóng đỏ.
\textbf{Đáp án của câu hỏi là B}
}
\end{ex}

\begin{ex} %[2D13V66]
Trong túi có 3 viên bi xanh và 2 viên bi đỏ. Lấy ngẫu nhiên một viên bi. Khẳng định nào sau đây là ĐÚNG?
\choice
{Chắc chắn lấy được viên bi màu vàng}
{Không thể lấy được viên bi màu xanh}
{Chắc chắn lấy được bi đỏ}
{\True Có thể lấy được viên bi màu đỏ}
\loigiai{
Vì trong túi có cả bi xanh và đỏ nên có thể lấy được bi đỏ (hoặc xanh).
\textbf{Đáp án của câu hỏi là D}
}
\end{ex}

\begin{ex} %[2D13V66]
Gieo một con xúc xắc có các mặt từ 1 đến 6 chấm. Biến cố 'Xuất hiện mặt 8 chấm' là sự kiện
\choice
{\True Không thể xảy ra}
{Chắc chắn xảy ra}
{Có thể xảy ra}
{Luôn luôn xảy ra}
\loigiai{
Xúc xắc chỉ có từ 1 đến 6 chấm nên không thể xuất hiện mặt 8 chấm.
\textbf{Đáp án của câu hỏi là A}
}
\end{ex}

\begin{ex} %[2D13V66]
Hôm nay là ngày thứ Hai. Biến cố 'Ngày mai là ngày thứ Ba' là sự kiện
\choice
{Có thể xảy ra}
{Không thể xảy ra}
{\True Chắc chắn xảy ra}
{Không xác định}
\loigiai{
Thứ Hai luôn luôn liền kề trước thứ Ba nên chắc chắn xảy ra.
\textbf{Đáp án của câu hỏi là C}
}
\end{ex}

\begin{ex} %[2D13V66]
Một hộp có 5 chiếc bút chì màu đen. Bạn An thò tay vào hộp lấy ra 1 chiếc bút. Bạn An
\choice
{Có thể lấy được chiếc bút chì màu đỏ}
{\True Không thể lấy được chiếc bút chì màu đỏ}
{Chắc chắn lấy được bút chì màu đỏ}
{Chắc chắn lấy được tẩy}
\loigiai{
Trong hộp không có bút chì đỏ nên không thể lấy được.
\textbf{Đáp án của câu hỏi là B}
}
\end{ex}

\begin{ex} %[2D13V66]
Gieo một đồng xu có hai mặt Sấp và Ngửa. Kết quả xuất hiện mặt Ngửa là sự kiện
\choice
{Chắc chắn xảy ra}
{Không thể xảy ra}
{Luôn luôn sai}
{\True Có thể xảy ra}
\loigiai{
Đồng xu có thể rơi vào mặt sấp hoặc mặt ngửa.
\textbf{Đáp án của câu hỏi là D}
}
\end{ex}

\begin{ex} %[2D13V66]
Mặt trời mọc ở hướng Đông vào mỗi buổi sáng là sự kiện
\choice
{\True Chắc chắn xảy ra}
{Có thể xảy ra}
{Không thể xảy ra}
{Tùy từng ngày}
\loigiai{
Quy luật tự nhiên: Mặt trời luôn mọc ở hướng Đông nên chắc chắn xảy ra.
\textbf{Đáp án của câu hỏi là A}
}
\end{ex}

\begin{ex} %[2D13V66]
Một tuần có 8 ngày là sự kiện
\choice
{Có thể xảy ra}
{Chắc chắn xảy ra}
{\True Không thể xảy ra}
{Xảy ra vào năm nhuận}
\loigiai{
Một tuần luôn chỉ có đúng 7 ngày, do đó một tuần có 8 ngày là không thể.
\textbf{Đáp án của câu hỏi là C}
}
\end{ex}

\begin{ex} %[2D13V66]
Trong lọ có 4 viên kẹo cam và 4 viên kẹo dâu. Mai lấy ra 5 viên kẹo. Khẳng định nào sau đây là 'Chắc chắn'?
\choice
{Cả 5 viên kẹo đều là kẹo cam}
{\True Trong 5 viên kẹo lấy ra có ít nhất 1 viên kẹo cam}
{Cả 5 viên kẹo đều là kẹo dâu}
{Không có viên kẹo dâu nào}
\loigiai{
Vì chỉ có 4 viên dâu nên lấy 5 viên chắc chắn phải có ít nhất 1 viên kẹo cam.
\textbf{Đáp án của câu hỏi là B}
}
\end{ex}

\begin{ex} %[2D13V66]
Một em bé mới sinh biết đi xe đạp ngay lập tức là sự kiện
\choice
{Có thể xảy ra}
{Chắc chắn xảy ra}
{Thỉnh thoảng xảy ra}
{\True Không thể xảy ra}
\loigiai{
Em bé mới sinh không thể biết đi xe đạp.
\textbf{Đáp án của câu hỏi là D}
}
\end{ex}

\begin{ex} %[2D13V66]
Khi quay một chiếc kim trên vòng tròn chia làm hai nửa: nửa đỏ và nửa xanh. Kim dừng lại ở phần màu đỏ là sự kiện
\choice
{\True Có thể xảy ra}
{Chắc chắn xảy ra}
{Không thể xảy ra}
{Không bao giờ xảy ra}
\loigiai{
Kim có thể dừng ở màu đỏ hoặc màu xanh.
\textbf{Đáp án của câu hỏi là A}
}
\end{ex}

\begin{ex} %[2D13V66]
Tổng số chấm khi gieo hai con xúc xắc lớn hơn 1 là sự kiện
\choice
{Có thể xảy ra}
{Không thể xảy ra}
{\True Chắc chắn xảy ra}
{Tùy lúc}
\loigiai{
Mỗi con xúc xắc ít nhất 1 chấm, hai con ít nhất $1 + 1 = 2 > 1$, nên chắc chắn lớn hơn 1.
\textbf{Đáp án của câu hỏi là C}
}
\end{ex}

\textbf{Phần 2. Trắc nghiệm đúng sai}
\setcounterpageref{ex}{0}

\begin{ex} %[2D13V66]
Xét tính đúng sai của các nhận định về khả năng xảy ra của sự kiện:
\choiceTF
{\True Một sự kiện chắc chắn xảy ra nghĩa là nó luôn luôn diễn ra trong mọi trường hợp.}
{Một sự kiện có thể xảy ra nghĩa là nó không bao giờ xảy ra được.}
{\True Sự kiện con gà đẻ ra quả trứng vịt là sự kiện không thể xảy ra.}
{\True Lấy 1 viên kẹo từ túi có 10 viên kẹo sô-cô-la chắc chắn được viên kẹo sô-cô-la.}
\loigiai{
\begin{itemize}
    \item Ý a) Đúng: Đúng định nghĩa.
    \item Ý b) Sai: Có thể nghĩa là có khả năng xảy ra. Sai.
    \item Ý c) Đúng: Đúng thực tế sinh học.
    \item Ý d) Đúng: Đúng vì túi chỉ có kẹo sô-cô-la.
\end{itemize}
\begin{itemize}
    \item a) \textbf{Đúng}.
    \item b) \textbf{Sai}.
    \item c) \textbf{Đúng}.
    \item d) \textbf{Đúng}.
\end{itemize}
}
\end{ex}

\begin{ex} %[2D13V66]
Trong một hộp có 3 quả bóng xanh và 3 quả bóng đỏ:
\choiceTF
{\True Lấy ra 1 quả bóng thì có thể được bóng xanh.}
{\True Lấy ra 1 quả bóng thì không thể được bóng vàng.}
{Lấy ra 2 quả bóng thì chắc chắn cả 2 quả đều là bóng xanh.}
{\True Lấy ra 4 quả bóng thì chắc chắn có ít nhất 1 quả bóng màu xanh.}
\loigiai{
\begin{itemize}
    \item Ý a) Đúng: Đúng.
    \item Ý b) Đúng: Trong hộp không có bóng vàng. Đúng.
    \item Ý c) Sai: Có thể lấy được 1 xanh 1 đỏ hoặc 2 đỏ. Sai.
    \item Ý d) Đúng: Vì chỉ có 3 bóng đỏ nên 4 quả chắc chắn có ít nhất 1 xanh. Đúng.
\end{itemize}
\begin{itemize}
    \item a) \textbf{Đúng}.
    \item b) \textbf{Đúng}.
    \item c) \textbf{Sai}.
    \item d) \textbf{Đúng}.
\end{itemize}
}
\end{ex}

\textbf{Phần 3. Trả lời ngắn}
\setcounterpageref{ex}{0}

\begin{ex} %[2D13V66]
Trong hộp có 5 viên bi đỏ và 1 viên bi xanh. Hỏi cần lấy ra ít nhất bao nhiêu viên bi để chắc chắn có 1 viên bi màu đỏ?
\shortans{2}
\loigiai{
\textbf{Đáp số:} $2$\\
Trường hợp xấu nhất lấy phải 1 viên bi xanh, lấy thêm 1 viên nữa chắc chắn là bi đỏ. Cần lấy $1 + 1 = 2$ viên.
}
\end{ex}

\begin{ex} %[2D13V66]
Một con xúc xắc có 6 mặt từ 1 đến 6 chấm. Có bao nhiêu mặt có số chấm là số chẵn?
\shortans{3}
\loigiai{
\textbf{Đáp số:} $3$\\
Các mặt chẵn là 2, 4, 6. Có 3 mặt.
}
\end{ex}

\begin{ex} %[2D13V66]
Trong túi có 10 chiếc kẹo gồm hai loại: kẹo cam và kẹo dâu. Biết chắc chắn lấy được ít nhất 1 chiếc kẹo dâu khi lấy 6 chiếc kẹo bất kì. Có bao nhiêu chiếc kẹo cam trong túi?
\shortans{5}
\loigiai{
\textbf{Đáp số:} $5$\\
Để chắc chắn có kẹo dâu khi lấy 6 chiếc thì số kẹo cam nhiều nhất là 5 chiếc.
}
\end{ex}

\begin{ex} %[2D13V66]
Gieo 1 con xúc xắc, có mấy khả năng xuất hiện mặt có số chấm lớn hơn 4?
\shortans{2}
\loigiai{
\textbf{Đáp số:} $2$\\
Mặt 5 và mặt 6, có 2 khả năng.
}
\end{ex}

\begin{ex} %[2D13V66]
Trong hộp có 4 chiếc bút chì vàng. An thò tay lấy 2 chiếc. Có bao nhiêu chiếc bút chì màu khác vàng được lấy ra?
\shortans{0}
\loigiai{
\textbf{Đáp số:} $0$\\
Trong hộp chỉ có bút vàng nên lấy được 0 chiếc màu khác.
}
\end{ex}

\begin{ex} %[2D13V66]
Một tuần lễ chắc chắn có mấy ngày?
\shortans{7}
\loigiai{
\textbf{Đáp số:} $7$\\
Một tuần lễ chắc chắn có đúng 7 ngày.
}
\end{ex}

\newpage

\TaoPhieuKiemTra{Chủ đề 13 -- Làm quen với yếu tố thống kê, xác suất}{67}{Thực hành và trải nghiệm thu thập số liệu}{https://loptoanthaylinh.vn/thi/chuyen-de/2/toan2_cd_b67}
\textbf{Phần 1. Trắc nghiệm lựa chọn}
\setcounterpageref{ex}{0}

\begin{ex} %[2D13V67]
Trong giờ thực hành, bạn Nam ghi lại thời tiết trong 10 ngày: có 6 ngày nắng, 3 ngày mưa và 1 ngày râm mát. Ngày có thời tiết xuất hiện nhiều nhất là
\choice
{\True Ngày nắng}
{Ngày mưa}
{Ngày râm mát}
{Cả nắng và mưa bằng nhau}
\loigiai{
Ngày nắng xuất hiện 6 ngày, nhiều nhất.
\textbf{Đáp án của câu hỏi là A}
}
\end{ex}

\begin{ex} %[2D13V67]
Khảo sát loại quả yêu thích của 30 học sinh: 12 bạn thích xoài, 10 bạn thích cam, 8 bạn thích dưa hấu. Số bạn thích dưa hấu ít hơn số bạn thích xoài là
\choice
{$2$ bạn}
{$6$ bạn}
{\True $4$ bạn}
{$5$ bạn}
\loigiai{
$12 - 8 = 4$ (bạn).
\textbf{Đáp án của câu hỏi là C}
}
\end{ex}

\begin{ex} %[2D13V67]
Để đếm số chim bồ câu đậu trên mái nhà, cách kiểm đếm nào sau đây là nhanh và không bị nhầm lẫn?
\choice
{Đếm nhảy cóc tùy ý}
{\True Đếm lần lượt từng con từ trái sang phải rồi gạch ghi lại}
{Ước lượng đại khái không cần đếm}
{Nhắm mắt đoán số lượng}
\loigiai{
Đếm có trật tự từ trái sang phải và ghi lại là cách chuẩn khoa học.
\textbf{Đáp án của câu hỏi là B}
}
\end{ex}

\begin{ex} %[2D13V67]
Khi tung đồng xu 10 lần, bạn Mai ghi nhận được 6 lần xuất hiện mặt Sấp. Hỏi có bao nhiêu lần xuất hiện mặt Ngửa?
\choice
{$5$ lần}
{$3$ lần}
{$6$ lần}
{\True $4$ lần}
\loigiai{
$10 - 6 = 4$ (lần).
\textbf{Đáp án của câu hỏi là D}
}
\end{ex}

\begin{ex} %[2D13V67]
Bảng thống kê số cây mỗi bạn trồng được: An trồng 3 cây, Bình trồng 5 cây, Chi trồng 4 cây, Dũng trồng 2 cây. Bạn trồng được nhiều cây nhất là
\choice
{\True Bình}
{An}
{Chi}
{Dũng}
\loigiai{
Bình trồng 5 cây, nhiều nhất.
\textbf{Đáp án của câu hỏi là A}
}
\end{ex}

\begin{ex} %[2D13V67]
Từ bảng số liệu trên, tổng số cây cả bốn bạn An, Bình, Chi, Dũng trồng được là
\choice
{$12$ cây}
{$15$ cây}
{\True $14$ cây}
{$16$ cây}
\loigiai{
$3 + 5 + 4 + 2 = 14$ (cây).
\textbf{Đáp án của câu hỏi là C}
}
\end{ex}

\begin{ex} %[2D13V67]
Thực hành cân đo sức khỏe: Bạn Hùng nặng $28\text{ kg}$, bạn Tuấn nặng $25\text{ kg}$, bạn Dũng nặng $31\text{ kg}$. Bạn nặng cân nhất là
\choice
{Bạn Hùng}
{\True Bạn Dũng}
{Bạn Tuấn}
{Hai bạn nặng bằng nhau}
\loigiai{
Bạn Dũng nặng 31 kg, lớn nhất.
\textbf{Đáp án của câu hỏi là B}
}
\end{ex}

\begin{ex} %[2D13V67]
Trong hộp có 10 tấm thẻ ghi các số từ 1 đến 10. Rút ngẫu nhiên một tấm thẻ. Sự kiện rút được tấm thẻ ghi số có hai chữ số là
\choice
{Chắc chắn xảy ra}
{Không thể xảy ra}
{Luôn luôn xảy ra}
{\True Có thể xảy ra}
\loigiai{
Trong các số từ 1 đến 10 có số 10 là số có hai chữ số nên có thể xảy ra.
\textbf{Đáp án của câu hỏi là D}
}
\end{ex}

\begin{ex} %[2D13V67]
Bạn Nam gieo con xúc xắc 20 lần và ghi lại số lần xuất hiện mặt 6 chấm được 3 lần. Số lần xuất hiện các mặt khác là
\choice
{\True $17$ lần}
{$16$ lần}
{$18$ lần}
{$15$ lần}
\loigiai{
$20 - 3 = 17$ (lần).
\textbf{Đáp án của câu hỏi là A}
}
\end{ex}

\begin{ex} %[2D13V67]
Theo dõi số học sinh đi muộn trong tuần của lớp: Thứ Hai có 2 bạn, thứ Ba có 0 bạn, thứ Tư có 1 bạn, thứ Năm có 0 bạn, thứ Sáu có 3 bạn. Có bao nhiêu ngày trong tuần không có bạn nào đi muộn?
\choice
{$1$ ngày}
{$3$ ngày}
{\True $2$ ngày}
{$0$ ngày}
\loigiai{
Đó là thứ Ba và thứ Năm (2 ngày).
\textbf{Đáp án của câu hỏi là C}
}
\end{ex}

\begin{ex} %[2D13V67]
Tổng số lượt học sinh đi muộn trong cả tuần trên là
\choice
{$5$ lượt}
{\True $6$ lượt}
{$7$ lượt}
{$4$ lượt}
\loigiai{
$2 + 0 + 1 + 0 + 3 = 6$ (lượt).
\textbf{Đáp án của câu hỏi là B}
}
\end{ex}

\begin{ex} %[2D13V67]
Khi thu thập thông tin về chiều cao của các bạn trong tổ, dụng cụ đo thích hợp nhất là
\choice
{Cân đồng hồ}
{Đồng hồ bấm giây}
{Nhiệt kế y tế}
{\True Thước đo chiều cao (thước dây hoặc thước cuộn)}
\loigiai{
Đo chiều cao phải dùng thước dây hoặc thước cuộn.
\textbf{Đáp án của câu hỏi là D}
}
\end{ex}

\textbf{Phần 2. Trắc nghiệm đúng sai}
\setcounterpageref{ex}{0}

\begin{ex} %[2D13V67]
Xét tính đúng sai trong thực hành thu thập và xử lý số liệu:
\choiceTF
{\True Ghi chép số liệu cần trung thực, chính xác và rõ ràng.}
{Chỉ cần phán đoán theo cảm tính mà không cần đếm hay đo thực tế.}
{\True Bảng thống kê giúp ta dễ dàng so sánh các đại lượng và rút ra kết luận.}
{\True Sau khi thu thập, việc kiểm đếm lại giúp tránh nhầm lẫn.}
\loigiai{
\begin{itemize}
    \item Ý a) Đúng: Nguyên tắc khoa học. Đúng.
    \item Ý b) Sai: Khoa học thực nghiệm đòi hỏi đo đếm thực tế. Sai.
    \item Ý c) Đúng: Đúng vai trò của thống kê.
    \item Ý d) Đúng: Đúng.
\end{itemize}
\begin{itemize}
    \item a) \textbf{Đúng}.
    \item b) \textbf{Sai}.
    \item c) \textbf{Đúng}.
    \item d) \textbf{Đúng}.
\end{itemize}
}
\end{ex}

\begin{ex} %[2D13V67]
Khảo sát màu sắc yêu thích của 20 bạn học sinh: Màu xanh có 8 bạn, màu đỏ có 7 bạn, màu vàng có 5 bạn:
\choiceTF
{\True Màu xanh được nhiều bạn yêu thích nhất.}
{\True Số bạn thích màu vàng ít hơn màu đỏ là 2 bạn.}
{Tổng số bạn tham gia khảo sát là 25 bạn.}
{\True Số bạn thích màu xanh hoặc màu vàng là 13 bạn.}
\loigiai{
\begin{itemize}
    \item Ý a) Đúng: 8 > 7 > 5. Đúng.
    \item Ý b) Đúng: 7 - 5 = 2 bạn. Đúng.
    \item Ý c) Sai: 8 + 7 + 5 = 20 bạn, không phải 25. Sai.
    \item Ý d) Đúng: 8 + 5 = 13 bạn. Đúng.
\end{itemize}
\begin{itemize}
    \item a) \textbf{Đúng}.
    \item b) \textbf{Đúng}.
    \item c) \textbf{Sai}.
    \item d) \textbf{Đúng}.
\end{itemize}
}
\end{ex}

\textbf{Phần 3. Trả lời ngắn}
\setcounterpageref{ex}{0}

\begin{ex} %[2D13V67]
Khảo sát sở thích đọc sách của 40 học sinh: 18 bạn thích truyện cổ tích, 14 bạn thích truyện tranh, còn lại thích sách khoa học. Có bao nhiêu bạn thích sách khoa học?
\shortans{8}
\loigiai{
\textbf{Đáp số:} $8$\\
$40 - 18 - 14 = 8$ bạn.
}
\end{ex}

\begin{ex} %[2D13V67]
Tung đồng xu 15 lần, ghi nhận được 8 lần mặt Ngửa. Hỏi có bao nhiêu lần mặt Sấp xuất hiện?
\shortans{7}
\loigiai{
\textbf{Đáp số:} $7$\\
$15 - 8 = 7$ lần.
}
\end{ex}

\begin{ex} %[2D13V67]
Đếm số ô cửa sổ của trường học: tầng 1 có 12 ô, tầng 2 có 14 ô, tầng 3 có 14 ô. Có tất cả bao nhiêu ô cửa sổ?
\shortans{40}
\loigiai{
\textbf{Đáp số:} $40$\\
$12 + 14 + 14 = 40$ ô.
}
\end{ex}

\begin{ex} %[2D13V67]
Trong đợt thi đua, tổ Một đạt 28 bông hoa điểm tốt, tổ Hai đạt nhiều hơn tổ Một 6 bông hoa. Cả hai tổ đạt bao nhiêu bông hoa điểm tốt?
\shortans{62}
\loigiai{
\textbf{Đáp số:} $62$\\
Tổ 2: $28 + 6 = 34$. Cả hai tổ: $28 + 34 = 62$ bông hoa.
}
\end{ex}

\begin{ex} %[2D13V67]
Một nhóm gồm 4 bạn, mỗi bạn đếm được 5 chiếc xe buýt chạy qua. Nhóm đã đếm được tổng cộng bao nhiêu chiếc xe buýt?
\shortans{20}
\loigiai{
\textbf{Đáp số:} $20$\\
$4 \times 5 = 20$ chiếc xe buýt.
}
\end{ex}

\begin{ex} %[2D13V67]
Trong túi có 6 viên bi đỏ và 4 viên bi trắng. Cần lấy ít nhất mấy viên để chắc chắn có 1 viên bi đỏ?
\shortans{5}
\loigiai{
\textbf{Đáp số:} $5$\\
Trường hợp xấu nhất lấy hết 4 bi trắng, lấy thêm 1 viên nữa: $4 + 1 = 5$ viên.
}
\end{ex}

\newpage


% =========================================================================
% CHỦ ĐỀ 14 - ÔN TẬP CUỐI NĂM
% =========================================================================
\section{Ôn tập cuối năm}

\TaoPhieuKiemTra{Chủ đề 14 -- Ôn tập cuối năm}{68}{Ôn tập các số trong phạm vi 1 000}{https://loptoanthaylinh.vn/thi/chuyen-de/2/toan2_cd_b68}
\textbf{Phần 1. Trắc nghiệm lựa chọn}
\setcounterpageref{ex}{0}

\begin{ex} %[2D14V68]
Số liền trước của số lớn nhất có ba chữ số là
\choice
{$999$}
{$1 000$}
{\True $998$}
{$997$}
\loigiai{
Số lớn nhất có ba chữ số là 999. Số liền trước là $999 - 1 = 998$.
\textbf{Đáp án của câu hỏi là C}
}
\end{ex}

\begin{ex} %[2D14V68]
Số gồm $7$ trăm, $0$ chục và $4$ đơn vị viết là
\choice
{\True $704$}
{$740$}
{$74$}
{$407$}
\loigiai{
Viết là $704$.
\textbf{Đáp án của câu hỏi là A}
}
\end{ex}

\begin{ex} %[2D14V68]
Chữ số $5$ trong số $358$ có giá trị bằng bao nhiêu?
\choice
{$500$}
{$5$}
{$15$}
{\True $50$}
\loigiai{
Chữ số 5 nằm ở hàng chục nên có giá trị là 50.
\textbf{Đáp án của câu hỏi là D}
}
\end{ex}

\begin{ex} %[2D14V68]
Từ các chữ số $3, 0, 8$, số chẵn lớn nhất có ba chữ số khác nhau lập được là
\choice
{$803$}
{\True $830$}
{$380$}
{$308$}
\loigiai{
Số lớn nhất có chữ số hàng trăm lớn nhất (8), chục (3), đơn vị chẵn (0): số 830.
\textbf{Đáp án của câu hỏi là B}
}
\end{ex}

\begin{ex} %[2D14V68]
Số tròn trăm liền sau của số $450$ là
\choice
{$400$}
{$460$}
{\True $500$}
{$600$}
\loigiai{
Số tròn trăm lớn hơn và liền kề sau 450 là 500.
\textbf{Đáp án của câu hỏi là C}
}
\end{ex}

\begin{ex} %[2D14V68]
Cho $x$ là số có ba chữ số sao cho $789 < x < 792$. Có bao nhiêu giá trị của $x$?
\choice
{\True $2$ số}
{$3$ số}
{$1$ số}
{$4$ số}
\loigiai{
$x \in \{790, 791\}$, có đúng 2 số.
\textbf{Đáp án của câu hỏi là A}
}
\end{ex}

\begin{ex} %[2D14V68]
Tổng của chữ số hàng trăm và chữ số hàng chục của số $625$ là
\choice
{$7$}
{$11$}
{$13$}
{\True $8$}
\loigiai{
$6 + 2 = 8$.
\textbf{Đáp án của câu hỏi là D}
}
\end{ex}

\begin{ex} %[2D14V68]
Điền dấu thích hợp: $500 + 40 + 9 \dots 500 + 50$.
\choice
{$>$}
{\True $<$}
{$=$}
{Không so sánh được}
\loigiai{
$500 + 40 + 9 = 549 < 550 = 500 + 50$.
\textbf{Đáp án của câu hỏi là B}
}
\end{ex}

\begin{ex} %[2D14V68]
Tìm số tự nhiên có ba chữ số, biết chữ số hàng trăm là số lớn nhất có một chữ số, hàng chục kém hàng trăm 3 đơn vị và hàng đơn vị bằng 0.
\choice
{$930$}
{$850$}
{\True $960$}
{$970$}
\loigiai{
Hàng trăm là 9, hàng chục là $9 - 3 = 6$, hàng đơn vị là 0: số 960.
\textbf{Đáp án của câu hỏi là C}
}
\end{ex}

\begin{ex} %[2D14V68]
Có bao nhiêu số tròn chục có ba chữ số có chữ số hàng trăm bằng $2$?
\choice
{\True $10$ số}
{$9$ số}
{$8$ số}
{$11$ số}
\loigiai{
Các số là: 200, 210, 220, ..., 290. Có tất cả 10 số.
\textbf{Đáp án của câu hỏi là A}
}
\end{ex}

\begin{ex} %[2D14V68]
Khi đổi chỗ chữ số hàng trăm và chữ số hàng đơn vị của số $437$, ta được số mới là
\choice
{$743$}
{$374$}
{$473$}
{\True $734$}
\loigiai{
Đổi vị trí 4 và 7 ta được 734.
\textbf{Đáp án của câu hỏi là D}
}
\end{ex}

\begin{ex} %[2D14V68]
Hiệu giữa số bé nhất có bốn chữ số và số lớn nhất có hai chữ số là
\choice
{$900$}
{\True $901$}
{$910$}
{$899$}
\loigiai{
$1 000 - 99 = 901$.
\textbf{Đáp án của câu hỏi là B}
}
\end{ex}

\textbf{Phần 2. Trắc nghiệm đúng sai}
\setcounterpageref{ex}{0}

\begin{ex} %[2D14V68]
Xét các khẳng định về số học phạm vi 1 000:
\choiceTF
{\True Số 1 000 gồm 10 trăm hoặc 100 chục.}
{Số 305 đọc là ba trăm năm mươi.}
{\True Dãy số 120, 140, 160, 180 là dãy các số chẵn tăng dần.}
{\True Tổng các chữ số của số 456 bằng 15.}
\loigiai{
\begin{itemize}
    \item Ý a) Đúng: Đúng.
    \item Ý b) Sai: 305 đọc là ba trăm linh năm.
    \item Ý c) Đúng: Đúng.
    \item Ý d) Đúng: 4 + 5 + 6 = 15. Đúng.
\end{itemize}
\begin{itemize}
    \item a) \textbf{Đúng}.
    \item b) \textbf{Sai}.
    \item c) \textbf{Đúng}.
    \item d) \textbf{Đúng}.
\end{itemize}
}
\end{ex}

\begin{ex} %[2D14V68]
Cho ba số tự nhiên A = 345, B = 354, C = 435:
\choiceTF
{\True A < B < C.}
{\True Số lớn nhất trong ba số là C.}
{Hiệu giữa B và A bằng 11.}
{\True Cả ba số đều có tổng các chữ số bằng 12.}
\loigiai{
\begin{itemize}
    \item Ý a) Đúng: 345 < 354 < 435. Đúng.
    \item Ý b) Đúng: 435 là lớn nhất. Đúng.
    \item Ý c) Sai: 354 - 345 = 9, không phải 11. Sai.
    \item Ý d) Đúng: 3 + 4 + 5 = 12. Đúng.
\end{itemize}
\begin{itemize}
    \item a) \textbf{Đúng}.
    \item b) \textbf{Đúng}.
    \item c) \textbf{Sai}.
    \item d) \textbf{Đúng}.
\end{itemize}
}
\end{ex}

\textbf{Phần 3. Trả lời ngắn}
\setcounterpageref{ex}{0}

\begin{ex} %[2D14V68]
Tìm số tự nhiên có ba chữ số khác nhau bé nhất.
\shortans{102}
\loigiai{
\textbf{Đáp số:} $102$\\
Hàng trăm chọn 1, chục chọn 0, đơn vị chọn 2: số 102.
}
\end{ex}

\begin{ex} %[2D14V68]
Tìm số liền trước của số tròn trăm bé nhất có ba chữ số.
\shortans{99}
\loigiai{
\textbf{Đáp số:} $99$\\
Số tròn trăm bé nhất có 3 chữ số là 100. Liền trước là 99.
}
\end{ex}

\begin{ex} %[2D14V68]
Có bao nhiêu số có ba chữ số có tận cùng là chữ số 5 và chữ số hàng trăm là 8?
\shortans{10}
\loigiai{
\textbf{Đáp số:} $10$\\
Dạng $8a5$ với $a \in \{0, 1, \dots, 9\}$. Có 10 số.
}
\end{ex}

\begin{ex} %[2D14V68]
Tìm một số biết rằng lấy số đó cộng với 250 thì được số tròn trăm lớn nhất có ba chữ số.
\shortans{650}
\loigiai{
\textbf{Đáp số:} $650$\\
Số tròn trăm lớn nhất có 3 chữ số là 900. $900 - 250 = 650$.
}
\end{ex}

\begin{ex} %[2D14V68]
Viết số gồm 5 trăm, 8 chục và 3 đơn vị.
\shortans{583}
\loigiai{
\textbf{Đáp số:} $583$\\
Viết là 583.
}
\end{ex}

\begin{ex} %[2D14V68]
Tìm hiệu của số lớn nhất có ba chữ số và số bé nhất có ba chữ số.
\shortans{899}
\loigiai{
\textbf{Đáp số:} $899$\\
$999 - 100 = 899$.
}
\end{ex}

\newpage

\TaoPhieuKiemTra{Chủ đề 14 -- Ôn tập cuối năm}{69}{Ôn tập phép cộng, phép trừ trong phạm vi 100}{https://loptoanthaylinh.vn/thi/chuyen-de/2/toan2_cd_b69}
\textbf{Phần 1. Trắc nghiệm lựa chọn}
\setcounterpageref{ex}{0}

\begin{ex} %[2D14V69]
Tính: $48 + 37$.
\choice
{$75$}
{\True $85$}
{$86$}
{$95$}
\loigiai{
$48 + 37 = 85$.
\textbf{Đáp án của câu hỏi là B}
}
\end{ex}

\begin{ex} %[2D14V69]
Tính: $92 - 46$.
\choice
{$56$}
{$48$}
{$36$}
{\True $46$}
\loigiai{
$92 - 46 = 46$.
\textbf{Đáp án của câu hỏi là D}
}
\end{ex}

\begin{ex} %[2D14V69]
Tìm $x$ biết: $x + 29 = 75$.
\choice
{\True $46$}
{$56$}
{$44$}
{$104$}
\loigiai{
$x = 75 - 29 = 46$.
\textbf{Đáp án của câu hỏi là A}
}
\end{ex}

\begin{ex} %[2D14V69]
Một lớp học có $35$ bạn học sinh, trong đó có $18$ bạn nữ. Lớp học có bao nhiêu bạn nam?
\choice
{$16$ bạn}
{$19$ bạn}
{\True $17$ bạn}
{$27$ bạn}
\loigiai{
$35 - 18 = 17$ (bạn).
\textbf{Đáp án của câu hỏi là C}
}
\end{ex}

\begin{ex} %[2D14V69]
Tổng của hai số là $80$, số hạng thứ nhất là $34$. Số hạng thứ hai là
\choice
{$56$}
{\True $46$}
{$44$}
{$36$}
\loigiai{
$80 - 34 = 46$.
\textbf{Đáp án của câu hỏi là B}
}
\end{ex}

\begin{ex} %[2D14V69]
Mai hái được $28$ bông hoa, Lan hái được nhiều hơn Mai $14$ bông hoa. Cả hai bạn hái được bao nhiêu bông hoa?
\choice
{$42$ bông}
{$68$ bông}
{$72$ bông}
{\True $70$ bông}
\loigiai{
Lan hái: $28 + 14 = 42$ bông. Cả hai bạn: $28 + 42 = 70$ bông hoa.
\textbf{Đáp án của câu hỏi là D}
}
\end{ex}

\begin{ex} %[2D14V69]
Hiệu của hai số là $25$. Nếu số bị trừ tăng $7$ đơn vị và số trừ giữ nguyên thì hiệu mới là
\choice
{\True $32$}
{$18$}
{$30$}
{$35$}
\loigiai{
Số bị trừ tăng thì hiệu tăng: $25 + 7 = 32$.
\textbf{Đáp án của câu hỏi là A}
}
\end{ex}

\begin{ex} %[2D14V69]
Giá trị của biểu thức $100 - 45 - 25$ là
\choice
{$40$}
{$20$}
{\True $30$}
{$35$}
\loigiai{
$100 - (45 + 25) = 100 - 70 = 30$.
\textbf{Đáp án của câu hỏi là C}
}
\end{ex}

\begin{ex} %[2D14V69]
Điền dấu thích hợp: $56 + 28 \dots 94 - 10$.
\choice
{$>$}
{\True $=$}
{$<$}
{Không so sánh được}
\loigiai{
$56 + 28 = 84$; $94 - 10 = 84$. Hai vế bằng nhau.
\textbf{Đáp án của câu hỏi là B}
}
\end{ex}

\begin{ex} %[2D14V69]
Bác Ba có $62$ con gà, bác bán đi $25$ con gà. Sau đó bác mua thêm $18$ con gà nữa. Bác Ba có tất cả bao nhiêu con gà?
\choice
{$45$ con}
{$65$ con}
{$57$ con}
{\True $55$ con}
\loigiai{
$62 - 25 + 18 = 37 + 18 = 55$ (con).
\textbf{Đáp án của câu hỏi là D}
}
\end{ex}

\begin{ex} %[2D14V69]
Tìm số tự nhiên có hai chữ số biết chữ số hàng chục hơn hàng đơn vị 4 đơn vị và tổng hai chữ số bằng 12.
\choice
{\True $84$}
{$73$}
{$95$}
{$62$}
\loigiai{
Hàng chục: $(12 + 4) : 2 = 8$. Hàng đơn vị: $8 - 4 = 4$. Số đó là 84.
\textbf{Đáp án của câu hỏi là A}
}
\end{ex}

\begin{ex} %[2D14V69]
Một cuộn dây dài $80\text{ m}$. Người ta cắt đi $35\text{ m}$ lần đầu và $20\text{ m}$ lần sau. Cuộn dây còn lại bao nhiêu mét?
\choice
{$35\text{ m}$}
{$15\text{ m}$}
{\True $25\text{ m}$}
{$30\text{ m}$}
\loigiai{
$80 - 35 - 20 = 25$ (m).
\textbf{Đáp án của câu hỏi là C}
}
\end{ex}

\textbf{Phần 2. Trắc nghiệm đúng sai}
\setcounterpageref{ex}{0}

\begin{ex} %[2D14V69]
Xét tính đúng sai của các phép toán cộng trừ trong phạm vi 100:
\choiceTF
{\True 37 + 63 = 100.}
{81 - 29 = 62.}
{\True Trong phép trừ có nhớ phạm vi 100, ta mượn 1 chục ở hàng chục của số bị trừ.}
{\True Nếu a + b = 90 thì b = 90 - a.}
\loigiai{
\begin{itemize}
    \item Ý a) Đúng: Đúng, tổng là 100.
    \item Ý b) Sai: 81 - 29 = 52, không phải 62.
    \item Ý c) Đúng: Đúng quy tắc.
    \item Ý d) Đúng: Đúng tính chất tìm số hạng.
\end{itemize}
\begin{itemize}
    \item a) \textbf{Đúng}.
    \item b) \textbf{Sai}.
    \item c) \textbf{Đúng}.
    \item d) \textbf{Đúng}.
\end{itemize}
}
\end{ex}

\begin{ex} %[2D14V69]
Một cửa hàng có 75 kg gạo, buổi sáng bán 28 kg, buổi chiều bán 32 kg:
\choiceTF
{\True Cả hai buổi bán được 60 kg gạo.}
{\True Cửa hàng còn lại 15 kg gạo.}
{Số gạo bán buổi chiều ít hơn buổi sáng.}
{\True Buổi chiều bán nhiều hơn buổi sáng 4 kg.}
\loigiai{
\begin{itemize}
    \item Ý a) Đúng: 28 + 32 = 60 kg. Đúng.
    \item Ý b) Đúng: 75 - 60 = 15 kg. Đúng.
    \item Ý c) Sai: 32 > 28 nên bán nhiều hơn. Sai.
    \item Ý d) Đúng: 32 - 28 = 4 kg. Đúng.
\end{itemize}
\begin{itemize}
    \item a) \textbf{Đúng}.
    \item b) \textbf{Đúng}.
    \item c) \textbf{Sai}.
    \item d) \textbf{Đúng}.
\end{itemize}
}
\end{ex}

\textbf{Phần 3. Trả lời ngắn}
\setcounterpageref{ex}{0}

\begin{ex} %[2D14V69]
Tính: $64 + 28 - 45$.
\shortans{47}
\loigiai{
\textbf{Đáp số:} $47$\\
$92 - 45 = 47$.
}
\end{ex}

\begin{ex} %[2D14V69]
Tìm $x$ biết $100 - x = 38$.
\shortans{62}
\loigiai{
\textbf{Đáp số:} $62$\\
$x = 100 - 38 = 62$.
}
\end{ex}

\begin{ex} %[2D14V69]
Nam có 45 viên bi, Việt có ít hơn Nam 12 viên bi. Cả hai bạn có bao nhiêu viên bi?
\shortans{78}
\loigiai{
\textbf{Đáp số:} $78$\\
Việt có: $45 - 12 = 33$. Cả hai bạn: $45 + 33 = 78$ viên.
}
\end{ex}

\begin{ex} %[2D14V69]
Tổng của hai số là 72, một số là 39. Tìm số còn lại.
\shortans{33}
\loigiai{
\textbf{Đáp số:} $33$\\
$72 - 39 = 33$.
}
\end{ex}

\begin{ex} %[2D14V69]
Một bến xe có 54 xe khách, có 18 xe rời bến và 9 xe vào bến. Bến xe có bao nhiêu xe?
\shortans{45}
\loigiai{
\textbf{Đáp số:} $45$\\
$54 - 18 + 9 = 45$ xe.
}
\end{ex}

\begin{ex} %[2D14V69]
Tính hiệu của số chẵn lớn nhất có hai chữ số và số 45.
\shortans{53}
\loigiai{
\textbf{Đáp số:} $53$\\
$98 - 45 = 53$.
}
\end{ex}

\newpage

\TaoPhieuKiemTra{Chủ đề 14 -- Ôn tập cuối năm}{70}{Ôn tập phép cộng, phép trừ trong phạm vi 1 000}{https://loptoanthaylinh.vn/thi/chuyen-de/2/toan2_cd_b70}
\textbf{Phần 1. Trắc nghiệm lựa chọn}
\setcounterpageref{ex}{0}

\begin{ex} %[2D14V70]
Kết quả của phép tính $458 + 364$ là
\choice
{\True $822$}
{$812$}
{$832$}
{$722$}
\loigiai{
$458 + 364 = 822$.
\textbf{Đáp án của câu hỏi là A}
}
\end{ex}

\begin{ex} %[2D14V70]
Kết quả của phép tính $725 - 389$ là
\choice
{$346$}
{$326$}
{\True $336$}
{$436$}
\loigiai{
$725 - 389 = 336$.
\textbf{Đáp án của câu hỏi là C}
}
\end{ex}

\begin{ex} %[2D14V70]
Tìm $x$ biết: $x + 276 = 615$.
\choice
{$349$}
{\True $339$}
{$329$}
{$891$}
\loigiai{
$x = 615 - 276 = 339$.
\textbf{Đáp án của câu hỏi là B}
}
\end{ex}

\begin{ex} %[2D14V70]
Một kho gạo có $850\text{ kg}$ gạo. Ngày đầu xuất $275\text{ kg}$, ngày thứ hai xuất $325\text{ kg}$. Kho còn lại bao nhiêu ki-lô-gam gạo?
\choice
{$240\text{ kg}$}
{$260\text{ kg}$}
{$350\text{ kg}$}
{\True $250\text{ kg}$}
\loigiai{
Đã xuất: $275 + 325 = 600$ kg. Còn lại: $850 - 600 = 250$ kg.
\textbf{Đáp án của câu hỏi là D}
}
\end{ex}

\begin{ex} %[2D14V70]
Tìm $x$ biết: $900 - x = 455$.
\choice
{\True $445$}
{$455$}
{$435$}
{$545$}
\loigiai{
$x = 900 - 455 = 445$.
\textbf{Đáp án của câu hỏi là A}
}
\end{ex}

\begin{ex} %[2D14V70]
Tính giá trị biểu thức: $340 + 260 - 150$.
\choice
{$440$}
{$460$}
{\True $450$}
{$500$}
\loigiai{
$340 + 260 = 600$; $600 - 150 = 450$.
\textbf{Đáp án của câu hỏi là C}
}
\end{ex}

\begin{ex} %[2D14V70]
Một đoạn đê dài $1\text{ km}$. Đội công nhân đã đắp gia cố được $650\text{ m}$. Còn lại bao nhiêu mét đê chưa đắp?
\choice
{$450\text{ m}$}
{\True $350\text{ m}$}
{$250\text{ m}$}
{$300\text{ m}$}
\loigiai{
$1\text{ km} = 1 000\text{ m}$. Còn lại: $1 000 - 650 = 350$ m.
\textbf{Đáp án của câu hỏi là B}
}
\end{ex}

\begin{ex} %[2D14V70]
Số hạng thứ nhất là $385$, số hạng thứ hai lớn hơn số thứ nhất $130$ đơn vị. Tổng hai số là
\choice
{$890$}
{$910$}
{$515$}
{\True $900$}
\loigiai{
Số thứ hai: $385 + 130 = 515$. Tổng: $385 + 515 = 900$.
\textbf{Đáp án của câu hỏi là D}
}
\end{ex}

\begin{ex} %[2D14V70]
Điền dấu thích hợp: $475 + 235 \dots 900 - 190$.
\choice
{\True $=$}
{$>$}
{$<$}
{Không so sánh được}
\loigiai{
$475 + 235 = 710$; $900 - 190 = 710$. Hai vế bằng nhau.
\textbf{Đáp án của câu hỏi là A}
}
\end{ex}

\begin{ex} %[2D14V70]
Một xưởng may ngày đầu may được $240$ chiếc áo, ngày thứ hai may được nhiều hơn ngày đầu $50$ chiếc. Cả hai ngày may được bao nhiêu chiếc áo?
\choice
{$490$ chiếc}
{$520$ chiếc}
{\True $530$ chiếc}
{$290$ chiếc}
\loigiai{
Ngày 2: $240 + 50 = 290$. Cả hai ngày: $240 + 290 = 530$ chiếc.
\textbf{Đáp án của câu hỏi là C}
}
\end{ex}

\begin{ex} %[2D14V70]
Tổng của ba số là $780$. Tổng của số thứ nhất và số thứ hai là $520$. Số thứ ba là
\choice
{$250$}
{\True $260$}
{$270$}
{$360$}
\loigiai{
Số thứ ba: $780 - 520 = 260$.
\textbf{Đáp án của câu hỏi là B}
}
\end{ex}

\begin{ex} %[2D14V70]
Cho phép tính $a - 350 = 420$. Giá trị của $a$ là
\choice
{$760$}
{$780$}
{$70$}
{\True $770$}
\loigiai{
$a = 420 + 350 = 770$.
\textbf{Đáp án của câu hỏi là D}
}
\end{ex}

\textbf{Phần 2. Trắc nghiệm đúng sai}
\setcounterpageref{ex}{0}

\begin{ex} %[2D14V70]
Xét tính đúng sai của các phép toán trong phạm vi 1 000:
\choiceTF
{\True 450 + 550 = 1 000.}
{1 000 - 245 = 855.}
{\True 325 + 175 = 500.}
{\True Nếu số bị trừ và số trừ cùng bớt đi 50 đơn vị thì hiệu không thay đổi.}
\loigiai{
\begin{itemize}
    \item Ý a) Đúng: Đúng.
    \item Ý b) Sai: 1 000 - 245 = 755, không phải 855. Sai.
    \item Ý c) Đúng: Đúng.
    \item Ý d) Đúng: Đúng tính chất phép trừ.
\end{itemize}
\begin{itemize}
    \item a) \textbf{Đúng}.
    \item b) \textbf{Sai}.
    \item c) \textbf{Đúng}.
    \item d) \textbf{Đúng}.
\end{itemize}
}
\end{ex}

\begin{ex} %[2D14V70]
Một trường có 450 học sinh khối sáng và 380 học sinh khối chiều:
\choiceTF
{\True Trường có tất cả 830 học sinh.}
{\True Khối sáng đông hơn khối chiều 70 học sinh.}
{Khối chiều đông hơn khối sáng.}
{\True Nếu thêm 170 học sinh nữa thì toàn trường có đúng 1 000 học sinh.}
\loigiai{
\begin{itemize}
    \item Ý a) Đúng: 450 + 380 = 830. Đúng.
    \item Ý b) Đúng: 450 - 380 = 70. Đúng.
    \item Ý c) Sai: Khối chiều ít hơn. Sai.
    \item Ý d) Đúng: 830 + 170 = 1 000. Đúng.
\end{itemize}
\begin{itemize}
    \item a) \textbf{Đúng}.
    \item b) \textbf{Đúng}.
    \item c) \textbf{Sai}.
    \item d) \textbf{Đúng}.
\end{itemize}
}
\end{ex}

\textbf{Phần 3. Trả lời ngắn}
\setcounterpageref{ex}{0}

\begin{ex} %[2D14V70]
Tính: $624 + 186 - 310$.
\shortans{500}
\loigiai{
\textbf{Đáp số:} $500$\\
$810 - 310 = 500$.
}
\end{ex}

\begin{ex} %[2D14V70]
Tìm $x$ biết: $x - 345 = 455$.
\shortans{800}
\loigiai{
\textbf{Đáp số:} $800$\\
$x = 455 + 345 = 800$.
}
\end{ex}

\begin{ex} %[2D14V70]
Một cửa hàng có 900 kg ngô. Sau khi bán hai lần thì còn lại 280 kg. Cửa hàng đã bán bao nhiêu ki-lô-gam ngô?
\shortans{620}
\loigiai{
\textbf{Đáp số:} $620$\\
$900 - 280 = 620$ kg.
}
\end{ex}

\begin{ex} %[2D14V70]
Tổng của hai số là 850, số lớn là 520. Tìm số bé.
\shortans{330}
\loigiai{
\textbf{Đáp số:} $330$\\
$850 - 520 = 330$.
}
\end{ex}

\begin{ex} %[2D14V70]
Một con tàu chở 750 tấn hàng, sau khi dỡ bớt 320 tấn thì nhận thêm 150 tấn nữa. Tàu chở bao nhiêu tấn hàng?
\shortans{580}
\loigiai{
\textbf{Đáp số:} $580$\\
$750 - 320 + 150 = 580$ tấn.
}
\end{ex}

\begin{ex} %[2D14V70]
Tính tổng: $125 + 250 + 375$.
\shortans{750}
\loigiai{
\textbf{Đáp số:} $750$\\
$125 + 250 + 375 = 750$.
}
\end{ex}

\newpage

\TaoPhieuKiemTra{Chủ đề 14 -- Ôn tập cuối năm}{71}{Ôn tập phép nhân, phép chia}{https://loptoanthaylinh.vn/thi/chuyen-de/2/toan2_cd_b71}
\textbf{Phần 1. Trắc nghiệm lựa chọn}
\setcounterpageref{ex}{0}

\begin{ex} %[2D14V71]
Tính giá trị của biểu thức: $5 \times 8 + 15$.
\choice
{$45$}
{$50$}
{$60$}
{\True $55$}
\loigiai{
$5 \times 8 = 40$; $40 + 15 = 55$.
\textbf{Đáp án của câu hỏi là D}
}
\end{ex}

\begin{ex} %[2D14V71]
Tính giá trị của biểu thức: $45 : 5 + 28$.
\choice
{$35$}
{\True $37$}
{$39$}
{$47$}
\loigiai{
$45 : 5 = 9$; $9 + 28 = 37$.
\textbf{Đáp án của câu hỏi là B}
}
\end{ex}

\begin{ex} %[2D14V71]
Tìm $x$ biết: $x \times 2 = 18$.
\choice
{$8$}
{$16$}
{\True $9$}
{$36$}
\loigiai{
$x = 18 : 2 = 9$.
\textbf{Đáp án của câu hỏi là C}
}
\end{ex}

\begin{ex} %[2D14V71]
Có $35$ chiếc kẹo chia đều cho $5$ bạn. Mỗi bạn nhận được bao nhiêu chiếc kẹo?
\choice
{\True $7$ chiếc}
{$6$ chiếc}
{$8$ chiếc}
{$5$ chiếc}
\loigiai{
$35 : 5 = 7$ (chiếc kẹo).
\textbf{Đáp án của câu hỏi là A}
}
\end{ex}

\begin{ex} %[2D14V71]
Một túi gạo cân nặng $5\text{ kg}$. Hỏi $9$ túi gạo như thế cân nặng bao nhiêu ki-lô-gam?
\choice
{$40\text{ kg}$}
{$50\text{ kg}$}
{$14\text{ kg}$}
{\True $45\text{ kg}$}
\loigiai{
$5 \times 9 = 45$ (kg).
\textbf{Đáp án của câu hỏi là D}
}
\end{ex}

\begin{ex} %[2D14V71]
Tìm $x$ biết: $x : 5 = 6$.
\choice
{$25$}
{\True $30$}
{$35$}
{$11$}
\loigiai{
$x = 6 \times 5 = 30$.
\textbf{Đáp án của câu hỏi là B}
}
\end{ex}

\begin{ex} %[2D14V71]
Tính: $2 \times 7 + 5 \times 4$.
\choice
{$32$}
{$36$}
{\True $34$}
{$38$}
\loigiai{
$2 \times 7 = 14$; $5 \times 4 = 20$; $14 + 20 = 34$.
\textbf{Đáp án của câu hỏi là C}
}
\end{ex}

\begin{ex} %[2D14V71]
Một cuộn dây dài $20\text{ dm}$ được cắt thành các đoạn dài $2\text{ dm}$. Cắt được bao nhiêu đoạn dây?
\choice
{\True $10$ đoạn}
{$8$ đoạn}
{$12$ đoạn}
{$5$ đoạn}
\loigiai{
$20 : 2 = 10$ (đoạn).
\textbf{Đáp án của câu hỏi là A}
}
\end{ex}

\begin{ex} %[2D14V71]
Điền dấu thích hợp: $5 \times 6 \dots 2 \times 10 + 10$.
\choice
{$>$}
{$<$}
{Không so sánh được}
{\True $=$}
\loigiai{
$5 \times 6 = 30$; $2 \times 10 + 10 = 30$. Hai vế bằng nhau.
\textbf{Đáp án của câu hỏi là D}
}
\end{ex}

\begin{ex} %[2D14V71]
Có $16$ học sinh xếp thành các hàng, mỗi hàng có $2$ bạn. Xếp được bao nhiêu hàng?
\choice
{$7$ hàng}
{\True $8$ hàng}
{$9$ hàng}
{$6$ hàng}
\loigiai{
$16 : 2 = 8$ (hàng).
\textbf{Đáp án của câu hỏi là B}
}
\end{ex}

\begin{ex} %[2D14V71]
Một người nuôi $5$ chuồng thỏ, mỗi chuồng có $4$ con thỏ. Người đó bán đi $6$ con thỏ. Còn lại bao nhiêu con thỏ?
\choice
{$12$ con}
{$16$ con}
{\True $14$ con}
{$20$ con}
\loigiai{
Tổng số thỏ: $5 \times 4 = 20$ con. Còn lại: $20 - 6 = 14$ con.
\textbf{Đáp án của câu hỏi là C}
}
\end{ex}

\begin{ex} %[2D14V71]
Tính: $0 : 5 + 2 \times 9$.
\choice
{\True $18$}
{$0$}
{$19$}
{$23$}
\loigiai{
$0 : 5 = 0$; $2 \times 9 = 18$; $0 + 18 = 18$.
\textbf{Đáp án của câu hỏi là A}
}
\end{ex}

\textbf{Phần 2. Trắc nghiệm đúng sai}
\setcounterpageref{ex}{0}

\begin{ex} %[2D14V71]
Xét tính đúng sai của các phép tính nhân, chia:
\choiceTF
{\True Số 0 nhân với bất kì số nào cũng bằng 0.}
{Số 0 chia cho số nào cũng bằng chính số đó.}
{\True 5 x 7 = 35 và 35 : 5 = 7.}
{\True Nếu tích là 20 và một thừa số là 2 thì thừa số kia là 10.}
\loigiai{
\begin{itemize}
    \item Ý a) Đúng: Đúng tính chất số 0 trong phép nhân.
    \item Ý b) Sai: 0 chia cho số khác 0 luôn bằng 0, không phải chính số đó.
    \item Ý c) Đúng: Quan hệ giữa phép nhân và phép chia. Đúng.
    \item Ý d) Đúng: 20 : 2 = 10. Đúng.
\end{itemize}
\begin{itemize}
    \item a) \textbf{Đúng}.
    \item b) \textbf{Sai}.
    \item c) \textbf{Đúng}.
    \item d) \textbf{Đúng}.
\end{itemize}
}
\end{ex}

\begin{ex} %[2D14V71]
Một nhóm gồm 6 bạn, mỗi bạn góp 5 quyển vở. Số vở này đem chia đều cho 3 tổ học sinh nghèo:
\choiceTF
{\True Tổng số vở các bạn góp là 30 quyển.}
{\True Mỗi tổ học sinh nghèo nhận được 10 quyển vở.}
{Nếu mỗi bạn góp thêm 1 quyển thì tổng số vở là 32 quyển.}
{\True Số vở của 2 bạn bằng 1/3 tổng số vở cả nhóm.}
\loigiai{
\begin{itemize}
    \item Ý a) Đúng: 6 x 5 = 30 quyển. Đúng.
    \item Ý b) Đúng: 30 : 3 = 10 quyển. Đúng.
    \item Ý c) Sai: 6 x 6 = 36 quyển, không phải 32. Sai.
    \item Ý d) Đúng: 2 x 5 = 10 quyển = 30 : 3. Đúng.
\end{itemize}
\begin{itemize}
    \item a) \textbf{Đúng}.
    \item b) \textbf{Đúng}.
    \item c) \textbf{Sai}.
    \item d) \textbf{Đúng}.
\end{itemize}
}
\end{ex}

\textbf{Phần 3. Trả lời ngắn}
\setcounterpageref{ex}{0}

\begin{ex} %[2D14V71]
Tính giá trị: $5 \times 9 - 18$.
\shortans{27}
\loigiai{
\textbf{Đáp số:} $27$\\
$45 - 18 = 27$.
}
\end{ex}

\begin{ex} %[2D14V71]
Tìm $x$ biết: $x : 2 = 8$.
\shortans{16}
\loigiai{
\textbf{Đáp số:} $16$\\
$x = 8 \times 2 = 16$.
}
\end{ex}

\begin{ex} %[2D14V71]
Có 40 quả cam xếp đều vào 5 đĩa. Hỏi mỗi đĩa có bao nhiêu quả cam?
\shortans{8}
\loigiai{
\textbf{Đáp số:} $8$\\
$40 : 5 = 8$ quả.
}
\end{ex}

\begin{ex} %[2D14V71]
Mỗi can chứa 5 lít dầu. Hỏi 7 can như thế chứa bao nhiêu lít dầu?
\shortans{35}
\loigiai{
\textbf{Đáp số:} $35$\\
$5 \times 7 = 35$ lít.
}
\end{ex}

\begin{ex} %[2D14V71]
Một sợi dây dài 18 dm cắt thành các đoạn dài 2 dm. Hỏi cắt được bao nhiêu đoạn?
\shortans{9}
\loigiai{
\textbf{Đáp số:} $9$\\
$18 : 2 = 9$ đoạn.
}
\end{ex}

\begin{ex} %[2D14V71]
Tính: $2 \times 8 : 4$.
\shortans{4}
\loigiai{
\textbf{Đáp số:} $4$\\
$16 : 4 = 4$.
}
\end{ex}

\newpage

\TaoPhieuKiemTra{Chủ đề 14 -- Ôn tập cuối năm}{72}{Ôn tập hình học}{https://loptoanthaylinh.vn/thi/chuyen-de/2/toan2_cd_b72}
\textbf{Phần 1. Trắc nghiệm lựa chọn}
\setcounterpageref{ex}{0}

\begin{ex} %[2D14V72]
Hình tam giác có độ dài ba cạnh lần lượt là $15\text{ cm}, 20\text{ cm}, 25\text{ cm}$. Chu vi hình tam giác đó là
\choice
{$50\text{ cm}$}
{$55\text{ cm}$}
{\True $60\text{ cm}$}
{$65\text{ cm}$}
\loigiai{
Chu vi: $15 + 20 + 25 = 60$ cm.
\textbf{Đáp án của câu hỏi là C}
}
\end{ex}

\begin{ex} %[2D14V72]
Hình tứ giác có độ dài bốn cạnh đều bằng $5\text{ cm}$. Chu vi của hình tứ giác đó là
\choice
{$15\text{ cm}$}
{\True $20\text{ cm}$}
{$25\text{ cm}$}
{$10\text{ cm}$}
\loigiai{
Chu vi: $5 \times 4 = 20$ cm.
\textbf{Đáp án của câu hỏi là B}
}
\end{ex}

\begin{ex} %[2D14V72]
Đường gấp khúc gồm $3$ đoạn thẳng có độ dài lần lượt là $12\text{ cm}, 18\text{ cm}, 30\text{ cm}$. Độ dài đường gấp khúc đó bằng bao nhiêu đề-xi-mét?
\choice
{$5\text{ dm}$}
{$60\text{ dm}$}
{$7\text{ dm}$}
{\True $6\text{ dm}$}
\loigiai{
Độ dài: $12 + 18 + 30 = 60\text{ cm} = 6\text{ dm}$.
\textbf{Đáp án của câu hỏi là D}
}
\end{ex}

\begin{ex} %[2D14V72]
Hình nào dưới đây có thể lăn được về mọi hướng?
\choice
{\True Khối cầu}
{Khối trụ}
{Khối lập phương}
{Khối hộp chữ nhật}
\loigiai{
Khối cầu có bề mặt cong hoàn toàn nên lăn được về mọi hướng.
\textbf{Đáp án của câu hỏi là A}
}
\end{ex}

\begin{ex} %[2D14V72]
Khối trụ có hai mặt đáy là hình gì?
\choice
{Hình vuông}
{Hình tam giác}
{\True Hình tròn}
{Hình chữ nhật}
\loigiai{
Hai mặt đáy của khối trụ là hai hình tròn bằng nhau.
\textbf{Đáp án của câu hỏi là C}
}
\end{ex}

\begin{ex} %[2D14V72]
Cho ba điểm $A, B, C$ thẳng hàng, điểm $B$ nằm giữa $A$ và $C$. Biết $AB = 8\text{ cm}, BC = 12\text{ cm}$. Độ dài đoạn thẳng $AC$ là
\choice
{$4\text{ cm}$}
{\True $20\text{ cm}$}
{$16\text{ cm}$}
{$22\text{ cm}$}
\loigiai{
$AC = AB + BC = 8 + 12 = 20$ cm.
\textbf{Đáp án của câu hỏi là B}
}
\end{ex}

\begin{ex} %[2D14V72]
Một hình chữ nhật có chiều dài $16\text{ cm}$, chiều rộng ngắn hơn chiều dài $6\text{ cm}$. Nửa chu vi hình chữ nhật đó là
\choice
{$22\text{ cm}$}
{$52\text{ cm}$}
{$20\text{ cm}$}
{\True $26\text{ cm}$}
\loigiai{
Chiều rộng: $16 - 6 = 10$ cm. Nửa chu vi: $16 + 10 = 26$ cm.
\textbf{Đáp án của câu hỏi là D}
}
\end{ex}

\begin{ex} %[2D14V72]
Một đoạn dây đồng dài $32\text{ cm}$ được uốn thành một hình tứ giác có bốn cạnh bằng nhau. Mỗi cạnh của hình tứ giác dài bao nhiêu xăng-ti-mét?
\choice
{\True $8\text{ cm}$}
{$6\text{ cm}$}
{$7\text{ cm}$}
{$9\text{ cm}$}
\loigiai{
$32 : 4 = 8$ cm.
\textbf{Đáp án của câu hỏi là A}
}
\end{ex}

\begin{ex} %[2D14V72]
Khối hộp chữ nhật có tất cả bao nhiêu mặt?
\choice
{$4$ mặt}
{$8$ mặt}
{\True $6$ mặt}
{$12$ mặt}
\loigiai{
Khối hộp chữ nhật có 6 mặt.
\textbf{Đáp án của câu hỏi là C}
}
\end{ex}

\begin{ex} %[2D14V72]
Một hình tam giác có chu vi $45\text{ cm}$, hai cạnh dài $15\text{ cm}$ và $18\text{ cm}$. Độ dài cạnh còn lại là
\choice
{$10\text{ cm}$}
{\True $12\text{ cm}$}
{$14\text{ cm}$}
{$13\text{ cm}$}
\loigiai{
Cạnh thứ ba: $45 - (15 + 18) = 12$ cm.
\textbf{Đáp án của câu hỏi là B}
}
\end{ex}

\begin{ex} %[2D14V72]
Vật nào sau đây có dạng khối trụ?
\choice
{Hộp diêm}
{Quả bóng đá}
{Viên gạch}
{\True Lon sữa bò}
\loigiai{
Lon sữa bò có dạng khối trụ.
\textbf{Đáp án của câu hỏi là D}
}
\end{ex}

\begin{ex} %[2D14V72]
Cho $4$ điểm phân biệt trong đó không có ba điểm nào thẳng hàng. Nối từng cặp điểm thì vẽ được bao nhiêu đoạn thẳng?
\choice
{\True $6$ đoạn thẳng}
{$4$ đoạn thẳng}
{$5$ đoạn thẳng}
{$8$ đoạn thẳng}
\loigiai{
Số đoạn thẳng: $4 \times 3 : 2 = 6$ đoạn thẳng.
\textbf{Đáp án của câu hỏi là A}
}
\end{ex}

\textbf{Phần 2. Trắc nghiệm đúng sai}
\setcounterpageref{ex}{0}

\begin{ex} %[2D14V72]
Xét tính đúng sai của các khẳng định về hình học:
\choiceTF
{\True Đường thẳng không bị giới hạn ở hai đầu.}
{Hình tứ giác là hình có 3 cạnh và 3 đỉnh.}
{\True Chu vi của một hình bằng tổng độ dài các cạnh bao quanh hình đó.}
{\True Khối trụ có thể lăn được khi đặt nằm ngang.}
\loigiai{
\begin{itemize}
    \item Ý a) Đúng: Đúng định nghĩa đường thẳng.
    \item Ý b) Sai: Hình tứ giác có 4 cạnh và 4 đỉnh. Sai.
    \item Ý c) Đúng: Đúng quy tắc tính chu vi.
    \item Ý d) Đúng: Đúng đặc điểm khối trụ.
\end{itemize}
\begin{itemize}
    \item a) \textbf{Đúng}.
    \item b) \textbf{Sai}.
    \item c) \textbf{Đúng}.
    \item d) \textbf{Đúng}.
\end{itemize}
}
\end{ex}

\begin{ex} %[2D14V72]
Một hình tam giác có độ dài các cạnh là 12 cm, 16 cm và 20 cm:
\choiceTF
{\True Chu vi của hình tam giác là 48 cm.}
{\True Cạnh dài nhất dài hơn cạnh ngắn nhất 8 cm.}
{Chu vi của tam giác lớn hơn 5 dm.}
{\True Nếu uốn một sợi dây thành tam giác này thì cần sợi dây dài ít nhất 48 cm.}
\loigiai{
\begin{itemize}
    \item Ý a) Đúng: 12 + 16 + 20 = 48 cm. Đúng.
    \item Ý b) Đúng: 20 - 12 = 8 cm. Đúng.
    \item Ý c) Sai: 48 cm < 50 cm = 5 dm. Sai.
    \item Ý d) Đúng: Đúng.
\end{itemize}
\begin{itemize}
    \item a) \textbf{Đúng}.
    \item b) \textbf{Đúng}.
    \item c) \textbf{Sai}.
    \item d) \textbf{Đúng}.
\end{itemize}
}
\end{ex}

\textbf{Phần 3. Trả lời ngắn}
\setcounterpageref{ex}{0}

\begin{ex} %[2D14V72]
Tính chu vi hình tam giác có độ dài ba cạnh đều bằng 9 cm.
\shortans{27}
\loigiai{
\textbf{Đáp số:} $27$\\
$9 \times 3 = 27$ cm.
}
\end{ex}

\begin{ex} %[2D14V72]
Một hình tứ giác có bốn cạnh dài lần lượt là 10 cm, 12 cm, 15 cm, 13 cm. Chu vi của tứ giác là bao nhiêu đề-xi-mét?
\shortans{5}
\loigiai{
\textbf{Đáp số:} $5$\\
$10 + 12 + 15 + 13 = 50\text{ cm} = 5$ dm.
}
\end{ex}

\begin{ex} %[2D14V72]
Một đường gấp khúc gồm 4 đoạn thẳng bằng nhau, mỗi đoạn dài 5 cm. Độ dài đường gấp khúc là bao nhiêu xăng-ti-mét?
\shortans{20}
\loigiai{
\textbf{Đáp số:} $20$\\
$5 \times 4 = 20$ cm.
}
\end{ex}

\begin{ex} %[2D14V72]
Khối lập phương có bao nhiêu đỉnh?
\shortans{8}
\loigiai{
\textbf{Đáp số:} $8$\\
Khối lập phương có 8 đỉnh.
}
\end{ex}

\begin{ex} %[2D14V72]
Một mảnh đất hình chữ nhật có chiều dài 25 m, chiều rộng 15 m. Chu vi mảnh đất là bao nhiêu mét?
\shortans{80}
\loigiai{
\textbf{Đáp số:} $80$\\
$(25 + 15) \times 2 = 80$ m.
}
\end{ex}

\begin{ex} %[2D14V72]
Đoạn thẳng AB dài 18 cm, điểm M là điểm chia đôi đoạn AB thành 2 đoạn bằng nhau. Độ dài đoạn AM là bao nhiêu xăng-ti-mét?
\shortans{9}
\loigiai{
\textbf{Đáp số:} $9$\\
$18 : 2 = 9$ cm.
}
\end{ex}

\newpage

\TaoPhieuKiemTra{Chủ đề 14 -- Ôn tập cuối năm}{73}{Ôn tập đo lường}{https://loptoanthaylinh.vn/thi/chuyen-de/2/toan2_cd_b73}
\textbf{Phần 1. Trắc nghiệm lựa chọn}
\setcounterpageref{ex}{0}

\begin{ex} %[2D14V73]
Một ngày có bao nhiêu giờ?
\choice
{$12$ giờ}
{\True $24$ giờ}
{$20$ giờ}
{$60$ giờ}
\loigiai{
Một ngày có đúng 24 giờ.
\textbf{Đáp án của câu hỏi là B}
}
\end{ex}

\begin{ex} %[2D14V73]
Tháng nào sau đây có đúng $30$ ngày?
\choice
{\True Tháng 4}
{Tháng 1}
{Tháng 3}
{Tháng 5}
\loigiai{
Các tháng có 30 ngày: Tháng 4, 6, 9, 11.
\textbf{Đáp án của câu hỏi là A}
}
\end{ex}

\begin{ex} %[2D14V73]
Đổi: $1\text{ m } 8\text{ dm} = \dots\text{ dm}$.
\choice
{$108\text{ dm}$}
{$81\text{ dm}$}
{\True $18\text{ dm}$}
{$9\text{ dm}$}
\loigiai{
$1\text{ m} = 10\text{ dm} \Rightarrow 1\text{ m } 8\text{ dm} = 18$ dm.
\textbf{Đáp án của câu hỏi là C}
}
\end{ex}

\begin{ex} %[2D14V73]
Mẹ mua một gói mì chính cân nặng $500\text{ g}$ và một gói đường cân nặng $500\text{ g}$. Cả hai gói cân nặng bao nhiêu ki-lô-gam?
\choice
{$2\text{ kg}$}
{$10\text{ kg}$}
{$5\text{ kg}$}
{\True $1\text{ kg}$}
\loigiai{
$500\text{ g} + 500\text{ g} = 1 000\text{ g} = 1\text{ kg}$.
\textbf{Đáp án của câu hỏi là D}
}
\end{ex}

\begin{ex} %[2D14V73]
Đồng hồ chỉ 8 giờ 15 phút. Kim phút đang chỉ vào số nào?
\choice
{Số 8}
{\True Số 3}
{Số 6}
{Số 12}
\loigiai{
Kim phút chỉ số 3 là 15 phút.
\textbf{Đáp án của câu hỏi là B}
}
\end{ex}

\begin{ex} %[2D14V73]
Có 3 tờ 200 đồng và 1 tờ 500 đồng. Tổng số tiền là
\choice
{\True $1 100\text{ đồng}$}
{$700\text{ đồng}$}
{$900\text{ đồng}$}
{$1 000\text{ đồng}$}
\loigiai{
$200 \times 3 + 500 = 600 + 500 = 1 100$ đồng.
\textbf{Đáp án của câu hỏi là A}
}
\end{ex}

\begin{ex} %[2D14V73]
Một can chứa 10 lít dầu. Người ta rót vào 2 can nhỏ, mỗi can 4 lít. Trong can còn lại bao nhiêu lít dầu?
\choice
{$3$ lít}
{$4$ lít}
{\True $2$ lít}
{$1$ lít}
\loigiai{
Đã rót: $4 \times 2 = 8$ l. Còn lại: $10 - 8 = 2$ lít.
\textbf{Đáp án của câu hỏi là C}
}
\end{ex}

\begin{ex} %[2D14V73]
Điền dấu thích hợp: $1\text{ km} \dots 950\text{ m} + 50\text{ m}$.
\choice
{$>$}
{$<$}
{Không so sánh được}
{\True $=$}
\loigiai{
$950\text{ m} + 50\text{ m} = 1 000\text{ m} = 1\text{ km}$. Hai vế bằng nhau.
\textbf{Đáp án của câu hỏi là D}
}
\end{ex}

\begin{ex} %[2D14V73]
Hôm nay là ngày 20 tháng 5. Còn bao nhiêu ngày nữa thì hết tháng 5?
\choice
{$10$ ngày}
{\True $11$ ngày}
{$12$ ngày}
{$9$ ngày}
\loigiai{
Tháng 5 có 31 ngày. Còn: $31 - 20 = 11$ ngày.
\textbf{Đáp án của câu hỏi là B}
}
\end{ex}

\begin{ex} %[2D14V73]
Nam bắt đầu học bài lúc 8 giờ tối và học xong lúc 8 giờ 45 phút tối. Nam đã học trong thời gian bao lâu?
\choice
{\True $45$ phút}
{$40$ phút}
{$50$ phút}
{$1$ giờ}
\loigiai{
Thời gian học là 45 phút.
\textbf{Đáp án của câu hỏi là A}
}
\end{ex}

\begin{ex} %[2D14V73]
Con gà nặng $2\text{ kg}$, con ngỗng nặng hơn con gà $3\text{ kg}$. Cả hai con cân nặng bao nhiêu ki-lô-gam?
\choice
{$5\text{ kg}$}
{$6\text{ kg}$}
{\True $7\text{ kg}$}
{$8\text{ kg}$}
\loigiai{
Con ngỗng: $2 + 3 = 5$ kg. Cả hai con: $2 + 5 = 7$ kg.
\textbf{Đáp án của câu hỏi là C}
}
\end{ex}

\begin{ex} %[2D14V73]
Để có đúng $1 000\text{ đồng}$ từ các tờ tiền $200\text{ đồng}$, ta cần bao nhiêu tờ tiền?
\choice
{$4$ tờ}
{$6$ tờ}
{$10$ tờ}
{\True $5$ tờ}
\loigiai{
$1 000 : 200 = 5$ tờ.
\textbf{Đáp án của câu hỏi là D}
}
\end{ex}

\textbf{Phần 2. Trắc nghiệm đúng sai}
\setcounterpageref{ex}{0}

\begin{ex} %[2D14V73]
Xét tính đúng sai của các đơn vị và số đo đại lượng:
\choiceTF
{\True 1 m = 10 dm = 100 cm.}
{Tháng 2 luôn luôn có 30 ngày.}
{\True 1 kg nặng hơn 900 g.}
{\True 5 tờ 100 đồng có giá trị bằng 1 tờ 500 đồng.}
\loigiai{
\begin{itemize}
    \item Ý a) Đúng: Đúng.
    \item Ý b) Sai: Tháng 2 chỉ có 28 hoặc 29 ngày. Sai.
    \item Ý c) Đúng: 1 kg = 1 000 g > 900 g. Đúng.
    \item Ý d) Đúng: Đúng, đều bằng 500 đồng.
\end{itemize}
\begin{itemize}
    \item a) \textbf{Đúng}.
    \item b) \textbf{Sai}.
    \item c) \textbf{Đúng}.
    \item d) \textbf{Đúng}.
\end{itemize}
}
\end{ex}

\begin{ex} %[2D14V73]
Một chuyến xe buýt khởi hành lúc 7 giờ sáng và đến nơi lúc 8 giờ 30 phút sáng cùng ngày:
\choiceTF
{\True Chuyến đi kéo dài 1 giờ 30 phút.}
{\True 1 giờ 30 phút bằng 90 phút.}
{Xe đến nơi vào buổi chiều.}
{\True Nếu xe khởi hành muộn 15 phút thì đến nơi lúc 8 giờ 45 phút.}
\loigiai{
\begin{itemize}
    \item Ý a) Đúng: Đúng.
    \item Ý b) Đúng: 60 + 30 = 90 phút. Đúng.
    \item Ý c) Sai: Xe đến lúc 8 giờ 30 sáng. Sai.
    \item Ý d) Đúng: Đúng.
\end{itemize}
\begin{itemize}
    \item a) \textbf{Đúng}.
    \item b) \textbf{Đúng}.
    \item c) \textbf{Sai}.
    \item d) \textbf{Đúng}.
\end{itemize}
}
\end{ex}

\textbf{Phần 3. Trả lời ngắn}
\setcounterpageref{ex}{0}

\begin{ex} %[2D14V73]
Một thùng sơn nặng 18 kg. Bác thợ dùng hết 5 kg. Thùng sơn còn lại nặng bao nhiêu ki-lô-gam?
\shortans{13}
\loigiai{
\textbf{Đáp số:} $13$\\
$18 - 5 = 13$ kg.
}
\end{ex}

\begin{ex} %[2D14V73]
Một bể nước chứa 50 lít nước. Sau khi dùng hết 32 lít, người ta đổ thêm vào 15 lít nữa. Trong bể có bao nhiêu lít nước?
\shortans{33}
\loigiai{
\textbf{Đáp số:} $33$\\
$50 - 32 + 15 = 33$ lít.
}
\end{ex}

\begin{ex} %[2D14V73]
Bác Nam có 1 tờ 1 000 đồng, bác mua vé xe buýt hết 700 đồng. Bác còn lại bao nhiêu đồng?
\shortans{300}
\loigiai{
\textbf{Đáp số:} $300$\\
$1 000 - 700 = 300$ đồng.
}
\end{ex}

\begin{ex} %[2D14V73]
Đoạn đường từ nhà đến trường dài 500 m, Nam đi học cả đi và về hết bao nhiêu ki-lô-mét?
\shortans{1}
\loigiai{
\textbf{Đáp số:} $1$\\
$500 + 500 = 1 000\text{ m} = 1$ km.
}
\end{ex}

\begin{ex} %[2D14V73]
Tháng 12 có bao nhiêu ngày?
\shortans{31}
\loigiai{
\textbf{Đáp số:} $31$\\
Tháng 12 có 31 ngày.
}
\end{ex}

\begin{ex} %[2D14V73]
Một hình vuông có cạnh dài 2 dm. Chu vi hình vuông đó là bao nhiêu xăng-ti-mét?
\shortans{80}
\loigiai{
\textbf{Đáp số:} $80$\\
Cạnh = 20 cm. Chu vi: $20 \times 4 = 80$ cm.
}
\end{ex}

\newpage

\TaoPhieuKiemTra{Chủ đề 14 -- Ôn tập cuối năm}{74}{Ôn tập yếu tố thống kê, xác suất}{https://loptoanthaylinh.vn/thi/chuyen-de/2/toan2_cd_b74}
\textbf{Phần 1. Trắc nghiệm lựa chọn}
\setcounterpageref{ex}{0}

\begin{ex} %[2D14V74]
Quan sát biểu đồ tranh số cây trồng được của ba tổ: Tổ 1 có 4 cây, tổ 2 có 6 cây, tổ 3 có 5 cây (mỗi hình vẽ đại diện cho 2 cây thật). Cả ba tổ trồng được bao nhiêu cây thật?
\choice
{$15$ cây}
{$25$ cây}
{\True $30$ cây}
{$20$ cây}
\loigiai{
Tổng số hình: $4 + 6 + 5 = 15$ hình. Số cây thật: $15 \times 2 = 30$ cây.
\textbf{Đáp án của câu hỏi là C}
}
\end{ex}

\begin{ex} %[2D14V74]
Trong một hộp có 5 viên bi xanh và 5 viên bi đỏ. Khẳng định nào sau đây là 'Có thể xảy ra' khi lấy ra 1 viên bi?
\choice
{\True Lấy được viên bi màu đỏ}
{Lấy được viên bi màu vàng}
{Không lấy được viên bi nào}
{Chắc chắn lấy được bi xanh}
\loigiai{
Có thể lấy được viên bi màu đỏ hoặc màu xanh.
\textbf{Đáp án của câu hỏi là A}
}
\end{ex}

\begin{ex} %[2D14V74]
Tung một con xúc xắc có 6 mặt. Khả năng xuất hiện mặt 7 chấm là
\choice
{Có thể xảy ra}
{Chắc chắn xảy ra}
{Thỉnh thoảng xảy ra}
{\True Không thể xảy ra}
\loigiai{
Xúc xắc chỉ có từ 1 đến 6 chấm nên không thể có mặt 7 chấm.
\textbf{Đáp án của câu hỏi là D}
}
\end{ex}

\begin{ex} %[2D14V74]
Khi đếm số lượng học sinh thích ăn các loại kem: Dâu có 12 bạn, sô-cô-la có 18 bạn, vani có 10 bạn. Loại kem nào được yêu thích nhất?
\choice
{Kem dâu}
{\True Kem sô-cô-la}
{Kem vani}
{Cả ba bằng nhau}
\loigiai{
Kem sô-cô-la có 18 bạn chọn, nhiều nhất.
\textbf{Đáp án của câu hỏi là B}
}
\end{ex}

\begin{ex} %[2D14V74]
Trong túi chỉ có các viên kẹo cam. Biến cố 'Lấy được một viên kẹo cam' là sự kiện
\choice
{Có thể xảy ra}
{Không thể xảy ra}
{\True Chắc chắn xảy ra}
{Chưa chắc xảy ra}
\loigiai{
Túi chỉ có kẹo cam nên chắc chắn lấy được kẹo cam.
\textbf{Đáp án của câu hỏi là C}
}
\end{ex}

\begin{ex} %[2D14V74]
Trên biểu đồ tranh, mỗi ngôi sao biểu thị 5 điểm tốt. Bạn Mai có 4 ngôi sao, bạn Lan có 6 ngôi sao. Lan nhiều hơn Mai bao nhiêu điểm tốt?
\choice
{\True $10$ điểm}
{$2$ điểm}
{$5$ điểm}
{$15$ điểm}
\loigiai{
Lan nhiều hơn: $(6 - 4) \times 5 = 10$ điểm tốt.
\textbf{Đáp án của câu hỏi là A}
}
\end{ex}

\begin{ex} %[2D14V74]
Để biết số lượng các loại xe trong bãi giữ xe, người ta tiến hành công việc nào sau đây?
\choice
{Cân tải trọng các xe}
{Đo chiều dài bãi xe}
{Sơn màu lại các xe}
{\True Kiểm đếm và phân loại}
\loigiai{
Công việc phù hợp là kiểm đếm và phân loại.
\textbf{Đáp án của câu hỏi là D}
}
\end{ex}

\begin{ex} %[2D14V74]
Gieo đồng xu 10 lần, có 4 lần xuất hiện mặt Sấp. Số lần xuất hiện mặt Ngửa là
\choice
{$5$ lần}
{\True $6$ lần}
{$7$ lần}
{$4$ lần}
\loigiai{
$10 - 4 = 6$ (lần).
\textbf{Đáp án của câu hỏi là B}
}
\end{ex}

\begin{ex} %[2D14V74]
Một năm có 13 tháng là sự kiện
\choice
{Có thể xảy ra}
{Chắc chắn xảy ra}
{\True Không thể xảy ra}
{Xảy ra vào năm nhuận}
\loigiai{
Một năm luôn chỉ có 12 tháng, không thể có 13 tháng.
\textbf{Đáp án của câu hỏi là C}
}
\end{ex}

\begin{ex} %[2D14V74]
Trong hộp có 3 quả bóng xanh và 1 quả bóng vàng. Cần lấy ít nhất bao nhiêu quả bóng để chắc chắn có 1 quả bóng màu xanh?
\choice
{\True $2$ quả}
{$1$ quả}
{$3$ quả}
{$4$ quả}
\loigiai{
Trường hợp xấu nhất lấy phải 1 bóng vàng, lấy thêm 1 quả nữa chắc chắn là bóng xanh. Cần $1 + 1 = 2$ quả.
\textbf{Đáp án của câu hỏi là A}
}
\end{ex}

\begin{ex} %[2D14V74]
Bảng theo dõi số ngày mưa trong tháng: Tuần 1 có 2 ngày, tuần 2 có 3 ngày, tuần 3 có 1 ngày, tuần 4 có 2 ngày. Cả tháng có bao nhiêu ngày mưa?
\choice
{$7$ ngày}
{$9$ ngày}
{$10$ ngày}
{\True $8$ ngày}
\loigiai{
$2 + 3 + 1 + 2 = 8$ (ngày).
\textbf{Đáp án của câu hỏi là D}
}
\end{ex}

\begin{ex} %[2D14V74]
Hôm nay là Chủ nhật. Ngày mai là thứ Bảy là sự kiện
\choice
{Có thể xảy ra}
{\True Không thể xảy ra}
{Chắc chắn xảy ra}
{Đôi khi xảy ra}
\loigiai{
Sau Chủ nhật là thứ Hai, không thể là thứ Bảy.
\textbf{Đáp án của câu hỏi là B}
}
\end{ex}

\textbf{Phần 2. Trắc nghiệm đúng sai}
\setcounterpageref{ex}{0}

\begin{ex} %[2D14V74]
Xét tính đúng sai của các nhận định về thống kê và khả năng xảy ra:
\choiceTF
{\True Biểu đồ tranh giúp người xem dễ dàng nhận biết và so sánh số liệu.}
{Một sự kiện chắc chắn xảy ra thì vẫn có lúc không xảy ra.}
{\True Trong một hộp có 2 viên bi đỏ, lấy ra 1 viên thì không thể được bi xanh.}
{\True Có thể dùng các vạch kiểm đếm để ghi chép số liệu nhanh chóng.}
\loigiai{
\begin{itemize}
    \item Ý a) Đúng: Đúng tác dụng của biểu đồ.
    \item Ý b) Sai: Chắc chắn nghĩa là luôn luôn xảy ra 100%. Sai.
    \item Ý c) Đúng: Đúng vì không có bi xanh.
    \item Ý d) Đúng: Đúng phương pháp.
\end{itemize}
\begin{itemize}
    \item a) \textbf{Đúng}.
    \item b) \textbf{Sai}.
    \item c) \textbf{Đúng}.
    \item d) \textbf{Đúng}.
\end{itemize}
}
\end{ex}

\begin{ex} %[2D14V74]
Kiểm đếm các loại hoa trong vườn: 15 hoa hồng, 20 hoa cúc, 10 hoa đồng tiền:
\choiceTF
{\True Tổng số hoa là 45 bông.}
{\True Hoa cúc có số lượng nhiều nhất.}
{Hoa hồng nhiều hơn hoa cúc.}
{\True Hoa đồng tiền ít hơn hoa cúc 10 bông.}
\loigiai{
\begin{itemize}
    \item Ý a) Đúng: 15 + 20 + 10 = 45. Đúng.
    \item Ý b) Đúng: 20 > 15 > 10. Đúng.
    \item Ý c) Sai: 15 < 20 nên ít hơn. Sai.
    \item Ý d) Đúng: 20 - 10 = 10. Đúng.
\end{itemize}
\begin{itemize}
    \item a) \textbf{Đúng}.
    \item b) \textbf{Đúng}.
    \item c) \textbf{Sai}.
    \item d) \textbf{Đúng}.
\end{itemize}
}
\end{ex}

\textbf{Phần 3. Trả lời ngắn}
\setcounterpageref{ex}{0}

\begin{ex} %[2D14V74]
Trên biểu đồ tranh, mỗi hình quả bóng đại diện cho 5 quả bóng thật. Cửa hàng vẽ 7 hình quả bóng. Cửa hàng có bao nhiêu quả bóng thật?
\shortans{35}
\loigiai{
\textbf{Đáp số:} $35$\\
$5 \times 7 = 35$ quả.
}
\end{ex}

\begin{ex} %[2D14V74]
Trong hộp có 8 viên bi trắng và 2 viên bi đen. Cần lấy ít nhất mấy viên bi để chắc chắn có 1 viên bi màu trắng?
\shortans{3}
\loigiai{
\textbf{Đáp số:} $3$\\
Lấy hết 2 bi đen, cần thêm 1 viên: $2 + 1 = 3$ viên.
}
\end{ex}

\begin{ex} %[2D14V74]
Khảo sát sở thích thể thao của 50 bạn: 25 bạn thích bóng đá, 15 bạn thích cầu lông, còn lại thích bơi. Có bao nhiêu bạn thích bơi?
\shortans{10}
\loigiai{
\textbf{Đáp số:} $10$\\
$50 - 25 - 15 = 10$ bạn.
}
\end{ex}

\begin{ex} %[2D14V74]
Tung đồng xu 20 lần ghi nhận 12 lần mặt Sấp. Có bao nhiêu lần mặt Ngửa?
\shortans{8}
\loigiai{
\textbf{Đáp số:} $8$\\
$20 - 12 = 8$ lần.
}
\end{ex}

\begin{ex} %[2D14V74]
Có 4 tổ học sinh, mỗi tổ trồng được 6 cây xanh. Cả 4 tổ trồng được bao nhiêu cây?
\shortans{24}
\loigiai{
\textbf{Đáp số:} $24$\\
$6 \times 4 = 24$ cây.
}
\end{ex}

\begin{ex} %[2D14V74]
Gieo 1 con xúc xắc, có mấy mặt có số chấm bé hơn 3?
\shortans{2}
\loigiai{
\textbf{Đáp số:} $2$\\
Mặt 1 và mặt 2, có 2 mặt.
}
\end{ex}

\newpage

\TaoPhieuKiemTra{Chủ đề 14 -- Ôn tập cuối năm}{75}{Ôn tập chung cuối năm}{https://loptoanthaylinh.vn/thi/chuyen-de/2/toan2_cd_b75}
\textbf{Phần 1. Trắc nghiệm lựa chọn}
\setcounterpageref{ex}{0}

\begin{ex} %[2D14V75]
Số gồm $8$ trăm, $5$ chục và $2$ đơn vị được viết là
\choice
{$825$}
{\True $852$}
{$582$}
{$258$}
\loigiai{
Viết là $852$.
\textbf{Đáp án của câu hỏi là B}
}
\end{ex}

\begin{ex} %[2D14V75]
Kết quả của biểu thức $345 + 255 - 180$ là
\choice
{$410$}
{$430$}
{$400$}
{\True $420$}
\loigiai{
$345 + 255 = 600$; $600 - 180 = 420$.
\textbf{Đáp án của câu hỏi là D}
}
\end{ex}

\begin{ex} %[2D14V75]
Tìm $x$ biết: $x \times 5 = 45$.
\choice
{\True $9$}
{$8$}
{$7$}
{$10$}
\loigiai{
$x = 45 : 5 = 9$.
\textbf{Đáp án của câu hỏi là A}
}
\end{ex}

\begin{ex} %[2D14V75]
Một mảnh vườn hình chữ nhật có chiều dài $24\text{ m}$, chiều rộng $16\text{ m}$. Chu vi mảnh vườn đó là
\choice
{$40\text{ m}$}
{$70\text{ m}$}
{\True $80\text{ m}$}
{$90\text{ m}$}
\loigiai{
Chu vi: $(24 + 16) \times 2 = 80$ m.
\textbf{Đáp án của câu hỏi là C}
}
\end{ex}

\begin{ex} %[2D14V75]
Số tròn trăm lớn nhất có ba chữ số đem trừ đi số chẵn bé nhất có ba chữ số thì được
\choice
{$798$}
{\True $800$}
{$899$}
{$802$}
\loigiai{
Số tròn trăm lớn nhất có 3 chữ số: 900. Số chẵn bé nhất có 3 chữ số: 100. Hiệu: $900 - 100 = 800$.
\textbf{Đáp án của câu hỏi là B}
}
\end{ex}

\begin{ex} %[2D14V75]
Một người có 1 tờ 1 000 đồng, mua 2 quyển vở giá 400 đồng một quyển. Người đó còn lại bao nhiêu đồng?
\choice
{$100\text{ đồng}$}
{$300\text{ đồng}$}
{$400\text{ đồng}$}
{\True $200\text{ đồng}$}
\loigiai{
Mua hết: $400 \times 2 = 800$ đồng. Còn lại: $1 000 - 800 = 200$ đồng.
\textbf{Đáp án của câu hỏi là D}
}
\end{ex}

\begin{ex} %[2D14V75]
Một sợi dây đồng dài $1\text{ m}$. Người ta cắt đi $35\text{ cm}$. Chiều dài đoạn dây còn lại là
\choice
{\True $65\text{ cm}$}
{$75\text{ cm}$}
{$55\text{ cm}$}
{$65\text{ dm}$}
\loigiai{
$1\text{ m} = 100\text{ cm}$. Còn lại: $100 - 35 = 65$ cm.
\textbf{Đáp án của câu hỏi là A}
}
\end{ex}

\begin{ex} %[2D14V75]
Điền dấu thích hợp: $5 \times 8 \dots 100 - 60$.
\choice
{$>$}
{$<$}
{\True $=$}
{Không so sánh được}
\loigiai{
$5 \times 8 = 40$; $100 - 60 = 40$. Hai vế bằng nhau.
\textbf{Đáp án của câu hỏi là C}
}
\end{ex}

\begin{ex} %[2D14V75]
Đường gấp khúc gồm 4 đoạn thẳng bằng nhau, mỗi đoạn dài 25 cm. Độ dài đường gấp khúc là bao nhiêu mét?
\choice
{$2\text{ m}$}
{\True $1\text{ m}$}
{$10\text{ m}$}
{$100\text{ m}$}
\loigiai{
$25 \times 4 = 100\text{ cm} = 1\text{ m}$.
\textbf{Đáp án của câu hỏi là B}
}
\end{ex}

\begin{ex} %[2D14V75]
Khối 2 có $280$ học sinh, khối 3 có nhiều hơn khối 2 là $45$ học sinh. Cả hai khối có tất cả bao nhiêu học sinh?
\choice
{$595$ học sinh}
{$615$ học sinh}
{$325$ học sinh}
{\True $605$ học sinh}
\loigiai{
Khối 3: $280 + 45 = 325$. Cả hai khối: $280 + 325 = 605$ học sinh.
\textbf{Đáp án của câu hỏi là D}
}
\end{ex}

\begin{ex} %[2D14V75]
Tìm một số biết rằng nếu lấy số đó chia cho 2 rồi cộng với 15 thì được 24.
\choice
{\True $18$}
{$16$}
{$20$}
{$22$}
\loigiai{
Số đó chia 2 được: $24 - 15 = 9$. Số đó là: $9 \times 2 = 18$.
\textbf{Đáp án của câu hỏi là A}
}
\end{ex}

\begin{ex} %[2D14V75]
Trong các khẳng định sau, khẳng định nào ĐÚNG?
\choice
{1 m = 10 cm}
{1 ngày = 12 giờ}
{\True 1 km = 1 000 m}
{100 đồng > 200 đồng}
\loigiai{
$1\text{ km} = 1 000\text{ m}$ là khẳng định chính xác.
\textbf{Đáp án của câu hỏi là C}
}
\end{ex}

\textbf{Phần 2. Trắc nghiệm đúng sai}
\setcounterpageref{ex}{0}

\begin{ex} %[2D14V75]
Xét các bài toán tổng hợp toàn diện chương trình Toán 2:
\choiceTF
{\True Số chẵn liền sau của 998 là 1 000.}
{Hình tam giác có 4 cạnh.}
{\True 500 ml + 500 ml = 1 l (hoặc 500 g + 500 g = 1 kg).}
{\True Phép nhân 2 x 9 có kết quả bằng phép nhân 9 x 2.}
\loigiai{
\begin{itemize}
    \item Ý a) Đúng: Đúng.
    \item Ý b) Sai: Hình tam giác có 3 cạnh. Sai.
    \item Ý c) Đúng: Đúng.
    \item Ý d) Đúng: Tính chất giao hoán. Đúng.
\end{itemize}
\begin{itemize}
    \item a) \textbf{Đúng}.
    \item b) \textbf{Sai}.
    \item c) \textbf{Đúng}.
    \item d) \textbf{Đúng}.
\end{itemize}
}
\end{ex}

\begin{ex} %[2D14V75]
Một kho hàng có 1 000 thùng sữa. Ngày đầu xuất đi 350 thùng, ngày thứ hai xuất đi 450 thùng:
\choiceTF
{\True Cả hai ngày xuất đi 800 thùng sữa.}
{\True Trong kho còn lại 200 thùng sữa.}
{Số sữa còn lại nhiều hơn số sữa đã xuất.}
{\True Nếu mỗi thùng xếp được 5 hộp sữa thì 200 thùng còn lại có 1 000 hộp sữa.}
\loigiai{
\begin{itemize}
    \item Ý a) Đúng: 350 + 450 = 800 thùng. Đúng.
    \item Ý b) Đúng: 1 000 - 800 = 200 thùng. Đúng.
    \item Ý c) Sai: 200 < 800 nên ít hơn. Sai.
    \item Ý d) Đúng: 200 x 5 = 1 000. Đúng.
\end{itemize}
\begin{itemize}
    \item a) \textbf{Đúng}.
    \item b) \textbf{Đúng}.
    \item c) \textbf{Sai}.
    \item d) \textbf{Đúng}.
\end{itemize}
}
\end{ex}

\textbf{Phần 3. Trả lời ngắn}
\setcounterpageref{ex}{0}

\begin{ex} %[2D14V75]
Tính giá trị của biểu thức: $425 + 375 - 200$.
\shortans{600}
\loigiai{
\textbf{Đáp số:} $600$\\
$800 - 200 = 600$.
}
\end{ex}

\begin{ex} %[2D14V75]
Tìm $x$ biết: $x : 5 + 12 = 18$.
\shortans{30}
\loigiai{
\textbf{Đáp số:} $30$\\
$x : 5 = 6 \Rightarrow x = 30$.
}
\end{ex}

\begin{ex} %[2D14V75]
Một hình tam giác có ba cạnh bằng nhau, chu vi bằng 24 cm. Mỗi cạnh dài bao nhiêu xăng-ti-mét?
\shortans{8}
\loigiai{
\textbf{Đáp số:} $8$\\
$24 : 3 = 8$ cm.
}
\end{ex}

\begin{ex} %[2D14V75]
An có 1 000 đồng, An mua một cây kem hết 800 đồng. An còn lại bao nhiêu đồng?
\shortans{200}
\loigiai{
\textbf{Đáp số:} $200$\\
$1 000 - 800 = 200$ đồng.
}
\end{ex}

\begin{ex} %[2D14V75]
Một xe khách chở 45 người. Đến trạm dừng có 12 người xuống và 7 người lên. Trên xe còn lại bao nhiêu người?
\shortans{40}
\loigiai{
\textbf{Đáp số:} $40$\\
$45 - 12 + 7 = 40$ người.
}
\end{ex}

\begin{ex} %[2D14V75]
Tính chu vi hình chữ nhật có chiều dài 30 cm và chiều rộng 20 cm.
\shortans{100}
\loigiai{
\textbf{Đáp số:} $100$\\
$(30 + 20) \times 2 = 100$ cm.
}
\end{ex}

\newpage

\end{document}
